% !TeX spellcheck = <none>

%\newcommand{\compileall}{def}


\documentclass[amstex,twoside]{ctexbook}


\usepackage{marginnote}

% 16开本
\usepackage[
		paperwidth=195mm,
		paperheight=271mm,
		%装订线宽10mm,页边距 28mm
		inner = 38mm,
		outer = 28mm,
		tmargin = 36mm,
		bmargin = 30mm,
		marginparwidth=2cm
	]{geometry}

\usepackage{needspace}
\usepackage{wrapfig}
\usepackage{float}
\usepackage[below,section]{placeins}
\usepackage{color}
\usepackage{tabularx}
\usepackage{longtable}
\usepackage{multirow}
\usepackage{rotating}
\usepackage{soul}
\usepackage{amsmath}
\usepackage{wasysym}
\usepackage{amsthm}
\usepackage{stmaryrd}
\usepackage{xltxtra}
% always load hyperref as last packages
\usepackage[colorlinks,linkcolor=black]{hyperref}
%\usepackage{todonotes}


%\usepackage{eso-pic}


\AddToShipoutPicture{%
\parbox[b][\paperheight]{\paperwidth}{%
	\vfill
\begin{center}
    \includegraphics[width=\textwidth]{pics/draft}\hfill
\end{center}
    \vfill
}}

%=====================================================================================
% font setting
%=====================================================================================
%=====================================================================================
% font setting
%=====================================================================================

\setCJKmonofont[Path="fonts/",BoldFont=AdobeKaitiStd-Regular.otf ]{AdobeFangsongStd-Regular.otf}

\setCJKfamilyfont{zhfs}[Path="fonts/", BoldFont=AdobeKaitiStd-Regular.otf]{AdobeFangsongStd-Regular.otf}


\setCJKmainfont[
		Path="fonts/",
		BoldFont = AdobeHeitiStd-Regular.otf,
		ItalicFont= AdobeKaitiStd-Regular.otf
	]{AdobeSongStd-Light.otf}

\setCJKfamilyfont{KaiTi}[Path="fonts/"]{AdobeKaitiStd-Regular.otf}
\newfontfamily\xingkai[Path="fonts/"]{xinkai.ttf}
\setCJKfamilyfont{xingkai}[Path="fonts/"]{xinkai.ttf}

%=====================================================================================
% define new command and environment 
%=====================================================================================
\renewcommand*{\marginfont}{\color{red}\scriptsize}
\newcommand*{\todo}[1]{\marginnote[%
\color{blue}\footnotesize%
TODO:#1%
]{%
\color{blue}\footnotesize%
TODO:#1}
}

\definecolor{gray}{cmyk}{0,0,0,0.2}

%=====================================================================================
% define new command and environment 
%=====================================================================================
\newenvironment{notice}{\tt}{}
\newenvironment{insertnote}{ \ttfamily\CJKfamily{KaiTi} }{\vskip 0.5cm }
\newenvironment{code}{\small\tt\begin{longtable}{p{0.8\textwidth}}}{\end{longtable}}

\newcommand{\email}[1]{\href{mailto:#1}{#1}}

\newcommand{\RTLpar}{% right-to-left paragraph alignment
  \leftskip=0pt plus .5fil%
  \rightskip=0pt plus -.5fil%
  \parfillskip=0pt plus .5fil%
}

\newenvironment{quotes}[2][0.55]{\pushQED{#2}%
\begin{flushright}%
\begin{minipage}{#1\textwidth}\begin{flushright}\noindent\it\RTLpar}{%
 \\------\popQED{}\end{flushright}\end{minipage}\end{flushright}\vskip 8mm }%

\newenvironment{filecontent}[1]{

\addtolength{\linewidth}{-\parindent}
\begin{minipage}{\linewidth}

\vskip 1ex

文件: #1

\vskip -2ex
\rule{\textwidth}{0.02pt}
\vskip -1ex

%\hline

}{	\end{minipage}}




\newcounter{lizi}[chapter]

\newenvironment{example}[1]{ \addtocounter{lizi}{1} \vskip 2.5ex \bf 例\hskip -0.05em\arabic{chapter}.\arabic{lizi}： #1%
\par %
 \tt\small\CJKfamily{zhfs}}{ \vskip 2.5ex  }

\newcommand{\theexample}{\arabic{chapter}.\arabic{lizi}}

\makeatletter\newcommand{\chatu}{\@ifstar%
                     \chatuStar%
                     \chatuNoStar%
}\makeatother

\newcommand{\chatuNoStar}[3][scale=0.35]{%
\begin{figure}[htbp]%
\noindent\centering%
\includegraphics[#1]{pics/#2}%
\caption{#3\label{fig:#2}}%
\end{figure}%
}


\newcommand{\chatuStar}[3][scale=0.35]{%
\begin{figure*}[H]%
\noindent\centering%
\includegraphics[#1]{pics/#2}%
\label{fig:#2}%
\end{figure*}%
}


\newcommand{\faqref}[1]{~附录 \ref{chap:FAQ}.\nameref{chap:FAQ}~“\nameref{#1}”}
\newcommand{\secref}[1]{ { \it节\ref{#1}~\nameref{#1}} }

\newcommand{\chapref}[1]{ { \it第\ref{#1}章~\nameref{#1}} }


\newcounter{faqs}

\makeatletter
\newcommand\@publisher{}
\newcommand{\publisher}[1]{
\renewcommand\@publisher{#1}
}

\newcommand\@company{}
\newcommand{\company}[1]{
\renewcommand\@company{#1}
}
\makeatother


%=====================================================================================
% style settings
%=====================================================================================

% 阿拉伯数字章节
\setcounter{chapter}{-1}
\CTEXsetup[number={\arabic{chapter}}]{chapter}
% 对齐设置
\tolerance=4500
\punctstyle{hangmobanjiao}
\CJKecglue{}
\CJKsetecglue{} %{\hskip 0.1em plus 0.05em minus 0.05em}

%=====================================================================================
% book info
%=====================================================================================
% book info
%=====================================================================================

\title{Linux枕边书}
\author{微菜}
\company{北京源智天下科技有限公司}
%\publisher{清华大学出版社}
\publisher{XX出版社}

%=====================================================================================


\begin{document}

%=========================================================
% 首页
%=========================================================
{
\thispagestyle{empty}
\makeatletter

\begin{center}


{\fontsize{2cm}{2.5cm}\selectfont \xingkai\CJKfamily{xingkai} \@title

 \vskip 2cm
 
\normalfont  \large  ------ Linux文化入门 
 
 }


\vskip 5cm

 {\LARGE
  \@author ~ 著 }


\end{center}

\null
\vfill



\begin{flushright}
本书由 \LaTeX + \tt{xeCJK} 排版。 \hfill \@publisher ~ 出版
\end{flushright}

\clearpage
\thispagestyle{empty}


\noindent\@author\\
\@title


\noindent Copyright~\textcopyright~2012-\the\year~All rights reversed.

\noindent 本书由 \@company 授权  \@publisher 出版发行。版权所有未经许可，不得已任何凡是复杂或者抄袭本书的任何部分。

\vfill

{
\noindent \bf版权所有，翻印必究
}

\vfill

\noindent 图书在版编目（CIP）数据

XXX

\vfill


\noindent
\hskip -1em
\begin{tabular*}{\textwidth}{>{\bf}ll>{\bf}ll}

出\hfill版\hfill者： &  \@publisher  & 地\hfill址: & 北京清华大学学究大厦  \\
&  \url{http://www.tup.com.cn} &  邮\hfill编: & 100084 \\
&  社总机：010-62770175 &  客户服务： & 010-62776969 \\

责任编辑： & \multicolumn{3}{l}{陈冠军} \\
印\hfill刷\hfill者： & \multicolumn{3}{l}{  XXX 印刷厂  } \\
装\hfill订\hfill者： & \multicolumn{3}{l}{XXX装订厂} \\
发\hfill行\hfill者： & \multicolumn{3}{l}{XXX} \\
开\hfill本： & \multicolumn{3}{l}{ 16开    } \\
版\hfill次: &  \multicolumn{3}{l}{  ~    } \\
书\hfill号: &  \multicolumn{3}{l}{ ISBN XXXX~    } \\
印\hfill数: &  \multicolumn{3}{l}{  ~    } \\
定\hfill价: &  \multicolumn{3}{l}{  ~~~~元    } \\
\hline
\end{tabular*}


\makeatother

%TODO 发行的时候解除这个
%\makeatletter
\thispagestyle{empty}


\noindent\@author\\
\@title


\noindent Copyright~\textcopyright~2012-\the\year~All rights reversed.

\noindent 本书由 \@company 授权  \@publisher 出版发行。版权所有未经许可，不得已任何凡是复杂或者抄袭本书的任何部分。

\vfill

{
\noindent \bf 版权所有，翻印必究
}

\vfill

\noindent 图书在版编目（CIP）数据

XXX

\vfill


\noindent
\hskip -1em
\begin{tabular*}{\textwidth}{>{\bf}ll>{\bf}ll}

出\hfill 版\hfill 者： &  \@publisher  & 地\hfill 址: & 北京清华大学学究大厦  \\
&  \url{http://www.tup.com.cn} &  邮\hfill 编: & 100084 \\
&  社总机：010-62770175 &  客户服务： & 010-62776969 \\

责任编辑： & \multicolumn{3}{l}{陈冠军} \\
印\hfill 刷\hfill 者： & \multicolumn{3}{l}{  XXX 印刷厂  } \\
装\hfill 订\hfill 者： & \multicolumn{3}{l}{XXX装订厂} \\
发\hfill 行\hfill 者： & \multicolumn{3}{l}{XXX} \\
开\hfill 本： & \multicolumn{3}{l}{ 16开    } \\
版\hfill 次： &  \multicolumn{3}{l}{  ~    } \\
书\hfill 号： &  \multicolumn{3}{l}{ ISBN XXXX~    } \\
印\hfill 数： & \multicolumn{3}{l}{  ~    } \\
定\hfill 价： &  \multicolumn{3}{l}{  ~~~~元    } \\
\hline
\end{tabular*}

\makeatother

\chapter*{致谢!}
\addcontentsline{toc}{chapter}{致谢}

\newpage\thispagestyle{empty}
}

%=========================================================
%目录
%=========================================================
\tableofcontents

%=========================================================
% 序言
%=========================================================
\chapter{序}
\begin{quotes}{George Santayana}
Those who cannot remember the past are condemned to repeat it.
\end{quotes}

\ifdefined\compileall


不知道从什么时候开始用上了Linux，已经不记得刚开始用Linux的时候那段辛苦并乐在其中的学习过程了。
有一天，我重新用起了曾经那么熟悉，那么“简单”的操作系统，结果大失所望。我看到的是一个完全反人类的系统。
它的操作似乎很简单，但是一直总有说不出的感觉，让我觉得它的操作非常别扭。我感觉我第一次接触到Linux的时候也是这种感觉。

有人说，这是一种习惯。而我却更乐意认为那是一种文化。Linux和Windows是一种截然不同的文化。
Linux技术在飞快的进化，各种新技术层出不穷。但是骨子深处的UNIX文化确没变。

有人说，Linux不是UNIX，BSD、Sun Solaris、IBM AIX这些才是UNIX。
没错，这些系统确实继承了AT\&T\footnote{第一个UNIX系统由贝尔实验室开发。贝尔实验室当时属于AT\&T公司。}
的代码\footnote{BSD和AT\&T打了几年官司，后经重写将AT\&T代码完全剔除了。}，还获得了
OpenGroup\footnote{UNIX走向分裂后，几个大公司成立了OpenGroup共同持有UNIX商标。}认证
。但是对我来说，只有Linux才真正在文化上继承了UNIX。
既然Linux的文化源自UNIX，要学习Linux将不得不首先了解UNIX，
了解UNIX从诞生到发展，发展到分裂，分裂到混乱，混乱中统一，统一却失去了UNIX的文化，再到借Linux重生的故事。

故事虽然不是技术，但是了解一个技术诞生的背景，能深刻的理解它。

融入UNIX文化后，再来学习Linux就得心应手了。

但是我不愿意独立一个章节介绍UNIX历史，因为这本书是带领大家学习Linux的，不是故事书，
所以我把整书的序部分用来介绍UNIX的历史文化。因为序比较长，干脆就叫第\thechapter章。

\section{为什么又写一本书}

我不记得自己为了学习Linux查阅了多少资料，翻烂了多少本书。最终翻了那么多书后发现，
现在讲解Linux入门书籍多数都讲求“速成”而采取填鸭式讲解。就算速成后成功学会Linux，最终也会因为不了解Linux的文化而处处碰壁。
然后到网上发牢骚抱怨Linux不人性化。

\begin{notice}
习惯于Windows的人无法习惯于Linux的操作。所有与Windows不一致的操作都被认为“不人性化”。
\end{notice}

必须首先要让读者了解到UNIX文化的精髓，而后开始学习Linux，抛开一切对Windows的比较。这样学会的Linux，遇事首先会以UNIX的思维模式思考，然后提出非常UNIXy\footnote{UNIX的，UNIX化的。}的解决办法。

所以按照人对事物的认知，我对Linux的教学过程的安排将不同于其他的书，我的重点在“文化”，介绍原理而不简单的介绍的用法。
确实我会在书中有专门的章节安排详细讲解个各种实用工具，包括服务器的搭建，但读者可以把它当作是一个手册，用到的时候再翻阅。

在对发行版的考量上，我采用了Gentoo作为本书的参考发行版。在市场上还没有哪本入门书是使用Gentoo作为发行版的。
对于使用人数最多的Ubuntu\footnote{非常流行的一个发行版}发行版，我想一定存在非常非常多的教程。但是没有哪个教程能让我满意。
不是介绍的太过于片面而没有考虑到其他发行版，就是过分介绍了发行版之间的不同点。
原因就是Ubuntu（还有Fedora\footnote{RedHat的社区发行版，其改动将回馈到RedHat企业版。}这些）发行版，并不是只将软件进行打包，还进行了大量的修改使之符合发行版的风格。

\begin{notice}
就拿 /etc/init.d/ 下面的服务启动脚本来说吧，没有哪个发行版的脚本是一致的。
都经过了大量的修改。只介绍一个发行版还不够，但是为了“全面”而大而全式的把所有的发行版都介绍一遍，对初学者来说简直就是个噩梦。他会发现他需要为每个发行版都学习一次init脚本。
\end{notice}

如果说哪个发行版最尊重上游软件\footnote{所谓上游，就是软件的作者。发行版将软件从原始站点下载后重新打包，如果称发行版为下游，那原始软件开发者就是上游。}，进行最小的改动的话，那就是Gentoo了。

Gentoo几乎所有的软件包都是直接从上游代码编译而来，中间很少有加自己的补丁\footnote{对一个软件的修改被形象的称呼为打补丁。}。
有限的几个补丁也是为了修复bug而不是对软件进行修改。

这种接近上游软件的状态非常适合初学者的学习。

\section{UNIX的史前文明}

当你看《天宫韵音》\footnote{86版本的西游记电视剧的片头曲}的时候，最先播放的画面就是悟空出世。第一个音符响起的时候，悟空从绷床\footnote{没错！仔细看，确实是绷床！后期处理居然没处理掉。BUG无处不在啊！}上腾空而起。悟空就这样开天辟地般高调出世了。

但是UNIX的出世却非常低调，低调到需要介绍他那高调的前身\textbf{MULTICS}。

Multics的全称是\textit{MULTiplexed Information and Computing System}\footnote{注音这里的大小写混合。这种混合方式旨在告诉读者缩写的由来。}。我们先来了解一下计算机早期的背景。



\begin{insertnote}
\subsection*{Multics背景}

\chatu{TwoWomenOperatingENIAC}{为世界上的第一台电脑编程\protect\footnotemark}%
	\footnotetext{图片来自维基百科~\url{http://en.wikipedia.org/wiki/File:Two_women_operating_ENIAC.gif}}


计算机诞生的时候并没有操作系统，甚至连程序都没有。操作员重新插拔系统电路连接来实现某个功能。
如图\thefigure那样两个人在重新插拔电路为世界上的第一台电脑编程。

而今天的电脑只要执行不同的程序就能有不同的功能。程序的发明是个非常重要的进步\footnote{通过执行程序而不是直接修改电路实现计算功能，程序和数据一同存储，执行的时候同时进入主存储器，这种结构的计算机最早由冯.诺曼提出，又叫冯.诺曼结构。我们现在的个人电脑就是冯.诺曼结构的。}。有了程序就可以预先编撰好程序，只要输入了程序就能执行了，而不需要费时的重新插拔电路。

那个时候计算机程序的输入靠的是穿孔纸\footnote{用纸啊！太不环保了！}。一长条纸穿进去，程序就输入了，然后等待执行完毕。执行完毕后将结果打印到纸上
\footnote{别想多了，是简单的机械打印机。只能打印0\textasciitilde{}9数字、26个字母和特殊符号。有时就是这点功能都不能保证。可能只能打大写的A-Z而已。UNIX少说多做的传统就来源于此，不仅是为了节约纸张，更是节约打印时间。}。
每次程序执行完毕后操作员要准备好下一叠纸，按下机器的复位按钮机器重新开始读取纸张并执行....周而复始。

如果程序的执行不需要人准备下一叠纸就好了。第一个操作系统就在这样的需求下诞生了。不用人的帮助计算机就能批量执行程序。尽管一次还是只执行一个程序。但是能做到前一个程序执行完毕后自动读取下一个程序执行。

有时候一个程序要执行很长的时间，执行完毕后要花更长的时间打印结果，这个时候处理器的资源就是闲置了\footnote{写小看这个闲置，那个时候的计算机比劳斯莱斯贵多了。}。如果前一个程序在打印结果的时候后一个程序能开始执行就好了，这样就能最大限度的利用处理器的资源。

于是最初的多任务操作系统应运而生。它能将处理器的时间分配给多个程序，虽然在任意瞬间还是只有一个程序在执行，但是由于快速的在不同的程序间进行切换\footnote{专业术语叫上下文切换（Context Switch），所谓上下文切换，就是将当前程序的状态保存起来，然后调入先前保持的另一个程序的状态。这样那个被调入的程序就可以从被切换打断的那个地方继续执行。}，看起来就好像多个程序同时执行了。如图 \ref{fig:ContextSwitch}  所示。

\chatu[scale=0.65]{ContextSwitch}{上下文切换}

但是那个时候的计算机，还不能进行实时交互。程序员写好纸条后，要交给计算机管理员\footnote{确切的说是操作员}，管理员将纸条输入计算机，然后等待计算机将执行结果输出。最后输出结果交给程序员，这样程序员才能知道程序的执行结果。


随着计算机性能的提高，出现了一种分时操作系统（Time Sharing Operating System）。 
它可以让大型主机提供数个终端，用户通过终端使用主机。
终端经一条很长很长的线缆\footnote{也可以通过电话线呢！}连接到主机上\footnote{那个时候主机很大，要放到专门的地方嘛。}。
当时的终端就是一个电传打字机和一个机械打印机，只具有输入/输出的功能，本身完全不具任何运算或者软件安装的能力。要使用计算机的人将程序或者命令通过终端输给主机，主机将执行结果传回终端。最后来屏幕代替了打印机，成为视频终端。

\chatu[scale=0.8]{ctss}{分时操作系统}

如此一来，无论主机在哪里，只要在终端機前面进行输入输出的作业，就可利用主机提供的功能了。

前面已经讲过，之前的系统只能由一个操作员去操作。而分时系统为每一个使用计算机的人提供了一个终端。比较先进的主机大概能同时为30个终端服务。 至此，计算机终于能实时交互\footnote{其实只是接近实时。不过已经是一种进步了，不是么？}了。

Multics就是在分时系统的基础上开发的。

\end{insertnote}

MIT\footnote{麻省理工大学}计算机研究中心研究出来的CTSS(Compatible Time Sharing System)是世界上最早的分时系统，它在1961年首次演示。
正如一切软件开发者所遵循的鉄律：来不及添加的功能就放到第二版做吧！Multics就是CTSS来不及添加更多功能后希望开发的全功能第二版。1964年，MIT、通用电器（General Electric）和贝尔实验室\footnote{根据Multics Myths网站（一个网站致力于破除有关Multics的谣言）的说法，贝尔实验室应该是在1965年的时候才参加的。
}通力合作，计划开发一个完美的CTSS第二代。

Multics计划的目的是想要让大型主机可以\textbf{}为超过300个以上的终端服务和一整套完整的信息处理系统。
Multics支持模块化的硬件配置：CPU、内存，磁盘等硬件都可以在需要的时候在系统运行的过程中被动态增加或移除。在当时的系统差不多都是由汇编语言写成的时候，Multics就大胆的提出了使用高级语言\footnote{当时还没C语言呢！}编写整个系统。

Multics先进的系统模型的结果就是：从1964年立项1969年，整整5年过去了，Multics的开发一再受挫。
Multics迟迟不能完成，预定的目标越来越难以达成，在投入了巨大的人力物力资源后，
贝尔实验室终于决定在这场注定失败的项目\footnote{一种说法是Multics最终达成了设定的目标，所以不能算是一个失败的项目。但是我且以为Multics没能在期望的时间内达到设定的目标，就是失败的。}中提前退出。

但是Multics激发的灵感------高级语言编写的内核、文件系统、终端实时交互、多任务、多用户等------却在今后几十年内被延续到各种各样的系统上。这些最终将在接下来的几十年时间里分别在不同的系统上得以实现。

Multics是个优秀的系统。在参加Multics项目的日子里，贝尔实验室的研究人员已经习惯了不用等待就能通过交互式终端使用大型电脑的功能。
但是退出Multics项目后，他们又得继续原始生活了。由于习惯了Multics带来的便利，贝尔实验室的研究人员将发明一个简化的Multics，这个发明就是\textbf{UNIX}。

\section{UNIX的故事}

很显然Ken Thompson不愿意放弃Multics提供的交互式终端，
而且为了能继续玩他在Multics项目做研究的时候编写的“星际旅行”游戏，他头脑风暴了一把，用实验室废弃的一台PDP-7写了一个简陋的操作系统。
Ken Thompson为PDP-7开发的小型系统被称呼为UNICS（UNiplexed Information and Computing System） 。 这个名字的由来是因为Ken Thompson认为Multics太复杂，而Multics的功能可以用更简单的方式去实现。Dennis Ritchie注意到了他的同事为PDP-7开发的简单系统。两个人亲密无间的合作就这样开始了。

\chatu[scale=0.2]{KandRwithPDP7}{Dennis Ritchie（站）和Ken Thompson（坐）在工作}

Dennis Ritchie和Dennis Ritchie的加入为Ken Thompson的系统带来的第一个变化就是高级语言。他们开始计划用B语言\footnote{C语言的前身，Ken Thompson创造的。}重写这个系统。

在重写UNICS的过程中，Dennis Ritchie改进了B语言，使之成为更加实用的C语言。Dennis Ritchie和Ken Thompson一起实现了C语言，并在1973年用C语言重写了UNICS。哦，不对，是UNIX。因为谐音的关系，不知道什么时候开始UNICS就被称呼为UNIX了。

UNIX很快就在贝尔实验室流行开来了。
1974年，Dennis和Ken第一次对外公布了UNIX。他们在《美国计算机通信》(Communications of the ACM)上发表了介绍UNIX的论文。
这篇论文使得外界对UNIX产生了极大的兴趣，纷纷开始向贝尔实验室索要UNIX。

老实说，AT\&T是个大公司。恩不对，还是个垄断公司。AT\&T垄断了美国的电话网。
根据1958年为解决反托拉斯案例达成的和解协议，AT\&T被禁止进入计算机相关的商业领域。
所以,Unix不能够成为一种商品。实际上，根据和解协议的规定，贝尔实验室必须将非电话业务的技术许可给任何提出要求的人。
Ken Thompson开始默默回应那些请求，将包含着UNIX源代码的磁带和磁盘一包包地寄送出去——据传说，每包里都有一张字条,写
着“love，ken”（爱你的，ken）。

UNIX的源代码在高校间广泛传播\footnote{UNIX因此获得了一个非常重要的传统“附带源代码”。正是这个传统催生了后来的开源运动。}，
其中有一份代码被寄到了加州大学伯克利分校。在那里，UNIX被孵化，催生了BSD（Berkeley Software Distribution）。

在伯克利，一些重要的UNIX功能被开发出来。最重要的莫过于socket了。在BSD开发出socket之前，UNIX的网络功能简陋得可怜\footnote{UNIX有个uucp命令可以通过电话线执行远程拷贝，仅此而已。}。


AT\&T对UNIX源代码的“慷慨\footnote{没那么好心。不过是AT\&T还没意识到UNIX的商业价值，同时暂时性的被反托拉斯和解协议束缚了。}”奉献促成了UNIX的繁荣。

大量的新功能被开发出来，各种功能补丁雪花般飘到贝尔实验室。各种版本的UNIX不断涌现，商业化的UNIX系统也出现了。
伯克利开发BSD的Bill Joy加入了新成立的Sun\footnote{是Stanford University Network的缩写。2011年被甲骨文（Oracle）收购。}公司开发基于BSD的SunOS\footnote{SunOS从 5.0版本以后就不再基于BSD开发了。}。
就连微软也曾开发名叫Xenix的UNIX系统\footnote{但是微软把UNIX作为一个产品的热情并没有持续多久。}。


\subsection{UNIX混乱和分裂}

UNIX世界百花争鸣的日子随着AT\&T公司的拆分一去不复返。1983年，美国司法部对在针对AT\&T的第二起反托拉斯诉讼中获胜，AT\&T和贝尔实验室被拆分。这次判决将AT\&T从1958年的禁止将UNIX产品化的和解协议中解脱了出来。AT\&T马上忙不迭地将UNIX商业化------这一举措差点扼杀了UNIX。

那时，没有人意识到，UNIX的产业化会破坏UNIX源码的自由交流，而恰是后者滋养了UNIX系统早期的活力。AT\&T收回了原先的UNIX源码的许可，挂上高价（而在这之前，AT\&T只对UNIX源码收取象征性的费用），对源码散发严加防护。来自高校的贡献开始枯竭。

刚刚进入UNIX市场的几家大公司马上犯下了最重大的战略错误------试图通过差异化寻求有力地位------这个策略导致了
UNIX的分歧，抛弃了UNIX的跨平台兼容性,造成了UNIX市场分割。

AT\&T 拆分后的数年内，UNIX社区忙着UNIX大战的第一阶段-----System V Unix（AT\&T的将UNIX商业化后推出的UNIX）和 BSD UNIX之间的内部争吵。争吵分成不同的层面，有些属于技术层面（socket vs stream,BSD tty 对 System V termio），有些则属于文化层面。
AT\&T对UNIX的商业化使得UNIX作为统一的操作系统环境分崩离析。市场上充斥着各种互不兼容的UNIX。
给软件开发和系统维护带来了沉重的负担。

\subsection{统一的UNIX和POSIX标准}

1983年/usr/group标准组开始了以调解System V UNIX和BSD UNIX的API为目标的严肃的标准化工作。
随后IEEE以/usr/group标准为基础制定了POSIX
\footnote{其正式称呼为IEEE 1003，而国际标准名称为ISO/IEC9945。POSIX这个名称是由理查德·斯托曼应IEEE的要求而提议的一个易于记忆的名称。}
（Portable Operating System Interface\footnote{最初这个X是 of UNIX最后的那个X。但是后来鉴于POSIX可以用于非UNIX（这些系统被称为UNIX-like，如Linux和Minix）系统，于是去掉了of UNIX但是保留了X。 }，可移植操作系统接口）标准。并于1985年发布了第一版。

POSIX标准定义了UNIX（和UNIX-like）系统的API\footnote{ 参考\faqref{API}}（Application Programing Interface）和shell命令所要遵循的标准。

后续开发的各种UNIX版本也严格遵循这个标准。

\begin{notice}
注意：POSIX标准是源码级兼容的标准。意思是在一个系统上编写的软件可以在另一个系统上“无修改的”通过编译并正常运行。
但是如果是编译好的可执行程序拿到另一个UNIX系统下“不保证”能执行。

但是对于shell脚本就另当别论了。POSIX定义了一个shell所必须具备的功能和行为。只使用POSIX定义的shell功能的脚本可以未经修改直接在另一个UNIX系统上执行。
\end{notice}

那么UNIX的故事到这里就告一段落了。Linux文化的的另一个重要来源GNU也要登上历史舞台了。

\section{GNU 的故事}

GNU也不是什么开天辟地的存在。早在GNU出现之前，类似GNU提倡的那种开放源代码的黑客精神就已经存在了。他们相互共享软件的代码，相互修改对方的软件。相信软件是一种自由的可以共享的东西。他们当中最杰出的一名，在黑客文化行将灭亡的时候，发出了著名的《GNU宣言》。这个人就是Richard Matthew Stallman\footnote{他的简称RMS比他的名字更出名}。

\chatu{rms-full-size}{Richard Matthew Stallman}

\subsection{认识自由的含义}

1971年RMS进入哈佛大学学习，同年受聘于麻省理工学院人工智能实验室（AI Laboratory），成为一名职业黑客。

RMS在人工智能实验室的最出名的发明就是Emacs编辑器。Emacs打从一开始就以源码的形式进行发布。Emacs出色的设计使得它非常的流行。许多人对Emacs进行了改进。这些改进最终都回到了RMS的手里。这种共享代码的方式深受黑客们的喜爱。对于RMS来说，自由的共享软件的源代码是多么自然的事情。

但是商业化的浪潮也侵袭到了MIT。RMS感觉到他熟悉的黑客圈开始消失。

\begin{insertnote}
\subsection*{该死的私有代码}
在软件私有化的浪潮开始前，所有的软件销售的时候都是附带源码的（不过修改代码重新发布的权利还是没有的）。软件通常是作为硬件的“附属品”。
人工智能实验室购买了许多的硬件设备。当硬件驱动有错误的时候，RMS和他的黑客同事们总是能修正这些硬件驱动的BUG。他们甚至能开发出新的驱动让硬件提供不同的功能。

后来开发商逐渐的都不再提供驱动软件的源代码了。这样驱动程序出了问题这些黑客们只能对着机器发呆等待厂家修正。
可想而知RMS心里有多么的痛恨这些私有代码。如果他们有这些糟糕的软件的代码，就能修改它们而不是被动的等待厂商的修复。
\end{insertnote}

进入80年代，他的同事们许多都离开了实验室开始创立公司。
RMS觉得自己被逐出了黑客的伊甸园,他把这一切都归咎于私有代码。

\subsection{为了自由}

RMS认为所有的软件都应该是自由的。自由的意思就是:
\begin{itemize}
\item 运行的自由。你可以自由的决定运行或者不运行这个软件。运行的话怎么运行这个软件，在何种环境下运行\footnote{可别小看这个自由，许多软件规定只能在家运行，不能在作为商业用途在公司运行。}。
\item 修改的自由。你可以获得程序的源代码，并进行修改。
\item 再发布的自由。你可以自由的传播不论是你修改过的还是未修改过的软件。
\end{itemize}

RMS可以编写更多的自由的软件来做抗争。但是只要一个系统的核心------操作系统------是私有的软件，写再多的自由应用程序又有何意义呢？
于是RMS要做的第一件事情就是创建一个自由的操作系统。UNIX本来是自由的，却被AT\&T变成了私有的软件。于是RMS想到编写一个自由的UNIX兼容系统。这个项目的名字就是GNU，是Gnu is Not Unix\footnote{这种递归缩写在很多项目里用到。}的递归缩写。

造一个自由的操作系统首先需要编译器。于是他开始开发C语言的编译器。也就是GCC。到了1985年，RMS发现靠一个人的力量做事情太慢了，他需要帮手。于是他成立了自由软件基金（Free Software Foundation 缩写FSF）。同年RMS撰写了著名的自由宣言《GNU宣言》（The GNU Manifesto）。

1986年，GDB\footnote{GNU Debuger，GNU调式器}发布。次年，第一版GCC\footnote{GNU C Compiler，GNU C编译器。后来因为GCC能支持其他语言，GCC被改名为 GNU Compiler Collection，GNU编译器套件} 发布。至此开发一个软件所需要的3大件------编辑器、调式器、编译器------齐备。

RMS面临的另外一个问题是，他需要保护自由软件的成果不至于被偷窃。UNIX就是一个活生生的由共享源码的自由的软件变成了AT\&T的商业软件的例子。于是在1989年，RMS编写了著名的GPL\footnote{GNU Public License，GNU公共许可协议}协议。

\subsection{以Copyright保护Copyleft}
%\subsection{只有自由软件的完整操作系统}

自由软件的目的是\textbf{去版权化},也被 称为Copyleft。
但是首先要用版权保护自己。软件的作者通常会随软件附带一份许可协议，这份许可协议由版权法保护，用户同意即生效。而且要使用该软件用户必须同意（当然前提是这个国家对版权的保护非常好。）。
用户可以同意并使用该软件或者不同意并退回所购买的软件\footnote{这有点像是霸王条款？}。通常这样的许可协议是用来保护开发者的：微软公司就在他的许可条款里堂而皇之的说因使用微软的软件导致的损失微软免责，这种免责条款被大多数公司所使用。

GPL却利用了这点来\textbf{保证用户的权利}。用户被赋予了和开发者同等的权利，只要保证一点：他同意继续保持软件为GPL授权。

\begin{itemize}
\item 当你下载了GCC，你有自由执行它的权利。GPL授权的GCC不会在许可协议里说，“嘿伙计，你只能在Windows下执行该软件\footnote{微软的Office软件有此条款。}”

\item 当你下载了GCC，你就拥有了GCC开发者相同的权利。你可以查看GCC的代码------FSF的网站上提供GCC源码下载------就和GCC开发者能查看GCC的代码一样。

\item  当你下载了GCC，你就拥有了GCC开发者相同的权利。你可以修改GCC的代码------就和GCC开发者能修改GCC的代码一样。

\item  当你下载了GCC，你就拥有了GCC开发者相同的权利。你可以把GCC修改后（或者没有修改）分享给任何一个人------就和GCC开发者能把这个软件发给你一样，而不会在许可协议里说，“嘿伙计，你不能把本软件拷贝给任何一个人，否则就是非法拷贝，是盗版”。

\end{itemize}

所有的这些权利你都能获得，只要你同意
\begin{itemize}
\item 当你再发行这个修改或者未修改后的软件的时候，你保证不修改它的许可协议。让它的GPL协议传播给下一代。
\item 当你在自己的软件里用到了一个GPL授权的软件的一部分或者全部的代码的时候------就算你没有将这个软件的代码直接包含到你的软件中，而只是使用里它提供的函数库------你保证将你的软件同样以GPL协议授权。
\end{itemize}

GPL就是通过如上的传染性保证自己能像病毒一样迅速传播。随着GPL授权的软件愈来愈多，最后几乎不可能编写一个完全不依赖GPL代码的软件了。这样所有的软件最终都将被迫GPL化\footnote{我不知道这个词语之前有没有人用过，我希望读者能理解这里个词语的含义，就是改变原先的许可协议为GPL。}。

到1992年，GNU项目开发出了一系列的软件，包括编辑器、开发工具链（编译器，调试器，C库，autotools\footnote{参考}）、shell和系统工具等等，再有一个内核就是完整的自由的操作系统了。但是GNU的Hurd内核一再跳票，距离发布的日子遥遥无期。这个时候GPL授权的Linux内核刚好填补了Hurd的位置，使GNU成了一个彻底完整的系统。至此UNIX的文化经过80年代的短暂中断后在GNU和Linux的组合系统中延续。

\fi


%=========================================================
% 正文开始
%=========================================================
\chapter{UNIX哲学}
\begin{quotes}{Dennis Ritchie}
UNIX is very simple, it just needs a genius to understand its simplicity.
\end{quotes}

\ifdefined\compileall


UNIX的哲学到底是什么？UNIX的哲学一言以蔽之，就是：\chatu*{kiss}{KISS}


UNIX无处不在，生命力顽强。从1969年诞生的那一刻起，没有人想到它会如此成功。得益于C语言，UNIX以其绝佳的可移植性轻松的被移植到各种类型的机器上。小到各种移动设备，达到超级计算机，UNIX都能在这些设备上畅快的运行。恐怕没有哪个操作系统能像UNIX那样拥有持续40多年的生命力。UNIX不仅存在于UNIX操作系统本身。许多非UNIX系统都能找到UNIX的影子。


%TODO:  主要的系统管理思想和一些小工具也在这里介绍，比如 ps grep 用户和组，用户权限的概念也在这里讲解

\section{一切都是文件}

这是最明显的一条UNIX设计哲学了。通过将一切设备抽象为文件，就可以使用统一的方法去操作设备。不管设备是磁带机也好，猫也罢，都使用同样文件操作。唯一不同的，也只有操作的设备文件名不一样而已。彻头彻尾的反UNIX，就是将设备和文件系统隔离。

\begin{example}{备份 MBR}

在 Linux下，第一个磁盘的设备文件是 /dev/sda , 第二个磁盘的设备文件是 /dev/sdb 
要备份第一个硬盘的 MBR ,可以使用  dd 命令。

\begin{code}
dd if=/dev/sda of=mbr.bak bs=512 count=1
\end{code}

使用这条命令就可以把硬盘的第一个扇区备份到 mbr.bak 文件了。对于其他的UNIX系统，只有对硬盘设备的命名不一致而已，没有其他区别。

而彻头彻尾的反UNIX，就需要使用特殊的备份MBR/分区表的“工具”。
\end{example}


\begin{example}{查看硬扇区}

直接查看硬盘的扇区，就和查看一般文件没有任何区别。使用二进制编辑器打开磁盘设备就可以了。
如果只是为了查看一下，还有个非常方便的工具：\textbf{hexdump}

\begin{code}
hexdump /dev/sda
\end{code}

这条命令就可以把第一块硬盘的扇区数据直接显示到屏幕上，按Ctrl-C终止。想屏幕不要快速滚动的话，请仔细阅读下面的章节，你会找到办法的。

而彻头彻尾的反UNIX，就需要使用特殊的查看磁盘扇区的“工具”。
\end{example}

从两个例子可以看出，UNIX将硬盘看作一个普通的文件（只不过文件大小固定为磁盘容量）。像这样的设备文件UNIX将他们组织起来统一放到/dev目录下面。
/dev 目录下的文件一度是由内核进行管理的，不过Linux后来放弃了内核管理/dev的办法而改为由一个叫udev\footnote{udev最后被systemd合并。}的程序进行管理。


\section{鼓励程序间协作}

UNIX鼓励短小精湛的程序，不鼓励编写大而全式的程序\footnote{据说是因为当年PDP-7内存有限}。正是因为UNIX鼓励短小的程序，并相互配合，所以创造了无穷的灵活性。
UNIX通过组合各种小程序，能发挥出各个小工具独立工作的时候不能完成的功能。

实现程序间协作的最伟大的发明就是\textbf{管道}。通过管道，一个程序的输出可以变成另一个程序的输入，将多个程序进行串街------每个程序变成这条处理链上的一个\textbf{过滤器}------组合的这些工具能产生工具的开发者都没有预料到的功能。

\chatu{Douglas-McIlroy}{管道发明人 Malcolm Douglas McIlroy}

图\thefigure 为管道发明人Malcolm Douglas McIlroy。 \chatu[width=\textwidth]{pipe-draft}{管道} 图\thefigure 为他发明管道的时候提出的建议的打字稿扫描图片。

管道能够让你用简单的程序去构 造更复杂的程序。管道能够以意想不到的方式去使用命令,管道使得你可以用更简单的方式去实现复杂的功能。
UNIX简单就是美的设计哲学，让管道这种强大的功能简单到只需要使用“|”就可以实现了。

\begin{example}{还是查看硬扇区}

前面说过，直接查看硬盘的扇区，只要执行如下的指令：
\begin{code}
hexdump /dev/sda
\end{code}

想屏幕不要快速滚动的话，并不需要修改 hexdump 程序添加滚屏功能。这个功能只需要交给一个专门为\textbf{滚屏}而开发的程序: less。

只要执行这样的命令就能实现一次查看一个屏幕的内容。还可以向上回滚。
\begin{code}
hexdump /dev/sda | less
\end{code}

less 的功能就是将自己的输入以一个屏幕的内容显示，超过屏幕能显示范围的内容要等用户使用滚屏命令（PageUp PageDown 空格键和回车键）让less滚动到用户想要观察的位置。

如果我不太喜欢 less 提供的功能，我还有more呢！
\begin{code}
hexdump /dev/sda | more
\end{code}

more和less在功能上大同小异。

而彻头彻尾的反UNIX，就需要给所有有需要的程序添加滚动屏幕的支持。

\end{example}

UNIX鼓励程序之间通过管道或者其他IPC（Inter-Process Communication 进程间通信）进行协作。这样一个程序就可以专注于本质工作，不需要为了“多功能”而开发庞大复杂的软件。这就引出了下一个我要讲的UNIX哲学。

\section{一个程序只做一件事情并把它做好}

UNIX鼓励程序间的协作，一个程序只需要专注于本职的工作，并做到最好。UNIX下有各种简单的工具。但是这些简单的工具却将自己的本质工作做到了最好。
所以UNIX下通常不会有功能繁多，大而全式的软件。但是你有非常丰富的功能单一却很强大的小工具，并拥有将工具进行组合的能力。

而彻头彻尾的反UNIX，让每个软件都倾向于包含你可能用到的一切功能，试图提供一个“家长式”的程序。

\begin{example}{下载工具wget}

wget 是一个命令行下的下载工具。wget功能非常明确，http/ftp 下载。使用起来也非常简单
\begin{code}
wget [url]
\end{code}
但是wget做到了一件事——下载——并将它做最好。
但是，如果给你的是一个下载文件列表，wget可不会去支持它。这是和下载无关的功能。但是，你能通过组合命令实现：
\begin{code}
cat urllist | xargs wget
\end{code}
在这里，我们又利用了一个名字叫xargs的程序，xargs的程序功能很简单，就是接受标准输入，为每个输入的字符串作为最后一个参数调用一次指定的命令。
cat 是个功能非常简单的程序，将文件打印到标准输出。这个标准输出被\underline{\textbf{管道}}重定向给了xargs做为标准输入。

这个命令组合后的结果就是：为urllist里列出的每个url调用一次wget进行下载。自然也就是为wget实现了“下载文件列表里的url”的功能了。
\end{example}

UNIX就是因为提供了管道这种强大到底协作工具，所以得以鼓励开发各种功能单一的精湛的小程序。而功能单一的程序大大简化了软件的开发，也使得UNIX下的软件质量非常高。

\section{使用文本流，因为这个是通用的协议}

如果说，UNIX鼓励编写功能简单的程序，必定是因为它提供了程序间协作的机制，而要利用这样的机制，程序间唯有使用统一的协议才能实现正确的协作。
这个统一的协议就是文本流。

UNIX反对使用专有的二进制格式。UNIX天然的支持流格式\footnote{一个接着一个的字节}，连文件都被看成是无格式的字节流。
因为输出不仅可以被管道送给另一个程序处理，更有可能会直接打印到终端。
所以对于输入输出，UNIX更是鼓励文本字节流，也就是只使用可打印的\footnote{26个大小写字母，0到9阿拉伯数字，标点符号，空格和换行}文本字母。并且将这种文本用于进程间的通信。

\begin{example}{文本协议HTTP}
目前广泛使用的HTTP协议就是一个文本协议的例子。

比如你访问 \url{http://www.example.com/example.html} ， 浏览器连接到 www.example.com 主机的80端口，发出一个
\begin{code}
GET /example.html HTTP/1.1
\end{code}
请求。服务器收到后，如果找到了example.html文件，就返回
\begin{code}
HTTP/1.1 200 OK
\end{code}
并空1行后紧接着输出example.html文件的内容。

要模拟浏览器的请求，可以使用\textbf{telnet}这个程序。
\begin{code}
\#telnet  www.example.com 80\\
Trying \st{***.***.***.***}...\\
Connected to www.example.com.\\
Escape character is '\textasciicircum'.{ <--  \it 以上为telnet输出 }\\
GET /example.html HTTP/1.1 { <--  \it 自己输入的，输入后敲两下回车 }

\\

HTTP/1.1 200 OK\\
Date: Wed, 19 Sep 2012 15:41:40 GMT\\
Content-Length: 31\\
Content-Type: text/html; {  <-- \it 下面有一行空行}

\\

<html>hello example.com</html>

Connection closed by foreign host. <-- \it telnet输出
\end{code}

可以看到，HTTP协议使用文本模式进行通信。服务器的回应也是文本流。而且是行模式，没条信息存放一行。如果有多个信息，可以多行存放。HTTP协议没有规HTTP头的大小，HTTP头和传输的内容之间使用空行隔开\footnote{这就要求HTTP头本身没有空行。}。

\end{example}

除了HTTP，互联网上常用的协议，比如FTP、SMTP\footnote{发送邮件用的}、POP3\footnote{用于客户端从服务器接收邮件}都是基于文本的协议。
HTTP这种只有HTTP头是文本，后续发送的文件内容（尤其是二进制文件来说）是字节流。但是SMTP全程连接使用的都是文本。对于发送不可打印的字符（比如以一张图片作为附件），SMTP需要将二进制字节流转化为文本格式。这是文本协议的一个缺点。

但是瑕不掩玉，文本协议的好处有很多，比如可以无限扩展。

还是拿 HTTP协议来说（方便起见，我没有设计HTTP请求的二进制版本。）。刚刚访问 www.example.com获了example.html。该文件有31个字节长度，HTTP头中的信息Content-Length: 31也告诉了我们文件大小。如果是二进制协议，一般会设计使用一个32bit的整数代表返回的文件大小。如果需要处理很大的文件，协议恐怕需要使用64bit长度的整数。但是对于64bit都无法表示的文件怎么办呢？ 二进制协议始终有扩展性不佳的问题。如果是HTTP协议，不管文件多大，只要这个数字还能表示出来，就没有问题。

我们可以看到， HTTP 返回的除了文件大小，还有文件类型等信息，如果是二进制的协议，大概会如图 \ref{fig:binhttp} 这么设计，其中每个方块代表4个字节。要传输的文件紧挨着HTTP头的后面。

\chatu[width=0.2\textwidth]{binhttp}{二进制的HTTP协议，服务器返回部分}

每个类型的信息都分配了固定的字段。解析起来非常快，似乎够用了

等等，我还要添加个信息，表示这个文件该不该被缓存。如果这个文件是临时文件，我大概不希望浏览器将它缓存，而是每次都到服务器上获取文件内容。
好吧，我再添加个字段。

\chatu[width=0.2\textwidth]{binhttp2}{二进制的HTTP协议，服务器返回部分，版本2}

等等，已经运行了的大批服务器和浏览器怎么办？由于文件紧跟着HTTP头，使用新的二进制HTTP协议的服务器碰上老版本的浏览器，缓存属性不仅是被忽略了，更被当做文件内容的一部分了。难道为了一个小修改需要为新的协议修改所有的浏览器和服务器么？

能无限的扩展并保持极佳的兼容性和人类可读性\footnote{测试问题的时候不需要特殊的工具去解析二进制流的各个字段，只需要打印出文本就可以了。}，是文本协议的最大优点。

\section{懒惰原则}

UNIX认为人的时间比机器的时间要宝贵的多。如果有什么事情是机器能代劳的，就绝对不要浪费人的时间。

这是 \textit{The Art of Unix Programming} 的《附录D》收录的无名师的传记，说明懒惰的人写脚本来解放自己的时间。我想没有比这个故事更能表达懒惰原则的含义了。

\begin{insertnote}
\subsection*{无名师与shell脚本}

无名师曾对来访的程序员说：“Unix传统上认为，一行 shell脚本胜过万行C程序。" 

这个程序员自以为对C极其精通，说：“这不可能。Unix内核正是用C实现的。” 

无名师回道：“确实如此。不过，Unix传统上认为，一行shell脚本胜过万行C程序。" 

程序员颇为沮丧：“但是在C中我们可领会到尊者Ritchie的智慧。我们与操作系统和机器合而为一，可以获取无与伦比的性能。“ 

无名师回道：“诚如你言。不过，Unix传统上认为，一行shell脚本胜过万行C程序。” 

程序员冷笑着想悲愤离去。无名师向学生Nabi颌首示意，Nubi在黑板上写下一行shell脚本，问道：“尊敬的程序员，看看这行管道。用纯C实现，是 不是要一万行C代码？“ 

程序员沉吟念诵。最终他承认如此。 

“你需要多长时间来实现和调试那个C程序?"Nubi问道。 

“很长”，来访程序员承认。“但傻子才会干这个而不去完成更有价值的任务。” 

“那么谁更了解Unix传统？”无名师问道。“是写一万行代码的，还是看到任务的无谓而不去编码的？” 

听到此，程序员眼中一亮。
\end{insertnote}

%这个原则我们将在 \ref{section:pkgmgr} \nameref{section:pkgmgr} 一节中讲解autotools的时候体会到。
无名师告诉我们，能用脚本完成就不编写程序。这是懒惰的原则，这是“宁花机器一分不花程序员一秒”的原则。这条原则对非程序员也适用的，就是\textbf{尽量编写脚本将自己的工作自动化。}

%而彻头彻尾的反UNIX，如果让每个软件都倾向于包含你可能用到的一切功能，试图提供一个“家长式”的程序。

\section{UNIX和非UNIX}
UNIX设计哲学并非只能在UNIX身上找到。只是有的系统只学了个表面。更有反UNIX而行之的。

UNIX管道自发明后，已经被各个操作系统借用，包括DOS和Windows都实现了管道。Windows/DOS依旧借鉴了使用|作为管道的办法。
\par\begin{notice}
事实上DOS是作为运行于8086芯片上的蹩脚UNIX开发的。树形目录结构（第二版开始支持）、命令行操作、管道，这些都带有深刻的UNIX影子。
DOS没学到的（受限于8086的能力也无法实现）：多用户多任务，丰富的IPC接口。
\end{notice}

%但是Windows下管道是如此的难用，以至于

UNIX“一切都是文件” 的思想也在许多的操作系统上实现。只是学的并不好。Windows下硬盘实际对应文件为 \texttt{\textbackslash\textbackslash.\textbackslash{}PhysicalDrive1} 第二块硬盘最后一个数字为2。但是这样的文件是无法在文件管理器里看到的。
而且也不能作为文件名使用一般的程序打开。

UNIX使用轻量级进程，进程的创建开销非常小。鼓励开发多个小程序并使用IPC相互协作。
其他系统也实现了多进程能力，但是进程的创建开销（相对UNIX来说）非常大，因而总倾向于使用（便宜的）进程内通信------将当个的程序做的非常大，使用复杂的API进行内部通信。而这些都是BUG的温床。

UNIX鼓励使用文本作为协议。就连配置文件都是（带格式的）文本文件。而有些系统使用专有的二进制格式数据文件保存配置。一旦配置文件损坏，系统就无法使用简单的文本工具进行修复。甚至因为系统本身无法启动而不能进行修复。

\fi

\chapter{黑客和Linux}
\begin{quotes}{Linus Benedict Torvalds}
If 386-BSD had been available when I started on Linux, Linux would probably never had happened.
\end{quotes}
\ifdefined\compileall

\begin{quotes}{Richard Matthew Stallman}
Writing non-free software is not an ethically legitimate activity, so if people who do this run into trouble, that's good! All businesses based on non-free software ought to fail, and the sooner the better.
\end{quotes}

在计算机发展的初期，软件\footnote{那个时候系统软件都是汇编语言开发的，所以也谈不上什么源代码。}被作为硬件的附属品销售。
购买一个机器的硬件也就连带的附送了一系列的软件。谁也不会觉得软件是可以被“保护起来”盈利的\footnote{机器码可以非常容易的反汇编重新生成源汇编文件。既然软件是汇编开发的，那也没有什么秘密可以隐藏。}。相互交换彼此的程序随处可见。
直到有一天，西雅图的一个软件天才\footnote{为了避免反应迟钝的人不知道，这位天才是比尔\textbullet盖茨啦}说，他花了很多时间写好的软件，不希望被人随意拷贝。每一个获得拷贝的人都应该给他付钱。
有了天才的带领，自此商业公司开始独立销售软件而不再作为硬件的附属品。

不过这股软件私有化的浪潮还不曾传播到社会的各个交流。
大学里，实验室里，人们还是保持着将软件自由传播的传统。虽然他们不再能自由的传播商业软件，但是自己编写的各种
稀奇古怪的软件可以自由的交流。在没有internet的年代他们就开始利用Usenet\footnote{%
\textbf{来自维基百科：}Usenet是一种分布式的互联网交流系统，它的发明是在1979年由杜克大学的研究生Tom Truscott与Jim Ellis所设想出来的。Usenet包含众多新闻组，它是新闻组（异于传统，新闻指交流、信息）及其消息的网络集合。 }相互交流了。
%
这样一群天才不满于商业化的软件，时常有人对商业公司的系统进行破解，故得名hacker，中译“黑客”\footnote{%
由于只有身怀绝技的聪明人物才有可能破解商业软件。故而黑客给人一种智商非常高的印象。久而久之智商很高（或者因为专业技能好，看起来很高）的程序员都被称作黑客。所以现在说的黑客一般指的是计算机方面非常厉害的人而不再是搞破解的人。现在再搞破解和破坏的人，就没有黑客头衔了，而是Cracker，骇客。
}。

Linux的诞生就是因为黑客门不满足于被大公司控制的操作系统市场。

\section{从UNIX到Linux}

早期的UNIX鼓励源代码的密切交流。这和黑客们的思想如出一辙。伯克利开发的BSD更是吸引了大量的黑客进行贡献。70年代到80年代早期，黑客门过着自由的日子。


但是一场突如其来的官司让AT\&T从反托拉斯的禁锢中挣开\footnote{参考前言的内容。}。
脱缰的AT\&T将UNIX推入了深渊。

但是就算UNIX被推入了深渊，BSD还是一个避风港。但是这最后一点希望都随着AT\&T对BSD的起诉而消失了。

1983年，AT\&T发布了Unix最新版system V UNIX，这是一个商业化版本，付费才能使用，并且不得传播源码。
为了减少纠纷，伯克利分校规定，BSD本身依然保持免费，但是只能提供给持有AT\&T源码许可的公司。
不过，与此同时，伯克利的师生也着手开始将AT\&T的代码从BSD中逐渐去除，以便发布一个不受AT\&T控制的版本。
这一举动惹恼了AT\&T。AT\&T起诉BSD侵犯了Unix的版权。
这场诉讼对BSD打击极大，所有的开发活动都被迫停止，用户人心惶惶，担心自己也遭到AT\&T的追究，因此BSD的使用急剧减少。

这场战争给BSD带来了毁灭性的打击。


黑客们突然间没有了玩具。

最后在1994年，双方达到和解，BSD才恢复开发。但是历史的已经不会再给BSD机会了。因为在BSD官司缠身的时候，上帝把爱心悄悄给了Linux。

1991年，芬兰的一个大学生在网上公布了Linux，立即吸引了大量黑客的目光。从此BSD神码的都成了浮云。黑客门团结在天才周围缔造了Linux的神话。
不用说，这位天才的大学生名字就叫做 Linus Benedict Torvalds。

\section{Linux诞生}

UNIX自诞生以来就是学术界研究的对象。UNIX被移植到大量的机器上，唯独没有染指PC市场。UNIX的人轻视了这个以8086为核心构造的个人电脑。
也难怪，UNIX的多任务能力需要CPU能提供一些保护机制，防止未经授权的访问。而这正是8086所缺乏的。

但是当80386发布的时候，情况有了一些变化。80386具备了“真正的”CPU的特性，包括特权级别区分和分页内存支持\footnote{关于分页技术，请参考 \faqref{FAQ:Paging}。}。 有了这些支持，PC终于具备了运行UNIX的条件。

\subsection{被UNIX忽略的PC}
%那么我们还是字为何PC
计算机最初是作为美国军方的项目研究制造的。美军向来有向全世界开放自己的研究的传统\footnote{Internet也源自美国军方的研究。还有现在每天都在用的GPS卫星定位哦~~。}，所以计算机得以被商业化。
一家造卡片统计机\footnote{恩，先进的自动化的...算盘}的公司因卡片机的大量销售而积累雄厚的财力和强大的销售服务能力，决定进入未知的电子计算机世界。
并于1944年，开发出自动顺序控制计算机，也被称为Mark  I。

\chatu{ibm_mark1}{Mark I计算机}

随后在1951年它决定开始开发商用电脑，聘请冯·诺依曼担任公司的科学顾问，1952年12月研制出第一台存储程序计算机
\footnote{也就是第一台冯·诺依曼结构的电脑。别和第一台（插拔线路编程的）ENIAC混淆哦！}，也是通常意义上的电脑。

没错，这家传奇公司就是IBM。
它开发的第一台电脑就是IBM 701。这是IT历史上一个重要的里程碑。从此计算机进入了程序控制时代，也就是冯·诺依曼结构。

IBM在商用电脑上获得了巨大的成功。那个时候，IBM就是计算机的代名词，有人亲切的称呼IBM为“蓝色巨人”。IBM生产的都是商用电脑，而且都是巨无霸（参考上面的图片），没有一款是适合家庭使用的。
微处理器\footnote{将计算机核心运算和控制部件做到一个芯片上。由于核心非常小，故而得名微型处理器。典型的微型处理器就是Intel为PC生产的80x86系列CPU。}的诞生创造了Apple这样神话。IBM迫切需要一台使用“微处理器”的家用电脑以和Apple这样的新兴公司对抗。

时间紧迫，IBM来不及自己开发操作系统，于是将操作系统外包给了一个西雅图的软件公司\footnote{微软啦！}。CPU也一改自己生产的传统，采购了Intel公司的8088微处理器\footnote{和8086使用相同的指令集，只是在关键的性能指标上比8086次点，因此也有好处：便宜。}，最终的结果就是组合出了一台家用微型电脑。俗称PC（Persional Computer 个人电脑）。

PC就是装配着Intel生产的8088/8086处理器，运行着微软的MS-DOS操作系统的微型计算机。

%\subsection{GNU HURD 内核}
但是8086处理器实在赢弱，缺乏实现多任务所必须的功能。所以DOS是个单任务单用户的操作系统。

\begin{insertnote}
\subsection*{小知识：UNIX需要什么样的CPU}
至少要具备2种功能：特权级别和非特权级别划分、支持内存分页（也就是有MMU，Memory Management Unit内存管理单元）。

UNIX将操作系统内核放到特权级别执行，应用程序则使用非特权级别。当处理器处于非特权级别的时候，对内存和IO的访问是严格受限的：访问不允许的IO设备会产生一个异常，该异常被操作系统捕获\footnote{所谓异常捕获，就是处理器发生异常的时候自动跳转到内核提供的异常处理代码。}，然后中止试图越权访问的程序；访问不允许访问的内存——如内核使用的内存——会导致内存访问异常，操作系统捕获这个异常，中止试图越权访问的程序。内核自身运行于特权模式。
有了CPU的这种级别划分，UNIX能保护操作系统不受异常程序的影响，提高稳定性和安全性。

UNIX支持多进程运行。每个进程\footnote{参考 \faqref{FAQ:Process}}之间的内存必须进行隔离。也就是说，一个进程只能访问自己的内存，无法访问其他进程的内存。
事实上这个访问不是通过“权限检查”（试图访问其他进程的内存导致进程被中止）实现的，而是通过“完全的隔离”实现的。
完全的隔离，也就是任意进程拥有完整的内存地址空间。并不存在内存地址被划分成N个区域，每个进程只能访问属于自己的$ \frac{1}{N}$的区域。而是每个进程拥有整个的地址。
从 0 开始到 $2^{32}-1$ （对应于32位的系统，如果是64位的则是 $2^{64} - 1$），访问的通通都是属于该进程自身的内存。
不同进程可以使用相同的内存地址访问属于自己的不同的数据。这种内存隔离被称呼为“虚拟内存”。虚拟内存是支持多进程运行必备的功能。

虚拟内存必须依靠CPU的MMU提供的分页功能才能实现。
\end{insertnote}

由于8086/8088处理器并不拥有MMU，也没有实现特权级别划分，所以根本不具备运行UNIX的条件。

就这样PC市场被UNIX忽略，直到80386的出现。80386是Intel生产的第一个具备运行UNIX自个的CPU。
但是就算在功能上具备了运行UNIX的资质，80386从UNIX的观点来说，实在是太慢了。所以谁也没打算将UNIX移植到使用80386处理器的PC上。

PC机就这样继续运行着DOS，以及后来的Windows。

\subsection{386BSD、 Minix催生Linux}

要说PC真的没有一个UNIX可以运行那是假的。商业UNIX抛弃了PC，但是黑客们却像发现了新大陆一样为PC着迷。于是，各种黑客版本的UNIX被移植到80386芯片上。
其中最著名的莫过于386BSD（将BSD移植到80386处理器的项目）和Andrew S. Tanenbaum教授以教学为目的开发的Minix。

Linus在大学里用过UNIX就为它着迷。于是决定为自己买一台计算机运行UNIX，这样就不要跑去学校的机房了。
由于PC性能不佳，跑着UNIX的工作站又太贵了。Linus决定购买一台PC。时值BSD正好处于和AT\&T的长期诉讼中。法律风险是大大的。于是Linus装起了Minix。
至于其他版本的UNIX，就没入Linus的法眼，呵呵。

Minix是Andrew S. Tanenbaum教授以教学为目的开发的UNIX，但是并不是一个“自由”的软件。而且Linus对于Minix有相当多的吐槽。不仅仅是Linus，许多人都为他们的吐槽写了补丁。
但是Andrew以“只用于教学目的的系统不想变的太复杂”为由而拒绝添加这些改进。一个拒绝改进的开发者+虽然宽松但是不自由的授权\footnote{
购买Andrew的《操作系统设计与实现》这本书就能获得Minix代码并授权个人使用。仅限于非商业用途。}\textrightarrow ~Linus~决定\textbf{为自己}开发一个UNIX兼容的操作系统。


\subsection{以GPL开源}

Linus说自己在Linux上做的最正确的决定就是用GPL保护了Linux。1911年发布的第一版内核并不是使用GPL协议发布的。
他使用了自定的一个协议，内容大概就是：{\it Linux自由免费，随便用。如果你对Linux做了修改，麻烦你把修改后的代码给我一份，谢谢。}

在黑客界，有一个古老的传统：\emph{尽量去发现工具，实在没有就自己写一个。} 这条传统似乎对软件授权协议也有效。
%
Linus大概觉得自己跳过了第一步的发现，于是他决定重新开始依照老传统做一次——这次他发现了GPL，于是他不再发明，将Linux置于GPL的保护之下。

Linus为第一版Linux编写的协议和GPL想要表达的意思差不多，但是起草过程中有律师帮忙的GPL更加的严谨。

从头开发的Linux源码和GPL授权为Linux的发展构建了坚实的后盾。即没有侵权的后顾之忧也没有不怕被侵权后缺乏法律的保护。Linux就这样替代了386BSD，Minix等等等等
成为了PC上最受欢迎的类UNIX系统。同时由于Linux使用GPL授权，填补了GNU缺失的内核 GNU Hurd。

全世界的黑客开始行动起来，推动着Linux快速的发展。Linux成了分布式开发的典范。

\fi

\chapter{初识Linux}

\begin{quotes}[0.63]{The Gospel of Tux}
Now may you code in the power of the Source; may the Kernel, the Libraries and the Utilities be with you, throughout all Distributions, until the end of the Epoch. Amen. 
\end{quotes}

\ifdefined\compileall

Linus发布他的内核的时候，Linux就只是一个内核而已。一个操作系统不当当依靠一个内核就能工作。还需要各种程序供用户使用。最少的，用户需要一个shell进行人机交互。
用户需要管理磁盘，管理他的文件。等等操作都需要依靠各种各样的程序进行。如果用户只是将Linux内核运行起来，那他将什么都做不了。所以用户需要的，是一个Linux内核和一套工具软件。这样的组合才能构成用户操作机器的环境——也就是说，一个操作系统。于是人们有时候简单的称呼这样的操作系统为Linux。但是我们得知道，这只是泛称，Linux的实际含义就是Linus和其他黑客开发的一个内核，并不是一个完整的操作系统。

由于Linux只是一个内核，所以用户需要的工具软件就需要向别的地方寻找了。所幸RMS宣誓的GNU已经开发的差不多了，就只差个名字叫GNU Hurd的内核。等等，我们不是有Linux内核了么？

没错！只要我们将GNU的系列软件和Linux内核搭配起来，不就有一个完整的操作系统了么？我们需要一个shell, GNU bash就是我们要找的。我们需要各种命令行工具， GNU coreutils就是我们要找的，我们需要编辑器，GNU emacs就行。我们需要开发工具编写自己的程序，GCC GDB GNU make等等就是我们需要的。 我们要的各种构成UNIX系统的要素GNU通通提供了。额，除了内核。所以RMS提议使用Linux内核+GNU userland\footnote{
前文提到过，UNIX需要CPU提供特权和非特权两个代码级别。内核的代码运行在特权级别。用户的程序运行在非特权级别。故而用户程序统称userland。}的操作系统应该叫做GNU/Linux。

那还等什么，赶紧到GNU的网站上下载吧！

\subsection{使用Linux内核的操作系统——发行版}
且慢！我还未拥有一个UNIX环境，GNU网站上提供的都是源代码，Linux也是源代码，我如何使用他们？编译？
那么我首先需要一个UNIX系统，并有一个能工作的编译器，才能编译这些软件吧！

这是一个先有鸡还是先有蛋的问题。

再说，就算我已经安装了一个Minix或者BSD系统，想编译出一个使用Linux内核的完整操作系统，也是一个巨大的工程吧！

俗话说，有需求的地方就会有市场。

于是有那么些人开始将Linux内核以及各种软件搜集到一起，预先从源码编译成二进制的可执行程序然后打包放到光盘上出售。并编写了方便的安装程序。只要把光盘塞入PC的光驱，跟着光盘里安装程序的提示一步一步做下去就能安装好一个完整的系统。

这种以Linux为内核的操作系统和其他的软件一起打包组合而成的套装就是发行版。


\section{发行版的意义}
当然，发行版的意义并不是简单的将软件打包和降低初学者的门槛哦\textasciitilde

第一：通过将大量的程序打包到统一的仓库，使得很多默默无闻的软件被显示到可安装软件列表上。增加了这些低调的软件开发者的潜在用户。虽然说用户多了开发者不一定有好处，但是光是想到有那么多人会用自己开发的软件，想想都觉得来劲不是么？

第二：许多软件特定版本的\emph{组合}可能会导致一些潜在bug。这些是软件开发者无从测试的。一个特定版本的发行版将基础软件的版本全部确定下来。只要保证这些版本的配合不会产生bug就可以了。由于用户使用发行版提供的软件，就减少了各种没有预料到的组合。所以在生产环境，随机的升级某个软件是非常忌讳的哦\textasciitilde。

第三：许多软件开发商形容Linux是活动的靶子。每过6个月\footnote{Ubuntu每6个月发行新的版本。}，系统各个软件都版本大换血。原先测试好的软件活不过6个月——下一个版本出来就出各种问题了。但是开源软件因为是发行版负责打包编译，所以可以避免这儿个问题。发行版维护者会及时发行某些软件因为另一些软件的升级而不工作了，通常简单的问题维护者自己就可以将错误修正——同时将修改回馈上游。开源软件可以放心的躲过每6个月一次的劫数——除了拒绝开源的。

第四：第三条提过，发行版积极的向上游回馈修改\footnote{咳咳，并不是每个发行版都那么好心。Ubuntu就曾经因为回馈不足引发舌战过。}。
有时候发行版回馈的补丁不仅是躲过6月劫数那么简单。一些发行版直接参与上游软件的开发，特别是和发行版结合紧密的软件。Fedora是RedHat赞助的社区发行版，默认使用GNOME桌面，而GNOME有相当一部分开发者就是RedHat雇员。


\section{包管理和软件仓库}

是什么东西促使一堆软件的集合变成了一个发行版呢？答案是包管理机制和围绕它建立的一个软件仓库。

\subsection{包管理}

人要用电脑，就免不了要安装软件。浏览器，办公软件，游戏等等。软件安装了，免不了要升级和卸载。
%人发行版需要

所有的程序都带上升级卸载程序，是个不错的主意，也是个愚蠢的主意。不符合UNIX的KISS原则。

符合KISS原则:
\begin{itemize}
\item 安装程序应该只有一份。
\item 所有的软件使用统一的方式安装，不能区别对待

\item 使用方便。提供直观间接的命令行或者图形界面。

\item 干净卸载。卸载不能留下任何不该留下的东西。所谓燕过不留痕。

\item 解决依赖问题。一个软件或多或少的会依赖其他的软件。最明显的例子，如果系统里没有安装libc，所有的软件都不可能运行。一个软件要想正常工作，它依赖的软件必定也要安装。　包管理要最终软件之间的相互依赖关系。安装一个软件的时候能将它依赖的软件也安装，卸载软件的时候能将依赖的软件（如果没有其他的软件也依赖它的话）一并删除。

\item 打包方便，不能方便了最终用户而是发行版维护人员痛苦不堪。

\end{itemize}

各个发行版使出浑身解数，开发了各种软件包格式和对于包管理器。比较流行的也就只有两个而已：RPM和deb。

\subsection*{RPM}

RPM是Redhat Package Manager的缩写，不过后来RPM被Redhat以外的发行版采用的，似乎RPM成了另一种缩写：Rpm Package Manager\footnote{吐槽吧，学的GNU啊。递归缩写}。

RPM格式的扩展名是.rpm，分源码包和二进制包。源码包将rpm规格说明文件（.spec）和程序源码打包为一个.src.rpm包（或者.srpm） 。源码包可以用来构建二进制rpm包。
要安装rpm包，使用rpm命令进行。具体使用请参考man手册。比较常规的用法我会在
%\secref{软件管理}　介绍。
第\ref{软件管理}章~\nameref{软件管理}介绍。

\begin{insertnote}
\vskip 0.5cm
man手册的用法很简单啦，打man后面跟个要帮助的命令，rpm就是
\begin{code}
\$man rpm
\end{code}
按q退出查看。PgUp和PgDown可以翻页。
\end{insertnote}

RPM被广泛使用，包括RHEL、Fedora、OpenSuSE、RedFlag都是使用的rpm包格式。

\subsection*{deb}

包管理世界的另一大流派是deb。deb因其所使用的软件包扩展名为.deb而得名。实质上是Debian的缩写。因为它是Debian专用软件包格式。
因为Ubuntu是Debian修改而来，自然也采用了deb包管理。

和RPM一样，deb也有源码包和二进制包。从源码包可以编译得到二进制包。要安装deb包，使用dpkg命令进行。具体使用请参考man手册。比较常规的用法我会在第\ref{chap:软件管理}章~\nameref{chap:软件管理}介绍。

\subsection*{ebuild}

Gentoo的ebuild严格来说并不是软件包格式。rpm deb之类的，可以直接安装对应的二进制包。包管理器检查包的依赖，依赖检查后将包里的文件解压安装到系统里就可以了。

Gentoo的ebuild却并不是这样的。ebuild大致是一个BASH脚本（不能直接执行，只能由包管理器调用）。Gentoo的包管理器通过执行ebuild完成“自动下载源码，解压，编译，安装”这样的操作。

%TODO
\todo{更多ebuild介绍}

\subsection{软件仓库}

发行版搜集了各种软件，并编译成二进制包。通常发行版收集的软件包数量以万计。这么多包是不可能一次性全部装到一张光盘上的。也完全没有那种必要。用户拿到光盘的那一刻开始软件就已经过时了。

如果将所有的包都集中存放到服务器上，用户要安装的时候直接到服务器上获取就可以了。用户获得的软件总能保证是最新的。只需要比较本地安装的软件版本和服务器上存储的最新版本就可以知道是否有更新了。

这样在服务器上集中存储的软件包就集合成了一个仓库。发行版的工作就是维护这个仓库，时不时的向仓库更新软件。

软件仓库可是个伟大的发明。乔帮主从Linux社区学到了软件仓库创立了App Store赚了大发。微软也开始学软件仓库的思想，Windows 8也开设了应用程序商店。

限于仓库的格式问题，不同的发行版有不同的工具让用户访问软件仓库。

RPM系的发行版，使用yum\footnote{不要和rpm混淆哦！一个是安装软件包的，一个是访问软件仓库的。}访问仓库。用户可以下载　XXX.rpm 然后使用　rpm -i XXX.rpm 安装这个包，但是更便捷的做法是直接使用 {\tt yum install XXX} 安装名为XXX的软件。yum会到仓库里搜索XXX软件并将对应的rpm文件下载下来，如果有未安装的依赖，将依赖的软件也下载下来。然后调用rpm安装。

\begin{example}{Fedora下安装手写输入法}
手写，handwrite，搜索
\begin{code}
\#yum search handwrite
\end{code}

在结果页面，有个叫　ibus-handwrite的软件。ibus是一个输入法框架，自然ibus-handwrite就是手写输入法了。
开始安装（注意需要使用root帐号运行，参考\faqref{FAQ:su}）
\begin{code}
\#yum install ibus-handwrite
\end{code}
回答yes后ibus-handwrite就被安装到系统上了。
当然，前提是保证网络连接：）
\end{example}

要删除软件，使用 yum remove XXX.

对deb系的发行版，使用apt-get（俗称超级牛力）哦。基本用法和yum一致。apt-get install 安装软件，apt-get remove 卸载软件。
区别就是，yum是在线搜索，apt-get是离线搜索。你需要每天使用apt-get update（也不需要每天啦！
长久没用apg-get的时候先执行一下apt-get update就行了。）刷新缓存，
之后apt-get就使用本地缓存的软件包信息而不用每次都连接到服务器进行搜索。

另外，更新系统yum使用的操作是　yum update，而apt-get同样的操作用来更新本地软件包信息缓存了，故而更新系统在apt-get下使用的命令是apt-get upgrade。

而Gentoo这样的系统嘛，也有非常强大的工具：{\bf emerge}。emerge　XXX就可以安装名为XXX的软件。

\begin{example}{Gentoo下安装vim}
要安装vim非常简单，emerge就可以了
\begin{code}
\#emerge　vim
\end{code}
emerge会忠实的到vim的网站上下载vim的代码，解压，编译，然后安装。如果vim依赖的软件还未被安装的话，emerge会首先安装它们。

\end{example}

\subsection*{在仓库中搜索}

虽然仓库很好用，但是在用的时候大家都会有一个疑问：怎么知道一个软件对应的包的名字？

这就需要到仓库中搜索需要的软件了。搜索条件一般为软件的名字。通常软件包的名字和软件本身的名字一致。但也有例外。就拿glibc来说，在deb系里，它的包名是libc6。
这些问题都有历史渊源，说来话长，咱就不废话了。总之记住：软件包的名字并不总是和软件本身的名字一致。所以有时候需要搜索后才能安装。另外一个问题是，有时候我们只是希望安装一个“类型”的软件，比如想安装个画图软件，但是恰巧并不知道任何一款画图软件的名字。

针对这2个问题有两个办法。

第一个问题，属于特定发行版下的“软件包改名”问题。一般使用发行版的软件管理工具的搜索命令。以软件本身的名字作为关键词搜索即可。比如　yum search vim、apt-cache search sudo、emerge -s vnc。

第二个问题，命令行下无解。如果使用图形工具可以方便的浏览仓库中所有的软件。还是分门别类的哦！当然，最快的办法莫过于祭出伟大的Google。
\fi

\chapter{安装 Linux}

\begin{quotes}[0.5]{insomnia}
It only takes three commands to install Gentoo!
\end{quotes}

\begin{quotes}{Anonymous}
“Where there is a shell，there is a way!”
\end{quotes}

\ifdefined\compileall


Linux介绍了那么多，是时候将它安装到我们的电脑上了。接下来要边安装边学习哦。要安装一个Linux系统，我们总是选择一个发行版来安装。那么到底选择什么样的发行版呢? 理想的发行版通常是这样的一个发行版：它非常第容易安装，非常的容易适用；软件仓库里的软件非常丰富，我可以用包管理安装一切软件而不需要自己手动安装；对于一些软件我希望能容易的进行定制；快，非常快。在对各个发行版做里权衡后，我选择Gentoo作为本书中介绍Linux知识所使用的平台。


\begin{notice}
 注意：Gentoo只是我们选择的众多发行版中的一个，我并不打算把本书变成一个Gentoo入门教程。如果这样，读者还不如选择阅读Gentoo官网上的手册。我尽量将所涉及到第知识通用化。如果这样的知识点只适用于Gentoo, 我会进行说明并辅之以其他发行版的等位操作\footnote{不同发行版达到同样的目的采取的不同操作。}。
\end{notice}

本书并不打算引起发行版之争\footnote{指不同的人使用不同的发行版，并常常在网上争论自己使用的发行版要比别人的优秀}，
但是又必须解释为什么使用了Gentoo这个非常罕见的发行版而不是同其他书本那样采用最流行的发行版，比如Ubuntu或者Red Hat系。

最直接的原因也许就是作者本身使用的就是Gentoo，以自己最熟悉的发行版写作，可以避免不熟悉系统进行错误的讲解\footnote{这确实是一个理由，但是Gentoo必定有比其他的发行版更值得作为本书的平台。}。
当然，如果仅仅这一个理由未免太牵强。
{ \it Gentoo可以让读者将精力放到学习Linux本身，而不需要关注发行版的细节，这就是我采用Gentoo的原因。}
另一个能做到这种效果的发行版为LFS\footnote{严格来说，LFS 算不上发行版，只是一个教你如何手工编译出一个能用的系统的手册。}（Linux From Scratch）。
但是LFS太繁琐，恐怕还没有学习到Linux的知识就已经中途放弃了。


\section{安装准备}

当然是准备一台电脑啦！Gentoo虽然号称为本机编译优化，但是编译本身非常耗费系统资源，所以一台主流配置的电脑是必不可少的！
\footnote{Gentoo支持的硬件架构非常多，不是一定要用个人电脑的。
但是作为初学者，我们就用最常见的个人电脑，不折腾各种稀奇古怪的电脑，最重要的是，Gentoo对x86体系也是支持的最好的（这不废话么，用户量摆那里么）。}


除了一台电脑，还要准备个Fedora或者Ubuntu的LiveCD\footnote{参考附录A. FAQ 的“什么是LiveCD”。
}，这里不推荐Gentoo自己的LiveDVD\footnote{同LiveCD。}，因为前者更容易获取和使用。

当然，还有网络。Gentoo虽然可以脱离网络进行安装，但是需要将所可能需要用到的源码提前下载，这是非常麻烦的。所以最好准备一个安装的时候可以访问的网络，也就是说LiveCD 环境里必须能使用网络。Fedora 或者Ubuntu那样的LiveCD 通常提供了完整的桌面环境，使用这些桌面环境的网络设置，让LiveCD 能访问互联网就可以了。你可能需要使用 NetworkManager 提供的图形客户端进行配置，通常是在桌面右上角的系统通知区域。要使用 NetworkManager 配置网络，请参考第6章的内容。

嗯嗯，不过就算是x86还分为x86\_32（ 没有特殊指明的情况下x86特指x86\_32） 和x86\_64\footnote{因为历史原因，AMD率先推出支持64位运算的x86CPU，所以Gentoo管x86\_64叫amd64，这和x86又被称呼为 i386 是一样的原因。}呢。
由于x86\_64是在x86的基础上发展起来的，所以x86\_64向下兼容x86。 也就是说，64位的CPU可以运行32位的软件。但是反过来就不行了。

既然存在两种x86，选择起来就费脑筋了。怎么知道自己的CPU是x86\_64的还是x86\_32 的呢？答案是，现在生产的CPU都是x86\_64了。当然如果你想准确的知道CPU到底支持不支持64位指令，在Window下你可以选择执行CPU-Z 软件获得CPU信息，如果是Linux就简单的多，在终端下执行lscpu命令就可以了。

（1）  Linux下使用lscpu：

在终端下执行lscpu，该命令的第一行输出就是CPU的体系结构。在我的电脑上，其输出为：

\begin{code}%
Architecture:        x86\_64\\
CPUop-mode(s):        32-bit, 64-bit\\
Byte Order:            Little Endian\\
CPU(s):                8\\
On-lineCPU(s) list:   0-7\\
Thread(s) per core:    2\\
Core(s) per socket:    4\\
Socket(s):             1\\
Vendor ID:             GenuineIntel\\
CPUfamily:            6\\
Model:                 42\\
Stepping:              7\\
CPUMHz:               1600.000\footnote{如果CPU支持动态降频，这里显示的是当前频率而不是CPU的标称频率。}\\
BogoMIPS:              6420.96\\
Virtualization:        VT-x\\
L1d cache:             32K\\
L1i cache:             32K\\
L2 cache:              256K\\
L3 cache:              8192K\\
\end{code}

第一行的Architecture:        x86\_64就表明此CPU的结构是x86\_64。


\vskip 1em
（2）  Windows 下使用CPU-Z：

\chatu{cpuz}{cpuz截图}
如图\thefigure{}所示。我红圈标注的EM64T表示CPU支持64位指令。对于AMD系的CPU，您大可放心，全部支持64位指令。


\begin{insertnote}
\subsection*{小插曲  64位 vs 32位}

如果是x86\_64的CPU，应该选择amd64\footnote{AMD64 就是x86\_64。}
的Gentoo还是x86的呢？既然现在的CPU都支持 64位指令的，为何还有那么多人使用 32 位的系统？

x86\_64虽然能执行x86代码，但并不是没有代价的。Windows下大部分现有的程序都是32位的，如果安装的64位Windows，32位软件并不能100\%兼容。而且所有的64位软件同时有32位的版本，所以Windows用户非但没有迫切的升级到64位系统的意愿，放而有为了兼容性考虑不得不安装32位系统的处境。

64位系统如果不执行64的程序，那性能还不如直接使用32位的系统和32位的程序。所以在 Windows 统治的世界里，64 位CPU的地位非常尴尬。32位CPU通过PAE\footnote{Physical Address Extension，物理地址扩展。允许32位CPU使用36位地址从而突破4G内存限制。}同样能获得 >4G 内存的支持（非服务器版Windows人为阉割32位版本的PAE支持，故而32位桌面版Windows不能支持超过4G内存。）。

大概也就是因为这个原因，所以不明真相的群众就以为64位除了支持大内存外一无是处。继续固守32位。64Bit vs 32bit的战斗因此打响。

x86\_64对x86的提升不仅仅是更大的内存，同时增加的有

(1)  更大的寄存器 64位CPU自然拥有64位宽度的寄存器。

(2)  更多的寄存器x86\_64相比x86增加了8个通用寄存器,增加8个MMX寄存器。

(3)  默认支持SSE\footnote{SSE指令是Intel为x86添加的扩展，加快了浮点运算的速度。}指令集。

(4)  更有效的指令编码。

这些都无形中增加了x86\_64指令编译的软件的运行速度。所以，如果你的CPU支持x86\_64，请必定选择x86\_64。
Linux上没有Windows上遇到的问题。
Gentoo下几乎所有软件都是本机编译的，自然不存在兼容性问题。
而个别非开源软件也多数拥有 64位版本，更不需要操心兼容性问题。
\end{insertnote}

好，我们列一个清单：


\begin{itemize}
\item[ \checked] 一台电脑
\item[ \checked] 一张LiveCD
\item[ \checked] Internet
\item[ \checked] 本书或者Gentoo 安装手册。
\item[ \checked] 耐心，很多很多的耐心
\end{itemize}

你准备好了么？

准备好的话，将LiveCD塞入光驱，开始一步步安装属于你的Gentoo吧。有关如何使用LiveCD请参考\faqref{FAQ:UseLiveCD}。
启动LiveCD，首先配置好网络环境，有关网络环境的配置请参考 \secref{sec:ifconfig}。


\section{shell 快速入门}\label{sec:quickbash}

由于安装Gentoo需要执行shell命令，所以我需要讲解一下最基本的shell命令。高级的shell操作讲解要到

机器只能接受二进制指令，可是人需要工作在更高一层。人希望直接提供文本形式的指令由机器执行。shell 是操作系统提供给用户的解释器。shell把来自用户的命令转化为机器指令并执行。

命令虽然是文本，但也不是自由格式的。计算机是死的，自然不能接受人类的语言，只能是比较固定的几个格式。

一条shell命令的最常用格式为:

\begin{code}
命令 [选项] [目标]
\end{code}

选项是可选的，可以有一个或多个；多个选项间用空格隔开；有些命令没有选项；目标通常是一个文件的文件名，或者目录名。
一行可以写多个命令，多个命令用“;”隔开。一个命令也可以跨行，在最后行尾使用  “\textbackslash”  即可在下一行继续写。


请看一个命令做为例子： 

\begin{code}
root@gentoo \textasciitilde \# ls -l --all /  
\end{code}
在这里：

\begin{itemize}
\item root@gentoo \textasciitilde \# 是shell提示符。\#是root账户，普通账户是\$。root@gentoo 表示当前登录用户为root,当前登录机器是gentoo,后面的 \textasciitilde{}　表示当前目录为　\textasciitilde{}，　\textasciitilde{}　在shell中表示用户主目录，也就是　\$HOME。root的主目录为 /root,　普通用户为/home/用户名。
\item -l 是一个短选项。短选项 “-”开头，后面跟单个字母。在这里 -l 表示以“长格式”列出。所谓“长格式”就是相对与一般格式，除了要列出文件名，好要列出文件所有者，权限模式，文件大小
\item -{}-all 是一个长选项。长选项使用两个“-” 开头，后面跟一个单词。在这里 -{}-all 表示列出所有 文件，包括隐藏文件。
\item / 在这里为操作对象参数，在这里是一个目录。表示要列出 / 下的文件。
\item ls 的所有参数都是可选的，如果不使用任何参数，ls 列出当前目录下的文件。

\end{itemize}

shell命令分两种，一种是内部命令，另一直是外部命令。内部命令是由shell直接执行的\footnote{最常用的内部命令是cd。}。
外部命令则是独立的程序。比如ls就是独立程序，它的路径为/bin/ls。shell使用一个PATH环境变量\footnote{参考附录A FAQ“什么是环境变量”。}来查找外部命令。

现在我们只需要知道shell最常用的命令形式就可以了。在以后的学习过程中，我会像大家介绍 shell 的其他命令格式，包括各种循环和函数使用。现在还要一个比知道更多的shell命令更重要的知识需要知道：
\begin{center}\em
shell字符串扩展
\end{center}

还是 ls 这个命令作为例子，假设当前目录下有“a.txt”“ab.txt”两个文件。执行下面的命令：

\begin{code}
root@gentoo \textasciitilde \# ls *.txt
\end{code}

结果就是列出 a.txt ab.txt 两个文件。

在这里：
\begin{itemize}
\item *.txt 会被扩展成 a.txt ab.txt
\item a.txt 和 ab.txt 被作为2个参数传递给ls。
\item “*”是通配符，表示匹配任意多个字符；“?”表示匹配任意的一个字符。

\begin{insertnote}
注意：shell的字符串扩展不仅仅支持*和?这样的通配符，更高级的用法将在后续章节将到。
\end{insertnote}


\end{itemize}


比如你执行的是下面这个指令：

\begin{code}
root@gentoo \textasciitilde \# ls ?.txt
\end{code}

结果就只显示a.txt文件。要取消shell字符串扩展，加引号即可，如下面这条命令：

\begin{code}
root@gentoo \textasciitilde \# ls "*.txt"
\end{code}

执行结果为：

\begin{code}
ls: 无法访问*.txt: 没有那个文件或目录
\end{code}


我们知道，ls *.txt 相当于 ls a.txt ab.txt，可是如果执行下面的命令

\begin{code}
root@gentoo \textasciitilde \# ls "a.txt ab.txt" 

ls: 无法访问a.txt ab.txt: 没有那个文件或目录 

root@gentoo \textasciitilde \# 
\end{code}

shell将“a.txt ab.txt”整体（不包含引号）作为了一个参数传递给ls了。

要理解shell的这种行为，首先要明白shell是如何传递参数给程序的。

\begin{insertnote}
\subsection*{shell如何传递参数给命令}

编写过C程序的同学都知道，C程序的入口点在main，main的原型为

\begin{code}
int main(int argc, char * argv[]);
\end{code}

argc为参数列表中参数的个数，argv是个参数列表。argv[0]是程序自己，argv[1]为第一个参数，依次类推。
参数是直接由启动该程序的操作系统传递进来的，不需要还调用什么函数来获得参数列表
\footnote{Windows下确实是调用GetCommandLine()获得参数列表的，由C语言启动函数（真正的入口点）负责获得参数列表然后调用main()。}。
操作系统知道调用该程序的参数列表是因为程序是由execv()系统调用加载的。原型：
\begin{code}
int execv(const char *path, char *const argv[]);
\end{code}
path就是要执行的程序。
argv的最后一个由NULL结尾，故而内核会知道有几个参数，转化为main里的argc。

\end{insertnote}

在本例中，
shell
启动ls程序的参数应该是这样的
\begin{code}
argv[0]="ls";\\
argv[1]="a.txt ab.txt";\\
argv[2]=NULL;

execv("/bin/ls",argv);
\end{code}

所以ls会把"a.txt ab.txt"（不包含引号）作为一个参数。也就是说，参数和参数之间的分隔是由shell完成的\footnote{Windows下的程序，参数和参数之间的分隔是C库的启动函数完成的。GetCommandLine()后分割命令行参数再调用main()}。


那么，shell进行字符串扩展和命令行参数分割的时候遵循的规则是这样的：

\begin{itemize}
\item shell进行字符串扩展的时候，只对游离的*和?进行扩展。如果*和？处于引号之内则不进行扩展。
\item shell以空格为参数分割的界限。但是对于使用*和?扩展出来的字符串，以文件名为界限。比如文件名内的空格不会成为分割的界限。
\end{itemize}


如果执行的是下面的命令
\begin{code}
root@gentoo \textasciitilde \#ls　*.txt c.txt
\end{code}

shell对命令解析的步骤是：

首先进行命令行字符串扩展，*.txt 扩展为　a.txt ab.txt；接着进行参数分割，这样就是3个参数: a.txt ab.txt c.txt；最后将这3个参数传递给ls。也就是执行下面的操作

\begin{code}
argv[0]="ls";\\
argv[1]="a.txt";
argv[2]="ab.txt";\\
argv[3]=NULL;

execv("/bin/ls",argv);
\end{code}

\subsection*{转义序列}

既然说过，空格作为参数分格符。如果用引号将带有空格的字符串引起来，则两个引号之间的所有内容都作为一个字符串。也就是说，那之间的空格就失去了分割参数的功能。
引号还可以用于避免字符串展开。功能如此强大，结果就是字符串本身如果就打算使用引号怎么办？shell将搞不清楚引号到底是代表这是一个字符串还是引号本身是字符串的一部分。

shell引入了转义操作。使用在字符串内部使用{\textbackslash\tt{}"}就可以表示这是一个引号。

\begin{example}{显示带 引号的字符串}
\begin{code}
root@gentoo \textasciitilde \# echo "你好引号{\textbackslash\tt{}"}"

你好引号"

root@gentoo \textasciitilde \#
\end{code}
\end{example}

\subsection*{本节回顾}

目前shell知识讲解仅限于能让读者小小入门，并顺利完成Gentoo的安装。现在小结一下shell的简单用法。

\begin{itemize}
\item shell以文本模式命令行作为用户界面。
\item 命令行的规则是　命令+参数，以回车键代表输入完成，请立即执行命令。
\item 参数之间以空格隔开。
\item 文件名参数可以用通配符，shell会自动扩展后作为参数传递给要执行的命令。扩展后的文件名是每个文件名一个参数而不管文件名本身是不是包含空格。
\item 如果需要避免shell对通配符进行shell扩展，加引号。
\item 避免空格成为参数之间的分割的办法也是将带有空格的参数用引号引起来。
\end{itemize}

掌握这些shell知识应该用于完成安装Gentoo足够了。


\section{引导管理器}\label{sec:grub}
Linux的启动过程是一个令人着迷的过程，首先是硬件的初始化——也就是固件初始化——接着是固件加载并执行引导程序。
引导程序负责将内核启动起来。然后内核执行第一个程序/sbin/init，后续的启动就交给init完成了。

引导程序自身是如何实现的我们不必关心，倒是引导程序如何被固件找到载，以及引导程序如何找到内核是个重点话题。

IBM创造PC的时候给PC开发了一个名字叫BIOS（Basic Input Output System）的固件。该固件被存储在一个ROM\footnote{只读存储器，断电后内容不会消失，存储内容生产的时候固化，无法修改}上。 BIOS提供了基本的输入输出功能，负责机器的上电自检（POST，Power On Self Test）和操作系统加载（其实是加载MBR。MBR再加载DOS。）。操作系统运行后BIOS作为基本输入输出设备的驱动程序存在\footnote{Windows、Linux这样的操作系统使用自己的驱动，并不依赖BIOS提供的驱动。}。

也就是说，固件初始化完成后，就读取硬盘的第一个扇区（MBR）并执行这一小段程序。那么，固件并不关心你用的是什么操作系统，用的是什么引导程序，它只是遵照约定，执行MBR而已。

微软的操作系统，不管是DOS还是Windows，其MBR都是一样的。就是查找活动分区\footnote{分区的一个额外属性。活动或者非活动。值得一提的是DOS无法访问非活动主分区。}并执行分区的第一个扇区（PBR）。所以微软操作系统的真正引导程序在系统安装分区（C：）的第一个扇区里，也称呼为分区引导扇区，Partition Boot Record。这是GRUB所以能支持引导Windows的关键原因哦。我们下文马上会介绍到。

Linux下普遍使用的引导程序——或者说是事实标准——是GRUB（GRand Unified Bootloader）。其大名是如醍醐灌顶，如雷贯耳，谁人不知，谁人不晓啊。

BIOS这种固件使用简单的伎俩查找引导程序。不过MBR只有区区512个字节，还要留64个字节给分区表，用于引导程序就只有448个字节了。想要做功能多点的引导程序都不行啊！空间问题。

GRUB使用的办法就是接力加载。存在MBR上的那份代码只是完整GRUB的加载器。因为完成的GRUB需要支持ext2/ext3/ext4/fat/iso9660等等各种文件系统，以便读取存储在磁盘上的内核，还要支持配置文件，还要有命令行模式，故而GRUB的代码不可能少于448个字节。

GRUB完整启动后，读取配置 /grub/grub.conf\footnote{注意，此/非Linux启动后的/，而是安装了GRUB的那个分区。如果没有/grub.conf 会查找 /boot/grub/grub.conf。所以在Linux环境下配置文件路径就是/boot/grub/grub.conf而不管/boot分区是不是独立的。}，显示操作系统选择列表。

Linux下第一个硬盘的第一个分区是/dev/sda1，但是GRUB下却是从0开始的，(hd0,0)。这点需要注意。grub.conf的引导项配置以title开始，例如

\begin{code}
title Gentoo\\
kernel /vmlinuz ro root=LABEL=GENTOO\\
initrd /initrd.img
\end{code}

这里的 /　指的是GRUB安装分区。如果GRUB被安装到/boot分区，那么grub下的/对应的Linux下的路径就是/boot。
kernel　命令指示后面加载一个Linux内核，内核文件后面跟随的字符串将成为内核参数。好多参数可以控制内核的行为。参数的介绍参考　\faqref{FAQ：kernelparamter}

如果 vmlinuz不在 /boot分区，比如想在命令行模式启动其他内核的时候，可以用　root (hdX,Y) 的格式指定一个GRUB根。如下面的

\begin{code}
title Test-kernel\\
root (hd0,1)\\
kernel /usr/src/linux/arch/x86/boot/bzImage ro root=/dev/sda2
\end{code}

kernel用来加载一个内核，initrd 用来加载initramfs。对启动Linux来说，initrd并不是必须的。initrd的详细解释可以参考　\faqref{FAQ:initramfs}。

不过，这样的GRUB可是对UEFI无效的哦~。

\subsubsection{UEFI}

IBM发明PC的时候远见不足，BIOS的历史局限性逐渐显现。BIOS最大的禁锢就是运行于16位模式。在连32位都要过时的年代，BIOS却还是16位的。

彻底的革新，就是要从固件开始就让CPU处于64位模式运行，固件的代码再也不用遇到16位模式下的限制了。固件能集成更多的功能，甚至能直接访问文件系统，集成功能强大的引导程序。

UEFI就是Intel提出的新的BIOS。和BIOS相比，UEFI对于操作系统来说区别就是：
\begin{itemize}
\item 纯64位启动。从UEFI到GRUB（并不是必须的存在，UEFI可以直接启动Linux）到Linux，中间没有16位代码。
\item UEFI使用GPT分区表，支持访问FAT文件系统。
\item 引导程序是EFI系统分区（标记EFI系统分区属性，文件系统为FAT）里的文件，扩展名为.efi。
\item Linux可以由　grub.efi　引导，也可以将自身直接作为　vmlinuz.efi 被UEFI直接引导。
\item 在操作系统环境下可以直接控制UEFI的引导选项。BIOS下必须进入BIOS设置程序才能设置。
\end{itemize}

要使用UEFI启动，需要两个要素：
\begin{itemize}
\item 一个UEFI能识别的分区，里面存放引导程序。UEFI使用GPT分区表，支持访问FAT文件系统。需要将一个FAT分区的分区类型设置为EF00。
\item 告诉UEFI从哪个引导程序\footnote{扩展名.efi，也就是EFI可执行程序。}引导。可以在UEFI的设置界面设置，也可以使用efibootmgr命令行实用工具设置。
\end{itemize}

能用于UEFI的引导程序并是上文所提到的grub。您有三个选择：一是使用打了EFI补丁的GRUB；二是去使用GRUB2；三是使用内核自己\footnote{没错，编译好后的内核文件本身也是一个合法的EFI可执行文件！当然，前提是打开EFU stub选项。}。GRUB2的语法比较晦涩难懂，暂时没什么必要的话还是继续使用GRUB比较好。我会在后面的章节再次详解GRUB和GRUB2的配置文件语法，如果将来想换到GRUB2进行引导的话还是可以修炼的哦。

打过补丁的GRUB可以从Fedora\footnote{Fedora从18开始BOOTX64.EFI不再是GRUB了。默哀。追Fedora还真难。没准等读者买到我的书的时候，Fedora 18已经成了落后的过去式了。}的LiveCD上提取。光盘根目录下的 /EFI/Boot/BOOTX64.EFI 就是EFI版的GRUB。如果没有Fedora光盘，其他的支持EFI启动的LiveCD上也有。通常在光盘根目录的EFI子目录里。

也可以上网找grub的rpm包
\footnote{%\begin{longtable}{p{\textwidth}}  
http://mirrors.163.com/fedora/releases/17/Everything/x86\_64/os/Packages/g/grub-efi-0.97-93.fc17.x86\_64.rpm
%\end{longtable}
}
然后提取 grub.efi。

获得grub.efi（BOOTX64.EFI）后，将grub.efi拷贝到EFI分区。\footnote{其实我比较推荐直接将EFI系统分区挂载为/boot，但是有的人更喜欢的做法是将EFI系统分区挂载到/boot/efi/。}

拷贝完成后执行efibootmgr将其添加到UEFI的引导程序列表中。

\section{init 系统选择}

Linux内核被引导后，会自动加载 /sbin/init 程序，之后系统的启动就完全交给了init程序。由于init是第一个启动的程序，所以他的PID\footnote{Linux用一个数字代表一个进程。需要对进程操作的地方可以使用该数字。比如向进程发送信号，就是用PID标识发送的目标的。}是1。\footnote{PID 0是个特殊进程（或者是根本就不是进程吧），每当系统空闲下来的时候PID　0才会执行，PID 0只做一个事情，就是让CPU进入节能模式。而且PID　0的个数就是CPU的个数。每个CPU都有一个属于自己的PID 0。所以系统真正的进程号是从1开始的。}

\begin{insertnote}
\subsection*{Linux启动过程}
计算机的启动步骤为 固件初始化 \textrightarrow{} 引导程序　\textrightarrow{}　内核加载并初始化　\textrightarrow{}　执行首个程序负责剩余的初始化　\textrightarrow{}　检查并挂载文件系统　\textrightarrow{}　启动后台服务　\textrightarrow{}　开启登录　\textrightarrow{}　（用户登录或者自动登录）　\textrightarrow{}　执行shell\footnote{Linux下命令行的shell就是bash。但是图形界面也可以算是一个shell哦，Desktop Shell是也。}。

PID 1 的任务就是操作系统启动过程中的{\em 最后一英里（last mile）}。

其实PID本身也不做什么操作，它只是依次启动其他服务：

启动磁盘检查服务，启动文件系统挂载服务（这些都死一次性服务，系统启动过程中一次性执行）；

启动udev服务，启动Dbus服务，启动电源管理服务等等（持续性服务，系统运行过程中一直执行的）；

最后启动登录服务。

命令行登录服务为 agetty 。在文本终端上显示登录提示接受用户登录。通常开启六个。tty1-tty6。使用组合键Ctrl+Alt+(F1-F6)访问。
图形登录服务（显示管理器，Display Manager）为　GDM或者KDM或者XDM .... 好多，就看用户的选择了。通常发行版有一个默认的。一般和使用的默认桌面有关。
默认使用Gnome桌面的发行版默认的显示管理器是GDM，默认使用KDE桌面的发行版默认使用的显示管理器就是KDM。也可以使用简洁的XDM。各有特点，但是功能都是一样的，可以随意互换。GDM一样能启动KDE桌面，KDM启动GNOME也没问题。

\end{insertnote}

简单的来说，启动操作系统过程中PID　1的任务就是bring up userspace。此外，PID 1还减负者收养孤儿的重任。

\begin{insertnote}
UNIX进程有父子关系。比如说你在shell下执行了一个程序，执行这个程序的进程就成为shell的子进程。子进程还可以有子进程的子进程。形成一个树形的进程关系。
\end{insertnote}

如果已经进程在子进程退出前就退出，那这个子进程就会成为没有父进程的“孤儿”。UNIX同时又规定一个进程必须有父进程，所以孤儿就被 PID 1收养成为PID 1 的子进程。


\subsubsection{sysvinit+openrc}

多数Linux发行版使用传统的sysvinit作为PID　1。　
sysvinit有一个运行级别概念。每个级别用一个数字表示。预定义的运行级别如下：

\begin{longtable}{|l|p{0.7\textwidth}|}\hline
0　& 关机\\\hline
1 & 单用户（救援）模式　\\\hline
2 & 无网络模式　\\\hline
3 & 多用户模式　\\\hline
5 & 图形模式　\\\hline
6 & 重启　\\\hline
\end{longtable}

不同的运行级别对sysvinit本身并没有特殊含义。只是人为这样的运行级别设定而已。
sysvinit使用配置文件/etc/inittab决定启动过程。其中有一行
\begin{code}
id:3:initdefault:
\end{code}
决定默认进入运行级别3。3替换改为5就可以默认进入运行级别5。

/etc/inittab告诉了sysvinit每个运行级别下应该执行什么程序。例如inittab通常包含这么6行：
\begin{code}
c1:12345:respawn:/sbin/agetty 38400 tty1 linux\\
c2:2345:respawn:/sbin/agetty 38400 tty2 linux\\
c3:2345:respawn:/sbin/agetty 38400 tty3 linux\\
c4:2345:respawn:/sbin/agetty 38400 tty4 linux\\
c5:2345:respawn:/sbin/agetty 38400 tty5 linux\\
c6:2345:respawn:/sbin/agetty 38400 tty6 linux
\end{code}

表示：在运行级别1 2 3 4 5下，启动c1，启动c1所需要执行的命令为/sbin/agetty 38400 tty1 linux, respawn表示服务退出后立即重新开启。
c2 c3 c4 c5 c6 在运行级别 2 3 4 5下启动，　1 下不启动。因为1是单用户模式嘛。

图形登录界面是由下面的设置控制的（以早期Fedora为例，现在的Fedora不使用Sysvinit）
\begin{code}
x:5:once:/etc/X11/prefdm 
\end{code}
/etc/X11/prefdm 查找用户偏好的DM\footnote{一般是GDM或者KDM。}并启动之。
可以看出图形管理器是在运行级别5开启的。


{\tt c1:12345:respawn:/sbin/agetty 38400 tty1 linux} 这一行的上面还包含了一些配置用于启动　/etc/init.d/ 下的服务程序。各个发行版的做法都不一样。没有借鉴意义。

总之最后的结果是是\textbf{按循序}启动 /etc/init.d/　下的脚本
（因为运行级别的关系不是每个脚本都会被执行。通常如果运行于级别5的话，所有服务都开启，也就是都会被执行。）。

我在这里强调了一个次序。因为服务和服务之间是有明显的依赖关系的。举些例子吧：提供HTTP服务的服务apache必须于mysql之后启动。因为php页面需要和mysql数据库交互；NetworkManager依赖DBus,因为它需要和DBus提供通信功能；NetworkManger同时也依赖wpa\_suplicant，因为它需要wpa\_suplicant实现wifi加密认证；等等等等。

维护各种服务之间的依赖关系保证正确的启动次序是一个系统管理员非常非常重要的工作。
但是也不必过于担心，发行版通常已经将这种重要的工作做的相当好了。

对于Gentoo，/etc/init.d/　下的脚本并非bash脚本，而是openrc脚本。openrc主要的作用是“解析和缓存依赖关系”。
每个脚本都包含　depend() 函数，openrc执行该函数获得依赖关系以决定脚本的执行次序。

而对于Fedora/RHEL这样的系统（老Fedora，Fedora 从版本16开始不再使用sysvinit）， /etc/rc?.d/ （?为运行级别数字）下面有　S??XX K??XX 这样的符号链接到　/etc/init.d/XX 脚本。 %话说这也是Gentoo的启动脚本名为openrc的由来。
S表示　Start，?? 是一个两位数的数字，代表启动的次序。
S表示　Kill，?? 是一个两位数的数字，代表停止的次序（停止的次序和启动次序一样重要）。\footnote{
我表示我已经被Fedora这种模式弄晕了，Gentoo的openrc好歹是设置依赖关系自动决定启动次序的。这也是Fedora后来果断废除sysvinit的原因吧。}Fedora的 /etc/init.d/下是实实在在的shell脚本。没有依赖关系。引入Gentoo的openrc不如彻底革新到自动确定依靠关系或者是根本不需要依赖关系的新服务管理系统。

\subsubsection{systemd}
\sidenote{有关 systemd 的内容, 将在 \protect\chapref{chap:systemd} 详细讲解. }
sysvinit大量依靠脚本（/etc/init.d/* ）实现系统启动过程。脚本的一大特点就是慢。
另一个诟病的地方就是服务之间的关系复杂，很难搞懂。而且各个发行版并没有统一这些脚本。有的发行版为脚本加入了依赖关系，有的要依靠手工设置（是发行版打包的时候确定的）启动次序。管理上非常复杂。

那么，针对sysvinit的问题：慢、管理复杂，Redhat的一个牛人开发了systemd做为PID 1以克服sysvinit的缺陷。

那么systemd到底有哪些好处呢？



\begin{itemize}
\item 极度并行化启动。
		systemd开发者说过，要系统启动更快的，第一条需要做到的是单位时间内启动更多服务。不像sysvinit那样串行启动，systemd是高度并行的。systemd同时启动多个服务。

\item 无脚本启动。systemd开发者说过，要系统启动更快，你还需要做到的是整个系统启动过程启动更少的程序。systemd实现不依靠一行shell脚本启动完整的系统。

\item	按需启动。systemd开发者说过，要系统启动更快，按需启动是必不可少的。按需启动功能更进一步的使整个系统启动过程启动更少的程序。

\item	无需配置依赖性。
		按需启动还使得服务之间的依赖关系成了过去式。你不需要直接指定服务之间的依赖关系——只要一个服务用到了另一个服务，另一个服务会被按需启动，\em 就好像它本来就已经被启动了一样。

\item	支持未经修改在容器内启动。
		有关 systemd 容器的内容, 将在 \protect\chapref{chap:systemd} 详细讲解. 
\end{itemize}

systemd能实现按需启动和避开依赖性配置，原因就是实现了socket activation。

\begin{insertnote}
\subsection*{Socket Activation解释}
如果所有的服务可以同时启动而不考虑依赖性，服务就可以最大化的并行启动了。如果不考虑依赖性，管理员也可以大大的轻松。
那么服务和服务之间为何会相互依赖呢？答案是依赖它创建的Socket。{ \it Apache要在mysql之后启动，是因为网页需要数据库，需要mysql提供的数据库服务。它只所以需要这个服务，是因为它需要连接到mysql创建的socket。NetworkManager需要DBus服务，是因为NetworkManager需要连接到DBus，其实是连接到DBus创建的那个socket。}

如果，如果所有的Socket一开始就已经全部创建完成，所有的服务就可以同时启动了。因为他们连接到对方的时候不会再遇到“连接拒绝：地址不存在”错误。

如果没有Socket Activation, 同时启动NetworkManager和DBus的后果是，如果NetworkManager试图连接到DBus的时候，DBus还为创建它的socket,NetworkManager将出现连接错误而是服务启动失败。

有了Socket Activation, NetworkManager和DBus不再需要串行化启动，同时启动即可。NetworkManager总能连接上需要的Socket，因为它已经在那里了。
甚至DBus自身都不用启动，NetworkManager也能正常开启直到NetworkManager向这个socket发送了一些数据——这个时候DBus必须开起以便处理这个数据。这就是“安需启动”的实现方式。

Socket Activation 将创建socket的任务交给了systemd——利用了UNIX父进程可以向子进程遗传文件描述符\footnote{socket也是文件。也是使用的文件描述符标识的。}的特性——由systemd创建后再由其启动的服务继承。
\end{insertnote}

除了Socket Activation，还有类似的Path Activation，DBus Activation。 

DBus Activation好理解，有的服务提供的并不是Socket，而是在DBus下的一个名称接口。服务使用者并不连接到该服务，而是直接连接到DBus，然后向DBus发送请求，由DBus转发给对应的服务。这是一种RPC\footnote{远程过程调用。}实现的说。如果该服务不存在（比如还未启动），DBus会返回错误。DBus Activation将未启动的服务的接口提前标记为“可用”，等该服务真的被用到的时候再将对应的服务启动。

Path　Activation原理类似。比如邮件投递服务，只要在 /var/spool/mail 下有文件，该服务就应该进行邮件投递，将成功投递的邮件删除。Path Activation将监视该文件夹的任务交给了systemd，文件存在后systemd激活邮件投递服务。

那么，决定使用sysvinit还是systemd了么？话说Gentoo给的不就是选择么？
%如果决定使用sysvinit什么都不用做，stage3里已经带上sysvinit了。如果决定用systemd就继续哦，话说现在还不着急安装，内核代码下好了么？下好的话可以回到上面编译内核去了的说哦。


\section{分区规划}
任何系统都是安装到硬盘使用的
\footnote{好吧，先忽略掉LiveCD 和  WinPE 这类不需要安装的系统，咱讨论的是一般用途的桌面操作系统。}，
安装到硬盘之前，必须先划好家。在介绍如何分区前，首先得知道什么是分区，然后参考
FHS\footnote{FHS 是文件系统目录结构的一个标准。规定了根分区下各个子目录的名称和用途。}
指示结合自己的实际情况规划好分区。

\subsection{分区基础知识}
最初，计算机使用软盘，软盘是没有分区的。后来蓝色巨人IBM发明了硬盘。硬盘的容量一下子比软盘大出好多倍。当时的MS-DOS所使用的FAT12文件系统无法管理那么大的硬盘。于是西雅图的巨人和蓝色巨人想出来把硬盘划成逻辑上的几个区域，每个区域大小都在MS-DOS能管理的范围之内——这样的逻辑区域就是分区。人们发现将硬盘划分为逻辑上的几个区域后，更容易组织硬盘上的数据了；而且一个文件系统错误只会影响到一个分区的数据，其他分区不受影响，数据的安全性也得到的提升。因而后来的DOS虽然将FAT12进化到了FAT16，能管理当时的大容量硬盘了，但是分区这个功能却保留了下来\footnote{虽然是MS-DOS的发明，但是Linux可不会拒绝这样的发明。Linux还支持 BSD 发明的分区格式，总之，Linux决定支持越多的分区表格式越好，这极大的方便了用户，不是么？}。

既然用户划了分区，操作系统总得知道用户到底怎么划分的，描述分区的数据被称作分区表。既有分区表，必须有个地方存储，也必须知道到哪里去读取分区表。MS-DOS的把分区表和引导程序放入硬盘的第一个扇区\footnote{扇区是硬盘最小寻址单位。参考附录A. FAQ里的条目“什么是扇区”。}。

硬盘的第一个扇区又被称呼为MBR\footnote{MBR是硬盘的第一个扇区，具体解释请参考附录A. FAQ 里的条目“什么是 MBR”。}。
MBR既要存储引导程序，又要储存分区表，是个寸土寸金的地方，分区表大小受限，只有4个表项。也就是说，一个 MBR分区表最多只能有4个分区。要是只有四个，似乎并不够用。

MS-DOS将其中一个分区作为扩展分区，然后再扩展分区里再建分区表。扩展分区里面的分区就是逻辑分区，MBR 上的分区表就是主分区表。主分区表里划分的分区自然就是主分区了。所以一个MBR格式分区表最多允许3个主分区和1个扩展分区，或者4个主分区。扩展分区里再创建逻辑分区，没有数量限制。
\chatu{logicalpart}{链式逻辑分区表}

\begin{notice}
注意：扩展分区里的逻辑分区并不是表格形式存储的，而是“链式”存储。如图\thefigure{}所示。每个逻辑分区包含查找下一个逻辑分区的信息。因而一个逻辑分区的破坏有可能造成链式效应，将所有的逻辑分区全部摧毁。
\end{notice}

由于近来UEFI\footnote{参考附录A “UEFI 和 BIOS”。}  的兴起，
UEFI指定的分区表格式GPT也流行开来了。
和MBR不同的是UEFI不需要专门的引导扇区，引导程序由UEFI直接从文件系统上加载。所以GPT只是分区表，不需要和引导程序共存。

MBR 的分区表只有64个字节大小，只能包含最多4个分区的信息。
GPT包含128个分区表项，最多允许一个硬盘划分成128个分区，足够了。
能表示更多的分区并不是GPT唯一的优点，MBR分区表只能管理2.2TB
\footnote{MBR使用4个字节表示分区起始位置的偏移量。
偏移量以扇区为单位，一个扇区为512字节，那么MBR分区表能管理的最大硬盘大小为$2^{31}*512Byte = 2TB$。
又因为硬盘厂家以1000为进制而不是1024，故而大约为2.2TB。}%
以下大小的硬盘，GPT却可以管理容量超过9ZB\footnote{1ZB=1024PB 1PB=1024TB. 硬盘厂家的计算是 1ZB=1000PB 1PB=1000TB。}的硬盘。
在大容量硬盘越来越普遍的今天，MBR显得越来越力不从心。GPT正好接替MBR成为今后PC硬盘的主流分区表格式。

如果你电脑的固件是UEFI，不管硬盘实际是否大于2TB，建议最好使用GPT分区表。

在UNIX操作系统下，一个文件系统是要挂载到某个目录下才能访问其中的文件的。Windows和DOS下，操作系统会分配一个“盘符”，但在UNIX下，使用挂载命令将文件系统挂到一个目录下。被挂载的目录就成了那个挂载的文件系统的根目录，也称为挂载点。实现这个挂载操作的命令就是\textbf{mount}。mount将一个分区“挂载”到一个目录下，命令格式为

\begin{code}
mount 磁盘设备  挂载目录 [-t 文件系统] [选项]
\end{code}

选项是可选的，文件系统类型没有给出的话，mount会进行自动探测，所以一般也不需要给出。

\begin{notice}
注意：mount只有root用户才能运行。
\end{notice}

\begin{example}{挂载U盘}
\it 通常大家只有一个硬盘。按照第一个硬盘为/dev/sda，第二个硬盘为/dev/sdb这样的编号，通常U盘就是/dev/sdb（U盘也是硬盘的一种）。所以我这里假设用户插入的U盘被识别为/dev/sdb。U盘一般只有一个分区，所以正确的表示U盘的设备文件就是 /dev/sdb1了。

\normalfont

建立一个/media/usbdisk目录用于挂载，然后执行mount将U盘挂载到/media/usbdisk。

\begin{code}
mkdir -p /media/usbdisk

mount /dev/sdb1 /media/usbdisk
\end{code}

执行成功后，跑到（cd命令进去） /media/usbdisk目录下看看（ls命令列出），是不是U盘上的文件都跑到这个目录下了呢！

\end{example}


有mount操作自然也有撤销的操作。撤销mount的操作由命令unmount完成。试试看执行unmount /media/usbdisk，等等！命令未找到？

是的，UNIX就是这么矫情。看似应该使用unmount的地方，UNIX使用的却是umount命令。umount命令的语法不用说都猜到了：
\begin{code}
umount 磁盘设备
\end{code}
或者
\begin{code}
umount 挂载目录
\end{code}

用哪个形式都可以。要卸载一个文件系统，必须那个文件系统里所有的文件都被关闭。如果还有程序在使用上面的文件就无法卸载。可以使用 lsof 命令查看打开的文件。比如
\begin{code}
lsof /media/usbdisk
\end{code}
就可以查看/media/usbdisk目录下所有被打开的文件以及是哪个程序打开的。确保lsof的输出为空即可安全卸载了。

\begin{example}{卸载U盘}
U盘用完后直接拔出是不正确的\footnote{尽管这是个美好的传说。}。操作系统可能还有未来得及写入U盘的数据，直接拔出会导致数据丢失。所以在拔出U盘前必须先卸载U盘。
执行：
\begin{code}
umount /media/usbdisk
\end{code}

然后就可以拔出U盘了。如果卸载失败，可以用lsof查看是否还有正在使用的文件。

\end{example}

\subsection{为磁盘分区}

\chatu{gparted-gui}{gparted界面}

\normalfont

有很多种工具可以对硬盘进行分区。今天的主角是 gparted , 这个是一个非常简单易用的分区管理软件，带图形界面的哦~~

要打开gparted，在终端执行sudo gparted。如果 Fedora 的LiveCD提示没有这个命令，则用命令sudo yum install gparted或者apt-get install gparted进行安装。gparted界面如图\thefigure{}所示。

在gparted里所做的一切操作都不是立即执行的，只有应用操作后才被一次性执行。所以未应用前撤销操作是来得及的。
下面我把为硬盘分区所需要的几个主要操作图解一下。请读者根据此处的图示熟悉操作流程，并在学习下一个小节后自行分区。

\chatu{gparted-apply}{gparted应用操作}
\noindent{}要应用操作：\\\indent
1. 单击“编辑”|“应用全部操作”菜单。如图\thefigure{}所示。\\\indent
2. 在弹出的警告对话框中选择“是”。

\chatu{gparted-newpart-boot}{gparted新建分区}

\noindent{}建立一个分区：\\\indent
1. 选择磁盘的空白区域，右击后在菜单里选择“新建”。\\\indent
2. 在弹出的对话框里指定一个大小、分区类型和文件系统，如图\thefigure{}所示。\\\indent
\hskip 4em %
\begin{notice} %
  注意：分区表为GPT格式的只能选择主分区。
\end{notice} \\\indent
3. 单击“添加”按钮，新建的分区就编排在了任务列队里了\\\indent

\chatu{gparted-delpart}{gparted删除分区}
\noindent{}删除一个分区:\\\indent
1. 右击要删除的分区在菜单里选择“删除”。如图\thefigure{}所示。 %\\\indent

\chatu[scale=0.2]{gparted-newpt1.png}{新建分区表}

\noindent{}新建一个分区表：\\\indent
1. 单击“设备”｜“创建分区表”菜单，如图\thefigure{}所示。

\chatu[scale=0.2]{gparted-newpt2.png}{新建分区表}

\indent2.在打开的对话框里点“高级”前面的箭头，选择一个分区表类型。默认为MS-DOS分区表，也就是MBR分区表。
我需要在这里我选择gpt分区表。如图\thefigure{}所示。

好了，gparted的基本操作就完成了，接下来学习一下文件系统的结构后按照自己的意愿为自己的系统分区吧。

\FloatBarrier

\subsubsection*{UEFI模式注意}

值得注意的是，如果使用GPT分区表，必须有一个200M左右大小的FAT分区并设置为EFI系统分区，如图\ref{fig:EFI-part}所示，为该分区打上boot 标记即可（操作：右击分区选择“管理标志(A)”）。

EFI系统分区一般挂载到 /boot/efi下，但是个人喜欢EFI系统分区和/boot共用。自由选择的说。为了日后少麻烦，请务必让EFI系统分区为第一个分区。虽然说不一定非得是第一个分区，但是第一个分区是eifbootmgr的默认假定，嫌麻烦的研究手册的话，还是以第一个分区作为EFI系统分区为佳。

\chatu{EFI-part}{为分区设置EFI启动属性}

\FloatBarrier

\subsection{文件系统结构标准（FHS）}

Linux继承自UNIX树形目录。每个目录各司其职，并被FHS（Filesystem Hierarchy Standard）标准化。表3.1所列是一个不怎么完整的FHS标准，但是对于读者决定如何为他自己的电脑分区已经足够了。

\begin{longtable}{|l|p{0.7\textwidth}|}\hline
\caption{FHS标准参考}\\
\hline
/ & 根目录 \\\hline
/bin & 基本系统程序。如 ls cat grep \\\hline
/sbin & 只供管理员使用的基本系统程序 \\\hline
/boot & 内核和引导程序，通常为独立分区\\\hline
/dev & 设备文件目录 \\\hline
/etc & 存放系统配置文件。不允许存放为独立分区\\\hline
/home & 存放非root用户的家目录。 
比如foo用户的家目录就是/home/foo
一般用来存在个人文件和个人设置。
通常强烈建议独立为一个分区。\\\hline
/lib & 系统基本库，被 /bin /sbin 里的程序依赖的库 \\\hline
/media & 可移动媒体的挂载目录。比如cdrom和u盘。被/run/media代替\\\hline
/mnt & 临时挂载目录\\\hline
/opt & 通常用于安装非开源软件，如Adobe Reader \\\hline
/proc & 虚拟的文件系统，用来获得内核和运行中的程序的信息\\\hline
/root & root用户的家目录\\\hline
/srv & 做服务器使用的时候 ，用来存储服务的数据。如一个git服务器通常将仓库存放于/srv/git\\\hline
/tmp & 临时目录。其中文件关机丢失。建议挂载为tmpfs虚拟文件系统，对于tmpfs的介绍可以参考第7章。 \\\hline
/usr & 非基本系统程序根目录，下面的结构和 / 差不多 \\\hline
/usr/bin & 同 /bin。只是非基本程序。如xeye  elinks . \\\hline
/usr/sbin & 同 /sbin。只是非基本程序。如system-config-firewall \\\hline
/usr/include & 头文件目录。\\\hline
/usr/lib & 同 /lib。只是其包含的库是不被/sbin、/bin使用的。\\\hline
/usr/share &  保存架构无关的共享数据。如图标\\\hline
/var & 系统运行中会不停变化的文件。比如各种log，包管理器的数据库，等等。通常在服务器上是独立分区，个人电脑可以不独立划分。\\\hline
/var/cache & 缓存。该目录下的文件可以安全删除，要求使用它的程序必须能重建它。比如fontconfig的缓存。\\\hline
/run 和 /var/run & 临时目录。包含本次系统运行时的信息。\\\hline
/var/log & 存放各种日志文件。\\\hline
/var/tmp & 临时文件，可能重启后还在。注意和/tmp的区别。\\\hline
\end{longtable}


知道了系统各个目录的作用，就好为目录进行分区了。通常我们使用4个分区 /boot、/、/home 和一个交换分区。

\begin{notice}
注意：我在使用“/home分区”这个术语的时候，我指的是为“/home”分一个单独的分区，并且将这个分区挂载到“/home”目录下。我通常不使用“挂载到/home目录的分区”而直接简述为“/home分区”。同样的，“/分区”（或者根分区）指的就是挂载到根目录的分区。
\end{notice}

/boot独立为一个分区的好处是防止根文件系统的错误蔓延导致/boot/*文件破坏无法正常启动。另一个好处是 /分区可以使用引导程序（如GRUB）不支持的文件系统格式。例如GRUB不支持XFS文件系统格式，但是你可以让/boot分区使用ext2格式从而使 /分区可以使用XFS格式。如果/boot没有独立出一个分区。那 /分区就只能使用引导程序所支持的文件系统了。

/home独立一个分区的好处就是将用户数据彻底和系统数据分离。由于/home目录通常频繁读取写入，和 /分区隔离可以避免 /分区被碎片化。

提示：其他目录，比如“/var”，可以独立一个分区也可以不独立。就看你系统的用处了。作为桌面用途，完全没必要。如果是服务器，“/var/log”可能会很大。邮件服务器的“/var/mail”就可能会非常大，所以非常有必要使用独立分区。

\subsection{文件系统选择}


Linux支持的文件系统非常多，有自家系列 ext2/ext3/ext4、甲骨文公司开发的btrfs，来自硅谷图像的XFS等等等等。

不过选一个文件系统没那么纠结的。对于根分区\footnote{根目录所在分区}，既然大部分发行版默认使用ext4，那就ext4吧。

为了引导程序的兼容性，/boot分区仍旧使用ext2或者ext3。

对于/tmp这样的目录，强烈建议使用tmpfs\footnote{参考第7章“文件系统”。}。tmpfs是一种内存文件系统，所有tmpfs上创建的文件都是在内存里的，关机后就会消失。非常适合用作/tmp的文件系统。既能加快系统速度，又能保护硬盘。

对于/home目录，如果是个BT爱好者，会有很多大文件，我推荐使用XFS文件系统。其他情况下还是ext4各方面性能比较平衡。

如果希望为自己选择最佳的分区格式，请参考第7章“文件系统”的内容。

\section{开始安装}

废话了那么多，介绍了一些安装前需要知道的基础知识，现在终于要进入正题了，安装Gentoo。
Gentoo的安装非常简单，首先是分区。这个前面已经介绍过了。我现在假定读者已经在LiveCD 环境中完成分区了。下面开始真正的安装。


\subsection{正式开始安装过程}

现在，Gentoo的安装可以正式开始了！现在再检查一下

\begin{itemize}
\item	使用LiveCD启动，处于LiveCD环境中。
\item	LiveCD里已经配置好网络，能访问Internet。
\item	完成分区规划。已经为将要安装的Gentoo规划并划分好分区。
\item	现在处于LiveCD的命令提示符下。
\end{itemize}

一切准备就绪，可以开始安装了。

首先要做的事情就是创建一个　gentoo　目录。这个将要用来安装Gentoo的分区就暂时挂载到这个目录下。如不特殊说明，我使用 \#　代表命令提示符， \#后的命令表示需要手工输入并敲回车执行。

\begin{code}
\#mkdir　/gentoo\\
\#mount /dev/sdXX /gentoo
\end{code}

这里　/dev/sdXX　是读者为Gentoo划分的根目录。

系统目录准备好了，接下来下载Gentoo的stage3。打开　\url{http://mirrors.163.com/gentoo/releases/amd64/current-stage3/} 或者如果打算使用x86就打开\url{http://mirrors.163.com/gentoo/releases/x86/current-stage3/}下载以最近的日期结尾的 stage3 压缩包。

然后解压到 /gentoo
\begin{code}
\#tar stage3-*.tar.bz2 -C /gentoo 
\end{code}

接着下载　portage 快照
\begin{code}
\#wget http://mirrors.163.com/gentoo/snapshots/portage-latest.tar.bz2\\
\#tar portage-latest.tar.bz2 -C /gentoo/usr/portage
\end{code}

解压完成后，复制dns配置　
\begin{code}
\#cp /etc/resolv.conf /gentoo/etc/resolv.conf
\end{code}

如果还有划分了其他目录，也要一并挂载进来

挂载home分区，boot分区等（如果有划分独立分区的话）
\begin{code}
\#mount /dev/sdXX /gentoo/boot\\
\#mount /dev/sdXX /gentoo/home
\end{code}

各使用读者划分好的分区代替命令中的　/dev/sdXX。

下面是挂载/dev/ /proc /sys的。这里使用宿主系统的/dev /proc 和 /sys　目录。
\begin{code}
\#mount --rbind /dev/ /gentoo/dev\\
\#mount --rbind /sys /gentoo/sys\\
\#mount --rbind /proc /gentoo/proc
\end{code}

好了，最好这些操作，Gentoo最小环境（能用但是不能启动）的基本环境就搭建完毕了。
执行

\begin{code}
\#chroot /gentoo　/bin/bash
\end{code}
或者如果你喜欢简洁的操作
\begin{code}
\#chroot /gentoo
\end{code}

\begin{insertnote}
\subsection*{chroot命令解释}

chroot 命令能将某个目录变成根目录并在新的根目录里执行shell或者命令行第二个参数指定的程序。

执行 chroot /gentoo 后，chroot 将自己的根目录（不影响系统其他进程的根目录）切换到新的目录，然后执行新的根目录里的 \$SHELL（通常是/bin/bash）。
这样获得的新shell就继承了chroot的根目录设定，也就是 /gentoo。 试试在里面执行　mkdir /test。然后开个新的终端检查看看　/gentoo/test 目录是不是在了。
\end{insertnote}


好了，现在就已经进入Gentoo的环境了。正如我在本章开头提出的一个引用那样，“Where there is a shell, there is a way.” 安装Gentoo的第一步事实上就是获得一个shell，不管这个shell是用LiveCD获得的，还是使用一个已经安装的系统。只要有了shell，就有办法安装Gentoo。

\subsection{emerge用法}\label{sec:quickemergeguide}

说到Gentoo不得不提portage和emerge。portage是Gentoo的软件包管理器，emerge是portage的命令行工具。能使用portage的不止emerge一个，这个是有官方背景的工具，所谓正统是也。

emerge的用法非常简单，假设你要安装pidgin，执行下面的命令就可以了:
\begin{code}
emerge pidgin
\end{code}

很简单，是不是？emerge会自动的下载pidgin然后编译安装，全程不需要人工干预。如果pidgin依赖的一些包没有安装，emerge会先下载这些包然后安装。再一次，全程自动化。不需要手工干预。

如果认为emerge只是自动下载编译脚本你就大错特错了。执行
\begin{code}
\# -p参数表示不实际执行安装操作，只显示需要执行的操作

\#-v表示输出详细信息。

emerge -pv pidgin
\end{code}

看到什么了？

\begin{code}
[ebuild   N    ] net-im/pidgin-2.10.6 USE="dbus gnutls gstreamer gtk ncurses networkmanager nls xscreensaver (-aqua) -debug -doc -eds -gadu -groupwise -idn -meanwhile -perl -prediction -python -sasl -silc -spell -tcl -tk -zephyr -zeroconf"
\end{code}

注意看输出，ebuild 后面有个 N 代表新安装，如果是 R 表示重新安装，U表示升级。D表示降级（不常发生，但是出现新版本有重大错误，Gentoo会紧急撤销新版本，这样已经按照的人会降级）
之后为包名，带版本号。有时候会出现NS，表示安装新版本和老版本共存。然后就是USE=$\cdots$ 了。USE=表示该包所有可选的USE参数。-xx 的表示关闭对应的功能， xx表示打开对应的功能。这些功能都是编译期选项。编译的时候就确定下来的。没有被编译进的功能软件运行时就好像这个功能本来就不存在一样。从源头上禁用了。
有的功能会在外围打上括号，表示该功能对于本平台是强制启用或者禁用的。

那怎么调节USE参数呢？ 3个办法。

\begin{itemize}
\item 使用USE环境变量。比如{ \tt  USE="-networkmanager -dbus sasl" emerge -v pidgin} 。 这个设置只影响本次编译的pidgin。
\item 到 make.conf 设置 USE=。 比如在 /etc/portage/make.conf 文件中加入 {\tt USE="gtk networkmanager -dbus"}（USE变量要用引号。）。这里的设置会是全局设置，影响到所有后续编译的软件。
\item 到 /etc/portage/package.use 文件中写一行“ { \tt  net-im/pidgin dbus gtk -ncurses -nls } ”（不包括引号）。 这个设置只影响pidgin，但是所有后续安装pidgin都会使用这个设置。
\end{itemize}

\subsection{配置make.conf}

话说emerge虽然很智能，编译下载全自动，但不代表就不需要配置了。portage的主配置文件就是 /etc/portage/make.conf。

make.conf 里最重要的设置莫过于CFLAGS和CXXFLAGS了。 
CFLAGS 和 CXXFLAGS 设置 C 和 C++ 的编译参数。
emerge所编译的所有程序都会使用这里的编译参数设置。

要配置make.conf需要编辑器，请参考\secref{sec:nano}和\secref{sec:editors}
的内容。

CFLAGS是C编译器参数。在这里设置的参数会被所有软件作为C编译器参数。CXXFLAGS是c++编译器参数。通常两者设置为相同即可。

推荐的设置是 CFLAGS="-O2 -march=native -pipe"。
（CXXFLAGS 同 CFLAGS。）

\begin{itemize}
\item -pipe 表示使用管道代替临时文件把编译后的汇编文件交给汇编器。
\item -O2 表示优化级别为2。gcc 有5个优化级别， -O0 -O1 -Os -O2 -O3。-O0 表示不优化。-O1 到 -O3，数字越大，优化级别越高，生成的代码也越快。之所以不推荐使用 -O3 乃是保守。-O3 优化太激进，可能会导致一些软件错误。-Os是一个特殊的优化级别，表示按照大小优化而不是速度。
\item -march=native 指示编译器按照本地CPU的指令集优化。编译生成的二进制程序将不能在和本机CPU型号（比仅仅是CPU牌子比一致，代数比一致也不行。比如在i7上编译的到奔腾上将无法执行。）不一致的机器上使用。这不正是使用Gentoo的一个原因么？
\end{itemize}


GENTOO\_MIRRORS 指示emerge下载源代码的时候使用的镜像。ebuild里虽然给出了软件的原始下载链接，但通常只作为后备资源。emerge会优先从Gentoo镜像服务器下载同名文件。
Gentoo全球镜像站点非常多，所以需要设置离自己最近的镜像以加快下载速度。目前速度不错的两个国内镜像站点分别是搜狐和网易。\url{http://mirrors.sohu.com/gentoo/} \url{http://mirrors.163.com/gentoo/}。可以到 \url{http://www.gentoo.org/main/en/mirrors2.xml} 查看全部的镜像站点。

SYNC %="rsync://mirrors.163.com/gentoo-portage"
设置的是 portage 目录的更新源。emerge 依赖 /usr/portage 目录下的 ebuild 文件获得编译和依赖信息。所以及时的更新这个目录是很重要的。使用 emerge --sync 
即可更新此目录。同样的，更新这样的目录也要选择一个快点的镜像。到Gentoo的网站查询镜像列表后选择一个填入。\footnote{网易提供的 \url{rsync://mirrors.163.com/gentoo-portage} 在大陆有不错的速度。}~ 由于 portage 树是用 rsync 进行增量更新的，所以将算不设置一个镜像也可以。因为如果每天更新的话，rsync的传输量并不大。所以可以不设置，使用默认的。

ACCEPT\_KEYWORDS%="~amd64 amd64"
指示emerge接受指定关键字的ebuild。关键字为 ARCH 和 \textasciitilde{}ARCH , ARCH为处理器架构名，如x86、amd64。设置了{\tt ACCEPT\_KEYWORDS="amd64"} （当然，这在amd64平台是默认设置）后，emerge只接受标记了amd64关键字的ebuild。
这样就把没有标记此关键字的ebuild过滤掉了。这样的做法是因为有的软件只有某个平台可以用。而 \textasciitilde{}ARCH 则表示“不稳定的包”。
比如一个ebuild如果标记为 \textasciitilde{}amd64 就表示在 amd64平台可用，但是不稳定。如果标记为 \textasciitilde{}x86就表示在32位x86平台可用，但不稳定。
其实这个稳定是从服务器的角度说的。如果做桌面使用，\textasciitilde{}ARCH都是太稳定了的说。
通常桌面使用会设置{\tt ACCEPT\_KEYWORDS="\textasciitilde{}amd64 amd64"} ， 如果是32位的则使用ACCEPT\_KEYWORDS="\textasciitilde{}x86 x86"。 这样标记为不稳定的包也能使用。通常一个软件发布新版本后，portage中出现的都是 \textasciitilde{}ARCH 版的ebuild，被测试充分稳定后才会进入 ARCH 。
使用 \textasciitilde{}ARCH 是推荐设置。想尝新的朋友其实Gentoo有更不稳定的包，就是 masked 包。
\url{/usr/portage/profiles/package.mask} 文件描述了portage树中所有被 masked 的包，被 masked 的包不论 KEYWORDS设置成什么，都会被过滤无法安装。除非手工解除屏蔽状态，具体的办法参考 \secref{sec:emerge}

MAKEOPTS%="-j8"
通常的设置为 -jN，其中 N 为cpu的数目。\begin{notice}
注意，超线程被认为是2个CPU。比如双核超线程的i3，逻辑CPU数目就是4个。
\end{notice} 这个设置被 make 使用，表示同时开启的编译数目。比如 make -j8，则 make 会尽可能的同时开启8个编译器进行编译。


\subsection*{处理循环依赖}
软件之间相互依赖并不是很罕见的事情。既然要从头进行编译安装，过程中必然会出现“A依赖B，B依赖A”这样的循环依赖。
有时候循环依赖还是间接的：“A依赖B，B依赖C，C依赖A”。可以把C想象成多个软件。
最典型的循环依赖就是工具链了：要编译任何软件都需要GCC和GLIBC，但是编译GLIBC自己也需要GCC，编译GCC自己也需要GCC和GLIBC。正因为这样会发生循环依赖，
所以完全从头编译一个Linux环境是非常麻烦的。如果非要这么做，参考LFS。

我们下载的是Gentoo提供的stage3。这里必要提一下何为stage3，有stage3必然有stage2、stage1咯。所谓stage1就是只包含一个bootstrap.sh脚本和配置文件的压缩包。
如果从stage1开始安装，就是执行bootstrap脚本。这个脚本的做法就是和LFS一样，从0开始构筑stage2。stage2则意味着包含了一个最小的“工具链”环境。
这个环境有编译器，链接器，bash，必要的命令行工具。也就是stage2就是个能编译自己的“最小”环境。从stage2开始继续bootstrap就能获得stage3。stage3包含了一个能编译自己的最小“Gentoo”环境——包含了emerge程序。因为emerge是python编写的，故而不能放入stage2。

也就是说，stage1/2/3 就是处理循环依赖，每次将基础扩大。到stage3开始，基本就可以避免循环依赖了。但是并没有绝对避免。是否会遇到循环依赖取决于用户的配置和需要安装的软件。
这也就是Gentoo已经很久不再提倡从stage1开始安装Gentoo的原因：解决循环依赖是噩梦。

stage3开始，已经没有硬循环依赖了。这个时候的循环依赖通常是这样的“A依赖B，但是B有个可选的选项导致会依赖A”。这样只要将那个可选的选项关闭，编译好B，然后再编译A，如果那个功能用户需要的话，重新打开后再次编译B。

\begin{example}{Dbus和systemd相互依赖}

直接使用 emerge systemd 的时候，系统会提示dbus开了[systemd]选项后依赖systemd，而systemd也依赖dbus。同时systemd编译的时候并不依赖开启systemd选项的dbus，但是运行的时候却依赖。

直接的后果就是emerge出现错误提示后退出。emerge的关键提示信息为 systemd depends on by dbus, dbus depends on systemd, systemd depends on dbus.
随后提示Note that circular dependencies can often be avoided by temporarily disabling USE flags that trigger optional dependencies.

较高版本的portage还会提示关闭可能避开循环依赖需要关闭的USE选项。而低版本做法则为进一布提示：systemd依赖dbus, dbus[systemd]依赖systemd。注意dbus后面加了个[systemd]，说明USE包含systemd后dbus才依赖systemd。

\end{example}


接下来要做的就是编译一个可启动的内核，编译一些引导必备软件，安装一个引导程序Gentoo就能启动了。


\subsection{内核编译}
编译内核，对初学者来说是个艰巨的任务。不过这是通往资深用户和熟练管理员的必经之路。该做的迟早要做，不如现在就做个了断吧。

首先获取一个内核代码。推荐使用git获得内核的代码而不是下载tar压缩包。因为更新到下一个版本的时候git只会下载差异部分而不会重新下载。
% git的具体用法请参考　\secref{sec:git}。

必须把内核代码放到　/usr/src/linux　哦，这样依赖内核头文件的软件才能找到所需的内核代码（仅限Gentoo设定）。

{
\bf\noindent{} 提示：
}{
\em shell　下一行太长可以在后面输入``$\backslash$''并在下一行输入剩余的命令。shell会在下一行自动出现提示符 $>$　。
}


\begin{code}
\#cd /usr/src\\
\#git clone \textbackslash \\
>~ git://git.kernel.org/pub/scm/linux/kernel/git/torvalds/linux.git
\end{code}

内核代码非常大，这一步操作会比较费时。可以先放一边，重新打开个新的终端然后chroot /gentoo，进行下一节的操作。等内核的代码克隆完成了再回到这里继续。

内核克隆完成后，接下来的任务就是编译内核。 编译内核是安装 Gentoo 的过程中最难的步骤。新人只要克服了编译内核这个大关，就再也没有什么能难倒他了。

进入内核目录 \texttt{/usr/src/linux} 并执行 \texttt{make menuconfig} 会看到如下的界面。

\chatu{kernelconfig}{内核配置界面} 

在这个界面，可以使用键盘的  \textuparrow\textdownarrow 来移动选项，用 \textleftarrow\textrightarrow 来移动下方的 \verb|<Select>    < Exit >    < Help >    < Save >    < Load >| 按钮。移动到对应的选项上，使用空格键切换选项状态。Linux内核的选项，基本上分 Y M N 三种回答。在 menuconfig 配置界面，就是 [*] [M] [ ] 三种。空格键可以在这三个状态间切换。如果有子选项，按回车键进入。退出子选项，要将下面的选项移动到 <Exit> 然后按回车。
所谓 Y 的意思，就是编译进内核 vmlinuz 文件本身。M 的意思是编译为模块，动态加载。N 就是取消对应的选项。有的选项只有 Y N 两种状态。 另外还有一种选项，是按回车后可以进入一个输入文本的界面，输入完成再按回车返回。如机器名的配置。
光标位于一个选项上的时候，按 Shit + "/" 也就是按 "?" 键将显示对应选项的英文帮助。

一开始，初学者可能会对内核那么多选项摸不着头脑。内核里的选项按照子系统进行过分类。
硬是要再分类，我认为就2类，一是功能性选项，二是驱动。

\begin{notice}
Linux 和 Windows 在驱动这一点上显著的不同： Linux 认为驱动是内核的一部分，而 Windows 认为驱动是硬件的一部分。因此 Linux 下驱动随内核附带。 Windows 下驱动随硬件附带。买新硬件一定要安装驱动。驱动由硬件厂家提供。
\end{notice}

自行编译内核，就需要知道自己电脑的具体硬件，以便选择驱动。如果漏选了驱动，轻则硬件功能缺失（不能上网，没有声音），重则系统无法启动。好在Linux支持按需加载驱动。就算选择了某个驱动，如果内核启动的时候没有检测到相应的硬件，对应的驱动便不会加载。
因此，对于驱动，有一个诀窍，不知道的时候，就选 M。 但是对于功能性选项， 如果选择错误，可能导致一些系统功能不可用，轻则导致一些程序无法执行，重则无法启动系统。好在这些功能选项，内核都有默认值。第一次执行 make menuconfig 的时候，这些选项已经按照默认值设定好了，基本无需干预。

下面以 X86\_64 架构为例详解内核的编译。


第一条选项，“64bit-kernel”，允许在 64 位环境下交叉编译 32 位内核。Gentoo 要本机使用，自然这里选 Y 。

下图 \ref{fig:kernelconfig1} 是 General Setup 。

\chatu{kernelconfig1}{内核配置界面-General Setup} 

\begin{enumerate}
\item Cross-compiler tool prefix  交叉编译的时候，设定工具链前缀。本地编译留空，不设置。
\item Compile aslo drivers which will not load 编译用不到的驱动。只有测试和开发人员需要。一般选 N。
\item Local version - apend to kernel release 在内核的版本号后面添加自己的版本字符串。一般不写了吧。看个人喜好了。
\item Kernel compress mode (XZ) 这个是内核的压缩格式。linux 内核编译完成后是一个 vmlinux 文件， ELF 格式。接着对 vmlinux 进行压缩，然后在头部添加启动和解压代码，形成 bzImage 文件。bzImage 文件通常拷贝到 /boot/ 分区，名字修改为 vmlinuz-版本号。 这里选择的就是内核的压缩格式。压缩率越高，生成的 bzImage 文件越小。载入就越快。 括号里显示的是当前选择的压缩模式。默认的压缩格式是 gzip 。 注意，系统里需要安装有相应的压缩工具。

\item Default hostname 就是不使用 hostname 命令设定本机的机器名的时候，默认的名字。默认是 "(none)"。 其实开机后， init 系统会根据配置设置名字的，所以这里的设置其实无关紧要。
\item Support for paging of anonymous memory (swap) 这个选项的意义在与支持页面文件和交换分区。内存不够用的时候就把硬盘上的一部分空间拿来当内存。哪怕内存很大，最好也选上。
\item System V IPC UNIX 系统的 IPC 功能，禁用该选项会导致一些程序无法运行。
\item POSIX Message Queues ~  POSIX 消息列队, 也是 UNIX 系统的功能，不建议关闭。
\item open by fhandle syscalls ~ Linux 新增的一个系统调用。已经有一些新的软件开始采用这项功能了。
\item Auditing support ～ systemd 需要启用 auditing 支持，并且 ``Enable system-call auditing support'' 和``Make audit loginuid immutable'' 同样需要启用。
\item Timers subsystem ~ 定时器子系统。定时器有三种工作方式。在这里建议大家选择 Idle dynticks system (tickless idle)。 如果系统空虚下来，就会关闭时钟中断。减低系统功耗。 
就算硬件没有高精度时钟（不太可能），也没有坏的影响。而 完全 tickless 模式，是 3.12 版本以后的内核才实现的功能，目前来说还不太稳定。慎重选择。
注意把 High Resolution Timer Support 也启用，它可以提供比1毫秒更高的精度(因为它可以读取HPET/TSC等新型硬件时钟源),可以更好的支持音视频等对时间精度要求较高的应用，强烈建议选"Y".

\item CPU/Task time and stats accounting ~~ 子选项 Cputime accounting 可以在  “Simple tick based cputime accounting”    、   (X) ”Full dynticks CPU time accounting“      、 ”Fine granularity task level IRQ time accounting “ 三个之间选择。精度依次提高，当然，对 cpu 的占用也依次提高。对进程的 cpu 时间统计有准确要求的人，可以适当的使用第二个乃至第三个选项。完全没要求的话，用第一个选项也不错。
BSD Process Accounting 这个功能目前来说，只有 bootchart \footnote{一个用来统计开机过程中的各程序资源使用情况的软件，用于分析开机过程，进行有针对的优化。} 这个程序用到了，如果不打算使用这个程序，完全没有启用的必要。“Export task/process statistics through netlink” 同样也是用于类似目的。平时使用可以无需开启。


\item RCU Subsystem ~~ RCU 是一个高效的读写锁，特别适用于读多写少的情况。如果不确定，这里面的选项全部都选 N。

\item Kernel .config support  ~~ 这个选项启用后，允许用户通过 /proc/config.gz 访问当前运行的内核的 .config 配置文件。实际意义不大。浪费内存而已。

\item Kernel log buffer size ~~  当你使用命令 dmesg 的时候，就能查看到内核启动开始到目前为止的日志输出。但是，如果你的系统运行了很长时间，内核产生了很多输出， dmesg 的输出会无限长么？ 当然是否定的。 这个就是用来设定用于存储内核日志的空间限制的。超过空间限制的日志，将会覆盖掉老日志。默认值是 17 , 也就是 128KB。

\item Control Group support ~~ cgroup 是 Linux 特有的一项进程管理功能。 systemd 就采用了 cgroup 进行进管理，能正对一组进程进行资源管理配置。cgroup 启用后，还能再进入其子选项。进一步定制 cgroup 的功能。如果打算使用 systemd (并且是强烈的建议使用 systemd )作为 init 系统， cgroup 是必须启用的功能。而 cgroup 子选项的功能，当然是选的越多，功能越丰富啦。不过，一般桌面系统，其子选项可以全部不启用。无需那么强大的功能。

\item Namespaces support ~~ 名字空间支持。 这个也是 systemd 要求必须启用的功能。 可能你会看到该选项变成了 -*- 

\item Require conversions between uid/gids and their internal representation ~~ 不确定就选 N 。 这个是内核里更严格的检查 UID/GIDS 转换关系。

\item Automatic process group scheduling ~~ 如果你使用了 systemd , 那么这个是鸡肋功能。这个是在内核里实现了 systemd 的 cgroup 分组管理。允许内核根据进程所连接的控制台自动纳入不同的 cgroup 分组。加强系统稳定性。

\item Enable deprecated sysfs features to support old userspace tools ~~ 用 Gentoo 怎么还会有老旧的程序呢？ 当然选 N 节约系统资源。这个是让  /sys/class/block 的内容出现在 /sys/block 的选项。 只有当使用一些老旧的系统程序的时候才需要。 Gentoo 下无需启用该选项。一些新的程序可能会在启用该选项后无法工作。

\item Kernel->user space relay support (formerly relayfs)  ~~ 选 N 。 


\item  Initial RAM filesystem and RAM disk (initramfs/initrd) support ~~ 如果打算使用 initramfs 就启用，不打算使用 initramfs 可以不启用。有关 initramfs 请参考 \faqref{FAQ:initramfs} 启用该功能也可以不使用 initramfs , 该选项会增加内核大约15k的体积。启用该选项会有个 “Initramfs source file(s)” 选项，这个无需搭理。除非你打算把 initramfs 编译进内核文件本身，而不是使用 grub 的 initrd 命令加载。这个一般是用于引导程序不支持 initramfs 的情况下启动系统。

\item Optimize for size  ～～  选 Y ， 原因参考   \secref{sec:optsize}

\item Configure standard kernel features (expert users) ~~ 如选项后的 “expert users” 指示一样，除非你是专家，否则不要选 Y 。选了 Y 将展开子选项， 里面可以配置一些内核标准功能。调节后，你的内核就不标准了，不拥有一些其他系统上的标准功能，还能继续兼容市面上的程序么？

\item Embedded system ~~ 选 N 。 除非你是做嵌入式开发的，对标准内核功能都要定制。

\item Disable heap randomization ~~ 为了增加安全性， Linux 对堆的分配是要采取一定的随机化处理的。如果想提高内存分配的速度，可以牺牲一定的安全性，开启次选项。

\item Choose SLAB allocator ~~ 使用最新的  SLUB (Unqueued Allocator)， 性能更好。

\item SLUB per cpu partial cache ~~ 让 SLUB 分配出来的对象对 cpu 更本地化，以便利用 cpu 的高速缓存。缺点就是释放的时候不能立即释放。对实时系统选 N 。

\item Profiling support ~~ 收集内核性能数据的系统，非内核开发者选 N 。

\item Kprobes ~~ 非内核开发者选 N 。 用于追踪内核函数调用。

\item Optimize very unlikely/likely branches ~~ 对几乎不可能执行的条件判定指令的一个分支进行优化，使得条件成立的时候开销更小。
\end{enumerate}

图 \ref{fig:kernelconfig2} 是 Processor type and features 。

\chatu{kernelconfig2}{内核配置界面-Processor type and features} 


“Processor type and features” 的子选项用来配置和 CPU 相关的一些特性。

\begin{enumerate}

\item 
Symmetric multi-processing support ~~ 启用  SMP 支持。 关于 SMP 可以参考 \faqref{FAQ:SMP}。 对多核处理器，这个是必须开启的选项。否则 Linux 将只支持单CPU单核的硬件。一般 PC 上的多核处理器都是 SMP 结构的。因此这里需要选 Y 。

\item Enable MPS table ~~ 老式 SMP 使用  MPS table ，而 PC 上使用的是 ACPI ，因此这里选 N 。 选 Y 也没关系，系统可以正常启动，只是会增加一些无用的代码，增加内存占用而已。

\item Support for extended (non-PC) x86 platforms ~~ 使用 x86 处理器的并不只是 PC ， 这个选项用于支持非 PC 系统。 如果你运行 Gentoo 的是 PC , 那么就安心的取消这个选项吧，还节约资源。

\item Intel Low Power Subsystem Support ~~ 在启用 Intel Low Power Subsystem 支持。好吧，你要是用的硬件是 
Intel Lynxpoint PCH 就可以启用。

\item Single-depth WCHAN output ~~ 简化 /proc/<PID>/wchan 让 CPU 使用更低。 建议使用  Y 这里。

\item Linux guest support ~~ 选 N , 除非你打算在一些特定类型的虚拟机里运行 Linux。

\item Memtest ~~ 选 N 。 除非你想在内核启动的时候测试内存。

\item Processor family ~~ 下面的子选项有很多，建议的选项是 Generic-x86-64， 除非你很清楚的知道自己的  CPU 类型。 PS: i7-4770 可以选 Intel Core AVX2。 

\item IBM Calgary IOMMU support ~~ 不清楚的就选 Y。

\item Maximum number of CPUs ~~ 最大支持的 CPU 数目。 这个数字选小点，可以节约一些内存。比如在 4 核超线程的 i7 配置下， 输入 8 。

\item SMT (Hyperthreading) scheduler support ~~ 让内核的调度器针对超线程优化。比如说，把 4 个线程分配到 0 2 4 6 这 4 个 CPU 上，而不是分配到 0 1 2 3 这4个 CPU上。因为 0 1 这2个 CPU 实乃一个 CPU 核心超线程虚拟出来的而已。如果你的 CPU 支持超线程技术，务必把这个选项启用。

\item Multi-core scheduler support ~~ 多核 CPU 选 yes。

\item Preemption Model ~~ 抢占模式。 Linux内核同所有的其他多任务操作系统一样，都使用“抢占式多任务”。有关抢占式多任务，参考 \faqref{FAQ:preemptmultitasking}。 但是，只有当该任务执行的是用户代码的时候，才是可以任意抢占的。如果该任务正在执行一个系统调用，导致进入内核模式，则只能在内核里预先定义的几个点上进行抢占。Server 模式就是在内核模式下不抢占， Desktop 模式则是在内核里预先定义的几个点上进行抢占。而 Low-Latency Desktop 模式则是只在几个 “绝对禁止抢占” 的为止禁止抢占，其他情况下可以任意抢占。
通常选 Desktop 模式。 如果你对响应有要求，就选 Low-Latency Desktop 模式。而当你想最大化发挥系统的吞吐量的时候，就选择 Server 模式。

\item Reroute for broken boot IRQs ~~ 一些主板需要启用该选项，否则内核会变得不稳定甚至无法启动。如果不确定，那就选 Y。

\item Machine Check / overheating reporting  ~~ 启用该选项允许 CPU 向内核报告错误，包括数据出错和过热。应对措施可能是打印一个警告，也可能是直接停机。建议开启，并依据自己的 CPU 类型启用子选项，  Intel  的还是 AMD 的 MCE。

\item Dell laptop support ~~ 是 Dell 笔记本的选 Y , 其他的 N 。

\item CPU microcode loading support ~~ 微代码热更新支持。一般来说没有太大的启用必要。 如果需要这个功能，还需要安装 sys-apps/microcode-ctl 和 sys-apps/microcode-data。

\item /dev/cpu/*/msr - Model-specific register support ~~ MSR 寄存器是 CPU 的特殊寄存器。启用这个选项允许 内核把寄存器透过 /dev/cpu/*/msr 文件暴露出来。允许用户空间程序设置 MSR 寄存器。 选 N。
\item /dev/cpu/*/cpuid - CPU information support ~~ 理由同上，允许用户空间程序透过 /dev/cpu/*/cpuid 访问 CPUID 信息。选 N。

\item Numa Memory Allocation and Scheduler Support ~~ 在一些多 CPU 系统里，每个 CPU 访问连接到自己节点的速度快，访问其他 CPU节点的内存速度慢。与是分配内存尽量应该分配在本地节点的内存里。 这个就是 Numa 内存分配器。PC 嘛，没有启用的意义。

\item Sparse Memory virtual memmap ~~ 加速虚拟内存映射的操作，选 Y。

\item Allow for memory hot-add ~~ 允许开机情况下添加内存。 这 ... 还是不要了吧，PC 在硬件上不支持啊。
\item Enable KSM for page merging ~~ KSM 是一个将相同内容的页面进行合并的措施。这个技术在普通桌面应用上并没有实际意义。 但是对于跑了大量的虚拟机的虚拟服务器提供商来说，他能将多个虚拟机的内存进行合并从而大大减小虚拟机的内存使用。虚拟机通常都跑的相同的操作系统。他们的内存，虽然有大量的内容是重复的，但是都是独立分配的，因此传统上这些相同内容都要白白占用内存。KSM 通过定期扫描物理内存，找出内容一致的页面，并执行合并（物理页面共享）。

\item Enable recovery from hardware memory errors ~~ 使用 ECC 内存就选 Y 。 ECC 内存是带数据校验的内存。多用于服务器，价格昂贵。一般个人电脑不会使用。

\item Transparent Hugepage Support ~~ x86 系统使用分页技术。每个页面的大小是 4k 。 但是，4k 的页面在大地址空间下需要大量的页表。因此 CPU 提供了 大页面支持。 可以使用大于 4k 的页面。Transparent Hugepage 就是透明的合并一些连续的已分配内存的页表，改用大页面页表，以加快 CPU 的地址翻译速度同时减少页表占用的内存。
建议启用。 “Transparent Hugepage Support sysfs defaults” 有两个选项，一个是 madvise 一个 always。 madvise 指程序需要调用 madvise 系统才对指定的内存区域使用大页面。 一个是 always 指系统无需madvise的指示就可以使用大页面。

\item Cross Memory Support ~~ 这是 Linux 新增的一个功能，允许一个进程读写另一个进程的内存——只要有足够的权限。

\item Enable cleancache driver to cache clean pages if tmem is present ~~ 当启用 zcache 的时候，启用这个功能非常有用，能显著的减少磁盘 IO。 当然， 你需要执行 tmem 这个程序。 没有 tmem 就选 N 。

\item Enable frontswap to cache swap pages if tmem is present ~~ 同理，能显著的减少交换空间的磁盘访问。 同上，需要执行 tmem 。 没有 tmem 就选 N 。

\item Contiguous Memory Allocator ~~ 允许子系统分配连续的大块物理内存。不确定就选N。

\item Check for low memory corruption ~~ 一些有问题的 BIOS 会把低端内存弄坏。开启这个功能周期性的扫描低端内存以检查是不是中招。

\item MTRR (Memory Type Range Register) support ~~ MTRR 寄存器用于指示 CPU 对内存的访问。为性能考虑，这个最好启用。

\item EFI runtime service support ~~ 如果你使用的主板支持 UEFI 或 EFI , 那么就开启这个选项。这个允许内核从 EFI 系统上启动。如果不支持 EFI, 开启这个选项也是安全的。 有关 EFI 和 BIOS 相关的启动信息参考  \secref{sec:bootloader}

\item Timer frequency ~~ Linux 内核需要周期性的时钟中断。时钟中断给了内核完成很多事情的契机。更高的时钟中断也带来更高的时钟进度。建议保留默认的 1000HZ。

\item kexec system call ~~ 启用 kexec 支持。 kexec 允许用一个新内核 “替换“ 正在运行的内核。可以不\textit{重启}换内核。需要安装 sys-apps/kexec-tools。


\item Build a relocatable kernel ~~ 如果要编译一个能被 kexec 的内核，就需要把内核加载到非默认地址。因为默认地址上已经有一个正在跑的内核了。 与是，就需要编译一个可重定位的内核。 可重定位的内核能加载到非默认地址。一些 EFI 引导器也要求内核能重定位。

\item Compat VDSO support ~~ 运行 2.3.3 以前版本的 glibc 需要启用这个。否则 N 。

\item Built-in kernel command line ~~ 默认的内核引导参数。关于内核引导参数，参考  \faqref{FAQ:kernelparamter}。 一些不能传递内含引导参数的情况下，将引导参数写死到内核里也是个好办法。

\end{enumerate}

Processor type and features 选项就解释完毕了。下面是 Power management and ACPI options。 这个选项用来配置内核的电源管理和 ACPI 功能。

图 \ref{fig:kernelconfig3} 是 Power management and ACPI options 内子选项。

\chatu{kernelconfig3}{内核配置界面-Power management and ACPI option} 


\begin{enumerate}

\item Suspend to RAM and standby ~~ 待机功能。选 Y 。
\item Hibernation (aka \texttt{'} suspend to disk\texttt{'}) 休眠功能。需要注意的是，休眠功能需要使用交换分区。另外启动的时候内核要添加引导参数 resume=交换分区名。 有关交换分区，请参考 \faqref{FAQ:swap}。

\item Opportunistic sleep ~~ 开启这个选项将更节能。
\item User space wakeup sources interface ~~ 允许用户空间创建唤醒事件。不清楚就选 N。

\item Run-time PM core functionality ~~ 如果不可怜你的电费单可以选N。否则就 Y 啦！

\item Enable workqueue power-efficient mode by default ~~ workqueue 有性能模式和节能模式，选 Y 让 workqueue 模式使用节能模式。建议的选择是 N 。

\item ACPI (Advanced Configuration and Power Interface) Support

\begin{enumerate}
\item AC Adapter ~~ 选 Y。

\item Battery ~~ 有电池供电选 Y。通常笔记本选Y ， 台式机不选。
\item Button ~~ 选 Y。 这个可以支持主机上的关机按钮，按下后桌面环境会弹出关机对话框。
\item Fan ~~ 选 Y 可以允许程序控制风扇。选 N 则不能控制。靠主板自己调节。
\item Dock ~~ 一些型号的笔记本，有一个可移除的底座，底座上有附加的功能。选这个选项用以支持这样型号的笔记本。
\item Processor ~~ 选 Y。 空闲的时候可以让处理器进入低功耗模式。
\item Processor Aggregator ~~ ACPI 4.0 里定义的新功能，让 CPU 进一步进入低功耗模式。
\item Thermal Zone ~~ 关了你的 CPU 可能会被烤焦 。
\item  ACPI tables override via initrd ~~ 对一些不合格的主板，有时候需要用自己产的 ACPI 表替换掉 BIOS 提供的，选上这个就可以在 initrd 里带上自己准备的版本，吼吼。 具体参考内核源码里的 Documentation/acpi/initrd\_table\_override.txt  文件。
\item PCI slot detection driver ~~ 不确定就选 N。 这个是用于确定 PCI 设备的物理插槽位置的。
\item Smart Battery System ~~ 智能电池系统。如果你的笔记本支持，就选 Y。
\item Boottime Graphics Resource Table support ~~ 不确定就选N。这个允许操作系统访问冷启动画面。
\item ACPI Platform Error Interface ~~ 允许硬件向系统报告错误。
\end{enumerate}

\item SFI (Simple Firmware Interface) Support ~~ 不确定就选 N。 运行内核在使用 SFI 固件的硬件上启动。 详情参考 \url{http://simplefirmware.org} 。

\item CPU Frequency scaling ~~ 现在 CPU 越来越注重节能了，都可以在运行的时候变更自己的频率。这个选项就是让内核控制好 CPU 的频率变更功能。频率控制有五种模式： 性能、节能、用户控制、自动-按需、自动-保守。性能模式，就是一直让 CPU 工作在最高频率；节能模式就是一直工作在最低频率；用户控制，就是允许用户自己选择要运行的模式（需要安装调节频率的小程序配合）；按需模式就是让内核按照当前的工作负载自动调节；保守的自动模式，和按需差不多，但是在调节频率上比较保守，对频率的调节采取小步慢进的模式。用户可以选择几样自己用的到的模式，并设定其中一个模式为默认模式。推荐的设置是全部启用，然后设定按需为默认模式。 

\item x86 CPU frequency scaling drivers ~~ 既然要调节 CPU 的频率，就需要一个驱动完成这项工作。这个选项下面列出了几个驱动。有 Intel 牌的，AMD 牌的， 和 ACPI 通用牌的。如果不清楚的话，建议全部选上，让内核自动检测硬件自动使用一个合适的。

\item CPU Idle ~~ 当然是必须启用的功能。

\end{enumerate}

接下来是 Bus options (PCI etc.) 选项。用来设置 PCI、 PCI-E、 ISA 之类的总线支持。
PCI 也 PCI Express 总线自然是要启用的。 留意下 PCCard (PCMCIA/CardBus) support ， 笔记本通常还保留有这个总线，用于插上网卡之类的设备。但是新出产的笔记本也取消了这个总线。毕竟上网卡之类的，都 USB 接口化了。
ISA 总线在 PC 上已经淘汰了。二十年前的硬件上兴许还能看到。 最新的内核可能已经没有 ISA 的选项了。

接下来是  Executable file formats / Emulations 选项。 这个选项用以控制内核支持的可执行文件格式。

\chatu{kernelconfig4}{内核配置界面 - Executable file formats / Emulations}

如图 \ref{fig:kernelconfig4} 这里只有 7 个选项。 32 位内含没有 IA32 Emulation 选项。64位内核启用 IA32 Emulation 选项可以执行 32 位的程序，目前 32 位程序并没有完全消失，所以这个选项是必须启用的。建议就是开启到如笔者这里就可以了。毕竟没人希望眼看着好好的程序执行不起来吧 ;)


下面，进入 Internet 时代的选项 —— Networking support。

\begin{enumerate}
\item Networking options ~~ 先提个醒，这里的选项基本上都要启用，否则你的网络功能不完整，那要望洋（internet）兴叹咯。另外，一些网络功能其实是本地程序要用的，并不是说不联网就可以不启用哦～。


\begin{enumerate}

\renewcommand{\labelenumii}{\roman{enumii}.}

\item Packet socket ~~ 喂喂，这个用来支持 AF\_PACKET ， 你不想一些程序跑不起来吧。 当然是选 Y。  Packet: sockets monitoring interface 顺便也得启用。
\item Unix domain sockets ~~ 这个用来支持 AF\_LOCAL 。 别怪我没提醒你， 禁用这个基本上你的系统就废了。这个是非常非常重要的 IPC 功能。dbus , X11 都在用它。

\item Transformation user configuration interface ~~ 当然启用。这个被一些网络程序用到，如 IPSec。

\item Transformation sub policy support ~~ 不确定就选N。
\item Transformation migrate database ~~ 不确定就选N。
\item Transformation statistics ~~ 不确定就选N。
 
\item  PF\_KEY sockets ~~ 除非你知道你在干嘛，否则选 Y。 这个被 IPSec 之类的工具用到。
\item TCP/IP networking ~~ 好吧，这个是 IP/TCP 协议支持，我想没有人会拒绝选 Y 吧。 吼吼。下面那么多子选项，如果你不太清楚，就选 Y ，免得出现网络问题。

\begin{enumerate}

\item The IPv6 protocol ~~ 好吧， IPv6 ，目前在中国还基本上没有使用 IPv6 的，除了一些大学。如果你在一些使用校园网的学校，可以勇敢的启用哦。如果你运行 systemd , 那么 IPv6 是必须要启用的。
可以只启用  IPv6, 下面的子选项都不需要启用，以节约宝贵的资源。

\end{enumerate}

\item Security Marking ~~ 如果你对 nfmark 不熟悉，那就别选 Y。
\item Timestamping in PHY devices ~~ 为接受到的原始数据包添加时间戳。当然，前提是硬件支持。不知道的话选 N。

\item Network packet filtering framework (Netfilter) ~~ Netfilter 是一个处理和过滤数据包的框架，当然要启用咯。因为防火墙就建立在这个框架之上。

关于 Netfilter 可是有很长的研究的，如果只是尽快的想让内核跑起来，暂时不选也没什么，顶多就是没有防火墙功能。其子选项非常丰富，如果不确定的，可以全部选上。

\item Asynchronous Transfer Mode (ATM) ~~  ATM 传输，你需要一个支持 ATM 协议的卡才需要启用这个选项。 不知道就选 N。 

\item Layer Two Tunneling Protocol ~~ L2TP 是一个隧道协议，用于支持 VPN。 在一些学校和公司，需要执行一次 L2TP  拨号才能连上因特网。如果有这方面的需求，就选 Y。

\item 802.1d Ethernet Bridging ~~ 不确定就选 N。 打算将 Linux 用来跑路由器上的，可以启用这个。

\item 802.1Q/802.1ad VLAN Support ~~ 不确定就选 N。

\item DECnet Support ~~ 有支持 DECnet 协议的网卡？ 没有吧，选 N。 现在都什么时代了，非以太网的协议早淘汰了。

\item ANSI/IEEE 802.2 LLC type 2 Support ~~ 如果要支持 PF\_LLC 这里就选 Y，不知道就选 N。

\item The IPX protocol ~~ 不同于 IP/TCP 协议，这个协议是为局域网开发的，从前老的能局域网对战的游戏都用这个协议 ... 如果有这方面的需求，就选 Y 吧。

\item Appletalk protocol support ~~ 局域网里存在老苹果电脑的，可以用这个协议互连。

\item CCITT X.25 Packet Layer ~~  X.25 协议支持。有这种网卡的就选 Y， 没有就选 N。 
\item LAPB Data Link Driver ~~ 参考上条。

\item Phonet protocols family ~~ 这个协议是 诺基亚 开发的，用来连它的猫的。不清楚就选 N。

\item IEEE Std 802.15.4 Low-Rate Wireless Personal Area Networks support ~~ 好吧，这个是 LR-WPAN， 不是 WIFI。 不清楚的选 N。

\item QoS and/or fair queueing ~~ QoS 是决定数据包优先级的东西。不清楚的就选 N，你将获得标准的行为。除法你打算进行定制。

\item Data Center Bridging support ~~ 不确定就选 N。

\item DNS Resolver support ~~ DNS 解析支持。一般选 N。 注意这个不是用户程序做的DNS解析，那个是 glibc 完成的。这里的 DNS 解析一般用在 CIFS 这类网络文件系统里，需要内核解析域名。

\item BATMAN Advanced Meshing Protocol  ~~  选 N 。 
参考  \url{http://www.open-mesh.org} 获取更多有关 BATMAN 协议的信息。

\item Open vSwitch ~~ 选 N。 参考  \url{http://openvswitch.org}。
\item Virtual Socket protocol ~~ 在虚拟机和宿主机之间传递包裹用。不确定选 N。

\item NETLINK: mmaped IO ~~ 对 netlink 支持内存映射操作。不确定选 N。
\item NETLINK: socket monitoring interface ~~ 不确定就选 N。

\item MPLS: GSO support ~~ MPLS 即多协议标签交换，是一种可提供高性价比和多业务能力的交换技术。该选项让碎片化的非 MPLS GSO 变成 MPLS GSO 包。对家庭用户来说，选 N。
\item High-availability Seamless Redundancy (HSR) ~~ 不清楚的话选 N。

\item enable BPF Just In Time compiler ~~ BPF, Berkeley Packet Filter, 是一种数据包过滤技术。用来网络抓包和监控。BPF 用一种二进制代码来描述数据包的过滤过程，内核解释执行该代码完成数据过滤。开启此选项，内核能够将过滤代码编译为本机代码，这样就无需解释执行了。加快了过滤速度。建议选 Y。

\end{enumerate}

\item Amateur Radio support ~~ 用来支持业余无线电。相关问题咨询 \url{http://www.tldp.org/docs.html} 
\item CAN bus subsystem support ~~ CAN bus 是一个低速串行通信协议。广泛用在汽车和一些医疗产品中。如果要使用这个协议， 选 Y 并在子选项中启用需要的驱动。

\item IrDA (infrared) subsystem support ~~ 红外通信协议支持。旧式手机和笔记本通常都有红外通信接口，可以传文件神码的。现在被蓝牙取代了，吼吼。如果你还有，就选 Y。 并参阅 Documentation/networking/irda.txt 文件。

\item Bluetooth subsystem support ~~ 蓝牙支持。有蓝牙设备的话，就选 Y 吧。 KDE 和 GNOME 都集成了蓝牙支持，无需额外安装蓝牙套件。可以直接在桌面环境管理你的蓝牙设备。记得从 Bluetooth device drivers 子选项里启用需要的驱动。通常 USB 接口的蓝牙接收器选 HCI USB driver 即可。

\item RxRPC session sockets ~~ 选 N 。目前并不完善。

\item Wireless ~~ 无线网络。 当然是 WiFi 支持啦！笔记本都需要。一些使用了 PCI 无线网卡的台式机也需要。

\begin{enumerate}

\item cfg80211 - wireless configuration API ~~  必须启用啊，否则你无法配置无线网卡。

\item Generic IEEE 802.11 Networking Stack (mac80211) ~~ 选 Y。

\end{enumerate}

\item WiMAX Wireless Broadband support ~~ 4G 无线宽带。和 LTE 一样属于 4G 上网。不确定就选 N 。

\item RF switch subsystem support ~~ 笔记本通常提供软开关，用来关闭无线网卡。启用这个选项以支持这项功能。

\item Plan 9 Resource Sharing Support (9P2000) ~~ 不确定就选 N。

\item CAIF support ~~ 选 N。

\item Ceph core library~~ 选 N。
\item NFC subsystem support ~~ 近场通讯。 手机上用到的，把两个手机相互靠近就能通信的一个东西。如果不是在造手机，还是选 N 吧。 

\end{enumerate}

接下来就是硬件驱动和文件系统驱动了。硬件驱动，简单的原则是，有这个硬件就选 Y 没有就选 N。不清楚有没有就选 M。

文件系统则是，用到就选 Y 没用到就不选。不常用到，想节约内存，就选 M。

唯一要注意的是，一定要选上 AHCI 。 Serial ATA and Parallel ATA drivers 下面的，  AHCI SATA support。
如果你的电脑比较老旧，不支持 AHCI 或者还是 IDE 硬盘，那一定要将对应硬件选 Y 而不是 M。否则开机过程中会出现先有鸡还是先有鸡蛋的问题。没有磁盘驱动就无法访问模块文件，无法访问模块文件就无法载入磁盘驱动。

还有， SCSI 支持一定要选上。“SCSI device support” 和 “SCSI disk support” 是必选的。因为虽然没有 SCSI 设备，但是内核是将 SATA 硬盘模拟成 SCSI 硬盘的。

设备驱动的选择是编译内核最大的难题。需要非常熟悉自己的硬件。很难一次配置成功。所以建议新手使用 genkernel。

\subsubsection{新手使用 genkernel}

编译内核最难的地方就是配置内核了。 make menuconfig|xconfig|gconfig 最终的目的就是获得一个.config文件。其实后面的编译步骤都不难，难的就是一个能正常工作的 .config 配置文件。
似乎某些重要配置不能随意回答Yes或No，对于没有多少经验的人来说，还是有不小的难度么。

Gentoo开发者似乎注意到了不会自行编译内核的しんまい%ざこ%ばか
\footnote{内核对你来说和这个词是一样的。}%
，所以特意开发了しんまい的内核编译助手：genkernel。虽然genkernel能直接编译出一个工作内核，
但是以genkernel获得一个能工作的.config文件并在此基础上修改定制才是修炼之道。

这里就教大家用genkernel获得一个能启动的内核的.config配置文件好了。

\subsubsection{initrd}\label{sec:initrd}

Linux启动过程中内核主动做的最后一件事情就是执行/sbin/init，那么倒数第二件事情必然是挂载根目录。
挂载根目录有2个前提条件：找到根分区和支持根分区所使用的文件系统。
找到根分区，第一点，需要其所在设备的驱动。第二点就是内核需要知道哪个是根分区。这个可以通过内核参数 root= 设定。

Linux下驱动可以编译进内核，也可以编译为模块动态加载。如果恰巧磁盘设备的驱动被编译成了模块，就出现“鸡生蛋和蛋生鸡”的问题了：不访问磁盘就无法加载模块，而不加载模块有没有磁盘的驱动。这对于自己编译内核的童鞋不成问题。只要知道自己的SATA/IDE控制器的型号，将对应的驱动编译进内核就解决了。但是如果是打算制作“通用内核”的发行版呢？

文件系统也是同样的问题，没有文件系统驱动就无法挂载分区访问上面的文件，无法访问文件就无法加载模块。这对于自己编译内核的童鞋不成问题。只要将自己根分区所使用的文件系统编译进内核就解决了。但是如果是打算制作“通用内核”的发行版呢？将所有的文件系统都编译进内核么？

initramfs就是为解决这个“鸡生蛋和蛋生鸡”的问题诞生的。

GRUB加载内核的时候通过initrd 命令加载一个initramfs，该initramfs就会成为内核的根分区\footnote{事实上是Linux内核拥有一个不挂载就存在的rootfs文件系统，Linux启动过程中将initramfs的内容解压到rootfs中。作为理解，可以认为initramfs挂载成Linux的临时根目录。}。内核初始化完毕后直接执行initramfs中的/sbin/init而不理会root=命令行参数。由initramfs中的/sbin/init程序负责加载磁盘驱动，查找根分区，挂载根分区到 /sysroot 并调用pivot\_root()\footnote{使用命令　man 2 pivot\_root查看介绍。}将找到的根目录替换为真正的根目录。然后执行真正的/sbin/init程序。

initramfs是一个目录的压缩包。由dracut程序生成。initramfs用于寻找真正的根目录时的临时根目录，通常只包含一些必要的驱动和一些用于查找根分区的脚本，但有时候可以成为一个完成的系统的根目录，这样的系统可以脱离磁盘运行。比如Fedora的安装DVD,使用一个巨大initramfs作为安装环境的根目录。该initramfs包含的是一个真正的OS，带有完整的工具。


通常 Gentoo 用户不使用 initrd ，但是遇到下列情况，你需要 initrd：

\begin{enumerate}
\item 使用 genkernel 编译内核。
\item root 文件系统放在 LVM/RAID 上。或者其他无法由内核直接挂载的地方，如一些网络文件系统。
\item root 文件系统使用的驱动没有编译进内核。
\item / 和 /usr  不在一个分区。
\end{enumerate}

genkernel 会自动生成 initrd ,但是如果你不依靠 genkernel , 可以通过 dracut 这个实用工具生成。

\begin{code}
\# emerge dracut

\# dracut -H
\end{code}

使用 -H 参数，表示生成的是本机使用的 initrd 如果没有 -H ，dracut 默认会尝试生成一个通用的initrd , 通常用于制作发行版的 LiveCD。

dracut 默认会生成 initramfs-<kernel version>.img 并存放到 /boot 目录下。通过在命令行上直接给出文件名，可以覆盖默认的文件名。
如：

\begin{code}
\# dracut -H   initrd.img
\end{code}

有关 dracut 的进一步的信息，参考 man dracut。

\subsection{系统初步配置—重启前配置}
在chroot环境中我们需要做的就是进行必备系统的安装，配置好基本系统然后重启进入Gentoo环境中进行后续的安装。所以这里基础系统的配置非常关键，他关系到你第一次能否成功地进入Gentoo系统。

\subsubsection{nano 编辑器快速入门}\label{sec:nano}
现在才讲nano我都觉得有点罪过。由于Gentoo的stage3里并没有包含vim，所以一开始必须使用nano进行系统配置，直到portage配置好后emerge了vim才能使用vim。
在此之前还是得用nano。

nano 编辑器的使用非常简单，打开一个文件用“ nano 要编辑的文件” 这种简单的指令就可以了。编辑也非常简单，支持
\textuparrow{} \textdownarrow{} \textleftarrow{} \textrightarrow{} 上下左右移动。退出按Ctrl-X，提示是否保存，回答yes就可以了。


\subsubsection{时区设置}

类UNIX系统下，系统内核并不使用本地时间。内核使用的是被称为UTC\footnote{英文Coordinated Universal Time, 法文Temps Universel Coordonné。为了协调英文和法文用户，缩写为 UTC。}的时钟：文件系统上保存的文件最后修改时间即UTC时间而不是本地时间。
需要设置时区才能正确从UTC转换到本地时间。

同时，系统硬件存储的通常是本地时间——为了和Windows兼容，以及发挥定时开机等功能——内核启动的时候也需要将本地时间转化为UTC。
所以时区信息分2个，第一，硬件时间是UTC还是本地时间，第二，本地时区。
系统在启动过程中会使用 /etc/adjtime 的信息确定是否为本地时间。

\begin{filecontent}{/etc/adjtime}
0.0 0.0 0.0\\
0\\
LOCAL
\end{filecontent}

前两行忽略，后最后一行 LOCAL 即向系统传达硬件时钟为本地时间这个信息。如果不存在这个文件，就默认为UTC时间。赶快拿起编辑器生成这个文件吧。呵呵。

时区信息保存在/etc/localtime，但是这个文件不要手工编辑。在系统的 /usr/share/zoneinfo/ 下保存有各个时区的 localtime 文件。拷贝为 /etc/localtime 就可以了。
比如大陆使用东八区，城市为上海\footnote{居然不是北京的说}，所以拷贝 /usr/share/zoneinfo/Asia/Shanghai 为 /etc/localtime 就完成时区设置了。

系统启动的时候，会读取硬件时间，然后转化为 UTC 时间设置为内核时间，该操作由hwclock程序完成。systemd的做法是由内核直接读取硬件时间（配置内核的时候选上 “Set system time from RTC on startup and resume”）， 到systemd启动的时候，按照本地时间和UTC的差拨快或拨慢内核时间即可。

使用 date 命令打印的时间是从内核获得的 UTC 时间然后依据时区设置转化为本地时间再显示。

\subsubsection{fstab配置}
安装Gentoo，可以说，最重要的配置莫过于fstab了。fstab配置文件是Linux文件系统树的根基。正是fstab指明了文件系统的建筑结构。

fstab 每行配置一个文件系统。语法是:
\begin{small}
\tt 
\begin{longtable}{cccccc}
文件系统设备	& 挂载点 & 文件系统类型 & 挂载选项 & dump & pass
\end{longtable}
\end{small}

文件系统设备可以直接使用 /dev/sdXX 这样的设备，也可以使用 UUID= LABEL= 这样的形式。
挂载选项每个文件系统都有特定的选项。如果没有特殊的需要使用 defaults 即可。dump 位如果设置为 1 表示如果检查失败就不能继续启动系统。通常设置为 1。
pass 对于 / 设置为 1, 其他都设置为 2 。

下面给出的是挂载 / 的例子。
\begin{small}
\tt 
\begin{longtable}{cccccc}
/dev/sda2	& /	 & ext4	& defaults & 1 & 1
\end{longtable}
\end{small}

使用 LABEL也可以，假设 / 分区的 LABEL 是 Gentoo

\begin{small}
\tt 
\begin{longtable}{cccccc}
LABEL=Gentoo	& /	 & ext4	& defaults & 1 & 1
\end{longtable}
\end{small}

值得注意的是，swap分区也是在这里设置的。只不过， 挂载点设置为 none，文件系统类型为 swap , dump 和 pass 设置为 0 

\begin{small}
\tt 
\begin{longtable}{cccccc}
LABEL=SWAP	& none	 & swap	& defaults & 0 & 0
\end{longtable}
\end{small}

实际上Gentoo带的stage3已经提供了一个 fstab的模板，只要将自己的/boot分区、/分区、/home分区对应的设备写进去，调节一下文件系统类型就可以了。呵呵。


\subsubsection{配置帐号，设置密码}

我安装Gentoo的时候常常忘记为Gentoo设置root密码，结果重启后发现无法使用root登录。无密码和空密码不是一回事。Linux不允许登录无密码帐号。所以这个时候要做的事情就是设置一个root密码。

\begin{code}
\#passwd\\
New UNIX password: *输入密码，注意这里不回显*\\
Retype new UNIX password: *再次输入密码*\\
passwd: password updated successfully
\end{code}

这样就设置好root帐号的密码了。

\subsubsection{语言设置}\label{sec:lang}

很多程序都有本地化语言支持。支持本地化的程序，通过检查环境变量LANG的内容决定使用何种语言。

LANG环境变量的格式为 语言.编码。比如简体中文为 zh\_CN，使用UTF-8编码，则LANG设置为 zh\_CN.UTF-8（zh\_CN.utf8等价）。
编码影响的是程序向终端输出的时候使用的编码。绝对只建议使用utf-8编码。

\begin{code}
\#eselect locale list
\end{code}

列出支持的语言，使用

\begin{code}
\#eselect locale set zh\_CN.utf8
\end{code}

设置语言为中文，编码为utf8。
如果zh\_CN没有出现在列表里，应该考虑一下 /etc/locale.gen 文件里是不是有
``zh\_CN.UTF-8 UTF-8''这么一行。没有的话添加了然后执行。

\begin{code}
\#locale-gen
\end{code}

这样就能重新生成语言支持文件。

\subsection{安装基础系统}

除了内核，引导程序，和init程序，还有一些必须的基础程序需要安装。主要是：开机fsck需要用到的各种文件系统的辅助工具和网络工具。

\subsubsection{安装网络工具}

网络工具的具体用法请参考\chapref{chap:network}。

这里作为一个快速安装的指南，算是一个占位符吧。告诉你，现在就得把网络工具安装了，否则待会重启进的了系统却用不了网络，无法继续安装哦！

安装何种工具就看你使用的是何种网络环境了。
一般如果是静态网络，安装 net-tools 即可。它提供了ifconfig。
但是强烈的推荐使用 NetworkManager 这样的成熟的动态配置工具，尤其是还有好用的GUI可用。NetworkManager依赖较多，如果你的网络可以用 ifconfig 配置下来，可以等重启后再安装。

不过如果使用了无线网络，最好还是现在就把NetworkManager装了，免得重启后无法继续安装系统。

如果使用 PPPoE 进行拨号的（ADSL接入，在电脑上进行拨号。），除了安装NetworkManager还得安装 rp-pppoe。这样 NetworkManager就可以进行拨号了。
具体用法请参考 \chapref{chap:network}

因为 networkmanager会引入 bluetooth 依赖，现在还用不到，可以关闭。之后可以重新编译。

\begin{code}
\#emerge rp-pppoe\\
\#USE=-bluetooth emerge networkmanager
\end{code}

\subsubsection{文件系统工具}

众所周知，不正常的关机会导致磁盘数据丢失。所以每次linux开机都会进行一次磁盘检查。如果上次是正常关机，通常检查会很快结束，网络词语叫秒杀。
如果上次经历了非正常关机，就需要花比较长的时间进行检查。至于检查嘛，自然是使用各种文件系统检查工具了。init系统会统一调用 /sbin/fsck 来检查文件系统，不过 fsck 只做一件事情（UNIX哲学嘛）：检查文件系统类型，调用 fsck.??? 进行检查。而fsck.??? 自然就是属于各个文件系统专用工具集合里的了。

\begin{code}
\#emerge e2fsprogs	\#安装 ext2/ext3/ext4 系列的文件系统工具\\
\#emerge xfsprogs	\#安装 xfs 文件系统工具\\
\#emerge reiserfsprogs	\#安装 reiserfs 文件系统工具\\
\#emerge btrfs-progs	\#安装 btrfs 文件系统工具\\
\end{code}

\subsubsection{安装内核}

去去去，查看一下内核编译完成了没。内核编译完成后，就剩下安装了。

\begin{code}
\#cd /usr/src/linux\\
\#make install modules\_install
\end{code}

\subsubsection{安装引导程序}

引导程序前面已经介绍过了，现在我们要做的就是安装一个。对于BIOS机器，执行下面的命令安装的老版本的grub:

\begin{code}
\#emerge grub-static
\end{code}

这个时候　grub　会被安装到 /sbin/grub，　16位模式的GRUB模块会被安装到/usr/lib/grub。拷贝/usr/lib/grub/一份到 /boot/grub。

\begin{code}
\#cp -a /usr/lib/grub /boot/
\end{code}

/sbin/grub 这个命令可以在Linux环境下执行，除了不能启动Linux内核，GRUB能用的其他命令可以在这个命令行下使用。
执行：
\begin{code}
\#grub
\end{code}

会出现grub的交互式命令行模式，这个和BIOS启动GRUB，按C进入的命令行模式是一样的。唯一的区别就是这个模式不能真正的启动Linux。
现在我们要做的就是执行　root (hdX,Y)命令定位到 /boot　分区（如果/boot分区列位有划分独立的分区的话，否则就是定位到/分区。），然后执行{\tt setup (hd0)} (或者 hd1,看你要安装到哪个硬盘的MBR上。)。
这么做的意思就是将 grub 的“接头人”装到MBR上，由这份接头人性质的代码加载 /boot/grub/stage2实现GRUB的启动。

也可以不进入grub命令行模式，直接使用 grub-install 命令进行安装。但是熟悉grub交互式命令行模式对以后使用grub修复受损的MBR也是非常有用的。所以推荐使用grub的交互式命令行模式。

如果是UEFI的机器，则需要安装 EFI 版本的grub。不过linux内核本身也是可以直接被UEFI启动的。但是使用引导程序提供了一些灵活性，还是值得的。
grub-efi 在 gentoo-zh overlay里。介绍和使用overlay请参考 \secref{sec:overlay}

\begin{code}
\#emerge grub-efi
\end{code}

如果还没添加 gentoo-zh overlay 将找不到对应的包，可以使用下面的命令添加 gentoo-zh overlay。具体请务必查阅\secref{sec:overlay}。

\begin{code}
\#USE=git emerge layman \#管理 overlay 的工具\\

\#layman -a gentoo-zh \# 添加 gentoo-zh overlay，如果还没添加的话 

\#echo~source~/var/lib/layman/make.conf~>>~/etc/portage/make.conf
\end{code}

按照提示，拷贝 /usr/lib/grub/grub.efi 到 EFI 分区。

\subsubsection*{使用efibootmgr将grub.efi加入引导}

上文提到过UEFI支持在操作系统环境下管理启动选项。这个管理操作由efibootmgr工具实现。当然，前提是系统是由EFI模式启动的。目前大部分EFI主板都提供了BIOS模拟，所以很有可能您的LiveCD是用BIOS模拟模式启动的，这样就会失去EFI管理能力。所以确保您是用EFI模式启动的LiveCD，可以参考　\faqref{FAQ:UseLiveCD}。

当然第一件事情是安装它。
\begin{code}
\#emerge efibootmgr
\end{code}

首先显示出现有的引导选项
\begin{code}
\#efibootmgr
\end{code}

添加GRUB引导选项，将　\textbackslash{}grub.efi 添加到引导选项中去。
\begin{code}
\#efibootmgr　-c -L GRUB -l '\textbackslash{}grub.efi'
\end{code}

确认结果
\begin{code}
\#efibootmgr
\end{code}

记住　Boot????*GRUB　那4个问号代表的数字，表示新添加的GRUB处于何种次序。将GRUB设为第一引导次序，也就是默认引导项，如下文的命令。????就是刚刚找到的新添加的GRUB引导项的数字。

\begin{code}
\#efibootmgr　-o ????
\end{code}

这样就可以了。重启后系统会自动运行 grub.efi　进入grub模式，grub就可以将内核加载了。

\subsection{配置引导程序}\label{sec:bootloader}

grub 的配置文件就是 grub.conf。 UEFI 模式下，是和 grub.efi 在同一个目录下。 BIOS模式下，是在 /boot/grub/grub.conf。（果把grub安装到boot分区的话）

在 \secref{sec:grub} 我们已经大概熟悉了一下grub.conf文件的语法。
这里要做的就是将内核加入 grub 的引导选项了。

\begin{code}
title Gentoo\\
kernel /vmlinuz ro root=LABEL=GENTOO\\
\end{code}

主要，这里的 /vmlinuz 是按照 grub 的路径来算的。比如 grub 安装到/boot分区，而内核文件（通过 make install 安装）通常为 /boot/vmlinuz-内核版本。相对 grub 的路径就是 /vmlinuz-内核版本。如果 grub 安装到 /boot/efi/ (/boot/efi 为独立的 EFI 分区) ，则必须使用 root (hdX,Y) 这样的语句选择 /boot 分区或者 / 分区（如果/boot不是独立分区的话）。
然后再使用 /vmlinuz-内核版本 或者 /boot/vmlinuz-内核版本 （/boot不是独立分区的时候） 这样的路径。

如果使用了 initrd 的话（强烈的不推荐Gentoo系统使用initrd），还需要加上 initrd 命令加载 initrd文件。initrd文件可以由dracut程序生成。参考 \secref{sec:initrd}

\section{重启后安装}

前面的安装都是在liveCD环境下进行的。现在要进行重启并进入Gentoo自己的环境中进行后续安装。这次重启万一失败就得重新进入liveCD进行修缮。
通常失败无外忽 grub 没有正确安装，或者内核没有配置正确。

好了，退出 liveCD 下的那么多 shell , 重启电脑吧。

运气好的话，大概你的电脑已经出现登录提示了。以 root 登录进行安装。既然一个基本的Gentoo已经能自我启动了（不一定进入可用状态），Gentoo的自我救赎式的安装就已经完成了。
剩下的无非就是按照自己的需求将一些用到的软件进行安装以让系统进入可用状态。让我们随着使用Linux的过程中随着安装缺失软件吧！

\fi

\chapter{Linux的日常}

\begin{quotes}[0.52]{Jef Raskin}
A computer shall not waste your time or require you to do more work than is strictly necessary.
\end{quotes}

\ifdefined\compileall

\chapter{Linux的日常}

\begin{quotes}[0.52]{Jef Raskin}
A computer shall not waste your time or require you to do more work than is strictly necessary.
\end{quotes}

其实很重要的一个问题：我们为什么要学习Linux？须要意识到操作系统只是一个工具，而要使用工具是必须经过学习的。越是强大的工具越是需要学习后方能使用。
得心应手的工具，必定经历常年累月的使用。Linux不仅仅是一个用来学习的系统，还是一个可以经年累月的使用的系统。否则，一个不能使用的系统，也没有存在的价值。

我们需要一个称心如意的系统，必定是需要悉心调教和了解的系统。

\section{某日常的man-pages}

UNIX命令繁多，载籍浩瀚。一个一个介绍纯属不可能之任务。所以一个随时可以借用的手册（manual）是比不可少的。
不过且不论manual是怎么被缩写成man的，man已经成了UNIX管理员的必备工具了。正所谓有问题，man 一下。

要使用man手册请首先按照它：
\begin{code}
\#emerge man man-pages man-pages-zh\_CN
\end{code}

man-pages-zh\_CN带来了部分man手册的中文汉化，按照后有部分手册会使用中文页面。当然，前提是安装支持UTF-8编码的groff-utf8并配置man程序使用groff-utf8替代groff。具体做法请参阅 \faqref{FAQ:zhman}。

上千条命令和系统API的帮助页面被收集起来，形成了man pages，需要查阅的时候只需要使用 ~man + 命令名/API名称~即可方便调阅。


实战演练：
\begin{code}
\$man man
\end{code}

进入阅读模式后，可以使用上下方向键翻动。按“Q”键退出。按“/”键，终端底行会由行数显示改为显示一个/,然后键入字符串，按回车即可查找该字符串。
man本质上是调用的less进行最终的显示。所以其实这里使用的就是less的命令。使用参数 -P 还可以制定其他的分页器。

man pages将系统手册分成8个类别。
\begin{enumerate}
       \item      命令行命令
       \item      系统调用
       \item      标准C库和POSIX定义的API
       \item      设备文件
       \item      文件格式
       \item      游戏
       \item      杂项
       \item      系统管理工具和后台服务
\end{enumerate}

遇到有重名的时候就可以指定一个数字表示显示指定类别的手册。否则显示序号最小的那个类别的手册。

\begin{example}{printf}
shell提供了printf命令用于打印格式化的字符串，同时printf()也是标准C库的组成部分。

直接使用 \texttt{\$man printf} 得到的就是printf命令的手册。要显示 printf() 函数的手册， 必须使用 \texttt{\$man 3 printf}。
\end{example}

%由于man手册为手册划分了类别，所以我们提到参阅某某某的时候，后面都会打上括号，指出是哪个类别。比如会有“参考man手册 printf(3)”这样的说法，指的就是参考第3类目的手册，也就是说，使用 \texttt{ \$man 3 printf } 查看。有时也会省略“man手册”，直接使用“参考printf(3)”这样的用法。

{
\bf 
由于man手册为手册划分了类别，所以我们提到参阅某某某的时候，后面都会打上括号，指出是哪个类别。比如会有“参考man手册 printf(3)”这样的说法，指的就是参考第3类目的手册，也就是说，使用 \texttt{ \$man 3 printf } 查看。有时也会省略“man手册”，直接使用“参考printf(3)”这样的用法。
也就是说，“参考XXX(1-8之间的数字) ” 已经成为“让读者使用man命令参考相应内容，并指出该内容隶属于哪个类别”的固定句式。英文里，通常是 “see printf(3)” 这样的用法。

\it 
%由于man手册为手册划分了类别，所以我们提到参阅某某某的时候，后面都会打上括号，指出是哪个类别。比如会有“参考man手册 printf(3)”这样的说法，指的就是参考第3类目的手册，也就是说，使用 \texttt{ \$man 3 printf } 查看。有时也会省略“man手册”，直接使用“参考printf(3)”这样的用法。

}


\section{UNIX权限基础}

UNIX是一个多用户系统，意味着UNIX支持同时登录多个帐号互不干扰的操作。要实现这一点就需要对帐号进行权限隔离。良好的权限设计也是杜绝恶意软件的基石。

什么是权限呢？说到底，权限就是一种信任。权限越高的程序，必然也是越受用户信任的程序。受用户信任的程序能从用户获得越多的权限。同时用户本身相互直接也有信任度问题，其中，管理员权限最大，因为他是系统的“拥有者”。管理员觉得用户（的程序）有“夺权”倾向，所以要对普通帐号进行限制。用户觉得自己执行的程序也有可能作出意想不到的行为，除非程序是自己编写的，他不太容易信任第三方程序。所以用户也会想到限制已经帐号下执行的程序的权限。权限是一种谨慎的信任。

最高一级的权限自然是系统所有者，也就是管理员。其次才是内核，然后是系统程序，应用程序，权限依次递减。
很多从非UNIX系统上过来的人都不能理解管理员的权限为何比系统本身要高。对于Windows这样的系统，最高权限是System，其次才是Administrator，有些操作连管理员都是被禁止的。
可是UNIX系统下，管理员的权限凌驾于系统之上。

为什么管理员的权限要凌驾于系统之上呢？ 因为操作系统是“被管理员授权来管理计算机”的。所以，UNIX系统本身严格遵循“不做管理员没有授权的事情”，和“不阻止管理员做任何事情”。
哪怕管理员要 \texttt{"rm -rf /"} UNIX都不会阻止管理员做这件事情。如果哪个系统阻止管理员做危险的事情，已经是“仆人成了主人”——本末倒置了。

\subsection{用户和组}\label{sec:usersandgroups}

UNIX权限是按照用户和组进行组织的。权限最高的是root用户，也称特权帐号，非root用户被称为普通用户。
在内核中，每个用户都对应一个UID，每个组也对应一个GID。root用户即UID=0,GID=0。
每个程序都会继承父进程的UID和GID。也就是说，登录哪个用户，执行的程序就是以哪个用户的权限执行的。登录程序通过改变自己的UID从而影响它启动的shell的UID，不过
只有UID=0的进程才有权更改自己的UID\footnote{使用setuid()系统调用}，这就避免了非特权进程将自己改为特权进程。

只有以root用户执行的进程才能进行的操作有
\begin{enumerate}
\item TCP/IP网络程序，允许绑定 < 1024 的端口。非特权进程无法绑定<1024的端口。
\item 修改ip地址。
\item 加载内核模块。
\item 挂载文件系统。
\item 使用chroot。
\item 执行关机、挂起操作。
\item 更改iptables。
\item 向任意进程发送信号。非特权进程只能向相同用户的进程发送信号。
\item 创建原始套接字。用于网络抓包和发送原始数据包。

等等。
\end{enumerate}

这些都是在系统调用层面进行UID检查保证的。内核对 UID=0 和 UID>0 做了区别对待。

\subsubsection{su和sudo}

有时候我们需要临时的切换到root帐号，而不用退出当前shell并重新登录。这个时候你就需要祭出sudo/su了。

在任何命令前加sudo即可让命令以root账户执行。

\begin{example}{普通用户显示/root目录内容}
\$ls /root 

ls: 无法打开目录/root: 权限不够

\$sudo ls /root

密码： (输入当前用户的密码，不显示星号)

ab.txt  a.txt  Desktop
\end{example}

当然，使用 sudo 的前提是将自己加入到 /etc/sudoers 。 通常 /etc/sudoers 会包含 \texttt{\%wheel ALL=(ALL) ALL} 这一行\footnote{如果被注释了，就解除注释状态。什么？不会，删除前面的 \# 就好了}。意味着将用户加入wheel组即可使用sudo命令。

su 命令用户开启一个新的shell，默认运行于 root 账户下。注：需要输入 root 密码。

su 命令也可用于切换到非root账户。比如 \texttt{su - smith} 就可以切换到 smith 账户。加 - 参数让 su 不要继承当前shell的一些环境变量，并且新开的shell切换到家目录而不是保持在当前目录。


记住：sudo 输入当前帐号的密码，su 输入目标帐号的密码。

\subsection{文件操作权限}

虽然在系统调用这一层面，由UID是不是0而确定不同的权限\footnote{Linux引入了新的机制，被称呼为capabilities。将原先root拥有全部特权和非root没有任何特权的二值判定改为细分的多个capabilities判断。比如进程拥有 CAP\_NET 权限就可以绑定<1024的端口而不必要求root权限。}。看上去似乎UID不为0就是普通权限了，那样还不如用个 bool 标记呢。用户和用户组的概念自然有其存在的必要。
那就是对文件的权限控制。

UNIX下，所有的文件（包括设备文件）都有所有者。文件被一个用户和组所拥有。
文件的所有者模式通常以 “用户名：组名” 的形式表示。
然后定义了文件的3种操作模式：读、写和执行\footnote{UNIX下一个文件是否可执行并不是看它的扩展名，而是看对要执行的用户来说是不是有可执行权限。}。对于目录，执行权限意味者可以切换到该目录作为工作目录\footnote{简单的说，就是可以 cd 进去}。
文件的权限很好理解，这里着重介绍一下目录的权限。
只读的目录意味着不能在目录下创建和删除文件——但是可以修改该目录下所存储的文件——只要拥有对改文件的写权限即可。没有读取权限的目录意味着无法列出目录内容。但是目录的权限并不会递归传递给子目录。子目录需要为自己设置权限。
\begin{example}{~}
所以虽然john对于/home目录是只读的，但是对/home/john是可写的。因为/home目录对john的权限并不能传递到子目录/home/john。
\end{example}

权限的设置是对3类用户设置的：
\begin{enumerate}
\item 拥有者。文件的所有者的权限。读、写和执行。
\item 所属组。文件的所属组的权限。也就是同组的其他用户的权限。因为UNIX下一个用户可以里加入多个组。读、写和执行。
\item 其他。即不是文件所有者，也不是同组的用户的权限。读、写和执行。
\end{enumerate}

系统其他方面的权限都是以组来奠定的。比如声音设备，设备的所有者是root,组是audio。对其他的权限是无。对audio组的权限是读写。那么只要将用户加入到audio组，就获得了声卡的访问权限。
因而系统并不需要在声卡方面添加新的系统调用进行权限验证。

通常一个桌面用户至少要加入的组有：users,video,audio,cdrom,cdrw,usb,plugdev。加入video组获得显卡的直接渲染支持，加入cdrom组以便访问光驱，如果有刻录机的话，加入cdrw组以便使用刻录机，等等。这些权限都是在文件权限上实施的。原理就是让对应的设备文件所有者设置为相应的用户和组
\footnote{内核自然没有UID,GID和账户对照表，默认创建的设备文件都是属于root:root的，文件所属权是由udev负责修改的。}。

文件权限的表示采用的是八进制表示法。

\begin{example}{一个 0664 权限模式的文件解析}

\end{example}

\begin{tabular}{|c|c|c|c|c|c|c|c|c|}
\hline \multicolumn{3}{|c}{所有者}&\multicolumn{3}{|c|}{组所有者}&\multicolumn{3}{c|}{其他}\\ 
\hline 读 & 写 & 执行 & 读 & 写 & 执行 & 读 & 写 & 执行 \\ 
\hline 1 & 1 & 0 & 1 & 1 & 0 & 1 & 0 & 0 \\ 
\hline 
\end{tabular} 

将 $ (110110100)_2 $ 使用八进制表示就是 $(0667)_8$。

\section{Shell 日常应用}

平常说学习shell神码的，其实是一个统称，实际上包括了一个shell和若干个shell能调用的工具。这些工具主要集中在 coreutils 和 linux-utils 包中。也就是说，学习shell不光是学个shell命令的语法格式，还包括一些常用工具的用法。没有这些工具的存在，shell能做到的事情是非常有限的。

shell是操纵电脑的接口。一切对电脑下达的指令都是shell代为执行的。要浏览网页？需要浏览器吧？可是浏览器不会自己动起来。需要下达执行浏览器的指令。shell理解了，然后执行了。
需要删除文件？可是文件不会自己消失。需要下达删除文件的指令。shell理解了，然后执行了。$\cdots$

shell的地位如此重要，以至于有人说，shell就是UNIX\footnote{Cygwin 提供了完整的UNIX命令，但并没有全部实现POSIX。尽管如此，cygwin虽然是一套Windows程序，还是被称呼为UNIX环境。}。

\subsection{bash名称的道听途说}


Bash是多数Linux发行版默认的shell。当然Gentoo也是不落俗套的默认为Bash了。Bash这个名字来源于\textbf{Bourne-Again SHell}。为什么说“又”（Again）呢？因为贝尔实验室发布的第七版UNIX——Unix Version 7搭载的正是\textbf{ Bourne SHell}， 这个shell由\textit{Stephen Bourne}编写。
FSF发扬自己借用UNIX名称的传统美德，给自己的shell起名为Bourne-Again SHell。

名字是借用的Bourne SHell，不过，如果认为它开发出来是为了替代Bourne SHell那就错了，bash诞生最初是为了替代 \textbf{Mashey shell 或者说 PWB shell} —— 这也是 Bourne SHell 所要替代的 。

\subsection{shell的字符串展开}

虽然已经在前面的 \secref{sec:quickbash} 介绍过，
不得不再次在这里重申，shell对字符串会进行分割后将其作为参数字符串数组传递给被调用者。而对于特殊的字符串会执行展开操作，更重要的是，展开操作是发生在字符串切分\textbf{之前}的。避免字符串展开的办法是加引号。同时引号也可以用于避免引号所包含的空格被用于字符串切割开。

会被执行字符串展开的除了带 * 和 ? 的通配符，还有 [] 和 \{\} 。 提醒一下，如果展开的时候没有匹配到适合的文件名，会自动取消展开而保留原字符串。

\subsubsection*{[ ] 的展开}

[~] 的用法是：[字母列表]。字母列表还可以使用a-z这样的简写形式表示连续的a-z而避免写全26个字母。

[~] 只展开为一个字符。

\begin{example}{~}

假设当前目录下有如下的文件:\\
a1 a2 a3 a4 a5 a6 a7 a8 a9 b1 b2 b3 b4 b5 b6 b7 b8 c1 c2 e1 e2

现在使用 ls [ab]1 这样的命令。[abef]1 会被扩展为 a1 b1 e1 三个文件名。

也就是说，传给 ls 的文件名参数有3个， a1 b1 e1。至于 f1 , 因为文件名不存在，也就不会扩展到。

如果使用的是 ls [gh]1 则会提示没有 [gh]1 这个文件。因为没有符合的文件名，所以表达式也就无从展开，以原样传给ls。

\end{example}

总而言之，可以认为 [~] 是一个过滤器，将符合表达式的文件过滤出来。如果没有一个文件名符合表达式的要求，则传递表达式本身。

\subsubsection*{ \{\}  的展开}

\{~ \} 的用法是：{字符串1,字符串2,字符串3,...} 空字符串也是可以接受的。但是连空字符串在内至少要包含两个字符串，否则不展开。

\begin{example}{~}

假设当前目录下有如下的文件:\\

m.c m.cpp m.o m.s m.h

现在使用 ls m.{c,cpp,h,lll} 这样的命令。m.{c,cpp,h,lll} 会被扩展为 m.c m.cpp m.h m.lll 四个文件名。

也就是说，即使文件名不存在，也会扩展到。因此除了显示m.c m.cpp m.h这3个文件信息外，还会提示 m.lll 没有这个文件。

如果使用的是 ls [gh]1 则会提示没有 [gh]1 这个文件。因为没有符合的文件名，所以表达式也就无从展开，以原样传给ls。

\end{example}

\subsection{shell变量和环境变量}
环境变量和shell变量都是字符串变量。使用的时候都通过 \$变量名 引用。shell会使用变量具体的值替换掉引用对象。
环境变量和shell变量在使用的时候没有区别。唯一的区别是环境变量会被子进程继承，而一般的shell变量只限在该shell中使用。
变量的赋值使用=号，不能两边有空格。

\begin{example}{变量展示}

下面的命令，将字符串“linuxbook”赋值给变量mybook,然后通过 \$mybook引用。shell会将\$mybook替换为“linuxbook”。

\begin{code}
mybook="linuxbook"
xelatex \$mybook.tex
xpdf	\$mybook.pdf
\end{code}
\end{example}

可是如果我已经有一个变量叫 my ， 则 my 和 mybook 两个变量会发生混淆：我实际想使用的是  \$my然后在后面接上book。但是shell会按照mybook变量来替换。

 取消这种混淆的办法是使用 \$\{变量名\} 进行引用。

\begin{example}{变量展示 -ver 2}

下面的命令，将字符串“linuxbook”赋值给变量texfile,然后通过 \$texfile引用。
将字符串"xelatex"赋值给变量tex,然后用 \$\{tex\}引用。

\begin{code}
\$~tex="xelatex" \\
\$~texfile="linuxbook" \\
\$~\$\{tex\} \$\{tex\}book.tex \\
\$~\$tex	\$texbook.tex
\end{code}


这里，“\$\{tex\} \$\{tex\}book.tex”时间上被替换后是执行的“xelatex xelatexbook.tex”。

可以把变量理解为一个宏，shell处理的时候会将\$和变量名本身 替换成变量的内容。

\end{example}

shell的变量不需要声明，给一个新变量赋值就是声明了。
一个变量使用 export 导出后就变成环境变量了，会被子进程继承。

比如这样设置环境变量USE给emerge：

\begin{code}
\$~ export USE=go  \\
\$~ emerge gcc
\end{code}

后续启动的任何子进程都将继承到USE环境变量。如果只想让本次的emerge继承到这个变量，可以用下面的形式。

\begin{code}
\$~ USE=go emerge gcc
\end{code}

这种情况下，USE会变成子进程 emerge 的环境变量，emerge结束后，USE变量就撤销了。

有关变量的另一个注意事项是，变量是在字符串展开之前替换的，也就是说，变量展开后，会继续对展开后的字符串再次进行字符串展开功能。使用单引号会禁止变量替换。

\begin{example}{字符串展开和shell变量}

\begin{code}
\$~ anyfile='*' \\
\$~ echo '\$anyfile'\\
\$~ echo "\$anyfile"\\
\$~ ls "\$anyfile"\\
\$~ echo \$anyfile\\
\$~ ls -d \$anyfile\\
\$~ ls \$anyfile\\
\end{code}

解释：
\begin{enumerate}
\item echo '\$anyfile' 变量没进行替换，所以显示结果为 \$anyfile

\item  echo "\$anyfile" , 变量被替换为 * , 所以显示的结果就是 * 
\item ls "\$anyfile" , 由上一条命令证实，变量被替换为 * ,显示结果却是 “ls: 无法访问*: 没有那个文件或目录”，表明引号禁止了字符串展开。
\item echo \$anyfile ，变量首先被替换为 * , 然后执行字符串展开，所以显示了当前文件夹的内容。
\item ls -d \$anyfile ， 变量首先被替换为 * , 然后执行字符串展开，所以执行结果和 ls -d * 一致。所以显示了当前文件夹的内容。不显示子文件夹内容，由 -d 参数保证\footnote{参考ls(1)}。
\item ls \$anyfile ， 变量首先被替换为 * , 然后执行字符串展开，所以执行结果和 ls * 一致。显示当前文件夹的文件和当前文件夹的子文件夹的文件。
\end{enumerate}


\end{example}



\subsection{shell和文件管理}

在任何系统下，用户都免不了要和“文件”打交到。印象中不需要和文件打交道的操作系统除了“Jarvis\footnote{Tony Stock的AI管家。也是钢铁侠衣服的辅助系统}”就不存在了。
以至于，有个操作系统认为操作系统的任务就是文件管理，把自己称之为磁盘操作系统（Disk Operating System）。

文件管理包括文件和文件夹的创建、重命名、删除和移动。多数情况下文件夹和文件是使用同样的命令进行操作，除了创建。

创建一个文件夹使用 mkdir 命令。

\subsubsection*{mkdir}

mkdir 是 make directiroy 的缩写。DOS/Windows 下的缩写为MD。

mkdir 接受一个或多个目录名作为参数。如果目录已经存在，mkdir会失败。加 -p 参数可以忽略已经存在的文件夹。

如果要想一次创建多个层次的目录，可以加 -p 参数。如 \texttt{mkdir -p a/b/c/d/e/f/g}。
文件夹创建后的访问模式默认为 0775。

mkdir的具体用法请参考mkdir(1)。


~ \newline

相比之下，创建文件的方法就多了去了。通常用户文件是由编辑器\footnote{不要光理解为文本编辑器啊，绘图软件也是图像文件的编辑器不是么。}创建。
另外创建空白文件的办法是使用touch命令。

\subsubsection*{touch}

touch 可不光是用来创建文件，还可以修改文件的最后修改时间——如果文件已经存在的话。

touch 的用法很简单，以文件名作为参数即可新建空白文件或者更新文件的最后修改时间。如果指定了 -d 或者 -t 还可以使用指定的时间作为文件的最后修改时间，默认使用当前时间。
-d 参数使用的格式同 date 命令的输出， -t 参数使用的格式为 [[CC]YY]MMDDhhmm[.ss]\footnote{[~] 表示可选。}。 
如:

\begin{code}
\noindent touch linuxbook.tex -d "2012-10-29 13:37:20"\\
\indent \quad 或者 \\
\noindent touch linuxbook.tex -t "201210291337.20"
\end{code}
touch的具体用法请参考touch(1)。


重命名文件和移动文件（夹）是同样的操作。使用的都是 mv 命令。

mv乃move的缩写，命令很简单：\texttt{mv OLDNAME NEWNAME}。具体用法请参考mv(1)。

复制的英文单词是copy，考虑到UNIX的神缩写传统，你一定不指望UNIX会使用copy作为复制文件的命令名称。没错，复制文件的命令名称也被缩写了，非常精简的cp。

直接使用cp的话并不能进行文件夹的复制。文件夹有子文件夹，子又有子，子又有孙，子子孙孙无穷尽。要拷贝一个文件夹，必然是一个递归的操作。所以拷贝文件夹需要添加 -R 参数。
至于符号链接\footnote{见 \faqref{faq:symlink}}，一般拷贝文件夹希望拷贝的是符号链接本身而不是符号链接指向的文件，所以加 -d 参数表示保留符号链接。这样拷贝文件夹通常使用法为“cp -dR 源文件夹 目标文件夹\footnote{-d -R 可以使用 -dR 这样的形式。对多数GNU软件适用。}”, 可以用 -a 代替 -dR 。

拷贝文件的时候提到了递归，提一下，删除文件也得递归。因为删除目录的系统调用rmdir(2)必须用在空目录上。所以对于一个非空的目录是不能删除的。
要删除一个非空文件夹，首先删除其下的所有文件和文件夹。对于非空的子文件夹，先删除子文件夹下的文件和文件夹。这就是一个递归操作了。
删除文件和文件夹使用rm(1)命令，rm默认不进行递归，所以对非空文件夹执行rm会失败。使用-r参数能让rm执行递归操作。


\subsection{管道和重定向}

UNIX管道在进程间架起了看不见的桥梁，联通了二者之间的输入和输出。而使用进程间的管道简单到只需要用“|”隔开两条命令。前一条命令的输出自动的就成为后一条命令的输入。

\begin{example}{统计portage拥有的软件包数目 - ver 1}
%ls /usr/portage/*-*/* -d | while read f ; do [ -d $f ] &&  ls -d $f  ; done | wc -l
\begin{code}
\$ls /usr/portage/*-*/* -d -1 | wc -l
\end{code}

通过前面的portage介绍可以知道，portage树的结构是  “分类名/包名/包名-版本.ebuild” 这样的2层目录。而“分类名”也是 “大类别-小类别” 这样的结构。
所以 “*-*”就可以匹配到类名而将 profiles eclass  licenses metadata scripts 这样的目录和 “header.txt skel.ChangeLog” 这样的文件排除掉了。连 virtual 目录都排除了。这样获得的文件夹的数目就是portage包含的包的数目。

等等，好像分类目录下，并不只是软件包名命名的目录。有的分类目录下有一个顶层的metadata.xml。描述一个分类里所有软件的公共信息。比如 \url{/usr/portage/games-fps/metadata.xml}。 这条命令似乎会把这样的 metadata.xml 算成一个目录而进入统计。
\end{example}


\begin{example}{统计portage拥有的软件包数目 - ver 2}

让我们进行改进，将第一次 ls 输出的文件名中，把非目录过滤掉：

\begin{code}
\$ls /usr/portage/*-*/* -d -1 | while read filename ;  \textbackslash \\
	\qquad do [ -d \$filename ] \&\&  ls -d -1 \$filename  ; done | wc -l
\end{code}

通过中间引入一层过滤层，把非目录过滤，这样 wc 获得的行数才是真正的包的数目。while这样的用法是shell循环语句，在下一个小节会详细介绍。

\end{example}

通过这2个例子相信已经能体会到管道的魅力了。

其实一个程序的输出何止是可以作为另一个程序的输入，本身也可以直接打印到文件中的嘛。在UNIX“一切都是文件”的哲学下，控制台对程序来说也是文件。既然都是文件，那么程序的输出除了可以输出到控制台或者通过管道输出给其他程序，自然也可以直接输出到文件中去。

不需要对程序进行任何修改，只要在shell里使用简单的技巧就能做到了。越是强大的功能用起来越是简单直接。重定向输出居然只是简单的使用 “>”就能做到了。

\begin{code}
\$ls /usr/portage/*-*/* -d -1 > portage\_ls
\end{code}

打开 portage\_ls 文件，看是不是输出被直接写入这个文件了？
多执行几次，你会发现文件并没有变大。因为shell在使用重定向的时候，会将原来的文件内容清空。使用 “>>” 即告诉shell使用追加模式。不清空原有文件内容，向文件末端追加。

既然输出可以重定向，那么输入也可以。否则怎么体现UNIX的灵活强大呢！

\begin{code}
\$ ( while read filename ; \textbackslash \\
	\qquad do [ -d \$filename ] \&\&  ls -d -1 \$filename  ; \textbackslash \\
	\qquad done ) <  portage\_ls > portage\_lsdironly
\end{code}

打开 portage\_lsdironly 文件，非目录是不是已经被过滤了？

最后一条命令：

\begin{code}
wc -l portage\_lsdironly
\end{code}

就获得了统计结果了。

这就不使用管道，完全使用重定向完成了同样的功能。不过多了一些中间文件，对吧？


\subsection{bash是门编程语言}

Bash作为一个shell是尽职的。有了各种实用小工具\footnote{已经在system集合里，默认安装，自行安装使用\#emerge coreutils）} 和管道的配合，shell能发挥出非常强大的作用。

不过，要想完全发挥shell的作用，就需要将shell作为一门编程语言，使用shell提供的判断和循环语句。

\subsubsection{条件判断}

shell之条件判断分两种，一种基于 if-then-else-fi ， 另一种是 cmd1 \&\& cmd2 和 cmd1 || cmd2。

我们首先介绍 cmd1 \&\& cmd2 这种形式的命令。

写过 "hello world" 程序的人都知道，main() 函数的最后都要有一个 “return 0;”。 这个返回值到底放回给了谁？答案是调用它的父进程。如果是shell调用的，就是返回给shell咯。
执行过一条命令后，通过 \$? 变量就能获得刚刚执行过的命令的返回值。

做个测试吧：

\begin{code}
\$ ls \\
\$ echo \$?
\$ ls /root \\
\$ echo \$? \\
\end{code}

第一次调用 ls 显示当前目录的内容，返回值很明显是 0 。第二次显示 /root ，很明显没有权限，返回值是 2 。

shell有时候必须知道一条命令执行成功还是执行失败，shell就是通过查看命令退出的时候的返回值判断命令是否成功执行的。像“ls /root”这样的命令如果不是以root权限执行的话，一定会因为没有权限而失败。 shell规定了，一条命令执行后返回0表示成功，非0表示失败。

这个成功和失败就用在了 cmd1 \&\& cmd2 和 cmd1 || cmd2 语法上。
简单的说，“cmd1 \&\& cmd2” 表示cmd1命令执行成功就执行cmd2。而“cmd1 || cmd2”表示cmd1命令执行失败就执行cmd2。

拿先前的例子来说，我在while循环里使用了一个判断 
\begin{code}
  [ -d \$filename ] \&\&  ls -d -1 \$filename
\end{code}

意思是，如果“ [ -d \$filename ]” 成功，执行“ls -d -1 \$filename”。

在shell中， “ [~] ” 的意思就是执行条件判断。
如“\texttt{ [ "a" = "b" ] }” 就表示判断字符串 "a" 和 "b" 是否相等。
 [] 对于表达式成立会退出码0,不成立为1。再依据shell对退出码所采取的行为，就能使用 [] 编写表达式是否成立的条件而执行的命令了。
 
在本利中， [ -d \$filename ] 的意思是判断 \$filename 是否为存在且为目录。
事实上。 [] 是 test 的别名。“test -d \$filename” 等价于 “ [ -d \$filename ]” 。 所以想知道 [] 的详细用法可以查看man手册 test(1) 。

\subsubsection{while循环}

例子中还有一个用到的shell功能就是循环。 bash shell支持while循环和for循环。

while循环的语法为 
\begin{code}

while cmd of condition ; do \\
	cmd1\\
	cmd2\\
	...\\
done\\

\end{code}

命令可以和 do 写在一行上也可以不写在一行上。
do如果不和while写在同一行，while后面的那个命令的“ ;” 就可以不写。

如果“cmd of condition”执行成功就执行 do 和 done 之间的命令。重复执行直到命令执行失败。也就是返回值不为0。


\subsubsection{until循环}

until循环和while循环语法一致，不过是用until替换while。


\begin{code}

until cmd of condition ; do \\
	cmd1\\
	cmd2\\
	...\\
done\\

\end{code}

不一样的是，until循环是条件不成立才执行的循环。而while循环是条件成立时执行。


\subsubsection{for循环}

while循环用户知道循环条件的地方，for循环用于知道循环的内容的地方。

for的语法为：

\begin{code}

for var in array ; do \\
cmd1 \\
cmd2 \\
... \\
done 

\end{code}


array 是一个字符串的集合。比如 “/usr/portage/*-*/*” 会导致 shell 将其扩展为数万个文件名的集合。for为集合中的每个元素运行一次循环。
每次循环，都可以用“\$var”引用集合中的一个元素。


例中命令使用 for 循环可以修改成这样，于是出现版本3：

\begin{example}{统计portage拥有的软件包数目 - ver 3}

\begin{code}
for f in /usr/portage/*-*/* ; do  \textbackslash \\
	\qquad [ -d \$f ] \&\&  ls -d \$f ; done | wc -l
\end{code}
\end{example}

除了用通配符制造的文件名，for 循环还可以使用数组。数组将在下面的章节中节介绍。


\subsubsection{字符串}

字符串是shell的基本数据类型。有多种方式可以进行字符串的操作。所以到现在才介绍，是因为日常使用很少需要用到字符串操作。

\begin{enumerate}
\item 合并字符串	

合并两个字符串的办法很简单，就是将连个变量“贴”到一起，如变量 A 和 B ， 获得 A 和 B 连接在一起的字符串的办法就是 {\tt "\$A\$B"}。

\item 字符串替换

\$\{var/查找字符/替换字符\}

	这个操作中除了第一个参数是变量外其它两个都是字符;还有一点就是这个操作并不是把变量“var”中的字符替换了，而是返回一个替换后的字符串。原字符串并没有被替换。
这个查找是正向查找，就是从左到右执行查找，把找到的第一个字符串替换到。执行反向查找使用的是这个语句：


\$\{var/\%查找字符/替换字符\}

上面两个查找替换都是找到第一个匹配的字符串进行替换，而下面这个可以执行全部替换。

\$\{var//查找字符/替换字符\}

以上所有的操作都不会修改原始变量var而是返回一个新的字符串。如果要修改原始变量可以把新的字符串重新赋值给var嘛！

\item 取子字符串

从一个字符串变量获取子字符串的语法

\$\{var:位置\}

从“位置”开始取子串到最后。默认是从左边开始，如果“位置”为负数，则是从右边的第“位置”个字符开始，并且第一个位置为0。 %
如果不取到结尾而是限定一个长度，可以使用下面的语法：
%例：
%str1=abcABCabc123ABC
%echo ${str1:(-3)}#会输出ABC

\$\{var:开始位置:结束位置\}


%\$\{字串\#匹配字串\}

%（说明一下，这个是从左边第一个开始匹配，剥去最短“匹配字串”） %
%例：
%str1=abcABCabc123
%echo ${str1#a*c}#输出ABCabc123
%(2)${字串##匹配字串}
%（说明一下，这个是从左边第一个开始匹配，剥去最长“匹配字串”）
%str1=abcABCabc123
%echo ${str1#a*c}#输出123
%echo ${str1#b*c}#输出abcABCabc123，因为没有从第一个开始匹配
%(3)${字串%匹配字串}
%(4)${字串%%匹配字串}
%（说明一下，这与上面的（1）（2）是正好相反的，是从最后一个开始匹配的）

\item  字符串长度

\$\{\#字串\} %

例： str=abcdefg

echo \$\{\#str\} 输出7

\end{enumerate}






\subsubsection{数组}

shell支持字符串数组。 %不过好像日常使用没什么大的用处。在编写shell程序的时候会用到。所以还是有必要知道一下。
数组的语法有三种

\begin{code}

(1) array=(var1 var2 var3 ... varN) \\
(2) array=([0]=var1 [1]=var2 [2]=var3 ... [n]=varN) \\
(3) array[0]=var1 \\
\leftskip=2em
    arrya[1]=var2 
    
    ... 
    
    array[n]=varN 
\end{code}

计算数组元素个数：
\begin{code}
\$\{array[@]\}  或者  \$\{array[*]\}
\end{code}

引用数组的语法是 （n为下标）：
\begin{code}
 \$\{array[n]\}
\end{code}

可以使用 for 循环来遍历数组。

\begin{example}{统计portage拥有的软件包数目 - ver 4}

\begin{code}
dirlist=/usr/portage/*-*/*

for f in \$dirlis  ; do  \textbackslash \\
	\qquad [ -d \$f ] \&\&  ls -d \$f ; done | wc -l
\end{code}

顺便一提，其实shell进行字符串扩展的时候生成的就是数组，数组再就地展开为命令参数。所以{}“dirlist=/usr/portage/*-*/*”{}这样的用法即让dirlist变量变成数组。


\end{example}

\subsubsection{读取输入}
在 \textbf{统计portage拥有的软件包数目} 例子中看到过 \texttt{while read f} 这样的用法。读者猜想这是读取输入到变量f的意思。
没错，正是这个意思。
准确的说， read 的与否是这样的：

\begin{code}
read var1 var2 var3 ...
\end{code}

read 从标准输入读取一行，并将一行中的字符串分别赋值给 var1 var2 var3 等等变量。字符串的拆分规则和命令行一样，使用空格分隔。如果拆分出来的字符串个数比命令行提供的变量多，最后一个变量将包含余下的所有字符串。

如果读取输入的时候遇到文件结束（在例子中，管道的输出方结束输出），退出状态就是非0，导致例子中while循环终止。

要读取文件可以使用重定向，比如
\begin{code}

读一行

read line < sometextfile.txt

或者读整个文件

while read line ; do

...

done < somefile.txt
\end{code}


\section{文件编辑器}\label{sec:editors}

日常使用Linux的过程中，无法避免的一件事就是要编辑文本。可以说，计算机除了“计算”之外的一大用途就是“编辑\footnote{不论是编辑文本还是编辑视频，通通都是编辑的说。}”。
Linux世界有好多非常优秀的文本编辑器。由于现在还未安装X服务，所以只能使用命令行模式的编辑器。在开源界，有两大编辑器派系。一个是vi和之后的改进版vim，被称为编辑器之神；
另一个是被称作神之编辑器的emacs。


\subsection{编辑器之神vim}

vim是vi的改进，全称是\textbf{Vi IMproved}。	%
%
\newcommand{\keyboradparty}[1]{
#1是最好的编辑器，没有之一。


对于键盘党\footnote{对惯用键盘操作电脑，提倡抛弃鼠标的人的称呼。}来说，让双手离开键盘是一种罪恶。
没有比让双手保持在键盘上更能发挥打字的速度了。不仅仅如此，#1 还罪恶的提出，双手不仅不能离开键盘，甚至连主键盘区都不能离开。
对光标位置的操纵，明明有直观的四个箭头按键， #1 还是采用的定义主键盘区的快捷方式来控制。双手不离开主键盘区是  #1 的设计宗旨之一。

#1 被初学者说成复杂的原因大概也缘于此吧！ 毕竟使用鼠标直接定位是最直观方便的操作方式了。再不济用键盘上的方向键也顶用多了。好在 #1 仍然支持使用方向键，只不过不推荐用户这么做。

“除了#1，其他编辑器神码的都弱爆炸了”，资深#1用户如是说，“这个世界上只有两种人，一种人使用#1,另一种使用其他。” 这大概就是{}#1{}在用户心目中的评价吧。

先安装#1

\begin{code}
\# emerge #1 -av \footnote{在安装任何软件之前，使用“-av”这样的参数让emerge先打印将要进行的操作，然后回答yes/no继续安装永远都是一个好习惯。}
\end{code}

}
%
\keyboradparty{vim}

%\marginpar{
可能依赖很多，不过可以去掉{}X{}，这样vim就不会引入X相关的依赖。
%}
%可能依赖很多，不过可以去掉X，这样vim就不会引入X相关的依赖。

\begin{code}
\# USE=-X emerge vim -av
\end{code}

耐心的等待片刻，新鲜出炉的vim就准备好了。

vim为了支持前面提到的“纯键盘”操作，把编辑器设定成两种工作模式：命令模式和输入模式。在命令模式下，按键盘的结果是每个按键都被解释成相应的命令。
输入模式下，按下键盘按钮会在当前输入位置（光标所处位置）输入对应的字符。\footnote{这里还是要废话一下，以免误导大家。只有能显示的字符才能输入。比如即便是输入模式下，按{\tt "Del"}键的结果也是删除字符而不是输入“Del”字符。因为“Del”并不是能显示的字符。还有方向键在任何模式下都是移动光标位置。能输入的，只有英文字母大小写A-Z和数字0-9以及0-9上方的字符，还有标点符号。}

命令模式下按“i”进入输入模式。

任何时候按“Esc”键都可以退出输入模式转到命令模式下。

在命令模式下，先输入“:”，会开启命令行模式，终端最下一行出现“:”提示并开始回显输入的命令。命令行模式下，输入完整的一条命令然后按回车，vim就会执行该命令。
比如保存并关闭文件的命令为 “:wq”。意思就是先退出到命令模式，然后输入“：”进入命令行模式，接着依次输入‘w’、‘q’并按回车。

在命令行模式下，按N次某命令，和先用数字输入次数然后再按下某个命令字符，效果是一样的。
比如：按10次‘j’向下移动10行，和依次按‘1’、‘0’、‘j’的效果是一样的\footnote{马上下文就讲到光标移动了}。

特別提示︰在浏览本教程時，不要強行记忆。记住一点︰在使用中学习。\footnote{如果用不到，就不用学了。:）}

\subsubsection{打开关闭和保存文件}


打开文件的办法有两个，一是在命令行上给出文件名，就像这样

\begin{code}
\$ vim 文件名
\end{code}

二是先打开vim然后使用命令模式 “:o 文件名” 打开文件。

保存文件使用命令“:w”

退出使用命令：“:q”。如果文件已经修改，vim会提示文件被修改，不能退出。使用“:q!”命令忽略修改强制退出，或者使用“:wq”保持并退出。



\subsubsection{移动光标}

除了使用键盘上的方向键，还可以使用 ‘h’、‘j’、‘k’、‘l’
这4个键代替。

{
\vspace{-3ex}
\begin{wrapfigure}{l}{0.4\textwidth}

 \tt \mbox{

\begin{tabular}{ccccccc}
 &  &  & \UParrow  &  &  &  \\ 
 &  &  & k &  &  &  \\ 
 &  &  &  &  &  &  \\ 
\LEFTarrow & h &  &  &  & l & \RIGHTarrow \\ 
 &  &  &  &  &  &  \\ 
 &  &  & j &  &  &  \\ 
 &  &  & \DOWNarrow &  &  & 
\end{tabular} 
}
\end{wrapfigure}{
%
~ \\
~ \\
\noindent
\leftskip=-2em
提示︰ \\
\leftskip=2em
		\hbox to 1em {h} 每次按下就会向左移动。\\
		\hbox to 1em {l} 每次按下就会向右移动。 \\
		\hbox to 1em {j} 每次按下就会向下移动。 \\
		\hbox to 1em {k} 每次按下就会向上移动。 \\
\vskip 3ex
}

}


\subsubsection{文本的编辑操作}

\begin{enumerate}
\item 删除:

	输入模式下，使用退格键和删除键删除光标前后的字符；命令模式下，使用x删除光标处的字符。“dd”（按两下d）删除光标所在行。
	“dw”（依次按下d和c键，下同）删除光标处的一个英文单词。“de”删除到行尾。

\item 插入：

	将光标定位到要插入的地方，按‘i’进入插入模式。按‘Esc’退出。\\
	将光标定位到要插入的地方，按‘a’进入光标后移动一个字符，插入模式。按‘Esc’退出。\\
	将光标定位到要插入的行，按‘A’就自动定位到行尾巴，并进入插入模式。按‘Esc’退出。

\item 区域选择：

	选择一个区块通常是“复制”和“剪切”操作的前奏。
	命令模式下，按‘v’就进入区块选择模式，并定义区块的起始位置。移动光标选择终止位置。被选中的文本会反白显示。

\item 复制：

	选择好区块后，按‘y’。选中的文本会被拷贝到vim自己的剪切板。

\item 剪切：

	选择好区块后，按‘d’。选中的文本会被剪切到vim自己的剪切板。
	
\item 粘贴：

	在命令模式下，按 p 就可以把vim自己的剪切板的内容粘贴到光标位置。vim不是X程序，不能使用X界面的剪切板哦。

\item 撤销：

	按‘u’撤消最後執行的命令，按‘U’來修正整行。

\end{enumerate}


\subsubsection{从vim那获得帮助}

由于本书定位是“入门”，所以高级的vim话题就不探讨了，有兴趣的同学可以参阅其他书籍。通过本小节的阅读，使用vim编辑配置文件之类的工作应该能胜任了。

惊愕的是，vim居然自带了教程！使用



\begin{code}
\$ vimtutor
\end{code}

就能打开这个教程了。默认会打开本地语言版本的教程！如果你的本地语言设置不是中文\footnote{到\secref{sec:lang}脑补。}，使用 

\begin{code}
\$ vimtutor zh
\end{code}

即可打开。虽然是繁体中文，不过阅读上并无大碍。



\subsection{神之编辑器emacs}

emacs是GNU出品的正牌编辑器，出自RMS大神之手。故称之为神之编辑器。
%
\keyboradparty{emacs}


可能依赖很多，不过可以去掉X和gtk，这样vim就不会引入X相关的依赖。

\marginnote{可能依赖很多，不过可以去掉{}X{}和{}gtk{}，这样vim就不会引入X相关的依赖。}

\begin{code}
\# USE="-X -gtk -gtk3" emerge emacs -av
\end{code}

耐心的等待片刻，新鲜出炉的emacs就准备好了。

和vim不一样，emacs并不使用命令模式和输入模式这样的区分\footnote{因为这意味着在码字的过程中要不停的按‘Esc’和‘i’切换模式。}。
emacs下所有的命令都是立即可用的，不会要求你时刻记住当前工作模式并随时准备切换。


%TODO
\todo{写更多的 emacs 入门}


~


\section{系统管理基础}

如果一个系统能自发的自我维护就好了。不过这种系统目前只存在于Tony Stock的家里。想好好的用Linux干活，一些基本系统的维护知识还是必要的\footnote{汗！用惯了Windows的小白表示从来不进行系统维护。定期重装或者安装个还原系统。一般品牌笔记本不都会带个系统恢复功能么？}。
它可以在让你永远不会遇到系统瘫痪而必须重装的问题，或者就算瘫痪了都能起死回生。


\subsection{进程管理}
\subsection{系统服务管理}

%\section{shell 进阶应用}
\fi



\ifdefined\compileall

\section{软件的分发}

除了写给自己用的软件，大部分软件都需要向最终用户进行分发。软件的分发不仅仅是一次下载，一份拷贝。如何发行软件也是很考究的活。
安装发布内容的形式可以分为二进制发布和源码发布。

闭源软件没得说，只有二进制发布。开源软件通常使用源码发布，并为流行的发行版提供预编译的二进制发布。

商业软件使用二进制发布，通常还会带上一个安装程序。用户只要执行这个安装程序就可以自动的将软件安装到指定位置，调节好系统配置，建立快捷方式等等操作。
安装程序的傻瓜方式就是 “下一步”。使用Windows的人，不知不觉的养成了一路Next的习惯。有恶意程序就专门针对一路Next的人的习惯偷偷的给安装上一些广告插件。

Linux下的软件通常以源码形式发布。源码常常以tar打包并用gz或者bzip2压缩，扩展名 .tar.gz 或者 .tar.bz2。发布的文件名通常为 “软件名-版本.tar.bz2”。
用户将源码下载后，就进入了编译阶段。软件会在压缩包内带一份文档告诉用户如何编译，或者在官方网站上提供编译指导。

\section{软件的编译}
源码发布，不仅仅是把源代码打包发布完事，还需要带上一套构建脚本。但凡一个真正工作的现实软件，都不会只有一个 .c 文件完事。有的大型软件包含上万个文件。
如此众多的文件，一个一个手工调用gcc命令去编译显然不人性化，也非常烦琐。所以但凡一个软件项目，都有一套编译系统。

编译系统不仅仅用于从源码从头构建可执行文件，还包括追踪文件的修改，再次编译的时候只重新编译改动的部分。这样可以极大的加快开发的进程。因为开发者做一个小修改不需要重新编译所有的文件。
基于这样的目的，UNIX界发明了Makefile。

\subsection{Makefile}

Makefile是一个基于时间戳的依赖追踪器。Makefile 定义多个目标和依赖。并在目标的时间戳比依赖旧的情况下调用定义的命令行更新目标。
比如下面的一个Makefile文件：

\begin{code}
helloworld: helloworld.c\\
\qquad gcc helloworld.c -o helloworld
\end{code}

意思是当helloworld的最后修改时间比helloworld.c旧，或者helloworld文件不存在的时候，调用`gcc helloworld.c -o helloworld'命令更新helloworld文件。

\begin{notice}
但是直接执行 make , 则不会更新 helloworld, 必须执行 make helloworld。 因为 helloworld 不是默认更新目标。直接执行 make 只会更新默认的目标。默认更新目标是 all， 所以需要这样写

\begin{code}
all: helloworld
\\
\qquad
\\

helloworld: helloworld.c\\
\qquad gcc helloworld.c -o helloworld
\end{code}

all依赖helloworld，所以更新 all 前要更新 helloworld， 这样直接执行 make 即可更新helloworld了。

\end{notice}

通常发布的Makefile还定义一个特殊的目标,  install： 直接执行 make install 就可以把编译好的可执行文件拷贝到系统的 bin 目录下。

一旦项目变大，手写的Makefile一起膨胀。逐个的追踪文件的变化，并手工更新相互间的依赖关系变得非常的头疼。所以，最好的办法是有个工具能自动的生成Makefile。
这样很多软件都可以有统一的make行为了，比如支持 make test,  make install, make uninstall, make dist-bzip2 自动打包为源码发布包，等等。

利用更简单的规则生成符合标准的Makefile的工具有很多，最常用的莫过于autotools了，这个是软件发布的事实标准。除了autotools，还有scons, cmake, qmake 等。
这里着重介绍常用的autotools和cmake的使用。

因为本书并不是开发手册，所以并不打算教大家如何编写Makefile, 如何编写autotools的规则，如何编写CMakeLists.txt。需要知道这些详细规则的同学可以参考各自工具的手册。

\subsection{autotools}

autotools 并不是一个工具，而是一组工具的集合。autotools 在开发阶段需要 m4, automake , autoconf 三个工具的配合，最终生成的是一个名为 configre 的 shell 脚本和一系列的Makefile.in文件。生成的configure 脚本和Makefile.in被打包到最终发布的源码中。因此\textbf{用户并不需要安装autotools工具集\footnote{这是autotools非常重要的优点，最终发布的源码包只依赖shell脚本和make工具。并不依赖autotools本身。}}。用户需要执行configure脚本从Makefile.in生成Makefile。然后使用make编译即可。

典型的使用 autotools 编译工具的源码包的编译流程为：

\begin{code}
\#tar -xvf somepkg-ver.tar.bz2

\#cd somepkg-ver

\#./configure --prefix=/usr % [--prefix=/installpath --enable-XXX --disable-XXX --with-XXX --without-XXX] 

\#make

\#sudo make install
\end{code}

configure、make、make install 又被称为软件编译三部曲。括弧笑。

代码的编译期配置都发生在调用 ./configure 脚本的时候。所有的 ./configure 脚本都支持 --prefix 指定安装路径。通常将程序安装到 /usr 下。
autotools 的规则是这样的

\begin{enumerate}
\item 可执行文件安装到 \$\{prefix\}/bin 下
\item 库文件安装到 \$\{prefix\}/lib \footnote{如果指定了 --libdir=XXX , 则是安装到 \$\{libdir\} 下 } 下
\item 头文件安装到 \$\{prefix\}/include 下
\item 数据文件安装到 \$\{prefix\}/share 下
\end{enumerate}

所以指定了 --prefix 就把其他相对路径都确定下来了。如果不指定 prefix , 默认则是 /usr/local。

autotools 天生支持交叉编译。要交叉编译，只需要使用 --host=XXX 参数即可。
GNU为一个平台定义了一个叫 host triplets 的东西，一个平台使用 cpuarch-platform-kernel 这样的三组来唯一确定。比如64位PC上跑的linux系统就是 x86\_64-pc-linux-gnu。 x86\_64 是 CPU 架构， pc-linux 指的是操作系统内核\footnote{一些发行版使用 unknow-linux}，gnu指的是libc的牌子，表示使用glibc\footnote{嵌入式系统通常是uclibc}。这样处理这样一个平台的编译工具都使用 x86\_64-pc-linux-gnu- 前缀，比如 x86\_64-pc-linux-gnu-gcc , x86\_64-pc-linux-gnu-g++,  x86\_64-pc-linux-gnu-ld , x86\_64-pc-linux-gnu-ar . 如果指定了 --host=x86\_64-pc-linux-gnu-,  则 configure 生成的 Makefile 了默认就把各种编译工具设置为 x86\_64-pc-linux-gnu- 前缀的版本。


\subsection{cmake}\label{section:cmake}

autotools 对最终用户非常友好，可是开发者需要做大量的工作。显然搞明白 m4 automake autoconf (甚至还有 libtool) 相互之间的复杂关系非常令人头痛。所以 KDE 果断迁移到了 cmake 构建工具。连 KDevelop 这种重量级的IDE都放弃了原生autotools支持，转而支持cmake。

cmake 使用一个 CMakeLists.txt 文件来记录项目的编译配置。开发者通过编写 CMakeLists.txt 及其简单的规则， 然后由 cmake 处理成 Makefile。 cmake 除了能将 CMakeLists.txt 转化为 Makefile， 还可以转化为 Visual Studio 的工程文件，生成 Eclipse 的工程文件，CodeBlocks的项目文件。是个对 IDE 比较友好的 工具。

cmake 的通常用法是 

\begin{code}
\#tar -xvf somepkg-ver.tar.bz2

\#mkdir build \&\& cd build

\#cmake ../somepkg-ver -DCMAKE\_INSTALL\_PREFIX=/usr

\#make

\#sudo make install
\end{code}

cmake 推荐的做法是使用 out-of-tree build， 也就是不要在源码目录编译，而是新建一个目录用于编译。例子中 cmake 的第一个参数是源码路径。
cmake 在当前目录生成 Makefile，所以需要 cd 到 新建的 build 目录。

cmake 对交叉编译支持并不佳。通常仅用于执行本地编译。

不像 autotools 只需要开发者安装 autotools , cmake 需要用户也安装有 cmake ， 所以不太适合那种需要尽量少的外部依赖的软件。一些人有洁癖，多一个依赖都要吵吵闹闹的。所以才会有 Gentoo 这种系统诞生的嘛（笑）。

\section{发行版的包管理}\label{section:pkgmgr}

从源码编译软件，怎么说呢， 好是好啊。问题是麻烦。虽然有 autotools 这种傻瓜式的编译工具，可是依赖问题呢？一个一个的到源网站下载，然后又一个一个编译。如果编译失败还要仔细检查是否因为依赖没有安装。
这可是非常麻烦的事情，如果有工具能自动化这些事情，自然就再好不过了。这种自动化的工具就是包管理器。

二进制发布除了像Windows上的程序那样使用一个安装程序，还能以发行版的包管理器使用的二进制包进行发布。二进制包管理最有名的就是 rpm 和 deb 了。所谓rpm系和deb系就是指使用 rpm 包管理器的一系列发行版和使用 deb 包管理器的一系列发行版。也有重复再发明一次轮子的发行版，使用自己发明的包管理器。比如 archlinux。

基本上，二进制包就是将软件编译，打包，压缩，然后配上一些相关信息，比如包名啊，依赖关系啊，等等。包管理器检查一下包的信息，然后把压缩的二进制文件解压到系统里，到数据库里记录一下就完成软件安装了。
因为包里包含了依赖关系，因此包管理器得以自动处理依赖关系：自动下载安装缺失的依赖。

因为包管理器的存在，系统目录下的所有文件都由包管理器严格追踪。任何一个文件都有所属的包。

\subsection{软件仓库}

为了包管理器能自动的安装依赖，需要提供一个软件仓库。正是软件仓库和包管理器定义了一个发行版。

软件仓库是一个伟大的发明。在软件仓库发明之前，操作系统和应用程序独立更新；每一个应用程序都要自带一个自动更新机制——每个程序都要开个后台进程\footnote{个别缺德程序还注册成系统服务，禁止用户中止。}检查更新。

而软件仓库发明后，操作系统和应用程序之间的界面就模糊了。操作系统和应用程序使用同一套机制进行管理和更新。想要安装什么软件只需要打开统一的仓库搜索安装即可。不需要到茫茫网络上去搜索下载，说不定一不小心就安装上了恶意程序或者流氓软件。

不管是应用程序还是操作系统，都统一在一个包管理程序的控制下被安装、更新和卸载。\footnote{如此伟大的创意最终被乔帮主借鉴，搞了个App Store。}

\subsection{portage和emerge}\label{sec:emerge}

portage是Gentoo的包管理工具。
\textit{ Portage可能是Gentoo在软件包管理方面最瞩目的创新。它所具有的高度灵活性以及其他大量的特性，使其经常被认为是Linux中最好用的软件包管理工具。\footnote{语出自Gentoo官方文档}}
emerge是portage的命令行前端。绝大部分用户通过emerge使用Portage。

/usr/portage下存放了数万个ebuild文件，这些文件按照预定的格式组织，形成了Gentoo的软件仓库，又称portage树。

ebuild 包含了Portage管理工具（emerge）维护软件（安装，搜索，查询， ...）时所需要的所有信息

本地portage树需要经常同远程仓库进行同步。使用rsync可以方便的将本地的portage树与Gentoo官方发布的portage树进行增量同步。
Gentoo为rsync提供了便捷的命令{\tt emerge --sync}。建议的同步频率为不要高于每天一次。

%在 \secref{sec:quickemergeguide} 已经进行了适当的 emerge 入门了。这里进行一些高级进阶。

\begin{enumerate}
\item 查询软件　

使用　emerge -s [软件名称] 查询软件,如果需要查找软件的描述，使用大写的 -S 参数。
\begin{example}{查找名字中包含mplayer}
\begin{code}
\#emerge -s mplayer
\end{code}
\end{example}

\item 安装软件

直接使用　emerge [软件名]　即可安装。portage会自动处理好软件间的相互依赖关系。

\begin{example}{安装mplayer}
\begin{code}
\#emerge mplayer
\end{code}
\end{example}

如果需要查看mplayer有哪些USE标记，有哪些依赖会被安装，可以加　-vp 参数。参数　-v　告诉emerge显示详细信息，而参数-p则告诉emerge执行“假装(pretent)”安装，并不真的执行安装，也就是只显示将要安装的软件包。

\begin{code}
\#emerge -vp mplayer　

[ebuild   N    ] media-video/mplayer-1.1-r1::gentoo-zh [1.1-r1::gentoo] USE="X a52 aalib alsa cdio dts dv dvd dvdnav enca encode faac faad fbcon gif iconv jpeg jpeg2k libass libcaca libmpeg2 live lzo mad md5sum mmx mng mp3 network opengl osdmenu oss png pulseaudio quicktime rar real rtc rtmp samba shm speex sse sse2 ssse3 theora toolame tremor truetype twolame unicode vdpau vorbis x264 xinerama xscreensaver xv xvid -3dnow -3dnowext (-altivec) (-aqua) -bidi -bindist -bl -bluray -bs2b -cddb -cdparanoia -cpudetection -debug -dga -directfb -doc -dvb -dxr3 -ftp -ggi -gsm -ipv6 -jack -joystick -ladspa -lirc -mmxext -nas -nut -openal -pnm -pvr -radio -sdl -tga -v4l -vaapi\% (-vidix) (-win32codecs) -xanim -xvmc -zoran" VIDEO\_CARDS="-mga -s3virge -tdfx"
\end{code}

[ebuild   N   ] 表示新安装。　可能会出现的标记有：　　R 表示重新安装，U 表示升级，UD　表示降级，NS 则表示安装到新的　SLOT\footnote{一些软件支持多版本共存。Gentoo对于这些软件的处理办法是不同的版本安装到不同的SLOT。}。

\item 卸载软件

使用　emerge -C [软件名称] 来卸载软件。如果软件SLOT安装，而你只想卸载某个版本，可以使用　emerge -C =软件名称+版本 的形式卸载特定的版本。


\begin{code}
\#emerge -C mplayer　

 * This action can remove important packages! In order to be safer, use\\
 * `emerge -pv --depclean <atom>` to check for reverse dependencies before\\
 * removing packages.\\
\\
 media-video/mplayer\\
    selected: 1.1-r1 \\
   protected: none \\
     omitted: none \\
\\
All selected packages: media-video/mplayer-1.1-r1\\
\\
>>> 'Selected' packages are slated for removal.\\
>>> 'Protected' and 'omitted' packages will not be removed.\\
\\
>>> Waiting 5 seconds before starting...\\
>>> (Control-C to abort)...\\
>>> Unmerging in:  5 4 3 2\\ 
\end{code}

不过，你需要特别注意的是：Portage将不会检查你要删除的包是否仍被其他的包依赖。但是当你要删除一个可能破坏你系统的重要的软件包时，它还是会给予警告。 

当从系统中移除一个软件包时，之前那些\textbf{为了满足其依赖关系而自动被一并安装的软件包}将会被保留。
如果想让Portage找出那些现在可以移除的相关软件包，可以使用emerge的--depclean功能。

\begin{example}{移除系统中所有孤立软件包\footnote{孤立软件包指的是没有被其他软件依赖并且不是用户显式指定安装的软件}}
\begin{code}
\#emerge --depclean 
\end{code}
\end{example}

所谓用户显式指定安装的软件，其实会被写入world集合。有关world/slot/virtual东西的详细概念稍后介绍。


\item Meta软件包

有时候，一些软件包自身什么也不安装，为的是通过依赖关系将很多很多其他的包一并安装。比如　kde-meta，　只要安装　kde-meta , 就能正确的将　KDE 的百八十个组件都安装了。就不用自己一个一个去emerge了。

如果你试图从系统中移除一个这样的软件包的集合体，只是单纯地使用emerge --unmerge命令并不能完成你的要求，原因在于这些包的依赖关系仍然保留在系统中。 

\begin{example}{移除孤立依赖的软件包}

\begin{code}
\# emerge --update --deep --newuse world\\
\# emerge --depclean\\

\# revdep-rebuild\footnote{revdep-rebuild工具由gentoolkit包提供；使用前别忘了首先emerge它哦。}\\
 或者如果用的　portage-2.2 以上的版本则用\\
\#　emerge @preserved-rebuild \\
\end{code}

\end{example}




\item 更新系统

更新系统可以获得新功能，可以修正bug，安全漏洞。总之更新的好处多多。对于Gentoo而言，更新的另一个原因是，长久不更新的系统将失去无痛更新的机会甚至不得不重新安装。
更新能保持系统最佳状态。

由于Portage只能检查本地Portage树中已有的ebuilds文件，因此你首先应该更新你的Portage树。使用　\#emerge　--sync　更新本地的portage树，如果使用了overlay\footnote{下节}, 使用　layman -S 更新本地的overlay。


\begin{example}{更新整个系统}
\begin{code}
\# emerge --update --deep world
\end{code}
\end{example}

因为有时开发者会为那些你没有明确要求安装（但被其他软件包依赖而装入系统中）的包推出安全更新，我们也推荐偶尔运行一下这个命令。

每当你改变了系统中任何的USE标记后，你可能需要加入--newuse选项。这样Portage将会检查这个USE标记的变动是否导致需要安装新的软件包或者将现有的包重新编译。

\begin{example}{执行完整更新}
\begin{code}
\# emerge --update --deep --newuse world
\end{code}
\end{example}

\end{enumerate}

\subsection{有关SLOT，Virtual，分支，体系结构和Profile\protect\footnotemark{}}\footnotetext{这段是Gentoo官方文档的复制粘贴版。}

通过使用Portage，一个软件的不同版本可以共存于一个系统中。其他发行版倾向于直接在软件包名字中包含版本号（例如freetype 和 freetype2），Gentoo的Portage使用一种称之为SLOT的技术来实现这种并存。

一个ebuild为它自身的版本声明了一个确切的SLOT。具有不同SLOT的同一软件的Ebuild可以共存于同一个系统中。例如，上例中那个freetype包就拥有不同的ebuilds，里面分别有SLOT="1"和SLOT="2"的标志。

这样软件的包的名字就不会被随机的修改为另外的名字\footnote{Gentoo认为包名加上个2是一个很2的行为}。其他包管理工具，因为缺乏SLOT机制，不得不将版本好作为包名的一部分形成如　freetype2-2.XXX.rpm　这样的包。

有一些不同的软件包提供了类似的功能。比如metalogd，sysklogd和syslog-ng都是系统日志记录工具。那些依赖于“系统日志记录工具”的程序并不能随便的依赖于其中之一，比如metalogd，因为其他的系统日志工具可能也是很好的选择。好在Portage允许使用虚拟包：每一个系统日志记录工具都可以提供virtual/syslog包，因此应用程序们可以设定成仅仅依赖于virtual/syslog即可。

Portage树中的软件可以存在于不同的分支中。系统默认只会接受那些Gentoo认为稳定的软件包。绝大多数新提交的软件会被添加到测试分支里。这意味着在此软件被标示为稳定版前需要进行更多的测试。尽管那些软件的ebuilds已经加入Portage树，在它们未被加入稳定分支前Portage将不会安装它们。

有些软件只在某几个体系结构上可用。或者在其他体系结构中还不能运行，或者仍需要对其进行更多的测试，或者将软件提交到Portage树中的开发者还不能确定这个软件能否运行于其他体系结构。

每一个Gentoo安装都依附于一个确定的profile，此文件里除了其他信息外还包含了一个正常工作的系统需要的软件包的列表。


\subsection{当Portage抱怨的时候......}

恩，使用　portage 的过程并不总是一帆风顺的。portage 也会出现一些乍一看无法理解的错误。

\begin{enumerate}
\item 被阻碍的软件包

\begin{example}{安装了mariadb又安装mysql}


\begin{code}
\# emerge dev-db/mysql -av \\
\\
These are the packages that would be merged, in order:\\
\\
Calculating dependencies     ... done!\\
【ebuild  N     】 dev-db/mysql-5.5.29-r1 USE="berkdb community ssl -cluster -debug -embedded -extraengine -jemalloc -latin1 -max-idx-128 -minimal -perl -profiling (-selinux) -static -systemtap -tcmalloc \{-test\}" 26,143 kB\\
【blocks B      】 dev-db/mysql ("dev-db/mysql" is blocking dev-db/mariadb-5.5.29)\\
【blocks B      】 dev-db/mariadb ("dev-db/mariadb" is blocking dev-db/mysql-5.5.29-r1)\\
\\
Total: 1 package (1 new), Size of downloads: 26,143 kB\\
Conflict: 2 blocks (2 unsatisfied)\\
\\
 * Error: The above package list contains packages which cannot be\\
 * installed at the same time on the same system.\\
\\
  (dev-db/mariadb-5.5.29::gentoo, installed) pulled in by
    =dev-db/mariadb-5.5*[embedded,-minimal,-static] required by (virtual/mysql-5.5::gentoo, installed)\\
\\
  (dev-db/mysql-5.5.29-r1::gentoo, ebuild scheduled for merge) pulled in by
    dev-db/mysql\\
\\
\\
For more information about Blocked Packages, please refer to the following
section of the Gentoo Linux x86 Handbook (architecture is irrelevant):
\\
http://www.gentoo.org/doc/en/handbook/handbook-x86.xml?full=1\#blocked
\end{code}

mysql 和　mariadb 是相互冲突的，不能同时安装到一个系统上。
为了使安装得以继续进行，你可以选择不安装这个软件包，或者先将发生冲突的包卸载。如先移除　mariadb 再安装 mysql。
\end{example}


如果　mariadb 和mysql不是用户手动安装的，而是作为依赖的一部分自动安装的。升级的时候Portage有可能能自动的决定卸载mariadb安装mysql。
比如在笔者一次例行升级中，出现了　

\begin{code}
\begin{flushleft}
[ebuild  N     ] kde-base/print-manager-4.10.0:4  USE="(-aqua) -debug" 87 kB

[uninstall     ] kde-base/printer-applet-4.9.5:4  USE="handbook (-aqua)"

[blocks b      ] kde-base/printer-applet:4 ("kde-base/printer-applet:4" is blocking kde-base/print-manager-4.10.0)

\end{flushleft}
\end{code}

KDE 4.10 使用　kde-base/print-manager 替换了原先使用的kde-base/printer-applet，而这2个都是作为　kde-meta 的依赖安装的，所以 kde-meta 从 4.9 升级到　4.10 的时候，Portage能正确的卸载掉和kde-base/print-manager冲突的kde-base/printer-applet。

\item 被屏蔽的包

\begin{example}{Portage关于被屏蔽的包的警告}

\begin{code}

These are the packages that would be merged, in order:\\
\\
Calculating dependencies               ... done!


[ebuild  N   \#] media-gfx/blender-2.64a  USE="dds elbeem ffmpeg game-engine jpeg2k nls openexr openmp player sdl sse -3dmouse -col
lada -cycles -debug -doc -fftw (-iconv) -jack -openal -redcode -sndfile -tweak-mode" LINGUAS="zh\_CN -ar -bg -ca -cs -de -el -en -es
 -es\_ES -fa -fi -fr -he -hr -hu -id -it -ja -ky -ne -nl -pl -pt -pt\_BR -ru -sr -sr@latin -sv -tr -uk -zh\_TW" 

The following mask changes are necessary to proceed:

\qquad (see "package.unmask" in the portage(5) man page for more details)

\#required by blender (argument)

\# /usr/portage/profiles/package.mask:

\# Diego Elio Pettenò <flameeyes@gentoo.org> (05 Feb 2013)

\# Needs a complete ebuild rewrite to use CMake, and a new patchset to

\# unbundle the bundled libraries. Use at your own risk; don't ask for

\# a bump unless you can provide the two needed items (see bug \#364291).

=media-gfx/blender-2.64a\\
\\
The following USE changes are necessary to proceed:\\
 (see "package.use" in the portage(5) man page for more details)\\
\#required by media-gfx/blender-2.64a, required by blender (argument)\\
>=dev-cpp/glog-0.3.2 gflags\\
\\
NOTE: The --autounmask-keep-masks option will prevent emerge
      from creating package.unmask or ** keyword changes.
\end{code}

\end{example}

当你想安装一个对于你系统不可用的软件包。你会收到类似这样的屏蔽错误提示。你应该试着安装那些对于你系统可用的程序或者等待那些不可用的包被置为可用的。通常一个软件包被屏蔽的原因在于：

\begin{enumerate}

\item ~arch keyword 意味着这个软件没有经过充分的测试，不能进入稳定分支，请等待一段时间后在尝试使用它。或者如果你打算bleeding edge可以在　make.conf 加入　ACCEPT\_KEYWORDS="~arch"，　这里arch是你所使用的体系结构名字，如amd64、x86、arm、mips。

\item -arch keyword或-* keyword 意味着这个软件不能工作在你机器的体系结构中。如果你确信它能工作那么请到我们的bugzilla网站提交一个bug报告。
\item missing keyword 意味着这个软件还没有在你机器的体系结构中进行过测试。你可以咨询相应体系结构移植小组是否能对它进行测试，或者你自己为他们进行这样的测试并将你得到的结论提交到我们的bugzilla网站。
\item package.mask 意味着这个软件被认为是损坏的，不稳定的或者有更严重的问题，它被故意标识为“不应使用”。
\item profile 意味着这个软件不适用于你的profile。安装这样的应用软件可能会破坏你的系统，或者只是不能与你使用的profile相兼容。

\end{enumerate}

显然　blender　被　package.mask 屏蔽了。一个包的ebuild还在portage里缺又要屏蔽掉禁止用户安装，是不是觉得有点不能理解？
一般来说，新加入的包会先 mask 以免用户在软件稳定之前安装。除非用户自动自己在做什么，他可以手工解除屏蔽。另外一些包被mask是因为portage即将移除它。


\item 依赖缺失

如果出现

\begin{code}
emerge: there are no ebuilds to satisfy ">=sys-devel/gcc-3.4.2-r4".

\qquad\\

!!! Problem with ebuild sys-devel/gcc-3.4.2-r2

!!! Possibly a DEPEND/*DEPEND problem. 

\end{code}

这表示正在尝试安装的应用程序依赖于你的系统不可用的另外一些软件包。请到bugzilla查看是否有此问题的记录，如果没有查找到相关信息的话请提交一个报告。除非你的系统混用了不同分支，否则这类问题不应该发生，若发生了那就是一个bug。

\item 意指不明的Ebuild名称

\begin{example}{Portage对于意指不明的Ebuild名称的警告}
\begin{code}
!!! The short ebuild name "aterm" is ambiguous.  Please specify

!!! one of the following fully-qualified ebuild names instead:

 \\

    dev-libs/aterm

    x11-terms/aterm

\end{code}

\end{example}

表明你要安装的应用程序对应有多个同名的包。Portage会列出所有可供选择的名称匹配的包。你需要同时指定类别的名称——如 emerge x11-terms/aterm。


\item 循环依赖

\begin{example}{Portage关于循环依赖问题的警告}
\begin{code}
!!! Error: circular dependencies: 

ebuild / net-print/cups-1.1.15-r2 depends on ebuild / app-text/ghostscript-7.05.3-r1

ebuild / app-text/ghostscript-7.05.3-r1 depends on ebuild / net-print/cups-1.1.15-r2 
\end{code}
\end{example}

循环依赖应该是最头疼的问题吧。幸运的是，几乎只有安装过程中会出现，平时使用基本不会遇到。安装过程中会遇到循环依赖，通常是USE没设置好。

两个（或多个）您想安装的包由于循环依赖而不能安装。这很可能源于Portage树中的bug。请等一段时间后重新sync再尝试安装。您也可以去bugzilla看看是否已经有此问题的报告，或者提交一个关于它的报告。

若是USE的关系，Portage通常会在错误信息中给出调节USE的建议，则可以尝试调节一下USE参数。

\item 下载失败

\begin{example}{Portage关于下载失败的警告}
\begin{code}
!!! Fetch failed for net-misc/avbot, continuing...

(...)

!!! Some fetch errors were encountered.  Please see above for details.
\end{code}
\end{example}

当Portage下载指定软件的源代码失败时，它会尝试继续安装其它（若适用）的应用程序。源代码下载失败可能源于镜像服务器没有正确同步，也可能因为ebuild文件给出了错误的下载地址。那些保存源代码的服务器也可能因为某些原因当机。也有可能是国家防火墙阻断了某些源码的下载，也有可能是当地ISP试图劫持HTTP访问插入广告。总之原因很多。

过段时间重试一次看看问题是否仍然存在。

\item Digest验证失败

有时当你试图emerge一个软件包时，它会失败并给出以下信息：

\begin{example}{Digest验证失败}
\begin{code}
>>> checking ebuild checksums

!!! Digest verification failed:
\end{code}
\end{example}

这是Portage树中出现了错误的迹象，通常这是由于开发者在向Portage树提交一个软件包时出错造成的。

当digest校验失败的时候，请不要尝试自己去为此软件包重新产生digest。使用ebuild foo manifest并不能修复问题，反而几乎肯定会使问题变得更糟。

取而代之的应该是等待一至两个小时以便让开发者来修复Portage树。一般来说错误很有可能马上就会被开发者注意到，但Portage树的修复也需要一点点时间。当你等待的时候，到Bugzilla看看是否已经有人报告了这个问题。如果没有，那就为那个损坏的包提交一个bug报告吧。

一旦在Bugzilla上看到此问题已经修复，你只需要重新sync就可以下载下来那些修复后的digest。

\begin{notice}
重要: 但值得注意的是：这并不意为着你可以短时间内多次重复sync你的portage树（对于同一个rsync服务器）。正如（当你运行emerge --sync时）rsync策略所指出的那样，那些短时间内过于频繁进行多次sync的用户将会被更新服务器禁止访问（一般是将你的IP添加到禁止名单并保留一段指定的时间后才解除）。实际上，最好等到你的下次计划更新的日子再sync，因为这样你不会使rsync服务器过载而影响其他用户的正常使用。
\end{notice}


\item 系统Profile保护

\begin{example}{Portage关于profile中保护的包的警告}
\begin{code}
!!! Trying to unmerge package(s) in system profile. 'sys-apps/portage'

!!! This could be damaging to your system.
\end{code}
\end{example}

您要求移除系统核心软件包中的一个。它是您的profile中所列出的必需的软件，因此不能从系统中移除。

\end{enumerate}

\section{理解portage}

\subsection{文件和目录}

要理解　portage 首先要从他的目录结构开始。

所有的Portage配置选项都是通过变量来控制的。

因为许多配置命令在不同的架构下是不同的,　portage通过profile来选择　make.defaults 文件以继承　profile 里的默认配置。如果你打算修改配置，不要变动　profile 里的　make.defaults ，　而应该修改 /etc/protage/make.conf\footnote{你会发现还有一个/usr/share/portage/config/make.conf.example文件。顾名思义，它仅仅是一个例子而已——Portage并不读取这个文件。}　。因为它有更高的优先级，它的设置会覆盖profile里的默认设置。

除了　make.conf ，　portage 还直接支持环境变量——并且环境变量拥有最高优先级——环境变量通常只用于临时改变设置。
当你需要变更Portage安装软件的行为时，你需要做的就是编辑/etc/portage 中的文件。我们强烈建议你使用/etc/portage中的文件而不是通过修改环境变量来变更这些行为！


在/etc/portage目录中，你可以创建下列文档：

\begin{enumerate}
\item package.mask它列出了你永远不希望Portage安装的软件包。
\item package.unmask它列出了本来Gentoo的开发者不建议安装的，但是你希望能安装的软件包。也就是将开发者　mask 的包强制解除　mask 状态。
\item package.keywords它列出了还未被确认适合你的系统或架构，但是你希望能安装的软件包。也就是强制为某些包加入KEYWORDS。
\item package.use它列出了你希望某些特定软件包使用的而不是整个系统使用的USE标记。为每个软件设定独立的USE标记。
\end{enumerate}

这些并不需要一定是文件；它们也可以是目录。例如你可以在　/etc/portage/package.use 这个文件中写入　\texttt{"app-editors/vim X -python -gpm vim-pager"}，　也可以建立
/etc/portage/package.use 目录，并在　/etc/portage/package.use　中建立一个文件（文件名任意），在文件中写入 \texttt{"app-editors/vim X -python -gpm vim-pager"}。

将　/etc/portage/package.use 创建为目录通常是更好的主意——他允许将软件包进行分组，在每个分组里设定各个软件包的USE参数。避免了单个文件变得过于庞大而变得难以编辑。

\vspace{3ex}

Portage树的默认位置是/usr/portage。这是PORTDIR变量定义的。当你把Portage树保存在其他位置（通过改变这个变量），不要忘了相应地改变符号链接/etc/make.profile指向的位置。

如果你改变了PORTDIR变量，你可能也需要改变下面几个变量，因为它们不会知道PORTDIR变量的改变。这是Portage处理这些变量的方式导致的：PKGDIR、DISTDIR、RPMDIR。

\begin{notice}
注意：如果要移动默认的portage树的位置，将PORTDIR的修改写入make.conf而不是export一个环境变量。
\end{notice}

程序的源代码默认保存于/usr/portage/distfiles。这个位置是通过DISTDIR变量定义的。

\section{什么是overlay，如何使用}\label{sec:overlay}

portage里虽然有数万个\footnote{截至到本书写作的时候有32225个ebuild共计15731个软件。}ebuild，但是总有软件是游离在portage之外的。
总要有机制可以添加第三方仓库吧？这样有的特殊软件可以放到第三方仓库中，想要使用的人添加第三方仓库即可使用emerge安装这些软件。
用户还可以自己维护一个给自己用的仓库，编写自己用的ebuild，这样就有可以更强大的定制能力了。
Gentoo使用overlay来支持这种需求。

所谓overlay就是结构和portage一样的一个目录，该目录里存放的ebuild也会被emerge使用。如果overlay里的版本和portage里的一样高，则overlay里的ebuild有更高的优先级。
overlay通常包含portage没有的软件，或者相同版本但是打上不同的补丁以提供增强功能。只要在 make.conf 中设定 PORTDIR\_OVERLAY 为一个overlay的路径即可使用该overlay。
overlay可以有多个，亦即PORTDIR\_OVERLAY可以设定多个路径，以空白隔开（通常是每行写一个overlay，参考 /var/lib/layman/make.conf），相互之间的优先级看它在PORTDIR\_OVERLAY变量中出现的次序。后出现的优先级高。

\subsection{用layman管理overlay}

现成的overlay被gentoo官方收录到了一个xml文件中：

\url{http://www.gentoo.org/proj/en/overlays/repositories.xml}

使用layman(8)工具即可自动获取这些overlay。这些overlay使用版本控制工具\footnote{参考 \chapref{chap:VCS}}，用户也使用版本控制工具获得这些overlay。
所以layman依赖一些版本控制工具，可以用USE控制layman依赖的版本控制工具。

\begin{code}
\#USE="git subversion" emerge layman
\end{code}

这样会将 git 和 subversion （大部分overlay使用git和svn）作为 layman 的依赖而被自动安装（如果还没安装的话）。没有安装的版本控制工具将无法添加使用该工具管理的overlay。

在 make.conf\footnote{ /etc/portage/make.conf 别忘记哦！} 里最后加入一行 “{\tt source /var/lib/layman/make.conf}”（不包括引号）， 然后就可以使用 layman -a overlay名字 添加指定名字的overlay了。

\begin{insertnote}
注意：/var/lib/layman/make.conf只有在添加了第一个overlay后才会生成。所以可以在添加一个overlay后再加入对该文件的source操作。
\end{insertnote}

\begin{example}{添加 gentoo-zh overlay。}
\begin{code}
\#layman -f \\
\#layman -a gentoo-zh
\end{code}

第一次使用layman需要layman -f，这一步到gentoo.org下载repositories.xml文件。如果已经下载过可以跳过这条命令。

\end{example}

\subsubsection*{列出所有的overlay}

\begin{code}
\#layman -L
\end{code}

\subsubsection*{列出已经添加的overlay}

\begin{code}
\#layman -l
\end{code}

\subsubsection*{添加一个overlay}

\begin{code}
\#layman -a <overlay名字>
\end{code}

\subsubsection*{删除一个overlay}

\begin{code}
\#layman -d <overlay名字>
\end{code}

\subsubsection*{更新一个overlay}

\begin{code}
\#layman -s <overlay名字>
\end{code}

\subsubsection*{一次性升级所有overlay}

\begin{code}
\#layman -S
\end{code}

\else
\addtocounter{chapter}{1}
\fi

\chapter{图形界面}

\begin{quotes}{Steven Paul Jobs}
We made the buttons on the screen look so good you'll want to lick them.
\end{quotes}

\ifdefined\compileall


对命令行是不是有点枯燥了？是时候来点新鲜的东西了 只所以现在才讲，因为图形实在是太复杂了。说Linux图形系统的复杂度已经超过所有其他子系统的总和都不为过。

所有的图形系统，都有一个“指点”设备。图形界面是一种“看得见”的界面，同时需要有种东西能模拟“摸得着”。模拟“摸得着”就需要“指点”设备。通常在桌面型电脑上是鼠标，而在平板电脑上则是触摸屏。“指点”设备告诉计算机，用户“指”的地方和“点”的地方。

除了“指点”设备，传统的键盘也是一个重要的组成部分\footnote{在平板上，键盘被取消了。取而代之使用屏幕键盘和手写输入法}。键盘和指点设备都是用户输入设备。

图形界面下，多个图形程序相互之间共享输入设备和输出设备。输出设备通常就是指的屏幕。

同控制台程序只能接收字符化的用户输入和输出字符串相比，图形程序能输出表意更直观的图像，接收的输入也更为复杂——有鼠标事件，有键盘事件，还有其他图形程序发来的消息，更有系统发送的重绘消息。

由于系统要进行多任务处理，输入和输出设备就不可避免的需要进行共享。完成这个共享的，就是操作系统的窗口系统。

\section{图形、字体和窗口}

窗口是一块有属主并能联接到输入设备的屏幕区域（通常这个也是一块矩形区域）。这块有属主的区域可以为其所有者所控制，接受输入设备的输入并用来输出图形信息。

在操控系统控制的屏幕上有且只有一个窗口能接受键盘输入。接受键盘输入的窗口被称为当前窗口。用户敲击键盘所产生的输入事件被发送给当前窗口，并由该窗口的所有者接受输入事件。

用户可以移动和切换窗口。某个窗口变成当前窗口的操作称为激活。被激活了可以接受到键盘输入的窗口被称为拥有了键盘焦点。
键盘事件总是发送给拥有键盘焦点的窗口。
窗口之间可以相互遮挡，被称为窗口Z轴层级。
拥有焦点的窗口需要处在顶层，所以被激活的窗口总是自动被提升到顶层。
鼠标被移动的时候需要向鼠标经过的窗口发送移动事件。鼠标点击的时候向被点击的鼠标发送点击事件。

从上面的描述可以推测，操作系统的窗口系统的主要任务有：

\begin{enumerate}

\item	提供输入设备的共享。鼠标消息总是发送给鼠标所在的窗口。键盘消息总是发送给拥有键盘焦点的窗口。键盘焦点可以切换。

\item	维护窗口的位置、大小等属性。

\item	管理窗口父子关系。

\item	维护窗口层级。因为运行窗口重叠，层级在前的窗口会覆盖掉层级在后的窗口的内容。

\end{enumerate}

我为什么没有在这里提到窗口绘制呢？严格来说，窗口内容的绘制并不是窗口系统本身的工作，尽管窗口绘制也是图形界面的基础功能之一。窗口内容的绘制，包括文字的渲染，都是被称作为图形渲染系统的工作。
图形渲染系统并不是窗口系统的一部分，而是作为一个基础系统被窗口系统所依赖——窗口系统利用图形渲染系统提供的功能进行窗口绘制。
甚至为了安全，图形渲染系统只能透过窗口系统间接使用\footnote{Linux下可以绕过X直接在屏幕绘图的办法有
/dev/fb*，DirectFB和KMS。}。

图形渲染系统需要提供的功能有：
\begin{enumerate}
\item	提供一个framebuffer。这个framebuffer就是显示器的输出。显示器的内容是由framebuffer直接呈现的。提供framebuffer，意味着可有查询和设定framebuffer的分辨率，
色深，刷新率。

\item	提供绘图API。最基本的API就是bitblt，也就是向framebuffer执行块拷贝操作。其他的一切API都是可选的\footnote{这也就是为何现代X11被用到的渲染API只有XRender（硬件加速的块拷贝操作）一个而已。我们在wayland一节里提到。}。

%\item	提供字体查找和绘制
\item	（可选）提供硬件加速。也就是可有访问GPU，利用GPU加速图形操作。最佳的方式就是提供OpenGL绘图API并使用GPU加速。
\end{enumerate}

正图\ref{fig:guisystem}描述的那样。窗口系统是建立在图形渲染系统基础之上的。
GUI程序使用窗口系统管理窗口，并通过窗口系统获得自己对应的framebuffer区域，然后时候图形渲染系统执行界面渲染向用户呈现图形化的界面。
因为安全性的原因，GUI程序是不能直接使用图形渲染系统，必须通过窗口系统进行授权以获得一块屏幕区域。否则一个程序能向不属于自己的窗口绘制，不是乱套了么？

\chatu{guisystem}{窗口系统}

其他的功能，如“字体渲染”是建立在原有的功能基础之上的。应用程序可以将文字自行绘制为位图然后bitblt到framebuffer，所以字体功能并不是必须的。
但是为了程序间字体的一致性，通常这个可选的功能仍然由图形系统统一提供\footnote{X提供了Xfont核心字体和Xft扩展，用于绘制X内建字体和freetype矢量字体。不过Gtk和Qt都不使用它，而是使用freetype2库自行绘制然后使用XRender绘制到窗口。这样freetype2+fontconfig而不是X提供的基础反而成了Linux下字体渲染的标准。Windows{}和{}Mac OS X{}也都在窗口系统上内建字体支持。}。由图形系统统一提供的好处是系统里只有一套字体管理机制，不容易带来混乱\footnote{在freetyp2和fontconfig接管X的字体功能之前，Linux下的字体管理曾经相当的混乱。安装一个字体需要到X、Xft和fontconfig下分别配置三次。而现在只需要利用fontconfig的自动配置能力即可实现即拷即用——拷贝字体到字体文件夹下即可使用。}。

\chatu{fontsystem}{字体系统}

可以看出的是，操作系统可以不提供字体字体系统，完全不影响。如图中所示，字体选择和光栅化\footnote{显示输出设备通常是像素构成的位图，而高质量的字体则是矢量图形。需要将矢量图形转化为位图，这个过程叫光栅化。}可以由应用程序自己\footnote{或自己写或拜托freetype2这样的库}进行。

\section{X窗口系统}

X的全称是X Window System。安装名称，是一个窗口系统。其实并不全面。X除了是一个窗口系统，还是一个图形渲染系统。X不仅提供窗口管理功能，也提供了图形绘制功能\footnote{而正是因为X提供了图形渲染，而导致了X成为性能低下的代名词。下面马上会讲到原因 !\_!}。

X协议下的Linux图形系统构造是如图\ref{fig:Xstruct} 这样的。

X是MIT在1984年的时候开发的，到1987年的时候出到了第11版本，也是当前正在使用的版本。当前的第11版X称为X11。

X实际上只是一个协议。就像HTTP是个统治互联网的超文本协议一样，X是统治UNIX图形界面的协议。
X是个协议就意味着必须有Server和Client。Server和Client相互使用X协议进行通信。所有的操作都是Server代为Client进行的。Client只需要通过X协议向Server发生请求即可。

也就是说, Client可以和Server处在不同的电脑上。这也是X的一个优势——网络透明。

\chatu{Xstruct}{X结构}

X的优势也是一个劣势：通过网络进行的操作不可避免的带来了大量的开销和延时。X通过缓存一部分请求然后封装成一个包进行批量发送在一定程度上减小了网络开销。但是延时是不可避免的。

\begin{notice}

想想用户点击一个按钮后，要从硬件到内核再到X Server，这都是很快的，慢速的地方在进入X Server后开始了。X Server同时连接了大量的 X Client， 任意时间都有大量的请求过来，非常忙碌。鼠标事件首先被X Server缓存，等处理完X Client的请求，开始把鼠标事件进行翻译，找到鼠标所在的窗口，将鼠标点击事件按照X协议格式化，再发送给X Client。X Client收到鼠标信息后，要进行“视觉回馈”。于是紧接着开始了大量的绘图操作，而所有的绘图操作，在X环境下，都是通过X协议透过网络转发到X Server上代为执行的。
想到一个简单的鼠标点击就产生多少次网络来回的数据包吧。 再想想，X同时为数十个X Client服务——极限环境下可能成百上千个——X Server 很有可能忙于处理X Client的请求而导致延时大大增加。

% 想如果X学Windows 7在入门版只允许开3个程序。

\chatu{XstructInteract}{一个按钮的交互}

\end{notice}


不过X也使用了很多的办法来提高速度。当前大部分环境下，基于X Server和X Client都是在同一台电脑上执行的事实，X还可以做很多的优化。比如使用共享内存代替网络传输。

\begin{notice}

\chatu{Xstructuselesscopy}{多余的拷贝}

这对视频播放软件特别有用。播放器在播放一个高清视频视频的时候，相当于每秒向X Server发送25幅全尺寸未压缩的高清图片。这是接近每秒198MB的数据量。想到这么多的数据，事实上发生了至少3次拷贝：一次是从播放器到内核（调用send(2)），一次是从内核到
X server（X Server使用recv(2)接收网络数据），一次是X Server拷贝到显存。放大这些拷贝因素后，得出实际上是每秒近600MB的数据量。这大大占用了系统的总线带宽。
如果这个时候共享内存介入，拷贝就只有一次了：播放器准备好的数据直接被X Server共享获得，然后拷贝到显存。



\end{notice}

\subsection{窗口管理}

用户移动窗口，改变大小，全屏/最大化窗口，最小化窗口，改变窗口次序，切换键盘焦点等这些操作都是所谓的窗口管理。有管理就会有策略。移动窗口，可以移出屏幕不？
改变大小，可以任意变小不？改变次序，以何种形式更改窗口堆叠次序呢？这些都是属于策略问题。不同的人有不同的喜好。一些人还有不允许窗口重叠的喜好，想让所有的窗口都平铺到屏幕上；不同的人喜好不同的窗口装饰\footnote{窗口的外框，最上门一排有窗口的标题和各种控制按钮——最大化，最小化和关闭按钮}风格$\cdots\cdots$

既然策略是一直非常众口难调的东西，MIT发明X的时候决定，X只提供实现窗口管理的机制，不实现窗口管理策略。窗口的管理策略连同窗口装饰，都由一个特殊的程序——窗口管理器——实现。

窗口管理器和其他X程序一样，都是X Client。只不过这个Client比较特殊，一个X Server上有且仅有一个Client能是窗口管理器——在同一个X Server下不能启动两个窗口管理器。

\chatu[width=0.8\textwidth]{kwinconfigfocuspolicy}{KWin配置焦点策略}

窗口管理器定义了窗口管理的规则。比如大部分窗口管理器都使用“点击获得焦点”的策略。点击一个窗口后，能将键盘焦点移动到被点击的窗口上，并将窗口提升到最前端。许多窗口管理器（通过可配置选项）还支持由鼠标所在的窗口自动获得焦点。如图 \ref{fig:kwinconfigfocuspolicy}所示，KWin\footnote{KDE默认的窗口管理器}允许4中焦点策略。默认是“点击获得焦点”。

\chatu[width=0.8\textwidth]{kwinconfigstyle1}{KWin配置装饰风格}
\chatu[width=0.8\textwidth]{kwinconfigstyle2}{KWin配置装饰风格}

窗口装饰是窗口管理器绘制的，通过修改窗口管理器的配置来修改窗口装饰的外观，如图 \ref{fig:kwinconfigstyle1} 和 \ref{fig:kwinconfigstyle2}所示。
图中我使用了2个装饰风格。


%\FloatBarrier

\subsection{混合窗口管理器}

混合窗口管理器——composited window manager——带合成功能的窗口管理器。
即是一个窗口管理器又是一个合成器。窗口管理我们了解了，那么混合是什么呢？

合成——composite——如名字所言，是用来混合的。那么混合的什么呢？
在介绍混合器（Compositor）前我们来看一张绚丽的桌面截图\ref{fig:3ddesktop}\footnote{当然，老样子，还是我的桌面，避免被窥视隐私，将系统托盘图标和浏览器书签栏和谐了。} ：

\chatu[width=\textwidth]{3ddesktop}{三维桌面}

从传统的X结构上，这是不可能实现的。窗口不可能脱离屏幕出现在一个立方体上。
显然这绝对不是X Server能做到的事情。实现这个就需要窗口管理器做点小小的手脚了。

通常，一个X程序的窗口位于屏幕，窗口直接连接到帧缓冲区，调用XDraw*绘制会直接修改窗口所覆盖的帧缓冲区的内容，也就是直接绘制到屏幕上。

在合成环境下，窗口管理器利用特权\footnote{窗口管理相对一般的X Client有特权，否则怎么做管理。}向X Server标记窗口为redirected状态。被redirected的窗口，其绘图操作不会反应到屏幕上，而是绘制到一个缓冲区（off-screen buffer，离屏缓冲区）， 之后窗口管理器把各个窗口的离屏缓冲区的像素数据合成到屏幕上显示。合成期间可以加入各种效果，比如半透明（alpha合成），窗口阴影等等。也就是说，窗口管理器独占了整个屏幕的绘制，
其他程序的窗口被迫被重定向到一个离屏缓冲区绘制，绘制完成后由窗口管理器将绘制好的窗口内容合成到屏幕上。整个系统结构如图 \ref{fig:compositing} 所示。

\chatu{compositing}{X合成管理器}

那么一个窗口被绘制的路径是 X程序-->X Server-->离屏缓存-->合成管理器-->帧缓存（屏幕）。而没有使用混合管理器的窗口绘制线路图是：X程序-->X Server-->帧缓存（屏幕）。
由于多了合成操作，所以合成管理器对系统的要求要比直接绘制稍微多了点，故而合成管理器通常要求显卡具有“硬件加速”能力以使屏幕不至于看“幻灯片”。

\begin{insertnote}
\subsubsection{硬件加速小常识}
在计算机的显卡上，有一块固定的内存区域，被称之为“帧缓冲区”，其实内容就是一幅位图图片。显卡的扫描电路不停的扫描帧缓冲区的内容并转化为显示器驱动信号。
计算机向屏幕绘图的即为将帧缓冲区的图片进行更新。计算机绘制图形就是将帧缓冲区的每一个像素修改为指定颜色的过程。
通常程序使用比较高级的绘图原语——绘图原语通常提供“点，线，多边形，曲线”的绘制操作\footnote{复杂的图形都是又数不清的多边形和曲线构成的}——由绘图引擎将其转化为对每一个具体像素的修改。这个步骤叫“光栅化”。

由于绘图原语中的每一个操作都是对批量数据的简单操作，所以有可能通过硬件直接实现。这样费时的光栅化操作就可以交给专门的硬件电路执行，释放了CPU的处理能力。这个就叫硬件加速。

由于显卡硬件复杂度的提升，越来越多的绘图指令被“硬件加速”了。能够专门利用硬件加速的绘图指令集有2个，一个是Direct3D一个OpenGL\footnote{我们将在 \secref{sec:OpenGL} 介绍OpenGL。}。一小部分X的绘图操作也是硬件加速的——当然，前提是驱动支持。

由于Direct3D为微软专有，所有Linux下的硬件加速无一例外的是采用的OpenGL，故大部分混合窗口管理器都要求显卡支持OpenGL以便发挥显卡的功效获得最佳体验。

\end{insertnote}

现代的桌面有很多地方依赖于“半透明”窗口来提供最佳的用户体验。而要实现半透明窗口就意味着：窗口要和位于其下的窗口执行alpha混合。这就必须要求使用一个混合窗口管理器了。
所幸运的是，并不是所有的混合窗口管理器都要求有OpenGL支持的显卡，仔细优化的纯软件混合一样能获得流畅的桌面。
KWin检测到用户电脑显卡无法开启OpenGL的时候——不论是因为硬件太老而不支持KWin所需要的OpenGL功能还是因为缺乏正确的驱动程序——就自动进入软件混合模式。
作为对比，GNOME使用的mutter窗口管理器拒绝在不支持OpenGL的机器上启动。

看来编写一个性能不错的软件混合管理器难度相当的大。为了节省开发者的精力，我们还是老实的把驱动给安装好吧！

\subsection{OpenGL和Mesa}\label{sec:OpenGL}

前面在讲X结构的时候，都没有说到硬件加速。引入了硬件加速后，X的结构要发生重大的变化，前述的图例已经不能表达X的结构了。

说到硬件加速就不得不说OpenGL。OpenGL从其首页上的描述\footnote{ \url{http://www.khronos.org/opengl} }来说，是2D和3D图形硬件的编程接口。
OpenGL是一个编程接口（API）的规范，应用程序只要使用这个通用的API编写就可以利用强大的GPU的处理能力。OpenGL的实现通常由GPU制造商提供，包含在显卡驱动里。
使用OpenGL的程序链接到 libGL.so.1 库，libGL.so.1 库又由显卡驱动提供。硬件开发商在库中将OpenGL的调用转换为显卡GPU能执行的命令然后交又GPU执行。这样大量的绘图指令就交给GPU而不是CPU执行了。大大提升了图形程序的执行效率。
通常需要如此执行效率的程序都需要绘制3D图形，如3D游戏和3D设计软件。

在Linux的世界，除了显卡厂家提供的闭源的libGL.so.1实现，还有开源的实现。实现开源OpenGL的程序就是Mesa。

在初期，Mesa将OpenGL绘图指令转化为Xlib的绘图指令。中间的转换纯粹是由软件执行的。Mesa又被称为软件渲染。之后 Mesa 加入了针对特定显卡的“驱动”。OpenGL 指令由这些“驱动”翻译为对应显卡GPU的控制指令，然后再发送给显卡执行。我这里“驱动” 加了引号，是因为这个驱动是纯粹的Mesa针对不同显卡写的不同的“GPU指令生成器”。并不是通常意义上载入内核的驱动模块。
由于向显卡发送指令需要root权限，所以这个时候Mesa需要由 X server调用，而 X client则需要将 OpenGL 调用包裹到 X11 协议中。

不过，在加入了DRI\footnote{Direct Rendering Infrastructure}后，X client 就能直接使用显卡的OpenGL功能了。

\begin{insertnote}

\subsubsection{名词解释：DRI}

DRI是Xorg提出的一种全新的渲染架构。结构如图 \ref{fig:dri} 表示。


\chatu[scale=0.7]{dri}{dri结构简图}

X Client 程序还是利用Xlib发出传统的绘图指令。如果使用了OpenGL，并使用直接渲染, 绘图指令将绕过X Server之间经过 libdrm 发送给内核一个叫 drm （Direct Rendering Manager） 的模块。
drm再将指令直接发送给显卡。 因为中间的绘图过程绕过了X，直接向显卡发送绘图指令将图像写入帧缓冲，故而得名直接渲染。 这里请注意的是，直接渲染和间接渲染的差别就是渲染过程 有没有经过X server， 而和是不是利用了显卡加速是没有关系的。间接渲染一样可以使用显卡加速。 直接渲染的优势是中间没有经过X server， 减小了延迟，去除了X11协议的网络开销。
极大的提升了 3D 渲染的效率。

%因为在间接渲染模式下，OpenGL 也被包裹到了X11协议中发送给Xserver由 X server代为调用 OpenGL。

闭源驱动并不使用DRI，闭源驱动有自己的直接渲染方式。因为drm模块是GPL协议的，闭源驱动若使用了drm岂不是要开源掉自己OpenGL实现。所以闭源驱动会自己带有libGL.so.1而覆盖掉Mesa提供的。Gentoo使用eselect机制，安装的时候将libGL.so.1安装到各自的目录，然后使用
\begin{code}
\#eselect opengl set [name] 
\end{code}
设定一个 OpenGL 实现，系统目录下的/usr/lib/libGL.so.1将链接到被选择的OpenGL实现。
通常是 nvidia,xorg,fglrx三者中选一个。xorg 就是使用的Mesa提供的 OpenGL实现。可以用
\begin{code}
\# eselect opengl list
\end{code}
 列出系统已经安装的 opengl 实现。没有安装的并不会列出。

\end{insertnote}

\subsection{内核模式设置（KMS）}

显示器都有一个最合适自己的分辨率。没有人觉得将显卡输出的分辨率固定是一件好事。但是因为历史也技术上的问题，BIOS启动的时候确确实实是以一个固定分辨率输出的。进入图形界面后，就需要根据用户的设置或则自动检测到的显示器规格，对显卡的输出模式进行调整。

传统上，这一次调整是由 X 执行的。但是 X 的调整内核是不知道的\footnote{X 通过自身的显卡驱动直接控制显卡}。
于是如果执行终端切换\footnote{在内核提供的几个虚拟控制台之间切换。}，便带来了问题。
X 本身也是运行在一个虚拟控制台上，但是 X 将控制台的模式切换到了图形显示模式，而内核对这一情况全然不知。
如果内核不顾 X 正在控制显卡的事实，强硬的切换到其他终端，轻则花屏，重则死机。
因此 Linux 不得不引入了其他烦琐的机制： 
\textit{
当 X 运行的时候就禁止了内核的终端切换能力。当用户按下 "Ctrl+Alt+F?" 组合的时候，X 将 显卡反初始化到 X 启动之前的状态， 然后通知内核执行终端切换。如果用户从文本模式切换回X, 则由内核发出信号，X 接收到后再将显卡初始化回切换出之前的状态。
}
如果在切换的过程中，X 进程出问题，不论是出现锁死还是进程崩溃，都将导致整个系统停止响应用户输入。
问题的根源就在于：显卡的显示模式是X配置的，内核无从知晓 X 对显卡执行的任何更改。

为了克服这种弊端，X 的人将模式设置的代码从 X 移入内核。除
了避免终端切换引起的问题，更让文字模式的控制台拥有和图形环境一样的分辨率！
这就意味着，终端切换并不会引起显示模式的切换 —— 用户可以享受到非常快速的终端切换。

内核模式设置的好处不仅仅体现在更好的终端切换体验上。
通过将模式设置代码移入内涵，实际上简化了内核和 X 双方的开发而并不是如人们想的那样， 加重了内核的维护负担。
KMS 让 X 和内核共享驱动，而不是 X 的驱动和内核里的驱动“竞争”显卡的访问。

KMS 也让人脱离 X 实现图形界面成为了可能\footnote{参见 \secref{sec:wayland} }。

\section{安装基础X环境}

%一个基础的X环境包括一个X Server和支持库。这些库包括 Xlib（以及一些扩展）和libGL。值得注意的是，X Server 依赖于libGL而不是反过来。

在Linux世界，曾经有一个开源的X11实现名字叫做XFree86\footnote{x86上用的自由（Free）的X}。
不过 XFree86 的幕后公司开始将 XFree86 闭源并做商业盈利的时候伤害到了社区的感情。
大量的开发者逃离 XFree86 项目，以最后一个公开的 XFree86 源码为基础重新组建了一个 X.org 组织。
所以 Gentoo 曾一度提供两个X Server： XFree86 和 X.org 。
不过 UNIX 世界有个经典名言：离开源越近越繁荣。
被社区抛弃的 XFree86 很快销声匿迹了。
所以现在的X11就是用的 X.org 了\footnote{后来MIT将X11协议本身的维护也交由X.org了。这都是后话了。}。


前面图片里可以看到，X Server分成两个部分：DIX（Device Independent X）和DDX（Device Dependent X）。 DIX实现X11协议，和X Client通讯，接受请求。
DDX实现的是和硬件的交流。硬件分两个大总类：Video和input。
之前在图 \ref{fig:Xstruct} 提到过，X对显卡的操作是不经过内核的。X有自己的显卡驱动。X必须运行在root权限的缘由也是基于此。X通过将显卡的IO映射到自己的内存空间来实现直接操作显卡。对内存映射IO可以参考 \faqref{FAQ:MMIO}。

X Server即利用了这一点。X Server首先通过内核的帮助将显存映射到自己的内存里，之后就可以绕过内核直接对显存进行访问。除了显存外，显卡还有IO端口，用于向GPU发送命令。X Server也将显卡的IO端口映射到自己的内存里，之后就可以直接对显卡进行操作了。故而X需要显卡驱动。

显卡驱动是DIX的重要组成部分。X.org 发布的X Server中自带有大量的显卡驱动（开源驱动）。还有一些驱动由硬件厂家提供（闭源的二进制驱动），需要自行下载安装。

Gentoo portage有个重要的变量VIDEO\_CARDS\footnote{到 /etc/portage/make.conf 中设置} 指示emerge使用的显卡驱动，一来安装X的时候能正确的自动安装好驱动，二来能自动选择好Mesa使用的驱动，三来可以略过全部无关驱动的安装。

VIDEO\_CARDS的取值如表 \ref{table:videocards} 所示。

\begin{longtable}{|l|p{0.7\textwidth}|}
\caption{VIDEO\_CARDS 驱动列表}\label{table:videocards}\\
\hline
apm & Alliance ProMotion显卡驱动。 \\\hline
ark & ark 显卡驱动。 \\\hline
ask & ask 显卡驱动。 \\\hline
chips & Chips and Technologies 显卡驱动。 \\\hline
cirrus & Cirrus Logic 显卡驱动。\\\hline
%cyrix & Cyrix 显卡驱动。 \\\hline
%displaylink & displaylink 显卡驱动。 \\\hline
dummy & 伪显卡驱动。只用于Xvfb。 \\\hline
fbdev & 通用帧缓存驱动 ， 使用 /dev/fb* 设备。 \\\hline
fglrx & AMD/ATI 闭源驱动。  \\\hline
%geode & AMD Geode GX 和 LX 显卡驱动。 \\\hline
glint & GLINT/Permedia 显卡驱动。 \\\hline
i128 & I128 显卡驱动。\\\hline
%i740\footnote{服务器芯片。家用电脑没有使用过。} & Intel 740 系列显卡驱动。 \\\hline
%impact & Impact 驱动 \\\hline
intel & Intel 显卡驱动。如无例外，Intel显卡全部适用此驱动。\\\hline
%ivtv & ivtvdev 显卡的 IV-out 功能驱动。 \\\hline
mach64 & ATI Mach64 显卡驱动。 \\\hline
mga & Matrox 显卡驱动。 \\\hline
modesetting & 通用显卡驱动，适用于KMS\footnote{见 \secref{FAQ:KMS} } 帧缓存设备。\\\hline
meomagic & Neomagic 显卡驱动。 \\\hline
%nested & 伪显卡驱动，用于 Xnested。 \\\hline
%newport & Newport 显卡驱动。 \\\hline
nouveau & 开源的NVIDIA显卡驱动。带3D加速。 \\\hline
nv & 开源的 NVIDIA 显卡驱动。无3D加速。可用于nouveau无法驱动的旧型号显卡。 \\\hline
nvidia & 二进制NVIDIA显卡驱动，由 NVIDIA 官方提供。\\\hline
%nsc & nsc显卡驱动\\\hline
%omap & ARM 芯片供应商德州仪器制造的OMAP系列芯片的片上集成显卡的驱动。 \\\hline
%omapfb & OMAP芯片提供的帧缓存的驱动。\\\hline
%psb &  一些上网本使用的 GMA500 显卡驱动，intel贴牌，实际由poulsbo制造。\\\hline
qxl & Qemu虚拟机 QXL paravirt 显卡驱动。\\\hline
r128 & ATI Rage128 系列显卡驱动。\\\hline
radeon & 开源 AMD/ATI 显卡的驱动。 \\\hline
rendition & Rendition 显卡驱动。\\\hline
s3 & S3 显卡驱动。\\\hline
s3virge & S3 ViRGE 显卡驱动。\\\hline
savage & S3 savage 显卡驱动。\\\hline
siliconmotion & Silicon Motion 显卡驱动。\\\hline
sis &  SiS 和 XGI 显卡驱动。\\\hline
sisusb & SiS的USB显卡驱动。\\\hline
tdfx & 3Dfx 显卡驱动。\\\hline
tga & tga 显卡驱动。 \\\hline
trident & Trident 显卡驱动。\\\hline
tsend & Tsend Libs显卡驱动。\\\hline
v4l & video4linux驱动。\\\hline
vesa & 通用显卡驱动。最后的救命稻草哦！\\\hline
via & VIA/S3G 集成显卡驱动。\\\hline
virtualbox & VirtualBox 虚拟机里安装Gentoo的话使用这个驱动。\\\hline
vmware & vmware虚拟机里使用Gentoo的话安装这驱动。\\\hline

\caption{VIDEO\_CARDS 驱动列表}\label{table:videocards}
\end{longtable}


\begin{insertnote}
注意：使用何种显卡驱动不决定于显卡板卡本身的品牌，而是看显卡所使用的GPU芯片。
\end{insertnote}


其实值得说明的是，虽然看上去驱动很多，其实现在显卡世界经过多年激烈的竞争后，已经只剩下 intel AMD NVIDIA三巨头了。
intel自不用说，有且仅有开源驱动，VIDEO\_CARDS安心的写上intel即可。AMD和NVIDIA则是开源驱动和闭源驱动共存。使用开源还是闭源驱动，第一要务是：必须能跑起来！如果两种驱动都合适，就看各位自己抉择了，是使用厂家提供的驱动还是使用开源驱动。AMD 的闭源驱动叫 fglrx 支持的是 AMD自家显卡和收购 ATI 之前年开发的显卡， 开源驱动实际上安装的是 xf86-video-ati 不过，VIDEO\_CARDS的设置是 radeon，支持的显卡比fglrx稍微在旧的批次上多点——如果fglrx放弃支持你的老显卡了，你只能用radeon驱动了。 最早一批的ATI显卡使用的是r128驱动,没有闭源驱动支持，不过我很怀疑现在是否还能看得到这些显卡。

NVIDIA显卡则有3种驱动可选 nv nouveau 和 nvidia。 在nouveau之前nv是nvidia唯一的开源驱动，nouveau出现后已经不建议使用了。nvidia的闭源驱动有多个版本，96.XX 版本 173.XX 版本和 3XX 版本。分别对应的支持最老旧的NVIDIA显卡，很老的NVIDIA显卡和现在还在使用的NVIDIA显卡。具体你的显卡适用哪个版本可以到官方网站的驱动下载页面查询。
一帮来说，除了非常老的显卡，最新驱动都可以支持。

接下来要设定的就是输入设备的驱动了。输入设备的驱动由 INPUT\_DEVICES\footnote{同样在 /etc/portage/make.conf 设置。} 选择。

INPUT\_DEVICES可用的取值如表 \ref{table:inputdevice} 所示。


\begin{longtable}{|l|p{0.7\textwidth}|}
\caption{INPUT\_DEVICES 输入设备驱动列表}\label{table:inputdevice}\\\hline
evdev & 标准Linux输入设备驱动。适用于键盘和鼠标。 \\\hline
acecad & Acecad Flair input driver。  \\\hline
aiptek & Aiptek  USB 数字平板输入仪驱动。 \\\hline
elograhpics & Elographics 触屏驱动。\\\hline
fpit & Fujitsu Stylistic平板的触屏驱动 \\\hline
hyperpen & hyperpen 手写笔驱动 \\\hline
joystick & 游戏手柄驱动 \\\hline
mouse & 鼠标驱动（已被evdev替代） \\\hline
keyborad & 键盘驱动（已经被evdev替代）\\\hline
synaptics & 触摸板驱动，不安装则触摸板只是一个鼠标，不能使用触摸板的高级功能 \\\hline
tslib & 利用 tslib 实现的触摸屏驱动。 \\\hline
vmmouse & VMWARE 虚拟机鼠标，使用此驱动代替evdev可以获得VMWARE和宿主机的集成功能。\\\hline
%evdev -acecad -aiptek -elographics -fpit -hyperpen -joystick -keyboard -mouse -mutouch -penmount -synaptics -tslib -vmmouse -void -wacom
wacom & wacom手写板输入 \\\hline

\caption{INPUT\_DEVICES 输入设备驱动列表}
\end{longtable}


evdev 是 Linux下的新设备驱动,代替之前分开的鼠标和键盘驱动。Linux内核提供统一的 /dev/input/event* 设备文件供用户空间的程序使用。

如果不是使用BSD的话，一般只需要使用evdev即可。对于笔记本，还需要启用 

设置好  INPUT\_DEVICES 和 VIDEO\_DEVICES 即可通过

\begin{code}
\#emerge xorg-server
\end{code}

安装 Xorg 的X Server了。安装Xorg需要较长的时间，除了本身Xorg比较大之外，Xorg引入了大量的依赖也是重要的原因。而且很有可能会出现循环依赖。需要依据提示暂时禁止某个USE标记
以打破循环依赖。

\section{桌面环境}

光有一个X什么也做不了，X是一套实现GUI的机制，我们需要一个“操作环境”。正如shell提供了文本模式的操作环境一样，在X下也需要一个“shell”来提供图形模式的操作环境。
做到这样一个图形操作环境，需要：窗口管理器、桌面、文件管理器、任务栏（提供窗口切换）、系统托盘区域、各种小程序等等东西，整合一些提供这样功能的程序集合就构成了一个“桌面环境”。
Linux有很多各种各样的桌面环境，最流行的还是KDE和GNOME。

\subsection{KDE}

传统的UNIX使用的桌面环境是CDE——Common Desktop Environment。但由于昂贵的授权费和对其他一些可选的替代方案的不满导致了德国人Matthias Ettrich\footnote{同时也是LyX的作者。}于1996年发起了KDE（Kool Desktop Environment）项目。也许看看他当时的宣言能更好的理解当时的状况。

\begin{insertnote}


\begin{center}
KDE 项目宣告

新项目: Kool Desktop Environment (KDE)

{\scriptsize 急需程序员！}

{\footnotesize 译：微蔡}
\end{center}

{\large  \noindent \bf 动机}

自由UNIX衍生版特别是Linux的出现让UNIX变得更流行了。但是仍然缺乏一个外表好看，界面统一的自由桌面环境。便宜甚至免费的程序已不在少数，只要能提供一个真正的GUI环境，Linux+X11基本上就能满足所有人的需求了。

这其中当然有这样的GUI环境，Commond Desktop Environment（太贵了）和Lookin Glas（相对便宜点），当然还有其他一些能算的上是“GUI环境”的免费程序：
比如 Moxfm 工作的就很好。可惜他是基于Motif设计的。当然，更重要的问题是：啥是一个GUI环境？它应该是怎么样的？


首先，我列出什么不是GUI环境：

\begin{enumerate}
\item [\textbullet] X窗口系统不是GUI环境。如名字所言，是一个窗口系统。
\item [\textbullet] Motif 不是一个GUI环境。Motif创立者曾希望通过Motif创建一个GUI环境，不幸的是他们意见没能统一，只好将Motif作为一个控件库发布并附带了一个文件管理器。

\item [\textbullet] 窗口管理器也不是GUI环境。他们只是一个管理窗口的小程序。把很多东西放进窗口管理器本身并不是个明智的做法。
\end{enumerate}

就我所知，一个GUI环境应该提供一个完整的图形环境。允许用户用它来做每天的工作，比如开启程序、阅读邮件、配置桌面、编辑文件、删除文件、看看图片等等。
所有的这些部件都应该整合到一起协调工作。带一个编辑器图标的按钮如果启动了xterm这个终端仿真程序在里头执行个vi就不能算是一个真正的GUI。
可能你很早就失望的发现，点击桌面上漂亮的帮助图标却等来的却只是一个终端仿真器，里面因为运行着man而正在显示man手册。

{\large \noindent \bf 面向最终用户的GUI}

我的主意并不是为一个完整的UNIX系统或者系统管理员创建GUI——这种情况下一个shell和上千个命令行工具可能是个更好的办法——而是面向最终用户。为那些希望上上网，写写文档，玩玩游戏的用户准备。

直到我配置我女朋友电脑前，我真心觉得Linux更能做到这点\footnote{这句话为fcitx作者csslayer帮我翻译的，原文是“I really believed that is even yet possible with Linux until I configured my girlfriends' Box. ”}。
就别提我要处理多少种形式和风格各易的菜单、滚动条和文本编辑框了。
我已经知道一些控件必须在鼠标下才能进行输入，一些页面需要用鼠标中建拖动，一些编辑框甚至只能使用和emacs一样的快捷键——连“End”按键都不能支持。
各个程序也都有不同的选择文字的办法。甚至连菜单和按钮都完全不具有一致的行为（比如Xaw, Fvwm, XForms, Motif）。

一句话来形容这些雅典娜\footnote{智慧与技艺的女神，这里指的是每个程序作者都发挥自己的智慧和技艺创造出各种形式的界面，风格及其难以统一。}型的控件：
尽管有很多使用这些雅典娜控件的程序，我们最好还是丢了他们吧。“雅典娜控件”和“X是GUI”应该是一样的错误想法。这种雅典娜式的控件应该视为如何使用Xlib和Xt构建一个控件库的反面例子。
X应该有一个控件库作者的在线帮助，而无论如何不应该成为应用程序作者的文档\footnote{指的是应用程序作者应该使用一套统一的方便的UI库来创建程序而不是直接使用Xlib。}。 不信的是，应用程序作者就是靠的X手册编写的。

所以，我提议一个GUI环境首先就应该为所有的程序提供一个一致的“外观”。这正是本项目所以和其他的项目最本质的不同点。

就在几个星期前，一个对自由软件开发而言即免费又自由的UI控件库出现了，\url{http://www.troll.no}\footnote{光阴似剑啊！Qt最原始的网站早已不可访问了。} 去看看吧。

这个库就是Qt，并且将是X界面编程的革命。它差不多已经非常完善了，完全的C++类库，提供了Motif相仿的外观，或者可以配置为类Windows95的外观。

它由一家公司控制的事实反而是一个大大的好处：我们有代码和这么好的库，他们有测试人员。他们也全力的提高这个库的质量。他们还提供技术支持。这就意味着，Qt也能用于开发商业软件。
它是Motif的真正替代品 ：）。  让我觉得Qt很专业的地方就是它的代码了。它真的是非常容易使用的C++库。

Qt还是可移植的，是的，兼容Windows。当然这你不用当心。在你代码里稍微加点UNIX/X专有的东西就跨不到Windows平台了 :)。


我强烈的建议大家看看这个库。它有能力成为自由软件开发的领军库。使用它就轻而易举的逃过了那个会吃掉所有内存和CPU的怪物级TCL/TK库$\cdots$

是时候为桌面设立一个标准了。同时往内存里放10个不同的控件库是在是一点意义都没有。想象一下这些程序和使用的库：

\begin{enumerate}
\item[\textbullet]    fvwm （自己的控件）
\item[\textbullet]    rxvt （自己的控件）
\item[\textbullet]    tgif （自己的控件）
\item[\textbullet]    xv （自己的控件）
\item[\textbullet]    ghostview （雅典娜控件）
\item[\textbullet]    lyx （xforms 空间库）
\item[\textbullet]    xftp （Motif 控件库）
\item[\textbullet]    textedit （xview控件库）
\item[\textbullet]    arena （自己的控件）
\end{enumerate}

可能有人会狡辩说，UNIX机器内存很大，往内存里放那么多UI控件库不是问题。即便是这样，真正的问题是：所有这些控件（菜单，按钮，滚动条等）库统统都没有一致的外观和行为。
这不是学院研究的例子，这是我真实见到过的 :-\}

我知道不可能一下子就清除掉它们，我也希望Motif和Qt能共存。


{\large \noindent \bf The Kool Desktop Environment (KDE)}

我没有时间自己一个人全部做下来（也有LyX\footnote{LyX是一个基于\LaTeX{}的图形化编辑器。}是我的主要项目这个原因。）。
像KDE那样的打项目能很容易的切成各种小部分。其中就有给你的部分！如果你要学一些X编程，那干嘛不开始给KDE写点小东西呢？
要是你知道有人能写点程序，请千万不要让他们再写第1004\footnote{作者这里的意思是，别浪费才能和时间在克隆已经写过无数遍的游戏了，参加到KDE项目中来吧！}个俄罗斯方块或者地768个扫雷了。;-) 我们已经有很多了 ... 

然后接下来就是我的KDE项目需要的程序的列表。也许还有更多，这是一个开发的项目，欢迎添加。

{
\leftskip=2.2em
\small 

{\noindent \tt 面板：}

    基本程序。开始是Fvwm的一个模块运行。提供类型Windows 95那样的底部任务栏和CDE那样的顶部面板。面板左边有菜单按钮用来开启程序，中间有4个用来切换虚拟桌面的按钮，然后右边有些快捷启动。
    比如有个邮件图标并且能通知新邮件的到来，一个垃圾桶指示回收站是否为空并能快速打开。或许还可以有一些电子时钟放在最右边。
    最好还有个漂亮的图标用来锁定屏幕、注销和关机。所有这些都是（在图形界面下）完全可配置的。
    我还在想一个解决办法，只让可用的GUI程序出现在桌面上，并且新安装的程序能自动的出现在桌面上。

    我已经着手开始写面板了，非常希望能获得一些帮助。还有很多事情要做呢！比如一个工具用来设置一个壁纸。
    
    一个好看的图标也是很有帮助的！

{ \noindent \tt 文件管理器：}

    KDE的另一个重要软件。创建一个功能强大的bash图形替代品, 并对简单的任务有一个易用的界面。简单的任务就是上传文件，拷贝文件，双击打开文件啊之类的。

    我还想到为每个目录打开一个漂亮的窗口，每个文件都是一个图标。图标可以拖动以复制或者移动文件——甚至可以在窗口直接拖动。
    还有一个重要的功能就是支持移动设备的自动挂载，这样对用户就是透明的了，无需关心mount/umount命令。

    拖动应该使用比较优雅的方式实现。拖动的图标应该是一个特殊的从窗口（参照 Qt 的 xshape 例子）而不像xfm或是xfilemanager那样是通过设定鼠标的图标来进行。

    所以把文件作为一个图标放到桌面上也是可行的了。这非常有意思，不是么？ 一些应用程序用面板启动，现在很清楚了，图标就是个数据对象，而fvwm1和 FvwmFileMgr都不能确定一个图标是文件还是最小化的窗口。    

    Qt程序之间实现拖动并不难。文件管理器实现起来并不难，并且也花不了多少时间，有谁能接手这个任务？
    
{ \noindent \tt 邮件客户端：}

    真正舒心的邮件客户端。现在我发现最好的邮件客户端（X环境下的）是XF-Mail。作者表示一旦KDE项目启动会将其移植到Qt下。不过他表示需要一些协助（比如写个简单的弹出菜单）。 
   
{ \noindent \tt 简单的编辑器：}

    很小但很重要的项目。一个月都编辑不了一次文件的人不可能乐意去学vi这样的软件或者等x-emacs\footnote{emacs的优良传统，启动非常慢。}启动甚至都不会掏钱去买个motif-static-nedit\footnote{一个基于Motif开发的文本编辑器}，更不可能乐意使用怪兽级别的 tkedit\footnote{作者说tk是怪兽，自然用tk开发的tkedit也是如此了。} ....

	不幸的是，Qt-1.0 还没有提供多行编辑控件。但是Troll-Tech公司已经在测试版里提供了，所以几个星期后应该就可以着手编写文本编辑器了。
	
{ \noindent \tt 终端模拟器：}

    和CDE的终端模拟器差不多。就是带菜单，能设置字体的 xterm。 很好的项目，拿xterm的代码修修，移植到Qt下就可以了。
    
{ \noindent \tt 图像查看器：}

    这个程序在文件管理器里将会被默认用于查看gif，jpeg等图片文件。额，xv是不错，不过是个共享软件。启动速度也不快。有个没有菜单没有按钮的纯Xlib软件叫xli启动速度倒是飞快。把它的代码拿来移植
    到Qt下作为起始点吧。 %. But there is a plain Xlib programm without any menues or buttons called "xli". Get the sources and make it userfriendly with Qt!
    
    
{ \noindent \tt 其他一些有用的工具：}

        xdvi 改成 Qt-Gui
        ghostview 改成 Qt-Gui
        xmag 改成 Qt-Gui
        其他

{ \noindent \tt 超文本格式的帮助系统：}

    一个完整的桌面需要一个友好的帮助系统。我认为一个很好的选择就是 HTML （>=2.0） 。 需要一个基于Qt的HTML查看器。
    可以借用 Arena 的代码，不过Arena 启动太慢了。最好还是从头写一个。Qt 提供了非常简便的接口处理字体。对于一个帮助系统来说， HTML 2.0 已经足够了。再添加一些搜索功能就可以了
    因为传统的 troff 格式的man手册也能转换为HTML格式，我们还能集成man手册到帮助系统里。

    顺便一提：Troll Tech公司开出了一个Qt竞赛（看看他们的首页）。最好的Qt程序（不仅仅是功能上，还包括设计上。仅仅将现有的程序移植到Qt下还不够的 :-( ）将赢得 2000美金和 Qt 的Windows版的授权（价值2000美金。）。他们还提到了一个浏览器作为例子。所以一个 Qt 的浏览器将值 4000 美金哦 （把用不到的 Windows版Qt授权卖掉） ;-)。


{ \noindent \tt 窗口管理器：}

    最开始，KDE面板将作为Fvwm的一个模块。等做好了，很多东西就可以从 fvwm 窗口管理器里剥离了。我们不需要fvwm的菜单、图标和难懂的配置文件。我们只需要一个很小的靠谱的窗口管理器。
    所以从 fvwm里删除一切不需要的代码，应该很快就能完成。


{ \noindent \tt 系统工具：}

    不管你是什么类型的用户，都会需要的。一个图形化的密码设置对话框出现在我脑海里。还有更多的吧！可能会导致一个小型系统管理工具集合也说不定。
    

{ \noindent \tt 游戏：}

    俄罗斯方块我们已经有了（作为Qt的一个示例）。我们还需要的是很多其他的小游戏，比如纸牌。有一些已经跑在X下的游戏，干嘛比从中汲取资源，然后自己写个真正的游戏呢？就像微软的纸牌那样。
    有时候我就为了玩纸牌安装个Wine。还有很多很多的游戏，多多益善。要么从现有的修改，要么从头写，咋样都行：）
    
{ \noindent \tt 图标：}

    一组漂亮的图标。可以从3D-pixmaps开始。（如果我们用的UI库有按钮的话干嘛还需要用图片模拟按钮呢？）
    
{ \noindent \tt 文档：}

	文档总是好东西。在这之前，我们的超文本格式的帮助系统应该先设计出来。然后我们就能开始写HTML格式的帮助页面了。喏，在这之前我们还得有程序可文档啊！

{ \noindent \tt 基础设施——页面/文件服务器/管理员：}

    我们需要一个服务器来保存这么多文件和页面，通告大家项目的状态。特别是当前我们在做什么，还需要人做什么。稍后我会在
    
    	\begin{center}
    	\url{http://www-pu.informatik.uni-tuebingen.de/users/ettrich}
    	\end{center}

设置一个首页。当然，只包含本文和一些链接。之后我会捣鼓出一个真正的首页——如果有人帮我做这个事情这样我只需呀关心编码我会非常高兴的。
要是有谁想帮KDE设计和维护首页，我很乐意让他这么做哦。:)

{ \noindent \tt 讨论：}

	最后，也是最重要的事情 :-) , 如果你们有兴趣就请加入我们的邮件列表  \email{kde@kde.org} 讨论吧！
	订阅邮件列表只需要发送一封如下内容的邮件到 \email{kde-request@kde.org} 就可以了。
	
\begin{center}
	{\tt subscribe 你的邮箱}
\end{center}


{ \noindent \tt 应用程序：}

	当KDE开始广泛使用的时候，新的程序有望基于Qt开发。获得和桌面环境协调一致的观感。

    我们很可能会把 LyX 移植到Qt，这样一个舒适的字处理软件就有了。不过这还在和LyX组成员讨论中。

    一个矢量绘图软件也很好。Xfig虽然丑了点，不过还是不错的。 tgif 也很强大。这些丑程序的作者并不喜欢添加 Qt 的界面，因为Qt是 C++ 的，而原作者只希望使用C——一些平台没有c++编译器。
    当然，KDE才不关系连c++编译器都没有的平台。（而且我认为 g++ 已经到处都有了。） 也许 tgif 作者赞成让别人来添加Qt界面 （反正源码在那里，不知道是不是GPL的。）。
    因为 tgif 没有使用第三方控件库而是全部自己实现，添加一个Qt界面应该不是什么难事。Qt里直接使用Xlib还是相当轻松的，所以并不需要重写很多代码。
    C++也兼容C，这样保护原有的C代码和是相当的容易。

    一个新闻阅读器怎么样？也能集成到KDE里。还有很多很多很多 ............

}

所以有趣的事情还是蛮多的。你要是感兴趣请加入我们吧！召集二三十人我们就可以开始了。可能平安夜前我们就能给自己一个礼物了。

KDE将使用GPL许可发布。

我承认这听起来有点疯狂。我可以告诉你，我是否非常认真的。每个和我讨论的人都说希望能有个简洁的桌面。Qt是实现这个目的的一个机遇。
让我们把业余时间利用起来开始干吧！

希望很有很多很多的回帖!
此致,\\
\qquad 敬礼！

Matthias Ettrich
(\email{ettrich@informatik.uni-tuebingen.de})

顺便说一下：通常这样的帖子会遭到这样的回复：“你要喜欢GUI就去用苹果好了”；“我就喜欢桌面程序上有各种不同风格的控件，你要是蠢的学习就去用Windows好了” ； “内存不够啊，我只用Motif”；
“你永远不会成功的，做点别的事情吧。”；“干嘛Qt? 我更喜欢用 schnurz-purz-widgets with xyz-lisp-shell, 还是 GPL 的！”； 等等。
非常感谢这样的回复，让我知道我还是有梦想的人 ......

好吧，又来了：你可能想问我干嘛那么讨厌 Tk。 恩，我不喜欢它的设计哲学： Tk 连个文本编辑控件都没有，却有一个缓慢的字处理程序。其他控件也差不多。
再同TCL语言结合起来，这个 Tk 程序就慢到家了（当然也有例外）。我还没有看到哪个 C/C++ 程序是用Tk库的，尽管Tk也包含C接口。TCL/TK 做点原型还行，写点内核图形配置工具倒不错，用来写应用程序就算了。
我压根就不想让哪个 Tk 程序待在桌面上。Qt就简单多了。试试看吧！

我确定是最后一段废话了：我和Troll Tech公司没有任何联系。我只是单纯的喜欢他们的产品（看看他们的代码，质量非常高！）和市场策略：对自由软件免授权费。

\end{insertnote}

耐心的看完{}Matthias Ettrich{}关于开始KDE项目的声明后能发现一些有趣的事情：
商业的桌面环境有，CDE , Looking Glas。 贵而且也不咋好用；多数程序都是直接使用Xlib开发，重复发明轮子。而且外观都无法统一；少数有几个帮助开发X程序的控件库，要么丑陋要么本身是个怪物；当时唯一可以扛担次重任的就只有一个Qt，虽然不是GPL授权的软件，好歹对开发自由软件是不收钱的。

KDE就在这样的环境下起步了。到本文写作的时间为止，KDE已经发布到 4.10了。

Gentoo下安装 KDE 环境非常简单。只要一条命令即可

\begin{code}
\#emerge kde-meta
\end{code}

这会将所有的KDE组件全部按照。依据提示可能需要开启一些 USE flag。将其添加到 make.conf 即可。更方便的莫过于直接启用 kde profile

\begin{code}
\#eselect profile  set default/linux/amd64/10.0/desktop/kde
\end{code}

这个 profile 会预先将 KDE 相关的 USE flags 全局启用。

如果不想把全部的 KDE 都安装（那意味着要编译数百个软件包），可以只安装一个基础的KDE运行环境。 

\begin{code}
\#emerge kdebase-meta
\end{code}

这样将只安装运行KDE所需的最小环境。不过{}Matthias Ettrich{}在声明里提到的邮件客户端，各种小游戏，新闻阅读器等等东西可全都没了。
不过，这不失为一个好办法，先折腾出KDE桌面，进去了再说，慢慢再编译全 :) 


\subsection{GNOME-Shell}

GNOME 是各个主流发行版的默认桌面环境。不过这个名声得来的有点不是那么厚道，起码不是靠实力获得的。GNOME无论是在易用性还是在可定制性上都离KDE差了一大截。GNOME能获得今天的成就，原因是
一开始就站在了道义的制高点牵制住了KDE。这是为何呢？

原因就在Qt上。{}Matthias Ettrich{}开发KDE选择了Qt。而当时的Qt并不是一个开源软件。Qt由一家叫做Troll Tech的公司掌握。Qt采取办法是的是对开源软件的开发使用免费授权，只有开发商业软件才需要购买商业授权。但是社区的一些人开始担心Qt的未来。
特别是RedHat这样的大公司，不希望Linux桌面被一家挪威公司掐住咽喉。于是 RedHat 伙同 IBM ，再拉上其他公司\textbf{资助}了另一个项目：GNOME。

没有大公司的资助GNOME从一开始就不会存在。所谓兵马未动粮草先行嘛。GNOME名义为社区项目，实际上参与GNOME开发的人有大量是Linux界有影响力的大公司的全职员工。
特别是RedHat派出了数量最多的人员参与了GNOME项目。因此在一些人当中流传GNOME是RedHat项目的说法。这当然是无稽之谈，不过GNOME获得大公司的资金和人员上的直接资助是事实。

GNOME为了避免KDE依赖于一个私有的Qt软件的弊端，选择了Gtk作为其GUI界面库。GTK是GIMP的副产品。全称为 Gimp Tool Kit。是为了开发GIMP而编写的一个UI库。同Qt一样是对Xlib的一次包装。
和Qt不同的是Gtk采用C语言开发。

相比采用不标准的C++语言, GNOME选择了同Gtk一样使用C语言进行开发。当时C++还不是一个标准语言，使用C++的开发者并没有C语言多。所以从一开始GNOME的选择就能吸引更多的社区开发者。

\begin{notice}
注意：C++语言的第一个标准是在 1998 年通过的，又叫C++98标准。第二个C++标准是2011年通过的，叫C++11。有人会提到2003年通过的C++03，那个是对标准库的一次扩充和修正，语言本身没有变化。
所以KDE创立于1996年，实际上是在使用一个没有标准的语言。标准要到2年后才有。
\end{notice}

由于没有协议的问题，GNOME获得了飞快的发展。从1997年立项开始没过多久，就迅速的变成各大发行版的默认桌面环境。

GNOME的目的是提供一个完全自由的桌面。通过使用完全GPL的Gtk进行开发，GNOME规避了KDE可能遇到的潜在风险。为了应对这件事，Troll Tech宣布Qt以GPL许可发布。Qt正式开始了GPL和商业授权的双授权模式。可以说，Qt采用GPL协议后，GNOME一下子失去了靶子（重新创建一个完全自由的桌面。）。而Qt最后被Nokia收购的时候，更是采用了LGPL协议发布，连开发商业软件都不再需要购买授权了。

重重危机下，也许是GNOME推出Gnome-Shell寻求差异化的原因吧！

GNOME-Shell 和其他桌面有一个显著的不同点：GNOME-Shell即是窗口管理器又是面板。传统上，窗口管理器和面板是2个独立的进程。但是GNOME认为窗口管理器的任务是管理窗口，面板的任务之一
就是启动和切换窗口。窗口管理和面板是桌面环境非常重要的两个部件，最好还是纳入一个进程统一指挥吧！于是就有了GONOME-Shell。

GNOME-Shell一开始就提出了一个其他桌面项目想都不敢想的事情：强制要求机器支持OpenGL加速！ 这可真够大胆的。GNOME基于如下考虑：

\begin{enumerate}
\item 为了用户体验，更为了让应用程序开发者能专注于一个平台，强制要求支持OpenGL加速，使得为Gnome设计的程序可以假定硬件加速总是存在的而无需担心没有硬件加速。简化应用程序的开发。
\item 混合管理器一直存在，是的窗口“半透明效果”总是能实现\footnote{窗口的透明度是一种指示给窗口管理器的属性。只有带混合功能的窗口管理器才能实现半透明效果。半透明效果无法在应用程序自身实现。}。类似的必须依赖混合管理器的功能还有“窗口阴影”、淡出淡入。可以保证一致的用户体验。
\item 启用用户体验更好的桌面效果。并保证任何Gnome桌面都有这些效果，带来用户体验的一致性。
\item 强迫上游开发者修复硬件驱动问题。其实这个也很重要。GNOME-Shell刚推出的时候把原本很多没有暴露的驱动问题都暴露出来了。而且驱动开发者因为GNOME-Shell也显著的提升了驱动的质量。
不管是开源驱动还是闭源驱动，都因为GNOME-Shell的到来大大的提升了驱动的质量。

\item 话说，现在硬件上不支持 OpenGL 的电脑已经是古董中的古董了吧。基本都是驱动问题导致的不支持硬件加速。所以说，GNOME让社区想明白，专注于眼下需要支持的硬件。减少无畏的人力资源浪费。
\end{enumerate}



\section{办公和打印}

PC 称为通用平台一点没错，装什么就能干什么。装套Office软件就能拿来办公了。
即便不办公 ，平时也有码字的需要吧，也有把码下来的字变成实体的需要吧。


\subsection{LibreOffice}

在商界，最鼎鼎有名的office软件当属 MS Office了。但是，MS Office 不开发的格式，劫持了很多人。你想想看，当你有一天不花钱就无法打开自己辛苦编写的文档的时候，是什么样子的心情？



\subsection{使用CUPS打印文档}

\section{声音的那些事}
声音是一种波，爱迪生发明了留声机，利用唱片表明的凹凸起伏记录声音声音。
贝尔发明了电话，利用电流的变化通过导线传递声音。

磁带则利用磁带上磁性物质被磁化后的磁场强弱来记录声音。

\chatu{soundwave}{声音是一种波}

但是这个时候，声音是模拟的。直到计算机的发明，声音开始被数字化存储。
数字信号抗干扰，无失真，可以被重复读取，而且每次读取的结果都是一模一样的。
不会像唱机那样，越放声音越失真。数字化的声音，在同样的回放效果下，大大降低了对器材的保真度的要求。

\subsection{声音的数字编码}

要数字化声音，首先要进行编码。声音通常使用 PCM 编码。

做法是：将模拟的声音的电信号每隔 (1/采样率) 秒进行一次采样，即测量当前电压值，然后把电压值用一个数字表示。
可见，采样率率越高，用来描述采样的点的位数越高，对原始波形的描述就越精确。
如图 \ref{fig:soundwavepcm} 示意，越高的采样率和量化精度，表示图中的格子越小，采样结果越逼近原始波形。

\chatu[width=9cm]{soundwavepcm}{声波的采样和量化}

PCM 编码有2个指标：采样率和量化深度。指标越高音质还原越好。CD 的规格是每秒进行 44100 次采样，每个采样点用16位的数字表示。这两个规格的来源是因为一个“采样定律”说采样率要比最高频率高两倍方可完整采样。而人耳只能听到最高 20000hz 的声音。至于采样率为啥是  44100 而不是  40000 主要是为了给系统留有一定的余量，并且 44100 能同时被 50 和 60 整除\footnote{欧洲的电网运行频率是 50Hz 而美国的电网运行频率是 60Hz。}。 而使用 16bit 量化而不是 8bit 的原因是 8bit 实在太粗糙， 32bit 又太浪费（制定 CD 标准的年代，存储设备并不像今天那样廉价）。于是折中方案就是 16bit 量化 44100 采样率。

随着技术的进步，人们开始不满足于 CD 的规格，于是提出 24bit 量化和 192000 采样率的规格。未来还会有 384000 甚至更高 采样率规格出现，这些越来越高的规格，使数字音频能越来越忠实的记录和还原声音。

除了 PCM 编码，另一种数字化声音的办法被称为 “DSD”, 不像 PCM 要读取一个完整的采样点，DSD 利用 0 和 1 的密度变化来表示声波。也就是在声波波形高的地方 ， 数字化音频流中1 出现的要比 0 多， 反之 0 出现的比 1 多。 DSD 的这种编码被称作 "脉冲密度调制"。

DSD 工作在高达 2.8224 MHz 的采样率上，是标准 CD 的采样率的 64 倍。以 DSD 技术制作的CD叫做 “SACD(Super Audio CD)” 意思是超级 CD。

DSD或者PCM编码方式之间孰优孰劣有许多争议。而 Sony 牢牢的控制着DSD的专利，对DSD的推广极为不利。

包括 PC 在内的许多设备，压根就不是为DSD准备的。USB Audio 2.0的声音定义只有 PCM ，许多操作系统也只能处理PCM编码的声音。在电脑上播放DSD编码的声音需要转换为PCM编码方能播放。


\subsection{ALSA高级Linux声音架构}

要用PC播放声音，软件硬件的支持少不了。硬件上，电脑需要具备一块 “声卡”。声卡的作用就是将数字化的声音信号逆向转化为模拟信号。旧时代，声卡是一块 PCI 板卡插在主板上，用户需要额外购置。随着电子技术的发展，如今声卡都是集成在主板上的。

有了硬件条件，就要谈到软件上了。 软件上，除了需要一个播放器程序，要想声卡正常工作，还需要内核的支持。

最早出现于  Linux 上的驱动框架是 OSS， Open Sound System。  但是多年前就已经被彻底废弃。
目前 Linux 使用的是 ALSA 声音架构。 ALSA 全称 Advanced Linux Sound Architecture, 高级Linux声音架构。ALSA 分内核驱动和一个用户空间的库两个部分组成。

\chatu{kernelconfig5}{内核编译选项ALSA}

要使用 ALSA 需要在内核里启用 ALSA 选项，并且启用对应的声卡驱动。对于集成声卡，和 插于 PCI 插口的声卡，统统都是 PCI 声卡，要启用 PCI sound devices 然后进一步选择需要的驱动。

一般主板集成声卡，都是使用的 Intel HD Audio 规范，也就是说，只要如图 \ref{fig:kernelconfig6} 所示启用 Intel HD Audio 即可。

\chatu{kernelconfig6}{启用　HD Audio 声卡驱动}

如果不清楚需要为 HD Audio 启用什么codec 支持，就把里面的所有 codec suport 选项启用。

对于 USB 声卡 （一些 USB 转同轴设备，其实就是一块 USB 声卡，只是没有模拟输出，而是直接通过同轴输出数字信号，通常需要配合一个带同轴输入的功放或者DAC。还有 USB 发射器的无线耳机也是一块USB声卡。）因为有 USB Audio 规范，几乎所有的USB声卡通用一个驱动的。所以只需要选上 USB sound devices 即可。启用该选项后，会出现多个子选项，如果通用驱动无法识别，则需要使用这些专门的驱动程序。

启用 ALSA 驱动后，还需要安装 ALSA 的库。

\begin{code}
\# emerge media-libs/alsa-lib
\end{code}

其实无需直接安装，通常 alsa-lib 会被作为其他程序的依赖被安装。

\subsection{PulseAudio优秀的声音服务器}

诶？ 有了 ALSA 为啥还有 PulseAudio 呢？ PulseAudio 是干嘛的呢？

说道 PulseAudio 不得不先说声卡占用问题。对于没有硬件混音功能的消费级声卡来说，硬件同时只能支持一个回放流。也就是说，两个程序无法同时发声。如果一个程序已经占用了声卡，另一个程序就无法播放声音。除非你的声卡硬件上支持多个回放流，而这通常是不可能出现在消费级产品上。

但是，使用Windows这样的系统我确实可以同时听歌和看电影！系统的一些提示声音也能同时播放！好吧，我见到别人的 Linux 系统(如Ubuntu)也支持多个程序同时播放声音。
这是怎么回事？ 他们用的是很专业的声卡么？

其实不然，要让这一切成为现实，只需要软件混音。操作系统提供软件混音功能，支持任意多个程序播放音频流，系统将他们混音成一条，然后再统一使用声卡输出。

而 Linux 下，这个混音任务则交给 PulseAudio 完成。

\begin{notice}
\subsubsection*{为什么不在内核里混音}

一些系统(FreeBSD)确实可以在内核里执行混音，但是 Linux 采取了将混音交给用户空间的策略。 因为首先混音并不需要内核权限，放到内核里只能增加内核的复杂度增加bug而没有明显好处，其次混音通常有多种算法可以进行，放到内核里也不便于用户更换混音算法。

另外，即使是　Windows 也并没有在内核实现混音。现实意义上说，在内核实现混音只有一个好处，就是保证混音功能随时可用，而无需额外启动一个守护进
程。实在无需为了一些人的 “进程列表洁癖” 大费周张。一些人会说到效率问题，首先，即使是未压缩的音频数据，以当下的PC处理能力来说其数据量也是非常小的，不值得一提。
即便如此，混音进程也可以使用很多避免拷贝声音数据的技术。

\end{notice}

使用 PulseAudio ， 需要应用程序把声音数据交给PulseAudio 而不是交给ALSA。 用户需要在程序里设置输出为 PulseAudio而不是ALSA。绝大多数在启用  pulseaudio (通过 USE flag pulseaudio 启用，建议全局启用) 后，
就将  PulseAudio 设置为默认输出。基本上无需配置。

而对于一些只支持 ALSA 输出不支持 PulseAudio 的程序呢？ PulseAudio 也提供了解决办法：通过劫持 alsa 将数据重新导回 PulseAudio\footnote{虽然我在这里使用了劫持这个词语，其实这是 ALSA 本身就提供的插件机制。}。 因此只支持 ALSA  的 Flash 插件也能使用 PulseAudio 输出。

有的人说没有 PulseAudio ，只靠 ALSA 的一个 dmix 插件，一样能实现混音。他们当然没有注意到 dmix 插件的混音质量非常低，而且已经停止开发，被 ALSA 开发者标记为不推荐使用了。

过去 PulseAudio 还不完善的时候， Ubuntu 就开始使用 PulseAudio,  导致用户面临很多声音方面的问题不得不通过卸载PulseAudio解决，
给人留下了 PulseAudio 非常不稳定的印象。实际上现如今 PulseAudio 非常稳定。而且使用 PulseAudio 用户可以在多个声卡间无缝切换。如将正在播放的声音从内置声卡转移到蓝牙耳机而不需要重新开始播放。顺带一提，PulseAudio 是让你使用蓝牙耳机的唯一途径。

启用 pulseaudio 这个 USE flag 后， pulseaudio 会被作为依赖安装。
pulseaudio 基本上无需配置， 开箱即用。

\subsection{对 PC HiFi 的一些建议}

PC HiFi 即用 PC 作为播放设备打造的 HiFi 级音乐欣赏设备。本节的开始介绍了一些声音编码方面的知识，就是为 HiFi 准备的\footnote{不HiFi，了解这些干嘛，能发声就可以了。（笑）}。

HiFi，意思是 High-Fidelity， 直译高保真。 HiFi 追求非常准确的还原声音。 一般录音阶段是在录音室完成，因此对普通人来说HiFi 就是准确的回放音源上记载的声音。

在 HiFi 领域，声音的回放从 PCM 音源开始，而不管此 PCM 音频流是从CD读的还是从网络下载的，还是硬盘上的，还是从其他格式解码来的。

\begin{list}{\textbullet}{}
\item PCM 音频流，首先进入 DAC 设备，转换为模拟信号。HiFi 要求，这个DAC设备能尽可能的准确还原 PCM 流表示的波形。这一步，会引起失真，主要是两个原因，一个是 DAC 设备的基准电压不稳定，另一个则是时钟不稳定。电压不稳定，就不能准确的还原采样点的电位，而时钟不稳定，就不能准确的还原采样点的时间。在 HiFi 里，时钟的不稳定被称为 jitter。

\item 经过 DAC 转换后的模拟信号，首先进入前级放大，将信号的电压大幅度提高，接着送入后级放大，最后输出高电压大电流的音频信号。这一步使用功放完成。HiFi要求输出信号能尽可能的保证和输入信号一致，除了电压和电流大点，波形上要尽可能的一致。在这一步引入失真，便无法称为 HiFi 功放。

\item 放大后的音频信号推动扬声器发声。HiFi 要求扬声器的动圈尽可能的遵循输入信号进行运动。

\end{list}

使用 PC 进行 HiFi 回放，主要就是利用 PC 作为音源，并继续使用传统 HiFi 领域的功放和扬声器设备。

和传统的 CD 或者黑胶唱机相比， PC 具有更大的电磁干扰。 因此，集成声卡的 DAC 通常会引入大量的电磁噪音，成为整个系统的短版。于是，PC HiFi 的人无一例外的，都使用外置 DAC 来避免机箱内部严重的电磁干扰。

使用外置 DAC 有两个主要途径，一是使用声卡的数字输出接口。通常是 S/PDIF 接口。 S/PDIF 接口有光纤和同轴两个版本。
使用光纤或者同轴将电脑同DAC设备连接起来。一些功放集成了S/PDIF输入功能，可以无需额外购买 DAC 设备。
另一个途径是使用 USB 或者 IEEE1394 火线接口的声卡。

软件上，PC-HiFi 要求 bit transparency， 即 DAC 设备收到的 PCM 数据就是原始音频媒体里存储的数据。 即要求播放器不对声音进行修饰，也要求操作系统不要对声音进行修饰。 操作系统的软件混音和音量调节功能是常见的破坏 bit transparency 的特性。简单做法是将音量调节到最大，使用功放上的音量控制。尽量保持播放器为唯一运行的声音程序从而避免混音。

\subsubsection*{S/PDIF输出设置}

设置方式一：

需要安装 pavucontrol 。 执行 pavucontrol 打开设置界面，点选  “配置” 选项卡
然后将声卡的工作模式设定在 Digital Stero (HDMI) Output。

设置方式二：

安装 alsa-utils。 执行 alsamixer -c  0 ， 注意直接使用  alsamixer 调节的是 pulseaudio 这个虚拟输出， 要使用 -c 0 来选定第一张声卡，如果有多个声卡，还可以 -c 1 -c 2 来选择其他声卡。

\chatu{alsamixer}{alsa配置S/PDIF}

然后光标移动到 S/PDIF 选项上，按键盘上的M键，把状态从 off 切到 on， 如图 \ref{fig:alsamixer} 所示。

\subsubsection*{使用纯数字功放}

晶体管具有电流放大作用。直接利用晶体管的电流放大作用可制成模拟功放。依据类型又分 A 类功放， B 类功放和C 类功放。

晶体管对直流电进行放大，而模拟声音信号为交流电。A类功放使用单个的晶体管进行放大，故需要一定的偏置电流。输入声音信号负半周的时候，输入给晶体管基极的是减小了的偏置电流，输入声音信号正半周的时候，输入给晶体管基极的是增大了的偏置电流，  在没有输入声音的时候，晶体管仍然在工作，对偏置电流进行放大。

故而 A 类功放工作效率低下。于是发明 B 类功放。
B 类功放使用一对晶体管进行放大。信号的正半周和负半周分别使用1块晶体管进行放大，无输入的信号时两个晶体管都不工作，大大提高了电源的使用效率。
B 类功放的特点就是需要使用2个放大倍数一致的“对管”，
并且是一个PNP型晶体管，另一个是NPN型。

然而晶体管的放大有一个最低工作电压，大约是 0.7v。 低于 0.7v 的输入电压无法放大。因此在 $\pm$ 0.7v 范围的信号将无法放大，产生“交越失真”。
为了解决这个问题，给2个晶体管各提供一个偏置电压消除交越失真即 AB 类功放。不过 C 类功放在无输入情况下也会消耗一个小电流，效率没有 B 类功放高。

不管是 A ~ B 还是 AB 类功放，都对晶体管的线性工作区间要求很高。
但是晶体管的放大是非线性的(对不同的输入电压放大倍数是有差异的。并非等比放大)， 制造线性度很高的晶体管成本很高(这些线性很高的晶体管专门用于放大，得了个专门的名字功放管，从几元一个到几十元一个不等。而普通晶体管的价格则是按照分厘来算的。)，外加这些功放的效率很低，这些年发展出了 D 类功放。

D 类功放的特点就是让晶体管工作于 开关 两种状态而不是工作在放大状态。因此对晶体管没有线性要求，任意晶体管都能实现放大。
成本非常低，而对电源的利用非常高，理论工作效率为 100 \%。

D 类功放首先将输入的音频信号转变为高频 PWM 信号。以 PWM 信号的占空比表示输入信号的电压。再用 PWM 信号控制开关管直接调制电源。于是就能获得放大了的 PWM 信号。放大后的 PWM 信号经过一个滤波器即输出放大后的音频信号。

影响 D 类功放的音质的主要2点，PWM 信号的生成和输出滤波器。而所选用的晶体管只影响到电源的利用率。晶体管导通电阻越小，在晶体管上消耗的能量就越少，效率越高。

PWM信号的频率越高，对输出滤波器的要求就越低，声音质量就越好。

而纯数字功放，直接使用PCM数字信号生成 PWM 信号。免去了DAC带来的损失，免去了模拟转PWM带来的损失。理论来说上是最优秀质量最好的功放。是近年来兴起的新兴模式。纯数字功放因为绕过了 DAC 和传统的功率放大，以数字信号调制电源推动扬声器发声，所以又得名功率DAC。

\section{Wayland}\label{sec:wayland}

X确实是一个优秀的架构。如果说他有什么缺点，那就是它老了。真的老了。单是第11版\footnote{1987年发布的第11版X协议。}X协议也已经拥有超过25年的历史了。很多当年没有考虑到的、
或者不能考虑的技术，现在都已经变成理所当然普遍应用了\footnote{动态链接库和廉价的图形加速芯片。}。

\subsection{X的局限和扩展}

随着Linux桌面的发展，X暴露出来的第一个局限性就是硬件加速。硬件加速通过DRI比较完善的被解决了。DRI绕过了X11协议让图形程序直接操作图形设备。给X11开了第一个洞。
（不过 DRI 只用来实现了 OpenGL， 并没有将一般的X绘图API交给DRI实现。）

Linux桌面的第二个发展是出现了“混合窗口管理器”。混合管理器和DRI天生就是一对矛盾体。

\begin{insertnote}

混合管理器的实现依赖于X的一个特性：窗口管理器可以指示X Server将一个窗口的绘图操作进行重定向。混合窗口管理器将普通程序的窗口操作重定向到位图上，然后将不同窗口的位图图像“合成”到屏幕上。因为一般窗口的绘制依赖于X协议中的绘图指令。
X Server很简单的就能实现绘图操作的重定向。可是遇到了DRI这个事情就范了难。一个使用OpenGL直接渲染（ 间接渲染依然使用X协议来传输OpenGL命令，所以和混合管理器不冲突的。）
的程序，其窗口的内容是由它自己直接绘制的，并没有经过X Server。 这样混合管理器按照道理来说是不能获得直接渲染的窗口的内容的。就算能获得，混合管理器也必须知道窗口内容“变动”，以便更新屏幕内容。
跟糟糕的是直接渲染的窗口并不能使用重定向功能！
\end{insertnote}

为了解决这个矛盾问题，X11加了一个Composite扩展。允许窗口管理器追踪使用直接渲染技术绘制的窗口变化。并且允许直接渲染的窗口使用重定向功能绘制到一个缓冲区上。

这就是DRI这个洞的补丁了。

随着Linux发普及，人们自然希望程序的“字体”能漂亮些。原有的X协议定义了一套核心字体。但是第一不支持UNICODE，第二不支持矢量字体。于是X又修补了这个漏洞：添加Xft扩展。

Xft允许在X中使用truetype字体。不过在X中配置truetype字体相当的麻烦。首先要在 /etc/X11/xorg.conf 中配置存放字体的文件夹路径，然后每个字体文件夹下都要 font.dir 和 font.scale文件。更要命的是，添加了字体需要\textbf{重启X}才能使用。所以这个补丁补的也不怎么好，导致大家更乐意在程序中自己管理和渲染字体，而不使用Xft扩展。

接下来X打的一个补丁理解起来就比较费力了。这个扩展名字叫 XRender。

XRender的出现源自两个事实：一是应用程序越来越不喜欢直接使用Xlib中“点”、“线”等基本图形的绘制API去绘制复杂图形了，而是采取直接“贴”图的办法。这个图可能是预先绘制好的图片（图标等资源），
也可能是其他软件生成的——比如FreeType2库渲染而成的文字，还可能是其他的绘图API生成的图片——比如cairo。二是贴图的时候希望能使用Alpha通道合成——半透明效果。

X在设计之初就考虑到了各种绘图需求。除了提供点线面绘图之外，还允许使用像素图。X程序可以使用自己任何最合适的方式生成象素图片，然后贴到窗口上\footnote{XPutImage(3)}。
这样单纯的点线面不足已绘制复杂图形的时候还有这一个API可用。
现代的X程序已经不使用X提供的 XDrawString(3) 绘制字体了，而是使用Freetype2这样的库，在内存中绘制要一个字符的像素图然后贴到窗口上。
不过问题就出在这里。
我们知道矢量字体要想绘制的好看必须开启“抗锯齿”功能。字体的边缘出现灰度化的抗锯齿效果。如图 \ref{fig:Rasterization-antialiasing}。

\chatu{Rasterization-antialiasing}{抗拒锯齿字体示例}

因为抗锯齿，线的边缘一些像素就要“灰度化”。灰度化本身并不会导致Xrender的催生。关键的原因是：这个灰度化并不是颜色灰度化，而是透明值。

一个只是灰度化的黑色字体显示在红色色背景上，效果就如图 \ref{fig:Rasterization-antialiasing-no-alpha-blending} 一样，灰度化的地方应该使用红色和黑色的混合色，但是这里却是白色。

\chatu{Rasterization-antialiasing-no-alpha-blending}{没有alpha混合示意}

在纯白的背景上以抗锯齿的方式绘制一条黑色线段。边缘像素会灰度化。
事实是，这个灰度的像素仍然是黑色的，只不过增加了透明度，只是在白色背景下看上去是灰色颜色。在非白色背景下，又怎么能是灰色呢？
使用透明值而不是像素本身个的颜色，在非白色背景下才能混合为正确的颜色。

\chatu{Rasterization-antialiasing-alpha-blending}{alpha混合的抗锯齿示意}

也就是说，绘制带透明通道的图片的时候，是需要执行“Alpha混合的\footnote{Alpha混合是指两个像素混合为一个的时候，被叠加的像素使用叠加像素的透明值混合成新的颜色。而非Alpha混合，就是使用新的像素覆盖掉老的像素。}”。老式的X协议并没有提供Alpha混合。XPutImage是不接受带Alpha通道\footnote{计算机图形学中，称保存像素透明度的存储空间。}的图片的（或者说忽略Alpha通道。）。

%那么，知道抗锯齿是在边缘像素使用透明值实现的。如何绘制呢？就拿字体渲染为例。字体通常是绘制于白色背景，但是也有例外。所以绘制好的字体图片是包含透明通道值的。背景部分全透明，边缘部分像素是半透明的。接着应用程序将渲染好的字体图片贴到窗口上（通常来说，背景是白色的）。

应用程序有一个绕过这个限制的办法，就是使用XGetImage(3)首先获得目的像素，然后自己执行Alpha混合，最后将混合结果交给X。

听上去是一个解决办法，这也正是Xft使用的办法。但是问题就在X是网络协议，这样会导致X程序和X Server之间来来回回的交互数据，在网络延迟很大的时候，有严重的性能问题。

解决这个问题的补丁就是XRender。XRender扩展给X Server添加了Alpha混合功能。

XRender出来是补XPutImage的缺陷的，而XPutImage本身就是用来补 X 原自80年代的绘图方式（点，线段，椭圆，圆形，矩形）的。
如今这种80年代绘图形式现如今已经没有哪个X程序还在用了。应用程序通常使用Cairo（或者相同功能的库）这样的矢量图形库。
Cairo提供非常丰富的API用来绘制图形。而Cario并不使用X的绘图API，而是自己进行光栅化，最后使用XRender绘制到屏幕上。

\subsection{渲染库和静态链接}

X是一种网络透明的协议。为了能实现远程显示，X除了定义一套窗口管理规则，还定义了一套绘图API。X的很多局限性都是从绘图这个功能层面开始暴露的。如果X只包含了窗口功能而没有绘图功能，事情会变得简单很多。

如果X没有渲染功能，渲染可以借助 OpenGL 这样的库进行。OpenGL的一个强大的优势就是它是硬件加速的。有人可能觉得简单的二维图形怎么使用OpenGL，那就大错特错了。OpenGL并不是3维专用的库。
2D可以认为是3D的一个子集。Clutter这样的一个库即是使用OpenGL构建的GUI库。Clutter本身就足以证明OpenGL绝对足够担当GUI的二维界面绘制的任务。

除了OpenGL，还有很多软件渲染库存在。cairo就是这样一个。所以，就算将图形绘制功能完全从X移除，GUI一样能跑的很好。事实上这件事情已经发生了。X目前为止还在被使用的绘图API只剩下了一个：XRender。
其他的API包括X提供的字体渲染都已经没有人再使用了。我更乐意将XRender定义为Buffer提交。由应用程序使用各种绘图库绘制界面，然后将绘制好的图形提交给X Server。

那么X就只有窗口管理功能还在被继续使用了。不过，如果将渲染功能从X中移除，是做不到的。移除了就不再是X了。要声称自己是X Server就必须支持已经没有人使用的功能。这给X Server的维护带来了严重的负担。我们不禁要问：

\begin{center}
X协议设计之初为何要包含绘图功能呢？
\end{center}

这个问题的回答还得回到80年代初。没错，当时的一个技术限制导致了X不得不设计成 C/S 结构并在核心协议中包括绘图操作。

仔细回想一下，一个渲染库要包含各种绘图指令，并且有可能的话还能利用使用硬件加速，必定是一个庞大的图形库。如 310.19 版本的NVIDIA驱动，其OpenGL实现\footnote{/usr/lib/libnvidia-glcore.so.310.19}庞大到有35MB之巨！ 
虽然这个例子有点极端，但是这正证明了一点，图形渲染事一件复杂的事情。复杂的事情就需要更多的代码实现。

虽然 X 是个窗口系统，但是从代码数量上看，管理窗口用的代码远远少于窗口渲染的代码。

那么回过头来，80年代初的一个技术限制，就是“代码共享\footnote{其实就是动态链接库啦！}”技术远未普及到今天这个地步。缺乏“代码共享”技术，意味着所有使用图形界面的程序执行期
都必须在内存中载入一份图形渲染代码——这就是静态链接的弊端。这是对内存和磁盘空间的极大浪费。其次，如果渲染库升级了，使用它的程序必须全部重新编译链接以享用新的库带来的改进——功能上的和性能上的。更糟糕的事情是，如果你的程序要放到不同的硬件上执行（使用了不同的显卡）就必须要将为不同显卡编写的代码全部包含进去！最糟糕的是已经编译好的程序将无法使用新的硬件！

想象一下：libGL 这个库是由显卡驱动提供的。如果这个库是静态链接的，意味着你要么为不同的显卡链接不同的库生成不同的程序为客户使用\footnote{话说当年DOS平台的一些游戏就是这么干的。要分不同的显卡插入不同的磁盘进行游戏。}，要么就得找到一直办法同时将所有的厂商提供的OpenGL库链接到同一个程序中。


在没有动态链接库的时代，要想共享渲染库就只有一个办法了：使用C/S结构。C/S结构下，应用程序链接到一个比较小的“转发库”，转发库将所有的绘图请求通过
网络协议发送到一个“渲染服务程序”，由渲染程序代为执行。只要协议统一了，可以为不同的硬件平台编写不同的渲染程序。而一般的图形应用程序就可以不需要重写。
这个转发程序比较小，所以静态链接也没什么问题。听起来是不是很熟悉？\textbf{这也正是X所使用的结构！}

\subsection{混合}

即便是说“X本身的渲染功能，不用就不用吧”，还有一个麻烦事情摆在X的面前。

\chatu[scale=0.7]{x-architecture}{一个X事件的处理}

顺着图\thefigure{}标记的序号，一个输入设备事件首先被内核的evdev驱动层捕，然后发送到X Server。接着X Server\textbf{判断这个事件应该属于哪个窗口}并将事件转发给相应的X Client，如图中\textcircled{\oldstylenums{2}} 所示。接着X Client对事件响应，绘制界面变化来做这一响应，\textcircled{\oldstylenums{3}}所示。

本来X Server这个时候只需要再将X命令解析了执行了，把绘制好的图形叫到显卡即可。
但是因为混合窗口管理器的缘故，这个时候X Server还不能这么做，它要给混合窗口管理器发送通知告诉它窗口内容变更了，如图中\textcircled{\oldstylenums{4}} 所示。
混合窗口管理器获得窗口内容更新事件，要重新向X索取新窗口的内容，然后执行混合（Composite）。
接将合成的整个屏幕的内容交给X Server，如图中\textcircled{\oldstylenums{5}} 所示。
最后X Server将混合窗口管理器绘制的屏幕内容呈现到屏幕上。

同一个内容被绘制了三次！第一次是在Client绘制的时候，该绘制被重定向到离屏缓冲区，接着X Server告诉混合窗口管理器代表该窗口的缓冲区内容变化了。混合窗口管理器读取缓冲区数据，将所有的窗口的内容执行一次混合，这是第二次绘制。
第三次绘制是窗口管理器将最终合成后的全屏内容交给X Server，X Server再次将该内容绘制到根窗口上\footnote{根窗口是一个全屏窗口，由X启动后自动创建。}。
%也就是说，X使用C/S结构并不是为了网络透明，而是没有动态链接的时候的选择。

在这里，还有一个潜在的问题暴露出来了。所有的窗口事实上都是隐藏的！你看到的窗口实际上并不是窗口，而是窗口管理器用实际窗口内容绘制的一个矩形区域。
窗口管理器创建了一个全屏的窗口，将所有其他窗口都重定向了。以致他们的内容是不会呈现于屏幕上的。一个窗口被用户看到的唯一途径就是被窗口管理器混合后绘制于它创建的全屏窗口上。

也就是说，如果窗口管理器对窗口做了一些“变形”处理...... 那么 X 将不会知道一个鼠标点击事件真正应该发送给谁！

所以一旦窗口管理器这么做了，就必须接管X的事件判断。窗口管理器向X Server申请将所有的鼠标事件都揽给自己，然后自己进行判断鼠标事件到底该属于哪个窗口。再将该鼠标时间发送给对应的窗口（XSendEvent(3)）。

等等，听上去感觉混合窗口管理器本身就已经变成了一个 X 了呀，除了它自己没有携带显卡驱动。
事实上也正是如此。现如今，混合窗口管理器接管了越来越多的X本身应该做的事情。
为了执行正确的混合，窗口管理器也在内部维护了窗口位置关系。
为了“被变形的窗口”正确收到鼠标事件，窗口管理器也接管了鼠标事件判断。
不用说键盘焦点本身就是窗口管理器维护的。X的输入功能已经沦落为一个“形式过程”，真正的维护都交给了窗口管理器。

顺便一提，窗口绘制（且不说窗口系统根本就不应该包含这个功能）这个功能X也早就保不住了。
就连最后“将多个窗口的内容正确的执行‘裁剪’最终呈现于屏幕上”这个功能，也早就被混合窗口管理器接管了。

X 剩下的唯一能做的事情,就是一个消息中转站：用来连接窗口管理器和X Client。

说到这里，有人觉得，既然X所有的事情越来越多的被窗口管理器接管，并且有的事情也通过“直接渲染”不再需要X Server进行而是X Client直接进行了，那么让X Client直接和窗口管理器通信，踢开X Server怎么样？

\subsection{窗口管理器做为Display Server怎么样}

这个“让窗口管理器成为X Server”的想法，最终的结果就是导致Wayland项目的诞生。

\chatu[scale=0.7]{wayland-architecture}{Wayland结构}

在这个架构中，因为移除了X Server，同时也意味着移除了X Client通常使用的渲染方式。但是Wayland Client可以使用在X时代就存在的一直渲染方式：直接渲染。
应用程序使用 OpenGL 或者它认为最合适的库进行直接渲染。
然后将渲染好的窗口内容——可共享的缓冲区——通过wayland协议提交给compositor。传输的是一个共享缓冲区的引用\footnote{需要显卡驱动的支持。目前开源的显卡去的都支持，不过闭源的NVIDIA和ATI驱动没有提供支持，只能可以使用共享内存来共享这个缓冲区，有性能损失。}，而不是数据本身，所以非常快，无需当心性能问题。compositor直接使用这个缓冲区执行合成然后将最终结果直接绘制到屏幕上。

Wayland是一个纯粹的窗口管理协议，没有渲染功能。应用程序需要使用其他的渲染库绘制窗口内容。然后将绘制好的窗口内容提交给Compositor。

Wayland被誉为Linux桌面的未来。Wayland移除了X Server同时也移除了X的C/S结构带来的性能损失。Wayland结构下，窗口管理依然是C/S结构。由Wayland Compositor负责管理窗口。
但是最影响性能的渲染部分，则完全使用了“直接渲染”技术，彻底摆脱X的性能禁锢。

Wayland在桌面上还是一个非常试验性的项目，但是在嵌入式领域，已经开始应用了。影响Wayland在桌面环境应用的主要原因，就是驱动和旧的X程序。
Wayland通过XWayland来兼容旧时代的X程序。而驱动问题则是任重道远。

\begin{insertnote}
\subsubsection*{XWayland}

没有规定X Server必须以真实显卡硬件作为输出吧，如果X Server本身成为Wayland Client呢？由于Wayland使用的是直接渲染，意味着X实际上还是直接控制显卡执行渲染的，所以旧有的X驱动只需要稍微进行修改即可支持Wayland。而且不存在性能损失。
对于驱动没有Wayland支持的情况，还有 “shmwl”这个伪X驱动。使用软件渲染然后使用共享内存提交给Wayland。

\end{insertnote}

Wayland需要驱动支持，因为旧时代X使用的是用户空间的驱动。而Wayland不能使用X的驱动。所以所有的X驱动都不能使用。幸好DRI将大部分驱动拆分成了DRM内核模块和Mesa两个部分，
所以Wayland目前只能使用有DRI支持的显卡：也就是 intel 显卡。NVIDIA和AMD的显卡必须使用开源驱动才能使用Wayland。以为前文说过闭源驱动虽然也支持直接渲染，但是并不是DRI结构的。

目前Qt 和 Gtk 的已经实现了Wayland支持。应用程序稍作修改甚至不做修改即可支持Wayland，只要驱动问题解决了，Wayland还是能够普及开来，替换掉过时的X的。

\begin{insertnote}

\subsection*{Mir vs Wayland}

鉴于 wayland 项目进度太慢\footnote{Wayland 的发起者不是 Canonical 雇员， Canonical 无法干预进度以到达他们的预期。} ，ubuntu 的幕后公司开始投资于另一个 X 替代品 —— Mir。

Mir 采取同 Wayland 类似的基本架构，不同的是协议改用 protobuf 实现，整个项目采用 C++ 开发。

Canonical 在 Mir 的页面上解释说， 是因为 Wayland 无法满足他们的要求他们才开始开发 Mir 的。
而事实上 Canonical 列出的 wayland 的缺点都不存在，只是他们没有正确理解 wayland。 
被指出后， Canonical 迅速的更正了对 wayland 的错误攻击。

Canonical 的做法引起了轩然大波。部分开发人员指责 Canonical 通过创建 Mir 而不是帮助 wayland， 将 Linux 的桌面环境进一步分裂。
KDE社区的开发人员 更是指责 Canonical 故意错误的攻击 wayland 的缺点误导群众。有的 KDE 开发人员甚至直接表示 KDE 不会支持 Mir。\footnote{Martin Gräßlin 后来向大家澄清，不支持 Mir 不是因为 Canonical 伤害社区的做法，而是因为 KDE 的政策就是不会提供只能为一个发行版工作的功能。如果 Mir 获得 ubuntu 以外的发行版的接受（显然不大可能），KDE 还是会支持 Mir 的。}


但不管怎么样，Mir 的兴起对 Wayland 产生了积极的影响。Wayland 的开发进度在 Mir 发布后明显加快 —— 虽然 Wayland 开发者拒绝承认这是因为 Mir 的出现。



\end{insertnote}\fi

\chapter{网络}\label{chap:network}

\begin{quotes}[0.5]{Bill Gates}
I see little commercial potential for the Internet for at least 10 years
\end{quotes}

%\ifdefined\compileall

\chapter{网络}\label{chap:network}

\begin{quotes}[0.5]{Bill Gates}
I see little commercial potential for the Internet for at least 10 years
\end{quotes}


%\begin{quotes}[0.5]{microcai}
%The protocol is the kernel of the net. The kernel is the protocol of the UNIX. The UNIX is the source of the net.The source is the kernnel of the UNIX.UNIX, the net , the protocol , you are one.
%\end{quotes}

若要问世纪末最伟大的发明是什么，一定有人说，是Internet。Internet让全世界的电脑连在了一起。

\section{网络基础}

说到网络不得不提到OSI的7层网络模型。

\begin{table}[h]
\begin{center}
\caption{OSI模型}\label{tabel:ISOmodule}
\begin{tabular}{|c|c|c|c|}
\hline
& 数据单元 & 层 & 功能  \\ \hline % \cline{2-4}
\multirow{4}{*}{主机层}   &  \multirow{3}{*}{块数据} & 应用层  &  网络应用程序  \\ \cline{3-4}
&						 & 表现层 & 数据表现，加密解密，数据格式转化  \\ \cline{3-4}
&						 & 会话层 & 管理会话 \\ \cline{2-4}
& 分片数据  & 传输层 & 点对点的可靠通讯，流控制 \\ \hline
\multirow{3}{*}{媒介层}  & 数据包   & 网络层 & 网络寻址和路由 \\ \cline{2-4}
& 帧 & 链路层 & 物理寻址和交换 \\ \cline{2-4}
& 比特 & 物理层 & 物理信号传递：光电等 \\\hline
\end{tabular}
\end{center}
\end{table}

OSI定义的7层模型是参考模型\footnote{实际上构成互联网基石的TCP/IP协议并没有完全按照7层实现。}。

第一层，物理传输层。典型的我们每天都能见到的就是以太网和Wi-Fi了。太网使用同轴双绞线传递差分信号，由以太网协议规范。网卡的接口电路负责实现。
使用同轴双绞线传递差分信号，可以认为能传递 0 和1。Wi-Fi则使无线信号进行传输。
%操作系统无法管理到。不属于TCP/IP管辖范围。

第二层，链路层。典型的链路层协议就是以太网了。什么？又是以太网？ 额，不是说了么，实际上没有完全参照ISO的模型。以太网即是物理传输层也是链路层。
另外一个重要的链路层协议是PPP。PPP是个链路层协议，也就是第二层，那么谁给PPP提供第一层呢？如果第一层是电话线，那么PPP就是构建了拨号网络，如果第一层是以太网，那就是大名鼎鼎的PPPoE（PPP over Ethernet）。


在这个层，以太网发生的数据以以太网帧的形式进行。最大一次可以发送 1500Byte的有效数据。在链路层，以太网协议负责实现数据包的传递。每个数据包由“前导同步序列+太网帧头+负载数据+尾同步序列”构成。
同步序列用来在01比特流中识别出以太网帧。并且网卡在发送同步序列的时候还负责检测“碰撞”——多个网卡同时向一条线路上发送以太网帧。
遇到碰撞的时候，网卡负责重发。以太网帧头包含了负责数据的大小、发送方地址和接受方地址。
每个网卡被分配了全球唯一的地址。这个地址被称为MAC地址。操作系统能控制的就是以太网头部中的“发送地址和目的地址”以及负载数据。
这一层依旧使用以太网协议。

第三层，网络层。在这个层，运行的就是IP协议了。IP协议发送的依然是数据包。每个数据包由 “IP协议头+负载数据”构成。一个IP数据包正好作为以太网的负载数据发送。IP协议为每一台主机分配了一个全球唯一的IP地址\footnote{实际上IP地址划分出了几个区段，不能用于因特网，只能用于自己组件的私有网络中。这类私有地址在全球范围内是可以重复的。但是在每个私有网络内部仍然是不能重复的。}。IP地址通常使用点分十进制表示法。

第四层，传输层。自然就是TCP协议了。TCP在IP包的基础上提供了可靠的流式传输协议。TCP将数据打散为一个一个的包，然后利用IP协议传输每一个包。
如果发生丢包，TCP协议还会重新传输丢弃的包。到目的地，TCP会将数据包重新组合为数据流。

第五层，会话层。还是TCP协议。TCP包含了连接建立，数据收发和连接断开。建立一个连接就是一次TCP会话。TCP协议使用“三次握手”建立连接。断开连接则是四次握手。

第六层，表现层。在TCP/IP中，这个层已经不是他们负责的啦。如果应用程序使用了SSL加密协议，比如使用 HTTPS 访问网站，则SSL可以认为是这一层的。
直接传递数据而不加密的应用中，这一层是缺失的。

第七层，应用层。这类协议众多。访问网站使用的HTTP协议，访问FTP使用的是FTP协议，等等。

在每一层协议中，都有一个“负载”用来承担上一层协议的数据。而实际上，层之间甚至是可以交叉越级的。比如各种VPN\footnote{后面会介绍VPN是什么}协议就是利用的应用层来传输IP数据包。还有利用应用层传输以太网帧的。

\subsection{物理传输层和链路层}

\subsubsection{以太网}

\subsubsection*{小插曲：永恒的以太网}

\subsubsection{无线局域网Wi-Fi}

\subsubsection{点对点协议}

\subsubsection{虚拟局域网-VPN}

看上去是个应用程序的虚拟网其实是链路层的。

\subsection{网络层：IP协议和IP地址}

IP协议从OSI的模型上看，属于网络层。使用IP协议的网络就是IP网络\footnote{除了IP网外，还有电话网，有线电视网。}。
IP网络按照规模分为局域网，城域网和广域网。而当今世界上最大的广域网就是因特网。

IP网的作用类似快递——传递数据包。IP网并不保证数据包总是能到达目的地——虽然网络建设的时候总是力求降低丢包率，但是IP本身并不100\%保证数据传达。

IP网工作在网络层，需要使用链路层协议传输数据包。在全球范围内，IP骨干网使用的是ATM和X.25这样的光纤网络进行传输的。在本地局域网，使用以太网协议。从ISP到家庭用户，使用的是
ADSL这样的拨号网络。甚至可以不是物理网络：使用各种应用层协议传输IP包——这就是VPN的原理了。

由于IP协议不提供可靠的数据传输。IP只能尽量保证将数据包传递到目的地。在这里不可靠的意思有两层：第一个意思是数据包可能丢失，第二个意思是丢失不丢失发送方无从知晓。如果发送方知道数据丢失，则仍然任务协议是可靠的——发送方可以选择重发或者标记网络不可达——这意味着任何协议都无法完整数据包的送达。

IP的另一个不可靠特性是接受方收到的数据包次序和发送方并不能保证是一致的。可能会出现乱序。因为IP数据包是经过“路由”的。我们会在下一个章节介绍路由。

\subsubsection{路由表和路由协议}

上节说到IP协议传递的是分组的数据包，每个数据包标记了目的地址和发送地址。数据包在发送到目标之前会经过层层的“路由”。

TODO 路由示意图

解说： IP网络由主机和路由器以及连接它们的通讯线路组成。图中示意的是一台主机发送数据包到多台目标主机的过程。

主机A分别发送数据包到B、C、D，向B是直接发送。向C是通过B的转发到达。而发送到D的数据则经过了E、F、H的转发。另外一条可能的路径是B->G->F->H->D，不过这是一条比较“远”的路径。
像 F 这样的主机连接了大量的网络，专门用来转发数据包，被称之为路由器。B这样的主机，本身也提供服务，但是也转发数据包。在转发数据包的时候它是路由器，本身发送和接收数据包的时候它是主机。

一台主机是不是路由器主要看它在网络上的作用。我们有时候会把专门制造的用于路由的计算机称作路由器。当是IP网络上任意一台主机都可以担当路由器——只要它扮演了这个角色。

我们可以看到，一个数据包从A到D可以选的路径不止一条——但是有的路径经过的节点少，有的则多。也就是说，这些路径的质量是不能等同看待的。
衡量一个路径的标准自然就是用户最关系的：延迟和带宽。经过的节电越少通常来说延迟就越少。带宽自然是越高越好——为了尽量避免拥挤的线路，不仅仅是要选择带宽高的路径，还要选择带宽没被占满的。

如此一来就带来两个问题：第一个问题，IP数据包在进行路由的时候，如何知道哪条路径是可以将数据包准确的送达目的地的——这属于路由发现问题；第二条，拥有多个路径可选的时候 ，如何选择一个最佳路径——这属于路由选择问题。

在操作系统的实现上，内核是依据一个“路由表”执行路由的。路由表可以修改，所以选择最佳路径归结为“在系统运行中动态的修改路由表以让内核总是使用最佳路径”这个问题了。

路由表是一张表格，详细的列举了每个IP（也可以是一个区间，这样可以减小表的大小）应该使用哪个路径。比如对A来说，向C和B发送的数据，使用的网关就是B。路由表上记载的应该就是 C->B , B->B 这样对应的条目。但是对于D，路由表包含的是 D->E 。使用 E 来转发。而事实上E并不能直接将数据转发给D，E 的路由表里包含 D->F 这样的条目，因此E将数据包交给 F 来转发。数据就是这样层层转发的，路由表并不记录完整路径，而是接力棒一样。
所以用来接下一棒的路由器被形象的称呼为网关。

路由表可以由管理员手工输入（使用route(8)命令）。管理员输入的路由表被称为静态路由表。但是面对复杂的网络——尤其是那些运行在互联网的骨干上的路由器来说——手工维护显得非常不合理。
所以就催生了各种路由协议，在路由器之间动态维护路由表。粗略的列一下，大概有边界网关协议（BGP）、外部网关协议（EGP）和内部网关协议（IGP）。
大协议下还细分了许多子协议，用来进行距离探测，链路状态探测等。路由协议运行在互联网的骨干路由器上。用于路由器之间维护完整的优化的路由表。普通用户的机器通常并不需要运行路由协议。设置一个默认网关就足够了。

\begin{insertnote}
\subsubsection{默认网关}

一般家庭局域网的机器上，路由表只有两条：一条告诉本机哪个IP范围的IP是直接发送的，其他的都使用默认网关。属于本地局域网的IP自然是不经过路由的。机器判定本地IP的范围使用的就是
DHCP\footnote{参见\secref{sec:DHCP}}的时候获得的掩码。而大部分IP地址都是用过默认网关进行路由的。毫无疑问这个默认网关就是家庭路由器:)。

默认网关告诉操作系统，所有条目都不匹配的时候使用这个网关。因为路由表总是从上而下进行匹配的，所以默认网关总是路由表的最后一条。

\end{insertnote}

那么，接下来整理一下，我们在这个想象中的网络具体的观察一下IP协议的运转。

\begin{enumerate}
\item [局域网（Hub集线器）] \item [局域网（带交换机）] 主机A发送到主机B，因为是他们连接到一个集线器上，所以主机A首先发送一个以太网广播（ARP协议）获得B的mac地址。
然后就可以直接向B发送数据了。Hub集线器实际上就是把网线加终端电阻连接到一起。所以数据在Hub的端口之间是直接广播的。目的地址是B的mac地址的以太网帧被B收到——也就是收到A发送的IP包了。

\item [局域网（带交换机）] 主机A发送到主机C，因为是他们连接到同一台交换机上，所以主机A首先发送一个以太网广播（ARP协议）获得C的mac地址。
然后就可以直接向C发送数据了。数据由交换机交换传递到C连接到的端口，这样主机C的网卡就能收到A发送的IP包了。

\item [局域网（无线网桥\footnote{简称AP。}）] 主机A发送数据包到D，中间通过了无线网桥，将有线以太网协议转化为802.11a/b/g协议，使用天线发送。802.11a/b/g协议属于第一层协议。第二层依然是以太网协议。网桥在这里是做了
第一层有线以太网协议到802.11a/b/g协议的转化。IP是第三层协议，所以这个时候依赖的第二层以太网协议本身没有变化。对IP层来说，无线和有线是透明的。都是以太网。所以这里依然使用ARP协议获得接收方的网卡地址，然后直接发送。这里网桥和交换机是一样的功能。

\item [局域网（无线到无线）] 主机A发送数据到D，这个时候断开网线，自动切换到使用无线网络。发送依然要经过网桥。仍然是以太网，依旧是ARP协议首先获得对方的mac地址，然后发送。数据由网桥转发到目的主机。

\end{enumerate}

\subsection{传输层：TCP协议}
为了提供可靠的通信，可以在IP协议的基础上加入数据接受确认、超时重发、和顺序重组功能就可以提供一个可靠的面向流的协议。该协议提供给上层可靠的面向流的协议——可靠的意思是，协议告诉你发送成功接收方就一定已经收到数据，协议告诉你发送失败接收方就是没接收到；面向流的意思是，数据在底层会被分包以使用IP协议传输，但是如何分包是不需要上层操心的，上层交给协议的 就是一连续的字节流，接收方收到后会重组为正确次序的字节流。除了可靠和流式这两个特点，还应该加入拥塞控制——发送速度不要超过网络能提供的最大带宽避免网络拥塞。整合这些需求后，诞生的就是TCP协议了。

由于主机通信双方并不可能只有一个程序。所以当一个TCP数据送达，只根据发送方地址来区分（这样的话一个主机只能向另一个发起一个链接了）太不合理了，所以TCP协议引入了“端口”的概念。
如此一来，一个主机就可以执行多个程序，使用不同的端口进行通信，相互之间不干扰。这样确定一条TCP链接的就是两个IP地址和2个端口了\footnote{发送方和接收方的端口，发送方和接收方的IP。}。如果没有端口，则两台主机之间只能建立一条TCP连接。

如果应用程序不需要TCP提供的可靠流式传输，想直接使用IP这样简单的“尽力送达”服务，另外有一个UDP协议可以使用。UDP协议在IP协议之上简单的加了一层“端口”。当然加入端口和TCP的原因是一样的——没有端口的话，一台主机岂不是只能有一个程序能接收数据？\footnote{操作系统将无法确定接收到的数据包应该传递给哪个程序。
有了端口，使用UDP协议的程序可以绑定到一个端口上，这样对应的端口接收到数据操作系统就可以知道将数据递交给哪个程序。}


\section{网络配置}\label{sec:ifconfig}
\subsection{图形环境下的工具}
\subsection{命令行配置工具}
\subsection{网络自动配置和DHCP}\label{sec:DHCP}

\subsection{虚拟局域网}

\section{网络攻击和防火墙}
\subsection{网络攻击类型和检测}
\subsection{iptables防火墙} %\fi

\chapter{文件系统}

\begin{quotes}{Ken Thompson}
We have persistant objects, they're called files.
\end{quotes}

\section{磁盘上的文件系统	}
\subsection{块设备和磁盘}
\subsection{超级块和inode}
\subsection{典型磁盘文件系统举例}
\section{虚拟文件系统}
\subsection{/proc内核信息窗口}
\subsection{/dev设备文件系统和udev}
\subsection{虚拟内存盘tmpfs}
\section{网络文件系统	}
\subsection{NFS}
\subsection{Windows 网络邻居 CIFS}
\section{  其他的文件系统	}
\subsection{  LiveCD的最爱——压缩文件系统squashfs}
\subsection{  为Flash芯片设计的文件系统}

\chapter{架设服务器}
\section{  搭建 HTTP 服务器}
\subsection{  apache 用的最多的服务器}
\subsection{ nginx 轻量级服务器}
\subsection{ lighttpd 超轻量级服务器}
\subsection{ squid 加速代理	}
\section{  数据库	}
\subsection{  最流行的开源数据库 MySQL}
\subsection{  最优秀的开源数据库 PostgreSQL 	}
\subsection{  商业霸主 OracleDB	}
\section{  加速 DNS ，在本机搭建 DNS	}
\section{  共享打印机	}
\subsection{ CUPS 打印服务	}
\subsection{ Samba 打印机共享	}

\ifdefined\compileall


这个世界上有许多人总是说，计算机处理能力已经过剩了。

不，从没有。计算机的处理能力从来就没有过剩过。如果说过剩，只能说，他从来就不需要计算机。计算机硬件的发展速度永远跟不上人对计算能力的爆发性增长的需求，
调和这个矛盾的唯一办法就是进行优化。所谓优化，就是不改变功能的前提下让程序执行的更快，占用的资源更少。

优化可以在好几个层面进行。从最底层的汇编优化再到最高的设计优化，每个层级都有优化的可能。

\begin{enumerate}
\item 设计优化	~	即使是同样功能的程序，也拥有不同的设计。不同的设计通常导致非常差异化的编码思路，从根本上影响到一个系统的效率。不同的设计导致关键逻辑上采取了不同的算法，很大完成同样功能的不同的算法有根本上的效率差距。

\item 代码优化	~	编写程序的时候，即便是相同的算法和设计，也有机会选择不同的代码描述形式。参数传递是采用指针还是对象？编写宏还是函数？等等。这些因素也会影响到最终代码的效率。
好的代码\footnote{那种一眼看上去就知道代码要做什么。}总是比写的烂的代码\footnote{让人看了半天没看明白，格式混乱，诸如此类的代码}更讨人喜欢。但是好的代码并不总是能生成优化的结果的代码。有时候为了最终编译后的程序的执行效率，会采取一些trick编写代码，一些trick能极大的提高代码的执行效率，但是大大降低了代码的可读性。现代程序普遍不再使用trick，只编写具有良好可读性的代码。把优化的工作交给优化编译器。

\item 构建优化	~	通常一个软件项目包含多个功能模块，每个模块使用多个编译单元\footnote{一个源代码文件就是一个编译单元}, 通过一套构建系统\footnote{autotools,cmake,qmake等等。}调用编译器生成最终结果。许多软件可以将其中一些功能列为可选，禁用的可选功能直接导致跳过编译对应的编译单元。
这被称为条件编译。条件编译也可以发生在编译单元内部\footnote{学过C语言的人都知道，\#ifdef ... \#endif 预处理指令就是干这个的}。
通过条件编译，使得一些用不到的功能代码被排除在最终可执行文件之外。减少了冗余代码，节约了程序运行时所占用的资源。

\item 编译优化	~	编译器能“读懂”源代码并将其转化为等价的二进制指令。在这个过程中，如果加入优化过程，编译器能删除\footnote{这里的删除并不是编译器将代码从源文件中删除，而是说是编译器并不将其转化为机器指令，而是直接丢弃。}冗余代码\footnote{源码中存在的但是实际上编译器发行并没有被调用的函数，会被删除。一些做无效操作的代码被删除。一些重复操作被编译器合并为一个操作.}，常量折叠\footnote{将相同的数据合并}，常量表达式\footnote{编译器在编译的时候发行一些表达式的值是固定可知的，而且在编译期就能求得，故在最终程序中使用一个常量代替直接计算表达式}优化，循环展开等等。

\item 特定CPU优化	~	其实这也是编译优化的一部分。这里单独列出来讲一下。在CISC\footnote{复杂指令计算机，参考 \faqref{FAQ:CISCandRISC}}指令的计算机上，由于指令繁多，编译器具有较大的选择余地。同时，一种指令在一款CPU上执行较快，而做另一款CPU上则执行的较慢。这使得编译器无法正确的作出最优的选择，因为对于某个型号的CPU的最优选择并不适用于另一个型号。汇编优化的意思就是指定为某个特定的CPU上选择生成对于该CPU而言最快的指令。例如同样是浮点操作，在x86平台可以使用x87指令集或者使用SSE指令。某些型号的处理器，x87指令的速度和SSE的速度不相上下，而x87的精度要比SSE高，就因该尽量使用x87指令进行浮点运算。而有的型号，x87速度比SSE指令慢好几个数量级，自然优先使用SSE指令\footnote{除非用户表明了使用双精度浮点数}。不仅仅为某个特定型号的CPU进行优化有可能再另一个CPU上起反效果，有些指令并不是所有的CPU都支持的。例如SSE指令只在奔腾3以上的CPU中才支持。avx指令更是只有在sandy bridge\footnote{intel 2011年发布的core2处理器的代号}中才支持。如果软件使用了 core2指令，放到非2011年前生产的CPU中可全部无法运行了。所有通常二进制发布程序都不为用户所使用的CPU进行专门优化，这也丧失了一大性能。

\end{enumerate}

使用Gentoo的时候，由于我们只是用户，并不是上游软件开发者，虽然无法在设计上改动上游的代码，但是我们手里能控制的还有编译优化和构建优化。将不需要的功能通过USE选项筛剪，并调节编译器的参数生成为本机的CPU优化的最优指令。这也是portage给我们的最强大的武器。

\section{编译优化}

%优化的理论依据是，
现代程序通常由高级语言编写，使用编译器将其编译为机器可执行的二进制指令。在这个转换的过程中，事实上存在多个可能的结果
\footnote{从源程序获得最优化的结果是个NP完全问题。}。正所谓条条大路通罗马。
人们对编译器的期望是，它总能够选择其中执行速度最快的那个作为最终结果。所谓执行速度最快，通常是完成同样的功能使用最少的指令，或者指令序列总执行时间最短。
这也是编译器能对程序进行优化的依据。

\subsection{编译期优化}
前面提到过，生成最优结果对编译器来说是个NP完全问题\footnote{所谓NP完全问题，就是说在“有限”的时间内无法找到最优解的问题}。编译器必须能在有限的时间内给出编译结果。所以编译器提供了“优化级别”概念。

所谓优化级别，就是给出编译器执行优化过程的时候所尽的“努力”程度。最高级别的优化就是尽最大的努力生成编译结果。通常生成的结果运行速度最快，但也意味着最费时的编译。
最低级别的优化就是不优化，编译速度为最快，生成的结果程序的运行速度也是最慢的。

对GCC编译器来说，优化级别由参数 -Ox 指定。x是一个  0 到 3 之间的数字。数字越大，优化级别越高。0表示不优化,如果不给出优化级别，默认就是0。
对用户来说，只要知道使用“-Ox”参数开启指定级别的优化即可。

emerge{}使用的编译优化参数由 make.conf(5) 文件中的{ }CFLAGS{ }和{ }CXXFLAGS{ }指定。通常指定\texttt{"-O2 -march -pipe"}，其中的-O2即指定优化级别为2。这是Gentoo推荐
的设置。虽然说-O3有可能带来更好的性能。但是激进的优化措施有时候会带来意想不到的问题。为保守起见，Gentoo只推荐全局启用-O2级别的优化。

GCC支持另一个特殊的优化级别：-Os。‘s’在这里的含义是size，也就是为大小进行优化。

\begin{insertnote}
\subsubsection{什么是为大小优化？}

编译器进行优化的时候，有时候会碰到代码大小和执行速度的矛盾。有的代码可以优化为更少的指令，但是每条指令都更复杂，导致总执行时间更长；或者优化为更多的指令，但是每条指令更简单，而且在超标量处理器上能做到指令级的并行执行，总执行时间反而更短。
虽然说大部分情况下，更少的代码总是意味着更快的执行速度，但是遇到某些特殊情况时，编译器就必须作出代码大小和执行速度的权衡。
-Os{}优化级别就是告诉编译器，在遇到代码大小和执行速度两难的境地，优先考虑代码大小，使用让代码体积更小的优化措施。

\end{insertnote}

一些人认为，只有为内存受限的特定环境编写程序时才考虑-Os优化。但是Linus在一次邮件中鲜明的指出，至少在x86平台，-Os优化生成的结果总是最快的。下面列出邮件原文，也可以从内核邮件列表的网上归档上查看 \url{https://lkml.org/lkml/2003/2/6/182}。

\begin{insertnote}

\subsubsection*{Linus讨论-Os和-O2的邮件}

\leftskip 2em
\rightskip 2em
\parindent 0pt

\colorbox{gray}{\hbox to 5em {Date} Thu, 6 Feb 2003 15:16:16 -0800 (PST)}

\colorbox{gray}{\hbox to 5em {From}	 Linus Torvalds <>   }

\colorbox{gray}{\hbox to 5em {Subject}	Re: gcc -O2 vs gcc -Os performance }

\footnotesize 

\parskip 4ex

\noindent{}On Thu, 6 Feb 2003, Martin J. Bligh wrote:\\
> \\
> The observation re low repeat rate is interesting ... might be amusing \\
> to do some really basic profile-guided optimisation on this grounds,\\
> take readprofile / oprofile output, and compile the files that don't\\
> get hammered at all with -Os rather than -O2. Given their low frequency\\
> (by definition), I'm not sure that improving their icache footprint will\\
> have a measureable effect though.

Icache footprint has nothing to do with repeat rates, which is exactly why 
repeat rates are interesting for -Os.

Icache footprint is directly proportional to the \_static\_ size of the code 
(ie exactly the thing that -Os is supposed to optimize for), while 
instruction-level performance measurement is only valid on the \_dynamic\_ 
code.

And with modern CPU's with big caches, a \_lot\_ of cache misses are the 
forced kind - the startup costs, not the actual runtime cost. That's not 
always true (if you touch big data sets, you'll have replacement misses 
too, of course), but it's not really false either.

So think of the I\$ (and TLB, and page load/map - all the same) cost as a 
fixed cost that will always be there, but that -Os tries to minimize. 
That's \_one\_ dimension in the total cost.

The "traditional" -O2 kind of "try to make the code run fast" 
optimizations tend to try to minimize a totally different dimension, 
namely the dynamic code speed.
And the time required for running the program is the sum of the static and 
dynamic factors. In other words, a \_good\_ optimization should try to 
minimize not one or the other, but the sum.

And low repeat rates means that the dynamic component is smaller, which 
clearly makes the static component more important.

For example, if you are doing mp3 encoding, the repeat rates for the core 
loop are huge, and the code is small, so clearly the static component is 
largely insignificant. Use -O2.

But if you're running a GUI program then just the loading time is often
quite noticeable, and if you can improve that by, say, 10%, then that can
\_more\_ than make up for almost any amount of stupidity in your code.  
Especially since a lot of the code isn't even all that loopy and tends to
have low repeat rates. You're almost guaranteed to be better off using -Os
than -O2.

If you've got performance counter data, check the I\$ and ITLB miss ratios, 
and if they are at all noticeable, think about the fact that a I\$ miss 
tends to cost a lot more than a few more dynamic instructions. 

I suspect the kernel I\$ behaviour is generally pretty good, and the ITLB 
behaviour is improved even further thanks to large pages etc. That said, a 
user app that blows the I\$ will blow the kernel out of the I\$ too, so 
small is always beautiful, even in the kernel.

\qquad Linus

--\\
To unsubscribe from this list: send the line "unsubscribe linux-kernel" in \\
the body of a message to majordomo@vger.kernel.org \\
More majordomo info at  \url{http://vger.kernel.org/majordomo-info.html}\\
Please read the FAQ at  \url{http://www.tux.org/lkml/}

\end{insertnote}

在邮件中，Linus指出，CPU的指令缓存是有限的，更小的代码意味着程序能将自身更多的部分放入指令缓存从而带来指令缓存命中率的提升。更大的代码导致CPU更多的发生 Icache miss（指令缓存不命中），指令缓冲一旦没有命中，CPU不得不暂停执行到内存读取指令并刷新指令缓存。众所周知，CPU之所以有缓存设计，就是因为内存无法和CPU同频运行。
CPU的一级缓存是单周期访问的，但是容量有限。通常只有数十KB；而二级和三级缓存虽然大很多\footnote{在i7级别的CPU上三级缓存有8MB。但是入门级别的CPU可能就没有三级缓存}，但是访问时间却要好几个时钟周期；但是内存的访问是更慢的了，达到数十个周期。也就是说，如果一次Icache miss发生，CPU将会浪费几十个周期到二级和三级缓冲刷新，但是如果三级缓存中都没有，必须到主内存刷新，浪费的时钟周期将达数百个\footnote{缓存刷新的时候会执行预读取，所以浪费的是数个内存访问周期。}。随着CPU主频的提升，访问一次内存导致的CPU时钟周期浪费也会越来越高。更大的指令除了导致更容易出现Icache miss，也导致可执行文件变的更大，也将是初次加载的
时候要花费更多的时间将代码从硬盘读取到内存。

不论是从运行时的速度上考虑还是从加载速度上考虑，更小的代码总是意味着更快的速度。更何况将算没有将缓冲填满，更小的代码也给其他程序留下了更多的缓存空间。所以整个系统以-Os参数进行优化似乎能获得从整体而言最快的系统。



\subsection{链接时优化（LTO，Link Time Optimization）}

Linux系统从内核到应用程序，几乎都是由C/C++语言写成的。C家族语言都需要经历“编译”$\Rightarrow$“链接”的过程才能最终转化为可执行文件。

C家族语言的源代码是分开编译的。一个程序通常会包含多个源代码文件，每个源码文件分别编译，最后执行链接合并成一个可执行文件。

大多数优化都是发生在编译的时候。但是编译的时候，由于源码被拆分成了独立的文件，导致编译器无法对一些变量和函数的使用收集到足够的信息。如果编译器能“看到”所以的代码，而不是编译的时候的那一个源码文件，就能更好的做出优化决定。

而在C家族语言中，工具链上工具中的能“看到”全局的程序只有链接器而已。如果把一些优化交给链接器执行，链接器就能针对全局进行优化了。这个优化 就是链接时优化,Link Time Optimizatio。开启链接时优化
只需要使用“-flto”参数即可。将“-flto”加到CFLAGS,CXXFLAGS和LDFLAGS就能全局开启链接时优化。

\begin{notice}
注意：LTO是非常新的优化，所以并不能保证所有的代码都能正确编译。在成熟之前强烈的不建议开启。
\end{notice}

\subsection{PGO}

现代CPU具有超长的流水线，流水线越长，越是害怕条件跳转。条件跳转一发生，流水线就要清。所以CPU都会预测调节跳转指令。只要预测落空才需要清空流水线。
居于此，编译器可以生成容易让CPU的超标量执行单元正确的预判条件分支，让CPU预测容易的预测到经常执行的分支。
GCC提供了 \_\_builtin\_expect 扩展，告诉gcc一个条件表达式是否是经常为真还是经常为假。

内核里经常会见到这样的代码：

\begin{code}
if(likely( 表达式 ))\{\\
	... \\
\}
\end{code}

这个 likely() 就是一个宏包装一下\_\_builtin\_expect。

\textbf{这是程序员在写代码的时候能知道的情况，如果程序员并不知道条件分支的概率呢？}

程序中经常能遇到许多循环。编译器常见的优化手段是执行循环展开。比如下面这个代码 


\begin{code}
for(i=0; i<4;i++)\{\\
\qquad 	f(i);\\
\}
\end{code}

这样的代码，编译器直接将它展开成4次f(i)调用，展开成如下的代码

\begin{code}
f(0);\\f(1);\\f(2);\\f(3);\\
\end{code}

这就避免了循环控制结构上的比较和条件跳转操作，以为条件跳转总是相对开销巨大的。但是，这增加了代码的体积。以空间换效率。也就是说，小循环值得展开，大循环最好不要展开。
这中间有个尺度（和CPU的cache有关，所以为本机CPU优化意味着要考虑CPU的cache大小哦）。但是有的循环在这尺度附近，但是程序生命周期中很少执行，那么展开增加了代码的体积，由于执行次数不多，所以带来的速度提升有限。这就最好不要展开了。最关键的问题来了：编译器如何知道循环执行的频度？
这这编译期是无法判定的。只有程序实际执行了才知道。也就是我们说的，优化热点代码\footnote{频繁执行的代码。}。不是热点的代码何必优化呢？\textbf{编译器如何知道一个循环的热点程度以更好的作出循环展开优化呢？}

编译期无法知晓的问题又何止这些！综合来看，编译器光看代码，做静态分析，不能知道代码这执行期间的热点程度，也就不能很好的结合代码的热度进行优化。
PGO就是为这个问题诞生的一种优化形式。PGO的意思就是编译器要收集程序运行时的数据（ Profile），然后利用（Guided）收集到的数据进行最终优化（Optimize）。

具体的做法就是先使用 -fprofile-generate 参数编译一次，接着执行程序。加了-fprofile-generate参数编译后的程序运行过程中会自动收集并生成性能数据。执行完毕后，修改编译参数使用-fprofile-use再次编译。
这个时候gcc就可以从收集到的数据回馈执行更好的优化。

目前这编译脚本中集成了GPO优化的只有firefox，USE标记中加入gpo再编译firefox即可。

\begin{code}
\#USE=gpo emerge firefox
\end{code}

这样会编译火狐两次，第一次编译完成后会开启一个火狐，里面运行测试页面。测试完成后进入第二次编译状态。最终的结果也许比没有开启PGO优化快吧\footnote{之所以用\textbf{也许}这个词，因为GPO优化的效果取决于试运行期收集到的数据。}。


%\section{有目的的优化 – 查找性能瓶颈}
%\subsection{查找性能瓶颈}
%\subsection{  优化瓶颈}

\else
\addtocounter{chapter}{1}
\fi


\chapter{时光机器-版本控制系统}\label{chap:VCS}
\section{   历史和后悔药	}
\section{   中心式版本控制仓库 CVS	}
\section{   CVS 后继 SVN	}
\section{GIT 划时代的分布式版本控制\label{sec:git}}
\subsection{  Bitkeeper	}
\subsection{  Linus 消失的一周	}
\subsection{  版本控制设计为一个文件系统	}
\subsection{ 去中心化	}
\subsection{  GIT典型工作流	}
\subsection{  GIT 简单使用	}


\appendix

{ \renewcommand{\thesection}{\arabic{section}}

\chapter{FAQ}\label{chap:FAQ}

\section{什么是 LiveCD ?}\label{FAQ:LiveCD}
一种让 Linux 系统脱离硬盘，只需要一张 DVD 或者 CD 运行的技术。使用者只需要插入一张 LiveCD 就可以使用 Linux  系统了。当然，由于 CD-ROM 是只读的，所以对 LiveCD 系统本身所做的更改重启后都会消失。由于 Linux 发行版的膨胀，一张 CD-ROM 无法装下，于是就有了 LiveDVD , 除了容量增加外，和LiveCD 是一致的。此外还有 LiveUSB , 也就是使用 USB 存储设备的 liveCD。用法和 LiveCD 并无二致。

\section{如何使用 LiveCD。}\label{FAQ:UseLiveCD}

\section{UEFI和BIOS是什么？}


\section{什么是API?。}\label{FAQ:API}

\section{Shell 是什么？}
操作系统是由内核和一系列的系统程序组成的。

\section{什么是进程和线程？}\label{FAQ:Process}

\section{什么是内存分页？}\label{FAQ:Paging}

\section{如何不重新登录切换到root帐号？}\label{FAQ:su}

正如 \secref{sec:usersandgroups} 所讲，UNIX支持多账户登录。如果当前使用普通帐号登录，而某些命令需要使用root帐号的时候，就需要用到su这个命令了。
使用su命令的前提是加入wheel组，否则不能使用su切换到root。 “su”切换到root帐号，但是没有切换环境变量和工作目录，亦没有执行/root/.bash\_profile。{}“\hbox{su -}” 能切换到 /root 目录，提供类似直接登录的环境。用户很可能需要的就是带 - 参数的su。

su 还可以切换到其他用户，如“su smith - ”可以切换到smith帐号上。

另外一个临时使用root权限的办法就是在执行的命令前加 sudo 。 sudo 需要用户加入wheel组并在 /etc/sudoers 加入 “\%wheel ALL=(ALL) ALL” 一行。
sudo 通常来说只应该授权给信任的普通帐号 —— 比如管理员自己给自己设定的日常登录帐号（如果是桌面用户，自然就是自己用来登录桌面的帐号） —— 并且为普通用户设定好强壮的密码。

\section{内核参数}\label{FAQ：kernelparamter}


\section{Initramfs是什么}　\label{FAQ:initramfs}

Linux启动过程中内核主动做的最后一件事情就是执行/sbin/init，那么倒数第二件事情必然是挂载根目录。
挂载根目录有2个前提条件：找到根分区和支持根分区所使用的文件系统。
找到根分区，第一点，需要其所在设备的驱动。第二点就是内核需要知道哪个是根分区。这个可以通过内核参数 root= 设定。

Linux下驱动可以编译进内核，也可以编译为模块动态加载。如果恰巧磁盘设备的驱动被编译成了模块，就出现“鸡生蛋和蛋生鸡”的问题了：不访问磁盘就无法加载模块，而不加载模块有没有磁盘的驱动。这对于自己编译内核的童鞋不成问题。只要知道自己的SATA/IDE控制器的型号，将对应的驱动编译进内核就解决了。但是如果是打算制作“通用内核”的发行版呢？

文件系统也是同样的问题，没有文件系统驱动就无法挂载分区访问上面的文件，无法访问文件就无法加载模块。这对于自己编译内核的童鞋不成问题。只要将自己根分区所使用的文件系统编译进内核就解决了。但是如果是打算制作“通用内核”的发行版呢？将所有的文件系统都编译进内核么？

initramfs就是为解决这个“鸡生蛋和蛋生鸡”的问题诞生的。

GRUB加载内核的时候通过initrd 命令加载一个initramfs，该initramfs就会成为内核的根分区\footnote{事实上是Linux内核拥有一个不挂载就存在的rootfs文件系统，Linux启动过程中将initramfs的内容解压到rootfs中。作为理解，可以认为initramfs挂载成Linux的临时根目录。}。内核初始化完毕后直接执行initramfs中的/sbin/init而不理会root=命令行参数。由initramfs中的/sbin/init程序负责加载磁盘驱动，查找根分区，挂载根分区到 /sysroot 并调用pivot\_root()\footnote{使用命令　man 2 pivot\_root查看介绍。}将找到的根目录替换为真正的根目录。然后执行真正的/sbin/init程序。

initramfs是一个目录的压缩包。由dracut程序生成。initramfs用于寻找真正的根目录时的临时根目录，通常只包含一些必要的驱动和一些用于查找根分区的脚本，但有时候可以成为一个完成的系统的根目录，这样的系统可以脱离磁盘运行。比如Fedora的安装DVD,使用一个巨大initramfs作为安装环境的根目录。该initramfs包含的是一个真正的OS，带有完整的工具。


\section{汉化man手册}\label{FAQ:zhman}

man使用程序使用nroff程序格式化man手册，然后使用less向屏幕分页输出。为什么呢？因为man手册的显示会随着终端的宽度不同而变化。所以man手册并不是使用原始的ascii文本文件编写的。
man手册源文件使用非常隐晦难懂的一些标记，这些标记借助nroff程序格式化后就获得了适合特定宽度的终端显示的文本。正因为如此，所以man手册能在各种大小的终端上显示而不破坏风格。

GNU使用的nroff是groff（GNU出品，带个g是常有的事情），但是groff并不支持UTF-8编码的文本。所以导致man手册中不能出现非英文字符。所以有人对groff进行了修改，添加了UTF-8支持，这就是groff-utf8。要想显示中文man手册，按照man-pages-zh\_CN这个中文手册之外，还得安装groff-utf8并修改/etc/man.conf，让man调用groff-utf8而不是原来的groff进行格式化操作。

\begin{code}
\#emerge man-pages-zh\_CN groff-utf8
\end{code}

编辑： /etc/man.conf 

找到这一行

{
\tt
NROFF           /usr/bin/groff -mandoc
}

替换为

{ \tt
NROFF           /usr/bin/groff-utf8 -Tutf8 -c -mandoc
}

然后一些汉化了的手册就不会再显示乱码了。
gentoo-zh overlay 已经集成了这一变化，groff-utf8 安装完成后会自动修改/etc/man.conf，如果修改没完成可以手工进行。

\section{符号链接}\label{faq:symlink}

从需求上看，有时候我们希望能使用不同的文件名访问同一个文件。符号链接就是干这个活的。可以理解成一个快捷方式。符合链接分软链接和硬链接两种。硬链接只能用在文件上而无法对目录使用硬链接。而且只有在同一个分区上才能使用硬链接。软链接没有这方面的限制。

硬链接就是在文件系统上出现了使用同一个磁盘区域的多个文件。所以会有上述限制也不奇怪。软链接实际上是个特殊的文件，文件的内容是链接的目标的文件名。打开软链接，实际上会打开链接到的文件。使用readlink(2)系统调用才可以获得链接目标的文件名。

符号链接使用 ln(1) 工具创建。软链接要加“-s”参数。

创建链接的语法是

\begin{code}
ln [-s] 链接到的目标 链接名称
\end{code}


“链接到的目标”可以使用相对路径。如果是相对路径，并不是相对执行命令的时的当前目录，而是相对于“链接名称”的路径。

\section{内存映射IO-MMIO}\label{FAQ:MMIO}

CPU对外设的操作有两种形式：专门的 IO 指令和 MMIO。 和内存拥有地址意义，外设也占有IO地址。IO地址和内存地址是独立编排的。对内存的操作使用的都是一般的内存操作指令，如 mov 。
而对IO地址的操作使用的则是专用指令 IN 和 OUT。IO指令是特权指令，只有CPU处于特权级别的时候（执行内核代码的时候）才能使用。

除了专用的IO地址和IO指令，CPU对外设还可以采取内存映射的IO。在“内存控制器”\footnote{注意和MMU的区别}的配合下，将一些PCI设备的IO地址映射到上内存地址。CPU即可使用一般的内存操作指令对IO进行操作。由于CPU对内存操作的指令非常丰富，这就方便了设备驱动的开发。当然，还有一个原因：x86 CPU只给IO地址分配了0--65535共64k个地址。而PC外设越来越多，IO地址显然不够用了，借用内存地址吧！

当然，也有部分外设是只能MMIO的，比如显存。显卡带有一部分直接被GPU访问的内存，这部分内存称为显存。CPU要访问显存，需要通过AGP/PCI-E总线，以操作AGP/PCI-E控制器的办法间接访问显存，显然这样会太麻烦。不如直接采取MMIO的办法，把显存映射到某个物理地址区域，以对待内存那样对待显存。对CPU来说，显存不过是访问速度慢了点的内存。但是这带来了一致性。

\section{内核视频模式设置-KMS}\label{FAQ:KMS}

\section{CISC和RISC}\label{FAQ:CISCandRISC}

在开发CPU的人群中，诞生了两个派系：CISC和RISC。（人类从来不缺乏派系斗争啊！汗！）
CISC是{}Complex Instruction Set Computing{}的缩写。RISC全称为{}Reduced Instruction Set Computing{}。

CISC{}将指令设计为单条指令可以执行多个低级操作。如一条add指令即可完成访问内存获取数字，并和寄存器中的数字相加这两个操作。
CISC的指令灵活性比较强，完成同样的操作可以使用多种指令的组合方式做到。而且每种组合执行效率都不一样。
RISC则反过来，每个指令只提供一种低级操作，CISC中的典型指令通常需要多条等价的RISC指令。但是由于RISC指令比较简单，每条指令所做的操作非常简单，所以多数指令可以做到单周期执行。
而CISC由于指令比较复杂，使得译码电路也大大的复杂化。执行每条指令所需要的时间也不相同。

简单的说，每条CISC指令其实可以分解为多个更简单的微操作，而每条RISC指令就是一条微操作。通常无法继续分解。另外，多数RISC指令都设计为等长格式，不论何种操作的指令，一条指令的
长度都是固定的。CISC通常使用变长编码。指令依据所要进行的操作不同而占据不同长度。

\section{RPC}\label{FAQ:RPC}

在计算机中，远程过程调用指的是程序调用一个子过程，子过程在另一个进程（可以是本机的另一个进程，也可以是网络上的另一台主机上的进程）里执行，但是编程上不需要关心网络传输细节。
只需要执行一次普通的函数调用即可。这样的过程就是远过程。远过程调用实现上有很多。桌面上常用的是DBus。Web编程有XML RPC 和 JSON RPC等可选。

\section{INI文件格式}\label{FAQ:INI}

INI 格式是 微软在 MS-DOS 开始沿用到 Windows ME 的系统配置文件的格式，通常使用扩展名.ini 因此得名。虽然后来被“注册表”所取代，但是并没有完全消亡。INI 文件因为其简洁的风格，被大量的软件用作其配置文件格式。

INI 文件的设置以 键值对为基础。如

\begin{verbatim}
key=value
\end{verbatim}

同时 ini 文件可以将这些键值对分节，节名用 中括号表示，如

\begin{verbatim}
[section1]
key1=value1
[section2]
key2=value2
\end{verbatim}

这样的文件就有2个节。key2 属于 节 section2 . 不同的节下允许同名的 key 。


另外 ini 文件也支持 以 “;” 号起始的行作为注释。如


\begin{verbatim}
[section1]
key1=value1
[section2]
; 注释哦
key2=value2
\end{verbatim}

 }

%{ \chapter{历史上的派系斗争}\label{chap:wars}

\section{RISC vs CISC}

\section{大端 vs 小端}

\section{vim vs emacs}

\section{KDE vs GNOME}


 }

{ 
\newcommand{\manpage}[1]{
\parindent=0pt
\leftskip=0pt
\subsection{#1}
\parindent=-3.5em
\leftskip=3.5em
}

\chapter{man手册}

这里附上本书中出现的命令的man手册。各条命令以本书中出现的先后次序出现。这里只包含了本书中提到的命令和文件，并将其翻译。供英文吃力的同学作参考。


\section{第一部分：UNIX用户手册}


\subsection{test2}

\subsection{test}

\parindent=-3em
\leftskip=3em

\subsection{test3}

\section{第二部分：系统调用}

\section{第三部分：库函数}

\section{第四部分：设备和特殊文件}

\parindent=-3em
\leftskip=3em

\subsection{null,zero}


NAME \\
	null,zero ~ - ~ data sink

DESCRIPTION\\
asdfasdfasf
 }


\chapter{名称缩写}

\chapter{图片索引}
%\cleardoublepage
{
\renewcommand{\cleardoublepage}{}

\renewcommand\listfigurename{}
\vskip -3cm
%\addtocounter{chapter}{1}
%\addcontentsline{toc}{chapter}{图片索引}
\listoffigures 
%\cleardoublepage
}


\ifdefined\compileall

\chapter*{附录 \Alph{chapter}  Gentoo 的前身今世}
\addtocounter{chapter}{1}

\begin{notice}

声明: 本文的原始版本最初发表于IBM developerWorks，现在所有权归属Westtech Information Services。

本文档的中文翻译由\textbf{范华}完成。经范华同意，将本文以附录形式列出。

\end{notice}

\vskip 10ex

{\huge\bf Gentoo的前世今生，第一部 }
\addcontentsline{toc}{chapter}{Gentoo的前世今生，第一部}
\vskip 10ex


\section{Gentoo Linux发行版的诞生}

\subsection{我和linux}

现今对每一个linux爱好者（linux geek）来说，linux不再只是一个字面上的名称，她所呈现的一切对很多开发人员来说已经超过了他们所接触过的任何东西，linux比它们更强大、更令人着迷和称赞。当我在新墨西哥大学担任系统管理员时便与linux结下了不解之缘。那时因为我们的NT服务器运行得非常棒，我的手头上也有了很多空余的时间可以加以利用，就这样我的第一个linux操作系统（Debian）在一台Pentium 166的主机服务器上安下了家，接下来的不断、不断地学习和深入理解的过程使我对linux越来越着迷了……

一开始我学习了linux下的很多细节的东西：网络访问、执行备份、搞定samba等等。接着我建了一个qmail和apache的服务器并学习了python编程和shell编程。接着我还搭建了一个小型局域网接着就把linux娶回了家，在尝试过很多发行版后我最终选择了Stampede Linux这个版本\footnote{该版本从2001起就没有再更新了。}。你知道学习linux的过程是怎么样的吗？

\begin{itemize}
\item[]  第一: 努力搞清楚linux基本的东西
\item[] 第二: 当你已经有了相当好的掌握程度之后，学习定制你的linux，知识的累积会和你深入的程度成正比。
\end{itemize}

由于linux并没有隐藏任何东西，当linux对你来说变得越来越得心应手之后就可以开始探究技术和那些实现这些技术的工具了。


\subsection{Linux的潜能}

Linux提供了很多以前我所没有见到过的东西，如果一定要我用一个词来形容这些不可思议的话，我选择“潜能”这个单词：用来维护、改变、提高事物的能力，这种能力甚至能够冲破一些固有规则的约束。当我把kernel升级到一个更新的版本时，简简单单的就把我眼前的这个linux的性能提升了很多，更为令人兴奋的是这种改变几乎每时每刻都在进行着。而我也正是这种进步的一份子，伴随着linux的前进而不断进步着，对我而言这种感觉真的很棒。

如果你和我是同一类人，在你进入开源世界和linux世界之前大概看过位于Redmond和Cupertino的那些大公司们准备的下一代操作系统，它们确实如你所愿般的完美，然而那些东西却始终都只是一个虚幻的影子而已。然后就在我们慢慢等待的过程中linux来到了我们面前。虽然等来的这个精灵并不如我们预料的那么完美，但是她却提供给了我们这些喜欢动手hack的帅哥和美女一个亲手改变她的机会。就这样我们一边期待着一个更强大的操作系统，一边津津有味的hack我们的linux。日子一天一天过去，直到某天我们才突然发现原来期待着的那个强大的操作系统其实就在我们自己的手中，大家不约而同的笑了起来，也决定了继续在linux这条路上走下去，并继续hack下去。

\subsection{Linux的人文艺术}

我学到的另一件事就是Linux对人们的影响，这个话题可能听上去还真有点新鲜，是吧？Linux不仅仅只是一堆源代码的，它其实就是一个“社区”，从一开始的依赖这个社区解决我们提出的问题到付出我们的时间和经验去帮助他人，渐渐的我们也将会成为了这个社区的一部分。

IRC（Internet relay chat）既是一个交朋友的好地方也是一个很打发时间的场所。irc.openprojects.net上的\#stampede频道已经成为了我在网络上正式的安乐窝
\verb ^_^  。

那是我解答自己疑问的地方，也是第一次回答朋友问题的地方。\#stampede频道需要很多有安装经验的用户去帮助那些新手解决他们刚刚开始安装后碰到的各种各样的问题。由于那些新手在安装过程遇到的问题在irc中越来越普遍，原来很多有经验的Stampede Linux用户渐渐失去了他们一开始的热情。但是我依然还是很兴奋，因为很多菜鸟的问题我都知道解决的办法，要我忍着不回答那些问题我可做不到！当然我也并不是唯一的那个对解决新手问题乐此不彼的人，同样的家伙也有不少。我也承认自己也有那么点私心，想从那些更有经验的家伙们（不是指Stampede的开发人员）身上学到更多的东西。

\subsection{如何起步}

当有朋友问我如何才能加入一个开源项目时，我告诉他们的是首先是找一个能为他人做些什么的地方，就算那里只是解答一些很基础的问题。一份诚挚的渴望帮助他人的愿望是通往Linux社区的通行证，因为这份诚挚的愿望同样也扎根在每一个开源项目开发人员的心中（不仅仅只对Linux项目来说），它也应该扎根在那里。

沿着这条路走下去不可避免的你会遇到比你更有经验的同志，你将会从他们身上学到更多的知识，就像以前新手从你身上学习时一样。另一方面，当你积累起了更多的经验后在碰到某些问题时你就会用一个新方法去解决它而不是用以前惯用的一套思路。你遇到的一些开发人员有时会提出一些建议，有时又或者会需要一些帮助，他们更可能会邀请你加入他们的开发队伍；如果你的助人为乐成为焦点时，他们可能会笑着从你身边经过；如果你帮助了很多很多人之后，你在社区内肯定会备受瞩目。在Stampede和我身上这些故事都曾经发生过。

渐渐的我在Stampede的开发越来越深入，不久以后我就成为了一个正式的Stampede开发人员。在受到了Stampede的领导者Matt Wood的鼓励后，我开始对用于Stampede Linux软件包的原有的.slp机制进行升级。当时，.slp软件包格式包含一个.tar.bz2的软件包和后面的一个包含软件描述及软件包创作者等等在内的一个定长的页脚。这种实现的方式有两个主要问题：页脚部分实际上包含的内容根本达不到定长所约定的字节数；该格式没有预留任何扩充余地（也就是说如果未来没有办法加入一些可能需要的额外信息）。显然这些问题需要动一次大手术了，hoho。

和那些老资格的Stampede开发人员工作一段时间后，我拟了一个解决上面那些问题的草案。过了一阵子我便开始用Python先编写了一些原始的实现方案，新的格式(代号slpv6）有些类似与Amiga世界的IFF格式。下一代的.slp格式包含了了2的32次方个字段，字段种类为2的32次方种，每个字段最大数据段同样为2的32次方字节。新的格式不仅具有良好的扩充性而且比纯文本更加紧凑和简洁并易于解析。二进制代码和文本都能存储在这样的格式当中，该架构对其本身在未来的进一步发展带来了无限的可能性。我的想法是把这个新版的动态header加入到打包文件的结尾部分，从而这个新版本的.slp格式未来可以为Stempede用户服务相当一段时间并且同时又能和标准的UNIX档案文件保持不错的兼容性。

\subsection{丑陋的一面}

slpv6的开发进展很顺利，所有的资深开发者看到我取得的成果后都很高兴。不幸的是，两名刚加入的Stampede开发者想要自己掌控slpv6项目。由于不欣赏我选择的开发方向，他们花了很大劲诋毁和打击这个新的slpv6系统，虽然我也用了大量时间一边继续我的开发一边加入讨论再一边回应他们的攻击，但是这样做也没从根本上解决问题。最后一切都变的很明了，他们只是很擅长辩论，并且显而易见的是除非是走他们自己的路子，不然他们是不会罢休的。幸运的是我的项目依然得到了资深开发人员的认可和支持。可是这些讨论渐渐地使我背上了一些包袱，同时对Stampede的开发也产生了一些不好地影响。唉……

可惜我没办法使这些家伙消失，原来还可以在\#stampede频道里和那些高级的开发者互相交谈，但是现在不得不退了出来。每次只要我一进入那个频道，他们就开始变得很不友好，总是在破坏我想要进行的工作。这些家伙会使用各种各样的方法：比如一个开发者会议（其实只是想当着其他资深开发者的面侮辱我）。他们还尝试用投票的方法控制Stempede，当然那种投票只在他们可以得到更多支持的时候才会举行。但是自始至终我在这样的情况下都没有放弃过我的slpv6的开发工作。不用多说，资深开发者们都喜欢我的开发项目，并且他们也都支持我继续做下去（没有他们的支持，我不可能克服那么多困难并坚持下去）。

\subsection{对这些异类的了解}

我习惯于把这两个家伙和这种类型的开发者称为“异类”。虽然我的开发工作因此变得很很不愉快，但是我还是学会了怎么样去对付他们。就这点我乐于提供给各位一个对这些“异类”的全方面的描绘：他们的品质、采用的方法以及当你作为一个项目领导者怎么样才能对抗这些”异类“或是尽可能的用最小的代价去改变他们。

为了消除情绪上可能存在的危险，你需要具备一个先决条件：意志力。如果你不能用一种既礼貌又态度坚决的方式回应你的对手，事情就会变得很糟糕。“异类”的目的就是尽可能多的在你的项目中取得控制权，这么做会使他或她感觉更具有力量。首先，他们会对某个项目或是项目的开发人员进行片面的指责和抱怨，同时他们也会阻挠那些对这个项目富有建设性的提议。当然这些家伙在他们获得项目管理人员位置之前是不会对这个项目伸出任何的援手。目的就是使你确信只有依靠他们的那些“独道的、富有素养”的眼光才能最终解决问题，这样你就不得不给他们足够的权限去实现这些。

如果指责和抱怨没起什么作用，这些“异类”就会要求举行一个开发者会议。这将会给他们一个可以分裂你开发团队的机会。在觉得本方这方面已经得到了大多数人的支持后，他们就会举行一次投票决定（当然他们知道赢的会是他们的情况下）。如果并没有赢得投票或是投票被驳回，那么下周他们还是会提出举行一次会议以便再一次的分裂你的团队，然后再是那种无休止的循环。

如果会议的方法行不通，“异类”们将会变成革新运动者。他们会用一种更民主（也就是更容易操纵）的办法来取代先前压迫性的和非公平的决策方案。这些办法常常包括令人信服的让你去为你的开发团队中的大部分人做任何事。异类比较偏爱这个办法，因为你没有办法弃大多数投票表决的结果于不顾（活活活）。你许可这些事情发生的时候就已经把那把通往你的“Lexus”的“钥匙”交到了他们的手里，这将使你感觉越来越无力。

“异类”们用的另一种方法是激怒你的主要开发人员并使他们离开，然后在你的开发团队混乱的时候尝试重新组织该项目的管理团队。如果所有的努力都没有成功的话，他们会聚集尽可能多的叛离者并把他们安插在你的项目中，痛啊！

\subsection{对付这些异类}

区分这些家伙还是相当容易的。他们不会写一行代码（也不愿意写），相反他们会花大量的时间讨论那些更重要的问题（对了，就是那些管理方面的问题）。
假设你是一个项目管理者，对付他们非常容易。只需要告诉他们，在没有看到高质量的代码之前你是不会考虑他们所谓的建议的。
或者在他们提出“建设性“的批评之前强调对于某个项目有建设性得帮助也包括服从项目的管理人员。
如果他们开始编制优质的代码并且越来越有易于这个项目，那么就太好了。如果没有，就告诫他们离开。
在你忽略这帮家伙一段时间后，他们会选择离开或是一边采取行动一边写一些代码，世界就这样清净了\verb ^_^ 。

不幸的是Stampede的那些资深开发人员对”异类“并没有采取更多的管理措施。换句话说，他们许可了这两个家伙对我（和其他人）的无休止的纠缠。虽然这些资深开发者总是赞赏我的项目，但是对那两个家伙他们却并没有做的更多。然后终于有一天我决定制作一个自己的发行版，因为我觉得这样做比忍受那两个家伙更容易些。我退出了Stampede的开发团队并开始制定自己发行版的一些计划和草案。

一段时间之内，我对自己因为两个低等级开发者而离开一个项目还是感到有些不可思议。其实他们没有涉及到的实际情况却真正显示出这个项目存在很严重的管理方面的问题，如果高等级的开发人员不能或者不愿意确认Stampede的开发成果是可喜的和有益的话，我想我不会愿意继续留在那里。

\subsection{新的开始}

离开Stampede后我做的第一件事就是长长的舒了口气。喔……，整个世界都清净了。现在我有了足够的时间来思考我自己的Linux发行版的轮廓和将给Linux发行版的布局带来什么新的贡献。对Stampede感兴趣的一件事是它所具有的原生的性能（这得感谢它使用的带有实验性质的、并针对Pentium处理器优化过的pgcc编译器），所以我决定首先我考虑的就是性能。除了更少的CPU占用率以外，我还希望它更精简。很多发行版本（特别是那些流行的热缩塑料封装的家伙）默认启动了太多的daemons以至于打开一个xterm（X环境下的终端）后系统所剩余的可用RAM已经所剩无几了。我希望自己的发行版能更小也更强，为此我把目光放到了最大限度的榨取让这个操作系统运行的硬件平台的性能上。为此我下决心进行一个整体测试并处理掉所有细节中的性能方面的问题。

但是我真的很缺乏对应的资源，因为我是这个发行版的唯一的一个开发人员！我该怎样做才能只靠自己就鼓捣出不逊色于Redhat或是Caldera这样的产品呢？解决办法是采用自动控制技术。我必须写一些脚本以便所有的事情都可以自动搞定，这样我就可以事半功倍了。毕竟，电脑们这些方面做得更好，对吧？

很快我发现光是写一些自动化的脚本还远远不够，需要设计的是一整套能从源代码产生一个完整Linux系统的机制。我实验性的把它称做ebuild系统并且开始了工作。ebuild系统可以自动的建立所有一个发行版所需要的二进制文件，包括从解压源代码并打好相应的patch再到编译、封包的一系列过程的自动化解决方案。在一个基本、原始的ebuild可以工作后，我开始为一个Linux发行版必要的一些关键组成部分（像是gcc、glibc、binutils、util-linux等软件）撰写ebuild脚本。通过重新撰写初始化脚本（基于以前我为Stampede设计的初始化脚本）把原先的Stampede开发系统逐渐的演变成一个我自己的系统，接着我开始测试每一个自己建立好的新的软件包。

几个月之后我有了一个完整的，自主的Linux版本。我给她起了个名字『Enoch』然后坐着满足的笑了起来。但是什么改变了Enoch、Gentoo的发展又是怎么样的？续篇将会告诉大家Enoch是怎么演变成Gentoo的，而我在这条路上将要面对的更多新的挑战。

\vskip 10ex
{\huge\bf Gentoo的前世今生，第二部 }
\addcontentsline{toc}{chapter}{Gentoo的前世今生，第二部}
\vskip 10ex

\section{从Echo到Gentoo}

\subsection{Enoch踏出的第一步}

我在先前的文章中告诉了大家那段和Stampede开发团队在一起的、曾经最兴旺的时光和最后为什么离开的原因（就是想离那些有低级政治目的的、想控制项目的那帮家伙远点）。因为这些爱管闲事的好事者的干涉，我才会觉得装配一个自己的Linux发行版比在那种恶劣条件下改进Stampede要简单的多。幸运的是，我离开Stampede时是带着满满当当的经验离开的，这些经验与在Stampede的工作（应该是实质性的吧？）是分不开的，维护一些软件包也好、设计初始化脚本也好或是领导slpv6（下一代软件包管理系统）都使我相关方面的知识和经验得到了极大的丰富。

Enoch是我开始工作的这个版本的一个代号，得益于为它开发的高智能的包管理和升级系统，它将会是一个速度飞快的版本。我不得不承认这套智能化的系统在整个版本中占据了很大一部分位置，因为对于我这个光杆司令来说在那种重复性的劳动中消耗时间是没法接受的，所以才会要求开发中的系统必须自动为我完成那些琐事。另一方面完全由源代码来构建整个发行版（比那些“spin off”的版本、例如RedHat要好）也需要把工作划分好并尽可能多的挤出空闲时间来做这些工做。

使最基本的Enoch系统启动和运行之后，我回到了irc.openprojects.net并开设了自己的＃enoch频道。在那里我逐渐聚集起了10个开发人员组成的团队。在早期的那段时间里我们整天都聚集在IRC里，用空下来的时间制作我们的发行版。在我们无私的付出和大家的齐心协力的hack下，在不断的消除bug和新的bug的过程中，Enoch每天都在变化着，不管是专业化的程度还是各方面的功能都变得越来越出色。

\subsection{Enoch的第一块绊脚石}

不可避免的一天，Enoch碰到了它的第一块绊脚石。在加入了Xfree86、glib、gtk＋之后，我决定把xmms（一个基于X11/gtk＋的MP3/CD播放软件）弄进我的发行版，因为也该到了用音乐来调剂调剂的时候了！但是在安装好xmms之后启动它时......X死锁了！最初我觉得是自己使用的编译器的优化参数造成的（"-O6 -mpentiumpro"，在你看来有点诧异吧？）。第一个想到的解决办法就是用标准的编译器选项来编译，但是问题依然没有解决。然后只好到处寻找解决方法，接下来整整几个星期的开发时间我都用来追踪这个错误。一天，我收到了一个叫Omegardan的Enoch使用者的电子邮件，他也同样碰到了xmms的这个死锁问题。

交流了一段时间然后历经了n个小时的检测后我们发现死锁的原因在于POSIX的线程描述符（POSIX threads-related issue）。因为一些原因，pthread\_mutex\_trylock()函数没有返回它应该返回的值。作为一个Linux版本的创始者，这种类型的bug是我真的不愿意碰见的家伙。我指望开发人员能能够释出完美的源代码以便我可以把精力放到提高Linux易用性上，而不是把时间花在修复别人源代码的bug上。当然很快我就发现这种希望仅仅只是一个美好的想法罢了，相同的错误有时还是会出现。

在找到问题后，我们发现它不是xmms本身的问题，也不时gtk＋或glib的问题，也不是Xfree86 3.3.5没有thread-safe和死锁的问题，而是令人惊异的存在于Linux 的POSIX的线程执行本身，具体来说就是版本2.1.2的GNU C库（glibc）的部分代码中存在bug。我很震惊的是在Linux如此核心的部分居然存在这样严重的bug（而且我们为Enoch使用的glibc的版本是它的release版本，并不是什么prerelease版本或是CVS版本！）。

那么怎么样才能解决这个问题呢？我们不可能马上就能拿出一个修补方案，但是在浏览了一堆glibc开发人员的邮件列表后，我偶然发现了还有一个人也碰到了相同的问题，然后在glibc开发人员在回复他的邮件里我们找到了那个附带的补丁，它为我们解决了那个线程问题。但我令我好奇的是为什么同样使用glibc 2.1.2的RedHat 6没有受这个bug的影响（当时RedHat 6的发布时间先于那个补丁的出现）。为了找到答案，我下载了RedHat里glibc的SRPM包（source RPM）想看一下他们使用的补丁是怎么样的。

RedHat有他们自己的glibc补丁来解决pthread\_mutex\_trylock()函数的问题。显而易见的是他们也碰到了同样的问题，然后自己进行了修补。但是由于RedHat没有把这个补丁回馈到glibc的开发社区，其他人们就没有办法分享这个补丁。但是也可能是RedHat把这个修补方案回馈到了glibc的开发社区，然儿glibc的开发人员并没有接受这个修补方案。 或者这个bug只会在特定版本的binutils和特定版本的编译器连用时才会触发，然而RedHat使用的binutils和编译器的版本并不是这两个特定的版本（虽然RedHat还是给出了这个补丁）。我猜测我们永远也不会知道究竟事情的真相是什么样的，但是我学会的一件事情是：RedHat的SRPM包里有很多定制的补丁和增强代码，而这些代码和补丁看来从来没有回馈到原始的开发人员那里。我将会为此来上一段激昂的演说。

\subsection{激情的演说}

当你将一大堆各种各样的源代码汇聚成一个Linux发行版时，把所有你做好的bug fix和补丁反馈给原始的某个软件包的开发人员是一件相当重要的事情，就如我了解到的那样，这是发行版的开发人员为Linux做贡献的很多途径中的一个。我们也恰好就是这样的一群人，为的就是把很多不同的程序和软件集合在一起，让它们工作起来就像是一个整体。将来我们也会把我们们对一些软件所做的修改和补丁反馈回原始软件的开发人员以便其他的用户和后来的发行版能从中受益。如果你只是把补丁留在你自己那里，这样做不会对任何人有什么帮助，很多人们将会为一些相同的问题浪费掉大量的时间。这种不顾别人的方式违背了整个开源世界的精神和宗旨，同时对Linux的发展也只是有害无益。或许我应该说这样的做法对我们来说就是一个大大的“BUG”。

不幸的是一些发行版（啊咳）（RedHat）并不如其他一些版本（Debian）那样对整个开源社区分享他们的成果。

\subsection{编译器的艺术}

在我们尝试解决glibc 线程问题的时候，我给Ulrich Drepper发了封email（他是Cygnus的一员并且在glibc的开发中举足轻重）。我在e-mail中提到了我们碰到的POSIX线程问题和我们在Enoch中使用pgcc来获得优化的性能。在他的回信中他这样提到（我解释一下）：“我们自己的包含在CodeFusion中的编译器制作的可执行代码比其他的一些编译器、比如pgcc编译出来的代码执行速度更快速。”显然我对测试测试Cygnus那帮家伙开发的神秘的“turbo”编译器非常有兴趣。

因此我申请拿到了一个Cygnus Codefusion 1.0的demo拷贝以便我可以对它的性能做一个测试。Omegadan和我对测试的结果很吃惊，它同Ulrich提到的那样出色。x86的后端提高了90％的有关cpu-intersive的可执行文件的执行效率（比如bzip2）。几乎每一个程序都能从中获得至少10％的真实世界的性能提升，而我们所作的仅仅是换了一个编译器。Enoch的速度也因此获得了30％－40％的提升。同时性能也提高了不少，提升的幅度超过了我们以前把编译器从gcc切换到pgcc时提高的幅度。显然，在对这个编译器的测试后，我们很希望把这个编译器包含在Enoch中，有点幸运的是CodeFusion CD中的包含的源代码遵循的是GPL，这样在Enoch中使用这个编译器已经可以算是已经得到了完全的认可了..........，至少我们是这么想的。

\subsection{异常事件的发生}

为了能在Enoch中使用这个编译器，我给Cygnus的市场部主管发了一封电子邮件，但是期望中的“哦，拿去用好了，感谢使用我们的编译器！”这样的回复并没有收到，取而代之的是一句“虽然在技术上我们许可使用Cygnus的编译器，但是我们强烈建议不要在在Enoch中使用该编译器或是包含它的源代码。接着在我的回复中我问了他们这样一个问题：“既然不愿意让别人使用它的源代码，为什么还在以GPL的许可条例来发布它的源代码？”作为一个猜测，我觉得他们事实上是不想以GPL的方式来发布他们的源代码的，但是由于这个编译器是源自egcs（以GPL方式发布的），他们除了以GPL方式发布之外别无选择。

这是当某一个公司想使用开源的代码来生产私有产品这样的情况时，GPL如何阻止这样的事情发生的一个很好的例子。我比较有根据的一个猜测是Cygnus担心我们使用这个编译器后将会打击到他们整个产品框架的销售，更加奇怪的是不管是他们的行销方案还是InfoWorld的预览中都没有提及包含在CodeFusion中的那个新的编译器，因为CodeFusion销售的是一套“development IDE”而不是一个编译器。

为了缓解一下他们那种偏执的态度，我提出了个建议，就是在我们的Enoch主页上放置上CodeFusion的签注文件并加上一个链接来刺激CodeFusion的销售。从我个人的观点来说，我不认为一个“turbo”的Enoch会影响到CodeFusion（虽然它是一个IDE产品）的销售情况。但是我还在想方设法的令到他们愉快，比如告诉他们这个IDE的组件是一个商业化的产品，我们也并没希望或者有什么意图用Enoch来发行它。

我把这个（大方的）请求用电子邮件的方式发给了Cygnus，但是收到的确实另一个奇怪的回复。他们想得到所有我们关于“市场元素”方面的具有权威的权利（显然，这也包括了我们网站上的内容），真是太令人震惊了。Cyguns的营销团队似乎对Linux社区和GPL的运作一无所知，事到如今我终于决定终止与Cygnus彼此间的联系，因为再这样下去事情会变得怎么样谁都不知道。与此同时，我们为Enoch准备了两个版本，一个是内部的“turbo”版，一个是公开的“non－turbo”版，其实就是把决定留在将来再去做。

但是几个月之后，他们就把CodeFusion x86的backend换成了gcc 2.95.2，现在不只是那些知道包含在CodeFusion CD中的“隐秘的GPL编译器”的这群人可以获益，几乎每一个人都可以从这个新的优秀的backend中获益了。然后我们还是决定继续前行，尽量使用gcc来替代CodeFusion的编译器。在gcc 2.95.2已经越来越成熟的情况下，我们已经可以放开Cygnus了（同时，RedHat却为购买这个CodeFusion而花费了比较冤的一笔钱了。）（注：新的x86版本gcc 2.95.2的backend为新的Linux发行版提供了一开始我们提到的很重要的速度提升，它也为FreeBSD 4.0相对3.3.6版本速度上提升做出了很大的贡献。你注意到这两个提升的不同点吗？）

\subsection{肥皂盒}

感谢这件事情和其他的一些经验，我从中对那些以开源为主要获利手段的企业有了很深的理解。虽然对个人来说，乐于生产私有闭源软件这件事并没有任何错误的地方，但是一个开源企业搅乱或是拒绝与其他的开源世界合作是没有任何意义的；同样，不支持GPL或是其他的等等也没有什么意义。这是一个实践性质的并具有现实意义的观点。

思想和代码上自由的交换才是开源企业得以获利的根本，这点他们应该充分的认识到。反过来，对立与GPL标准只会破坏这个他们依赖于发展与繁荣的环境。换句话说，开源的环境是你事业的土壤，保护这片土壤的纯净还是很有意义的。

我也懂得在短时期内保留一些代码上秘密来获得财富的累积是一个颇具诱惑性的东西，先进的代码和特别的技术提供给了人们一个在竞争中获得优势的绝好机会，由此可以获得增长的销售业绩和利益。但是当你的目的是成为一个唯一的产品提供者，而这个产品商业的成分大于开源的成分时，开源世界是不会许可这样排外性质地使用开源或是相关东西的，这就是开源的意义。

\subsection{回到Enoch}

现在，我从自己的肥皂盒中出来并继续我的故事。 

由于Enoch已经变得越来越出色，更名的计划也就这样列入了我们的议事日程当中，接着“Gentoo Linux”诞生了。然后就是朝Gentoo Linux 的1.0版本努力前进中。大约也是这个时候，我决定帮我那台Celeron 300M（超频到450M并且十分稳定）的老电脑升级一下，新平台是一块崭新的Abit BP6主板（从市场上找到的双Celeron接口的）。在卖掉了老主板后我把我两个Celeron 366的系统集中起来，然后把Celeron 366超到了500Mhz，然后开始工作了。但是我注意到我的新机器不是非常稳定。 

显然我第一个反应就是把频率改回没超之前的366Mhz，但是随之而来却遇到了一个更奇怪的问题：不管CPU全速运转多少时间，系统都不会死锁；但是一旦空闲下来过一夜的话，系统有很大的可能就会完全死锁掉。是的，这是一个idle bug－－－－噢！在作了一些调查之后，我发现在这块主板上也有其他用户碰到了这个相同的问题。原因是BP6主板上的一个芯片（可能是PCI控制器）与标准规格有点不同或是比较古怪，这个东西就是造成Linux在空闲时候死锁的主要原因。

我渐渐的心烦意乱起来，因为我没法再去采购另外的PC部件了，Gentoo的开发也只好被迫终止下来。我也开始对Linux越来越有些悲观的情绪了并决定转向FreeBSD。是的，的确是FreeBSD！这部分就此为止了，我们第三部再见了：）

\vskip 10ex
{\huge\bf Gentoo的前世今生，第三部 }
\addcontentsline{toc}{chapter}{Gentoo的前世今生，第三部}
\vskip 10ex

\section{离开Linux后再度回归}

在前一篇文章的结尾部分，我说到因为新升级的双Celeron主板（Abit BP6）存在一个古怪的空闲时死锁的问题导致Gentoo开发停止。虽然解决问题的办法就是更换主板，但是我已经没有重新更换主板的资金了，这件事也打击了我对Linux的信心并使我决定中断Gentoo的开发并转向了FreeBSD。我需要的是一个可以正常运转的系统，而Linux在这个时候的表现并不尽如人意（一天到晚的死锁），那个当口，我觉得是好好接触接触FreeBSD的时候了，便在机器上安装了FreeBSD后开始了又一次的捣腾，在接下去的几个月中，我也几乎没有再碰过Linux一个指头。

\subsection{FreeBSD之印象}

首先，我真的很喜欢FreeBSD。我感觉这个操作系统是一个组合的很完美的系统，它的几乎每一个部分都同样精巧，而这种精巧的在Linux世界中几乎不存在。我的满意实质上是来源于那些FreeBSD中非常充足的man page，这可不像Linux里那些只有GNU info文档的很多软件那样让人根本没法用。

最最重要的是我对FreeBSD中维护与升级系统的ports系统印象非常深刻。与Linux维护与升级的方法不同，ports使用的不是二进制的软件包而是直接去原始的软件站点下载所需要的源代码并编译。不管你是安装Samba或是升级核心系统都是在你的机器上用源代码编译而成。这样的实现方法和我在Gentoo Linux中建立的那套机制有着异曲同工之处。从这点和其他许多方面来说，FreeBSD的这种设计符合我作为一个开发人员和一个系统管理员所期望的那种感觉。就这样，FreeBSD为我营造了整整几个月舒适的工作环境，同样我也很乐意于花些时间在这个出色的操纵系统中探求与获取知识。

\subsection{FreeBSD的优点}

很多Linux和FreeBSD之间的不同点都是源自与它们本身开发架构的不同。Linux的开发架构非常松散，我们只是依靠不同的发行版把分散在Internet上呈离散状态的很多部分组合成一个完整的Linux，而FreeBSD和其他BSD系统（OpenBSD和NetBSD）都有一个唯一的核心小组来确保源代码的单一性和协调性，这样至少每一种BSD自身都拥有一套统一的源代码设置。这是一件挺棒的事情，也是FreeBSD感觉上和Linux那种“patch集合”有所不同的主要原因。

接下来，我们在纯技术方面再作个比较。很多FreeBSD的粉丝都声称FreeBSD比Linux更合适用作服务器上跑的操作系统，他们会告许你在高负载情况下FreeBSD表现得更好，而且它的TCP/IP栈相对出色一些（如果你用Linux 2.2或更早版本的内核和FreeBSD作比较，我同意这个说法）。FreeBSD确实是一个很好的服务器操作系统，这点勿庸置疑，但是这只是FreeBSD相对Linux 2.2或更早的内核版本时的情况。我作为一个新版本内核的粉丝，早就在我的电脑上用上了2.4测试版的内核，它确是也很棒，从出色的TCP/IP栈到整个重新设计的“netfilter”系统都是。我觉得在不久的将来，新的性能标准将会由Linux来定义，而“free UNIX”将会在商业领域面对Linux强有力的挑战。

\subsection{FreeBSD的不足}

与服务器领域的应用不同，在桌面应用上，Linux占有绝对份额上的优势（仅相对BSD来说，Linux不管是对Win还是对MAC都完全处于下风）。所有最新的桌面应用软件一定是先在Linux上出现、在3D加速和声卡的支持方面，Linux也比BSD走在了前面。随着2.4版本内核的临近，Linux在这块地盘上还是会继续保持它的优势地位。

我对FreeBSD采用的UFS文件系统并不喜欢，虽然UFS相对Linux的ext2文件系统来说更健壮，但是付出的代价是那个另人昏昏欲睡的龟速。现在也有一个UFS文件系统的扩展叫“soft update”，它是把小块的IO操作聚合成大的文件块后再写入物理硬盘以提高文件系统的速度，就算“soft update”这套机制大幅提高了UFS文件系统的性能，我也没法就说在所有方面的比较中UFS都比ext2优秀。当然，UFS和“soft update”更加可靠，FreeBSD也可能会在文件系统的战争中击败Linux，但是请不要忘记，输给FreeBSD的仅仅只是现在的2.2版本或者更旧版本的Linux，这不代表将来也会。

现在，我们把话题转变一下，我们比较的双方是现今的Linux 2.2版本、2.4版本和FreeBSD。Reiserfs（一个新的日志型文件系统）已经给我们带来了一阵惊喜，而Linux还有蓄势待发的ext3、IBM的JFS和XFS文件系统，这些文件系统都在提供高可靠性的同时提供了优秀的性能。Reiserfs给了Linux在文件系统上超越FreeBSD的一个契机，这也是我认为Linux 2.4版本会上演大逆转的原因，FreeBSD的传统强项在未来2.4内核面前可能会荡然无存。

\subsection{回到Gentoo的开发}

几个月之后决定重新回到Linux世界的我在一台新的机器上又装了Gentoo。首先，回到Gentoo的开发中来是一个计算后的决定－－我已经花费了很多时间使自己成为一个Linux的万事通，而现在怀抱着BSD就等于是把以前学到的知识都浪费掉了，这样做我觉得不是很值得。而且在更新Gentoo Linux后那么一段很短的时间内，我为“为什么再次回到Linux怀抱”找到了几个新的理由，也就是前面提到过的kernel以及文件系统的改进等等。FreeBSD是一个宁静的家园，但是这样的宁静太安静了点，这样的宁静也包含着困惑。相反Linux世界充满着活力，发展也是日新月异。如果你所寻找的是兴奋和创新的地方，那么毫无疑问Linux就是你所向往的世外桃源。

Linux从2.0进步到2.2给我的感觉就是满失望的，但是2.4时代是绝对值得去守候着的，为此Gentoo Linux重新回到了我们面前，那种兴奋的感觉也重新回到了我的心中。

Gentoo Linux重生的另一个关键因素是我们开发团队的领导者－－Achim Gottinger。我想花一点篇幅对他所给予的帮助（使我我重新开始了Gentoo Linux的开发）致以诚挚的感谢。我在回到Linux世界之前就开始与Achim Gottinger有了电子邮件上的往来，在几乎每一封他的电子邮件中，我都可以看到一些新的.ebuild或者是些迫切需要修复的bug。在我回到Linux世界并重新开始了Gentoo的开发之后，Achim继续贡献着他的时间和精力使这个发行版步入正轨。直到最近，Achim和我都是Gentoo Linux仅有的两个开发者，这也是出于选择的结果。因为我们都使用几乎相同的发行版，也因为Achim的技术，我们可以轻松的完成非常巨大的工作量以至于我觉得加入第三名开发者并不会对我们的进展有什么帮助。现在Achim是Gentoo Linux开发组的负责人，几乎每天Gentoo的都会有基础部分中主要的提高。我们已经走到了这里，也已经准备好了CVS树为后来者提供一个协同开发平台，小心翼翼的逐步扩大Gentoo开发队伍的工作也开始付诸实施。

\subsection{新的版本}

我没有觉得花在BSD上的时间是在浪费。实际上，它给了我一个很好的机会来反省一下整个Linux社区存在的问题和Gentoo Linux应该做点什么来改进这些短处。. 

在新版本的Gentoo Linux中，我下决定不再使用pgcc或者什么非常优化的参数来编译所有的软件包，因为稳定性还是要放在第一位的，我们默认将会使用合理的优化选项（"-O2 -mpentium"），但也同时向用户提供了可以简单自定义的优化选项来满足了一些同胞希望得到最“bleed edge”的系统（通过我们的自动化系统完成）这么个愿望。

FreeBSD给了我一个关于“自动化定制系统如何工作？”这个问句一个很好的提示。我决定在我们的自动化定制系统（现在叫做Portage）中加入一些FreeBSD的特性来制作一个新一代的ports系统。 

Portage 可以说是Gentoo Linux的心脏，它所具备的东西远远超过一个简单的包管理机制或是一个系统管理机制。Portage通过它包含的对制作工具的设置和制作脚本可以使你从源代码构建一个完整的发行版系统。但对我来说更重要的是，Portage给用户提供了一个可以完全接触Gentoo Linux构建智慧的途径。对我们开发者来说，这意味着当Gentoo Linux不断发展的同时我们也记录下了一个发行版制作的过程。Portage的易用性和可读性也为越来越多的人提供了一个窥探Linux内部的窗口，它也为后来者贡献他们的代码和脚本打开了方便之门。

Portage是我们为他人展示Linux技术和原理的一条途径，通货学习自动化制作脚本，你可以看到大量各不相同的包是怎么互相适应并结合成一个整体的。如果你需要，你也可以从我们的站点上攫取整个CVS树然后自己hack并制作个人的Linux发行版。我们坚信这是一件好事情－－我们希望把知识交给渴望这些知识的人们以便他们可以把Linux带入一个新的领域。

\subsection{商业上的关注}

起初，有许多拥有不同背景的人们加入了Gentoo的开发中来。因为这个，我们的开发人员对于如何最终在Gentoo上获得经济利益也有许多各不相同的打算，对此我并没有太多的诧异。基本上有这么两种类型的开发人员：一类群体反对用Gentoo来追名逐利，另一类群体则对使Gentoo Linux成为一个成功的商业产品非常感兴趣。这是一个预料中会存在分歧的地方，第一类群体认为商业化的运作包含着腐化等不良的影响，而第二类群体则认为没有这么多的负面因素。

在以前还是Enoch的那段时光中，我对商业成份究竟有利还是有弊这点也很难做个了断。我验证过的是像Debian这样的Linux发行版真正忠于“自由”这样的事实，我喜欢这样。对比其他商业化的发行版，他们给用户带来的易用性包括了在各自的网站上提供一份完整的安装说明，这也是一个我想去借鉴的好东西。

同样，我也真心希望Gentoo Linux能够成为一个成功的商业版本，为了这个目的，我努力想在商业和开源之间找到一个平衡点，可是直到最近我还是没有能够找到这么一个黄金分割点。

\subsection{该做些什么}

我们该怎么做才能在商业化和非商业化中取得平衡呢？关键的一点是一定不能忘记我们的基楚和根本－－－Gentoo Linux 作为一个开源软件的根本和基础。所有我们作出的努力都必须遵循这个基础，这不仅仅是肯定开源软件或只是使用开源软件，还是对开源软件和开源发行版开发的鼓励和支持，也不会发对用这样的一个对待开源姿态来获取商业回报。更重要的是，我们绝不会采用商业化的模型，因为这样做对于其他发行版使用我们的源代码有阻碍作用。我们的开发团队对所有人来说都会是开放的和可接近的，而Gentoo Linux这个自由发行版不仅仅可以被大家接受还会因为很多人的鼓励而继续走下去。我们必会成为开源运动的倡导者，一个把这个理念贯彻到行动中而不是停留在文字层面上的倡导者。

如果某公司需要为一个商业化的基于Linux技术的需求使用Gentoo Linux，他们可以从我们的CVS树中攫取这些代码并马上开始使用它们，因为所有我们的分散的工作都是基于GPL。我们在确信所有基于Gentoo Linux的衍生产物都遵循GNU Public License的前提下是不会在任何地方限制别人使用我们的代码的。

我们希望有尽可能多的人们从我们的工作中受益，但是我们也希望尽可能多的能从你对Gentoo Linux的提高中获益。如果你公司的产品有很大一部份是基于Gentoo Linux的话，希望你可以把所有可分类的修改和提高发送给我们以便加入到CVS树中使更多的人获益。继续保管和改进你提交的修改后，你也能从我们所做的修改中受益。我们也鼓励商业实体和非商业实体之间的合作，举个例子来说：不管是在他的ISP中使用Gentoo Linux的系统管理员还是用Gentoo Linux构建商业服务器的公司都能从彼此对Gentoo Linux的改进中获益。是时候来促进在人们之间的自由代码交换了，这也只有开源软件可以做到。

\subsection{将来要走的路}

现在离Gentoo Linux 1.0 的发布已经很近了
（在你在developerWorks上读这篇文章的时候它可能已经发布了，想想现在的2006.0是不是大家有种沧海桑田的感觉
\verb ^_^ ??
）\footnote{作者在IBM DeveloperWorks上最初发表的时候，距离Gentoo Linux 1.0的发布还有时日}。但是Gentoo Linux将来的方向会是怎么样的呢？

当我们逐步迈向2.0版本时，我希望继续提升Portage作为Gentoo Linux核心的性能，因为任何关于Gentoo Linux主要的进步都会从Portage的进步开始。主要代码从bash转换到python的过程我也会继续下去，因为这么做会使我们加入新的设计（比如为我们的全自动构造系统设计的面向对象的新东东）。

除了Portage的修改，我还希望小心谨慎的寻找技术出色并且和我们使用相同版本的开发者加入我们的开发团队。在扩大了开发团队之后，我们可以为Gentoo Linux的加入更多的自动化定制脚本。比这更重要的是，适当扩大的开发团队可以使Gentoo Linux站在Linux技术的尖锋之上，这才是乐趣所在嘛：）

我们也希望商业化的Linux技术公司可以把Gentoo Linux作为他们产品的基础。现在我们已经有了这样一个关系，将来也会更多的，而这样的协作承诺充满着乐趣并对于Gentoo Linux的用户非常有益。

最后我要说的是，我们主要的目标是为Linux社区提供有意义的贡献。虽然可选择的发行版很多，但是Gentoo Linux还是拥有许多其他版本所没有的东西。我们对未来Gentoo Linux发展充满着信心，我们希望你也有同样的感觉。

\vfill

\begin{center}
{\Huge 全文完}
\end{center}

\vfill
\else
\addtocounter{chapter}{1}
\fi

\end{document}
